\documentclass[11pt,reqno]{amsart}

\usepackage[T1]{fontenc}
\usepackage[utf8]{inputenc}
\usepackage{lmodern}
\usepackage{microtype}
\usepackage{amssymb,amsmath,amsfonts,amsthm,mathtools,mathrsfs}
\usepackage{aliascnt}
\usepackage[dvipsnames]{xcolor}
\usepackage[margin=1in]{geometry}
\usepackage[shortlabels]{enumitem}
\usepackage{booktabs,tabularx,array}
\usepackage[colorlinks=true]{hyperref}
\hypersetup{
  linkcolor=blue!45!black,
  citecolor=magenta,
  urlcolor=blue!45!black,
  pdftitle={Attractive log gas I: partition functions and equilibrium fluctuations},
  pdfauthor={Matias G. Delgadino, Rishabh Gvalani, Matthew Rosenzweig}
}
\usepackage[nameinlink,capitalise,noabbrev]{cleveref}

\raggedbottom
\allowdisplaybreaks
\numberwithin{equation}{section}
\emergencystretch=3em
\setlist{listparindent=\parindent,parsep=0pt}
\setlist[itemize]{topsep=4pt,itemsep=2pt,parsep=1pt,leftmargin=1.8em}
\setlist[enumerate]{topsep=4pt,itemsep=3pt,parsep=1pt,leftmargin=2.2em}

\theoremstyle{plain}
\newtheorem{thm}{Theorem}[section]
\newaliascnt{prop}{thm}
\newtheorem{prop}[prop]{Proposition}
\aliascntresetthe{prop}
\newaliascnt{lemma}{thm}
\newtheorem{lemma}[lemma]{Lemma}
\aliascntresetthe{lemma}
\newaliascnt{cor}{thm}
\newtheorem{cor}[cor]{Corollary}
\aliascntresetthe{cor}
\newaliascnt{assumption}{thm}
\newtheorem{assumption}[assumption]{Assumption}
\aliascntresetthe{assumption}
\theoremstyle{definition}
\newaliascnt{mydef}{thm}
\newtheorem{mydef}[mydef]{Definition}
\aliascntresetthe{mydef}
\newaliascnt{remark}{thm}
\newtheorem{remark}[remark]{Remark}
\aliascntresetthe{remark}
\newaliascnt{conj}{thm}
\newtheorem{conj}[conj]{Conjecture}
\aliascntresetthe{conj}

\Crefname{thm}{Theorem}{Theorems}
\Crefname{prop}{Proposition}{Propositions}
\Crefname{lemma}{Lemma}{Lemmas}
\Crefname{cor}{Corollary}{Corollaries}
\Crefname{assumption}{Assumption}{Assumptions}
\Crefname{mydef}{Definition}{Definitions}
\Crefname{remark}{Remark}{Remarks}
\Crefname{conj}{Conjecture}{Conjectures}
\crefformat{equation}{(#2#1#3)}

\DeclarePairedDelimiter{\pa}{\lparen}{\rparen}
\DeclarePairedDelimiter{\brak}{\lbrack}{\rbrack}
\DeclarePairedDelimiter{\brac}{\lbrace}{\rbrace}
\DeclarePairedDelimiter{\abs}{\lvert}{\rvert}
\DeclarePairedDelimiter{\norm}{\lVert}{\rVert}
\DeclarePairedDelimiter{\ip}{\langle}{\rangle}

\renewcommand{\d}{\mathsf d}
\newcommand{\dd}{\mathop{}\!\mathrm d}
\newcommand{\T}{\mathbb T}
\newcommand{\R}{\mathbb R}
\newcommand{\C}{\mathbb C}
\newcommand{\Z}{\mathbb Z}
\newcommand{\N}{\mathbb N}
\newcommand{\E}{\mathbb E}
\newcommand{\Pb}{\mathbb P}
\newcommand{\Q}{\mathbb Q}
\newcommand{\Id}{\operatorname{Id}}
\newcommand{\Tr}{\operatorname{Tr}}
\newcommand{\TV}{\operatorname{TV}}
\DeclareMathOperator{\sgn}{sgn}
\newcommand{\rank}{\operatorname{rank}}
\newcommand{\HS}{\mathrm{HS}}
\newcommand{\op}{\mathrm{op}}
\newcommand{\diag}{\Delta}
\newcommand{\one}{\mathbf 1}
\newcommand{\be}{\beta}
\newcommand{\al}{\alpha}
\newcommand{\vep}{\varepsilon}
\newcommand{\g}{\mathsf g}
\newcommand{\K}{\mathsf K}
\newcommand{\Fr}{\mathsf F}
\newcommand{\Hr}{\mathsf H}
\newcommand{\Er}{\mathsf E}
\newcommand{\Ec}{\mathcal E}
\newcommand{\Pc}{\mathcal P}
\newcommand{\Ent}{\operatorname{Ent}}
\newcommand{\Gcal}{\mathcal G}
\newcommand{\Ecal}{\mathcal E}
\newcommand{\pcspip}[2]{\left\langle #1,#2\right\rangle}
\newcommand{\tl}{\widetilde}
\newcommand{\f}{\mathsf f}
\newcommand{\wh}{\widehat}
\newcommand{\muunif}{\mu_{\mathrm{unif}}}
\newcommand{\gm}{\mathrm{gm}}
\newcommand{\pf}{\mathrm{pf}}
\newcommand{\spec}{\mathrm{s}}
\newcommand{\coll}{\mathrm{c}}
\newcommand{\tr}{\mathrm{tr}}
\newcommand{\XN}{X_N}
\newcommand{\Part}{\mathfrak P}
\DeclareMathOperator*{\esssup}{ess\,sup}
\DeclareMathOperator{\diver}{div}
\newcommand{\step}[1]{\par\medskip\noindent\textbf{#1}}

\title[Attractive log gas I]
{Attractive log gas I: partition functions and equilibrium fluctuations}

\author[M. G. Delgadino]{Matias G. Delgadino}
\address{Matias Delgadino, Department of Mathematics, The University of Texas at Austin, PMA 8.100, 2515 Speedway, Austin, TX 78712, USA}
\email{matias.delgadino@math.utexas.edu}
\urladdr{https://math.utexas.edu/directory/matias-delgadino}
\author[R. Gvalani]{Rishabh Gvalani}
\address{Rishabh Gvalani, School of Mathematics, University of Edinburgh, James Clerk Maxwell Building, Peter Guthrie Tait Road, Edinburgh EH9 3FD, United Kingdom}
\email{rishabh.gvalani@ed.ac.uk}
\urladdr{https://webhomes.maths.ed.ac.uk/~rgvalani/}
\author[M. Rosenzweig]{Matthew Rosenzweig}
\address{Matthew Rosenzweig, Department of Mathematical Sciences, Carnegie Mellon University, 7127 Wean Hall, 5000 Forbes Avenue, Pittsburgh, PA 15213, USA}
\email{mrosenz2@andrew.cmu.edu}
\urladdr{https://matthewrosenzweigwork-max.github.io/}
\thanks{M.G.D. was supported by NSF grant DMS-2205937. R.G. is partially funded by the Deutsche Forschungsgemeinschaft (DFG, German Research Foundation) through SPP 2410, Hyperbolic Balance Laws in Fluid Mechanics: Complexity, Scales, Randomness (CoScaRa). M.R. was supported by NSF grants DMS-2441170, DMS-2345533, and DMS-2342349.}
\subjclass[2020]{60K35, 60F05, 82B21, 26D10}
\keywords{attractive logarithmic gas, Beckner--Onofri inequality, logarithmic Hardy--Littlewood--Sobolev inequality, modulated partition function, Carleman--Fredholm determinant, marked marginal expansion, Hoeffding decomposition, connected correlation, replica asymptotics, local size of chaos, degenerate $U$-statistic, phase transition}

% Public working-draft presentation (2026-10-28).
\usepackage{draftwatermark}
\SetWatermarkText{DRAFT}
\SetWatermarkAngle{45}
\SetWatermarkScale{0.8}
\SetWatermarkColor[gray]{0.92}

\begin{document}
\begin{abstract}

We consider the attractive periodic logarithmic gas on the volume-one cubic torus in the fixed-temperature mean-field scaling. After subtracting the leading macroscopic contribution around the uniform measure, we identify the exact threshold for order-one control of the modulated partition function: it is the global-minimization threshold of the macroscopic free energy. Below this threshold, the partition function converges to the Carleman--Fredholm determinant of the logarithmic covariance operator with error of order $N^{-1}$ on compact inverse-temperature intervals. Above the global threshold and below the collision threshold $2\d$, the normalized logarithm has a strictly positive lower limit, while fixed-particle integrability holds exactly below $2\d$. We obtain the same estimate on explicit temperature intervals for bi-Lipschitz images of the uniform measure, including a fixed small H\"older neighborhood of it. Under a uniform partition bound at a strictly larger coupling, with the reference measure fixed, we obtain all-order prescribed-reference partition and fixed-marginal density expansions, identify the singleton and pair responses, and prove quantitative first-order profiles for growing marginals. For the homogeneous law, fixed-integer replica expansions retain the full dependence on a nonvanishing marginal fraction and give an explicit proportional-marginal chi-square limit. Further consequences include all-marginal relative-entropy bounds, a static central limit theorem, bounded-test fluctuation expansions, and a near-uniform modulated log-HLS inequality with additive error of order $N^{-1}$. For the planar gas with quadratic confinement, we prove a global equilibrium entropy--energy margin and a prescribed-reference moment bound, uniformly on compact inverse-temperature intervals below $4$. The main argument combines the periodic Beckner--Onofri inequality with an exact regression covariance decomposition and reduction of the local coupling; the higher-order analysis uses cutoff-uniform marked Wick expansions. Two appendices treat one-component repulsive and common-background signed logarithmic gases.

\end{abstract}
\maketitle

\section{Introduction}
\label{sec:introduction}

\subsection{The attractive logarithmic gas}
\label{subsec:threshold-introduction}

Fix an integer $\d\geq2$ and let
\begin{equation}
 \T^{\d}=\R^{\d}/\Z^{\d},
 \qquad
 \muunif(\T^{\d})=1,
 \label{eq:torus-intro}
\end{equation}
where $\muunif$ denotes normalized Lebesgue measure.  Let $\g=\g_{\d}$ be the mean-zero periodic logarithmic kernel specified in \Cref{subsec:functional-domains}.  Its diagonal singularity is
\begin{equation}
 \g(x)=-\log|x|+O_{\d}(1)
 \qquad\text{as }x\longrightarrow0.
 \label{eq:kernel-local-form}
\end{equation}
The attractive logarithmic gas has partition function
\begin{equation}
 Z_{N,\be}
 \coloneqq\int_{(\T^{\d})^N}
 \exp\left\{\frac{\be}{2N}
 \sum_{1\leq i\neq j\leq N}\g(x_i-x_j)\right\}
 \,\dd\muunif^{\otimes N}(\XN).
 \label{eq:homogeneous-partition-function}
\end{equation}

In dimension two, the normalization \eqref{eq:kernel-normalization} gives $\g_2=2\pi G$, where $G$ is the mean-zero Green function satisfying $-\Delta G=\delta_0-1$.  The same logarithmic energy appears in two classical dynamical settings.  Under overdamped Langevin dynamics, its mean-field equation is the periodic parabolic--elliptic Patlak--Keller--Segel model of chemotaxis \cite{ChodrondeCourcelRosenzweigSerfaty2025}.  Under Hamiltonian point-vortex dynamics, the corresponding attractive Gibbs weight is the canonical one-species vortex ensemble at negative inverse temperature in Onsager's theory \cite{Onsager1949,Grotto2026}.  These dynamics are distinct, but they lead to the same mean-field-scaled logarithmic Gibbs weight.  The partition-function asymptotics studied here therefore describe equilibrium fluctuations of the homogeneous phase in both interpretations.  In the point-vortex formulation, the limiting Carleman--Fredholm determinant normalizes the Gaussian energy--enstrophy ensemble obtained by tilting mean-zero white noise by its renormalized kinetic energy \cite{Grotto2026}.

\subsection{Partition functions and equilibrium fluctuations}

The macroscopic free energy selects a density around which to study fluctuations. After its leading contribution has been separated, the remaining partition function records the next-order normalization. In the Hilbert--Schmidt regime, the Gaussian prediction for this normalization is a Carleman--Fredholm determinant. For an attractive interaction, a competing density may have lower free energy even while the uniform state has a positive Hessian and every fixed-particle integral is finite. The persistence of the Gaussian normalization is therefore a global question.

Our first principal result identifies the threshold for order-one normalization with the threshold for global minimality of the uniform density. Below this common threshold, \Cref{thm:principal} gives determinant convergence with error of order $N^{-1}$ on compact inverse-temperature intervals, together with a transported-reference version. Above the global threshold and below the collision threshold, a product trial gives a strictly positive lower limit for the normalized logarithm of the partition function. The sharp periodic Beckner--Onofri inequality identifies the global threshold with the collision threshold for $2\leq\d\leq10$ and with the spectral threshold for $\d\geq12$; in dimension eleven the global transition occurs strictly before spectral instability. The precise statements and attribution appear in \Cref{thm:beckner-onofri-input}.

Our second principal result identifies the spatial correlations beyond the scalar normalization. We call the labeled coordinates held fixed in a marginal integral \emph{marked}; all other particle coordinates are integrated out. Under a strict partition bound around the same prescribed reference, \Cref{thm:prescribed-marked-expansion} expands each fixed marginal density to every fixed order in $N^{-1}$ in every finite $L^p$ space. Its first coefficient is a sum of singleton and pair responses; translation invariance makes the singleton term vanish for the uniform reference. \Cref{thm:prescribed-growing-marginals} then gives the first-order profile and sharp divergence constants when the marginal size grows but remains a vanishing fraction of the system. The fixed-marginal expansion also determines higher relative-entropy, correlation, and bounded-test fluctuation coefficients.

A marginal containing a positive fraction of the particles exhibits a different limit. In the replica representation of an integer moment, the copies share the observed coordinates and have separate remaining coordinates; the overlap fraction is the observed marginal size divided by $N$. Theorem~\ref{thm:all-order-replica-expansions} keeps the full dependence of each fixed integer moment on this fraction. In particular, at positive inverse temperature and a fixed positive fraction within the stated window, the chi-square divergence converges to an explicit finite positive constant. The quadratic term of the limiting determinant is the local-chaos constant; its higher powers describe collective correlations beyond the first-order local profile. See \eqref{eq:local-chaos-proportional-limit} and \eqref{eq:main-overlap-expansion}. For fixed marginal size, the strong marked expansion already determines these integer moments. The additional content of the replica theorem is its uniformity at nonvanishing marginal fraction.

The partition bound has consequences under several distinct sets of hypotheses: nonasymptotic relative-entropy bounds for every marginal size, a central limit theorem for square-integrable tests, and a near-uniform modulated log-HLS inequality with additive error of order $N^{-1}$. These estimates supply static inputs to Part II, which studies the mean-field dynamics. Obtaining sharp time-dependent marginal rates requires an additional evolution argument.

The periodic results are complemented by a planar whole-space theorem under quadratic confinement. On every fixed inverse-temperature interval strictly below $4$, \Cref{thm:planar-confined-static-package} gives a uniform global entropy--energy margin around the thermal equilibrium and a uniform partition bound at a larger coupling for each fixed reference in a bounded-ratio neighborhood. These are the static inputs used by the companion dynamics paper; the theorem does not assert a whole-space determinant limit.

Appendices~\ref{app:repulsive-marked-oscillatory} and~\ref{app:signed-log-heat-kernel} treat one-component repulsive and common-background signed logarithmic gases, including the neutral two-component model. They give all-order partition and fixed-marginal expansions. The oscillatory Hubbard--Stratonovich argument and the heat-kernel Sine--Gordon argument give two proofs of the one-component partition and marginal theorem; the heat-kernel argument also treats deterministic charge vectors and bounded linear-statistic sources. The signed theorem uses a common reference for both species and keeps the exact finite-particle charge imbalance. Its hypotheses and parameter ranges are stated in the appendices.

\subsection{Relation to previous work}
\label{subsec:related-work}

Chodron de Courcel, Rosenzweig, and Serfaty \cite{ChodrondeCourcelRosenzweigSerfaty2025} study the attractive macroscopic phase diagram and use the modulated-free-energy framework developed by Bresch, Jabin, and Wang \cite{BreschJabinWang2023}. Their partition estimate coarse grains the empirical measure so that entropy and a macroscopic variational inequality apply to a density. Here we also decompose the covariance operator. The exact regression identity \eqref{eq:intro-regression-decomposition} leaves the unresolved microscopic interaction as a blockwise-canonical quadratic statistic.

The sharp Euclidean and spherical logarithmic Hardy--Littlewood--Sobolev and Beckner--Onofri inequalities, together with their conformal equality families, were established by Carlen and Loss \cite{CarlenLoss1992} and Beckner \cite{Beckner1993}.  Okikiolu \cite{Okikiolu2008} developed the compact-manifold framework relating these inequalities to regularized spectral traces and concentration at the round-sphere level.  The periodic problem has the same local logarithmic singularity, while its discrete Fourier spectrum and absence of dilation symmetry create independent collision and spectral obstructions.  We use the resulting periodic inequality only through \Cref{thm:beckner-onofri-input}.

In our companion paper \cite{DelgadinoGvalaniRosenzweig2026}, we prove determinant asymptotics and quantitative rates for repulsive logarithmic and Riesz gases in the Hilbert--Schmidt regime, as well as attractive logarithmic estimates under perturbative hypotheses on the reference measure and inverse temperature.  The present paper specializes to the periodic logarithmic kernel and reaches the complete global-minimization regime through the exact finite-partition regression identity and the reduction of the local coupling in \eqref{eq:intro-local-temperature}.

In dimension two, Grotto \cite{Grotto2026} obtained qualitative determinant convergence at sufficiently small attractive inverse temperature in the equivalent negative-temperature point-vortex formulation.\footnote{Grotto writes $G$ for the mean-zero Green function satisfying $-\Delta G=\delta_0-1$.  Since $\g_2=2\pi G$, his positive parameter for negative inverse temperature equals $2\pi\be$ in our notation.  Consequently, his fixed-particle condition that this parameter be less than $8\pi$ becomes $\be<4=\beta_{\coll}$.}  Our result is quantitative and, since $\beta_{\gm}(2)=4$, applies on every compact inverse-temperature interval in the full collision-subcritical range.

In \cite[p.~9]{WangZhao2026}, Wang and Zhao announce forthcoming work \cite{WangZhaoAttractive2026} on the attractive planar Coulomb gas at the explicit logarithmic HLS equilibrium.  In their normalization, the announced work proves uniform exponential integrability, convergence to the corresponding Carleman--Fredholm determinant, and fixed-particle integrability throughout the sharp interval $0<\be<8\pi$.


The modulated log-HLS inequality supplies the periodic static coercivity used in the companion work \emph{The attractive logarithmic gas II: mean-field dynamics}. \Cref{thm:planar-confined-static-package} supplies the corresponding planar confined equilibrium margin and prescribed-reference partition bound at a larger coupling.


Quantitative local independence for mean-field Gibbs measures has been studied by Lacker \cite{Lacker2022}, who obtains a quadratic relative-entropy scale in a high-temperature regime using a hierarchy of functional-inequality estimates. Ren and Wang \cite{RenWang2025} use the gradient structure of conditional laws and concentration estimates to treat further mean-field Gibbs models. Here, the singular attractive interaction is controlled through a uniform partition bound at a strictly larger coupling, with the reference fixed. The first projection is estimated by the normalized conditional density, and higher projections by a partial-overlap replica moment. These are static estimates for the modulated Gibbs law; propagation estimates require the separate dynamical argument of Part II.

The microscopic residual belongs to the classical theory of canonical, or degenerate, degree-two $U$-statistics initiated by Hoeffding \cite{Hoeffding1948}.  We use the decoupling theorem of de la Pe\~na and Montgomery-Smith \cite[Theorem~1]{delaPenaMontgomerySmith1995} for the local exponential estimate.  Serfling's Hilbert--Schmidt theorem \cite[Section~5.5.2]{Serfling1980} supplies the qualitative joint linear/quadratic weak limit used in \Cref{cor:static-clt}.  A direct quantitative comparison follows from the Hubbard--Stratonovich estimate under a higher-moment bound and logarithmic cutoff analysis; the sharper rate stated in \Cref{thm:principal} follows from the all-order expansion at order zero.



The exact shared-label identities belong to the replica formulation of marginal moments, while the recovery of the local pair profile uses the Hoeffding decomposition of canonical $U$-statistics.  The growing-marginal total-variation constant is obtained by a rank-independent Lindeberg replacement in the spirit of Chatterjee \cite{Chatterjee2006}.  The all-order argument uses hypercontractivity for homogeneous Wiener chaoses \cite{NourdinPeccatiReinert2010} and is related to the Banach-space Laplace expansions of Bolthausen \cite{Bolthausen1986} and Albeverio and Liang \cite{AlbeverioLiang2005}.  In contrast with those abstract expansions, the present proof is organized around spatially integrated Wick powers and a cutoff-uniform partition bound at a larger coupling adapted to the logarithmic singularity.  Related high-temperature partition expansions for one-dimensional beta ensembles were obtained by Dworaczek Guera \cite{DworaczekGuera2024}.


The distinction between a density expansion and a moment expansion is relevant to these comparisons. For fixed marginal size, the prescribed-reference theorem identifies signed spatial coefficients, their Hoeffding projections, connected correlations, and relative-entropy coefficients. The homogeneous replica theorem instead treats nonvanishing overlap. Their different ranges and observables are specified in \Cref{thm:prescribed-marked-expansion,thm:all-order-replica-expansions}.

\subsection{Proof architecture}
\label{subsec:proof-architecture}

The proof has two parts. The macroscopic entropy--energy inequality first yields a particle-uniform partition bound. That bound then supplies the exponential integrability needed for the marked and replica expansions.

A direct approximation of the interaction at every pair would accumulate an error of order $N$ in the exponent.  We instead fix one finite spatial partition and decompose the logarithmic covariance operator $T$ exactly.  The resulting identity has the form
\begin{equation}
 T=A+R-D,
 \qquad
 0\leq A\leq T,
 \qquad
 D\geq0.
 \label{eq:intro-regression-decomposition}
\end{equation}
Here $A$ is the covariance explained by linear regression on the cell-average variables, and $R-D$ is the conditional covariance of the remaining component. The kernel of $R$ has an explicit double-centering formula on every cell pair. This formula is suited to the local logarithmic estimates, while the finite-rank subtraction $D$ has a favorable sign.

The macroscopic entropy--energy inequality controls the complete coarse covariance through the domination $A\leq T$.  After conditioning on the complete cell-label vector, the particle coordinates remain independent inside their assigned cells and the residual kernel is canonical in each variable on every cell pair.  For a near pair involving $m$ particles, the mean-field coefficient can be written as
\begin{equation}
 \frac{\be}{N}S
 =\frac{\be m/N}{m}S.
 \label{eq:intro-local-temperature}
\end{equation}
Thus, the local inverse temperature is the global inverse temperature multiplied by the local particle fraction.  After choosing the partition sufficiently fine, the good cell-count event bounds every local particle fraction by a small fixed constant, so every local canonical array lies in a perturbative regime even when the global inverse temperature is close to $\beta_{\gm}$.  Double centering removes the additive logarithmic scale created by rescaling neighboring cells, mixed second differences control separated cells, and the exponentially small probability of overoccupation pays for an independent linear-in-$N$ bound on the log partition function.

The subsequent Gaussian estimates use the partition bound at an arbitrarily small amplification $r\be$ of the desired inverse temperature. They therefore preserve the entire strict range below the global threshold. A Hubbard--Stratonovich estimate under a higher-moment bound gives the direct cutoff comparison, and the marked expansion gives the higher-order logarithmic asymptotics.



For prescribed-reference marginals, the partition bound at a larger coupling supports a strong marked expansion.  At a finite cutoff that preserves nonnegative definiteness, an exact Cameron--Martin shift moves the marked particles into a centered one-site source.  The decisive estimate is a moment bound with the marked coordinates fixed: after rounding a slightly amplified number of unmarked particles to an integer, the spatial $L^p$ norm is controlled by a partition function at effective inverse temperature $r\be+O_{k,p}(N^{-1})$.  Thus, taking a high spatial moment adds only finitely many shared particles and does not replace the bulk inverse temperature by $p\be$.  Conditional Gaussian estimates are completed before the spatial norm is taken, which preserves cutoff uniformity at logarithmic collisions.

The expansion is performed on spatially integrated Wick exponentials. Bounds for integrated Wick powers and a good-event/bad-event decomposition control the exact exponential uniformly in the logarithmic cutoff. Gaussian sign symmetry removes the odd half-powers from the partition and marginal expansions. Bell polynomials give the coefficients, and inclusion--exclusion over subsets identifies their Hoeffding levels and dependence on the marked particles. Bounded test insertions allow half-powers and give locally uniform joint transform expansions. Throughout, the expansion order, marginal size, and number of replicas are fixed unless a separate growing-marginal theorem is invoked.


The order-one estimate also extends uniformly to a neighborhood of reference densities. Entropy duality then gives the static modulated log-HLS inequality with an additive error of order $N^{-1}$, without using determinant convergence.

For homogeneous replica moments, exact shared-label Hubbard--Stratonovich identities reduce a $q$-replica moment with overlap fraction $\theta$ to scalar partition functions at an arbitrarily small strict amplification of
\begin{equation}
 \be\bigl[1+(q-1)\theta\bigr],
 \label{eq:intro-replica-effective-coupling}
\end{equation}
rather than at $q\be$.  The resulting coefficient functions retain the exact overlap fraction. The local homogeneous profile is the uniform-reference specialization of the prescribed-reference marked expansion.

The local-chaos argument uses the order-one estimate through two different consequences. A two-replica Gaussian identity bounds a proportional marginal by partition functions at effective inverse temperature close to $\be(1+k/N)$. Hoeffding orthogonality then controls all higher projections at scale $k^2/N^2$. Translation invariance removes the first projection for the homogeneous law. For a prescribed reference, a normalized one-particle conditional identity and the spectral margin forced by the same partition bound at a larger coupling control that projection. Reference-adapted conditional expectations preserve the full partition bound at every finite Gaussian cutoff; no limiting Gaussian multiplicative chaos construction or continuation in the number of replicas is used.

For growing marginals, the one- and two-particle coefficients identify the first profile. A uniformly bounded proportional marginal and Hoeffding orthogonality control the remaining projections. A scalar Taylor remainder then compares entropy with chi-square divergence, and a rank-independent Lindeberg replacement identifies the total-variation constant.

\subsection{Organization of the paper}
\label{subsec:organization}

Section~\ref{sec:main-results} gives the precise setup and statements. Section~\ref{sec:collisions-preliminaries} collects the preliminary estimates and marginal identities. Sections~\ref{sec:coarse-decomposition} and~\ref{sec:canonical-residual} prove the particle-uniform partition bound; Section~\ref{sec:thresholds-asymptotics} derives the threshold identity, reference estimates, determinant asymptotics, modulated log-HLS inequality, and static central limit theorem. Section~\ref{sec:local-chaos} establishes nonasymptotic local chaos. Section~\ref{sec:prescribed-marked-expansions} treats the marked expansions and their fixed- and growing-marginal consequences, followed by the homogeneous replica analysis in Subsection~\ref{sec:homogeneous-replica-chaos}. Section~\ref{sec:planar-confined-static-package} treats planar quadratic confinement, and Section~\ref{sec:perspectives} discusses open questions. Appendices A--D contain the Gaussian and Wick estimates, Appendices E--F treat repulsive and signed gases, and Appendix~\ref{app:planar-confined-static-package} contains the planar confined argument.

\subsection{Statement on AI use}

This working manuscript was generated with the assistance of generative AI tools, including OpenAI's ChatGPT and Codex and Anthropic's Claude, for literature identification, the development and testing of mathematical arguments, exposition, and consistency checks. The authors have not yet reviewed and independently verified all AI-generated content. They are in the process of formalizing the results in Lean and revising the natural-language presentation. The manuscript is therefore being shared as a draft, and its statements, proofs, references, and exposition remain subject to correction. Any completed formal verification applies only to the statements and versions explicitly identified in the accompanying verification records.

\section{Main results}
\label{sec:main-results}

\subsection{Modulated Gibbs normalization and macroscopic thresholds}

We use the torus, kernel, and homogeneous partition function in \eqref{eq:torus-intro}--\eqref{eq:homogeneous-partition-function}; the exact Fourier normalization is \eqref{eq:kernel-normalization}. Each reference measure is fixed unless a family is explicitly specified.


Write the physical attractive interaction as
\begin{equation}
 W(x,y)\coloneqq-\g(x-y).
 \label{eq:attractive-interaction}
\end{equation}
Let $\mu$ be a nonatomic probability measure on $\T^{\d}$ for which the expressions below are finite.  When $\mu\ll\muunif$, we use the same symbol $\mu$ for the measure and its density with respect to $\muunif$; thus, function-space norms and pointwise bounds on $\mu$ refer to this density.  Set
\begin{equation}
 h_{\mu}(x)\coloneqq\int_{\T^{\d}}W(x,y)\,\dd\mu(y),
 \qquad
 I_{\mu}\coloneqq\frac12\iint_{(\T^{\d})^2}W(x,y)\,\dd\mu(x)\dd\mu(y).
 \label{eq:hmu-imu-definition}
\end{equation}
For a two-body kernel $f$, we write
\begin{align}
 \tl f(x,y)
 \coloneqq{}&f(x,y)-\int_{\T^{\d}}f(x,z)\,\dd\mu(z)
 -\int_{\T^{\d}}f(z,y)\,\dd\mu(z)
 \notag\\
 &+\iint_{(\T^{\d})^2}f(z,w)\,\dd\mu(z)\dd\mu(w)
 \label{eq:hoeffding-centering}
\end{align}
for its Hoeffding centering relative to $\mu$.  The modulated energy $\Fr_N(\XN,\mu)$ is characterized by the exact splitting
\begin{equation}
 \frac1{2N}\sum_{1\leq i\neq j\leq N}W(x_i,x_j)
 =N\Fr_N(\XN,\mu)+\sum_{i=1}^Nh_{\mu}(x_i)-NI_{\mu}.
 \label{eq:intro-modulated-splitting}
\end{equation}
Equivalently,
\begin{equation}
 \Fr_N(\XN,\mu)
 =\frac1{2N^2}\sum_{1\leq i\neq j\leq N}\tl W(x_i,x_j)
 -\frac1{N^2}\sum_{i=1}^Nh_{\mu}(x_i)+\frac1N I_{\mu}.
 \label{eq:modulated-energy-general-reference}
\end{equation}
The associated modulated partition function and Gibbs law are
\begin{equation}
 \K_{N,\be}(\mu)
 \coloneqq\E_{\mu^{\otimes N}}\left[e^{-\be N\Fr_N(\XN,\mu)}\right]
 \label{eq:partition-general-reference}
\end{equation}
and
\begin{equation}
 \dd\Q_{N,\be}(\mu)(\XN)
 \coloneqq\frac{e^{-\be N\Fr_N(\XN,\mu)}}{\K_{N,\be}(\mu)}
 \,\dd\mu^{\otimes N}(\XN).
 \label{eq:modulated-gibbs-law}
\end{equation}



For probability measures $\nu\ll\eta$, we use the conventions
\begin{align}
 \chi^2(\nu\mid\eta)
 &=\int\left(\frac{\dd\nu}{\dd\eta}-1\right)^2\,\dd\eta,
 \label{eq:homogeneous-chi-square-convention}\\
 \TV(\nu,\eta)
 &=\frac12\int\left|\frac{\dd\nu}{\dd\eta}-1\right|\,\dd\eta.
 \label{eq:homogeneous-total-variation-convention}
\end{align}

For $\be>0$ and a probability measure $f_N$ on $(\T^{\d})^N$, the associated modulated free energy is
\begin{equation}
 \Er_N(f_N,\mu)
 \coloneqq\frac1{\be N}\Hr(f_N\mid\mu^{\otimes N})
 +\E_{\XN\sim f_N}[\Fr_N(\XN,\mu)],
 \label{eq:modulated-free-energy}
\end{equation}
and it satisfies
\begin{equation}
 \Er_N(f_N,\mu)
 =\frac1{\be N}\left(
 \Hr(f_N\mid\Q_{N,\be}(\mu))-\log\K_{N,\be}(\mu)
 \right).
 \label{eq:modulated-free-energy-rewriting}
\end{equation}

The centered attractive covariance is
\begin{equation}
 \tl{\g}_{\mu}(x,y)\coloneqq-\tl W(x,y)
 =\g(x-y)+h_{\mu}(x)+h_{\mu}(y)-2I_{\mu}.
 \label{eq:hoeffding-centered-kernel}
\end{equation}
Whenever $\tl{\g}_{\mu}\in L^2(\mu^{\otimes2})$, let $T_{\mu}$ be the nonnegative self-adjoint Hilbert--Schmidt operator on $L_0^2(\mu)$ defined by
\begin{equation}
 T_{\mu}f(x)\coloneqq\int_{\T^{\d}}\tl{\g}_{\mu}(x,y)f(y)\,\dd\mu(y).
 \label{eq:positive-covariance-operator}
\end{equation}

Since $\g$ has zero $\muunif$-mean,
\begin{equation}
 h_{\muunif}=0,
 \qquad
 I_{\muunif}=0,
 \qquad
 Z_{N,\be}=\K_{N,\be}(\muunif).
 \label{eq:uniform-specialization}
\end{equation}
We henceforth write
\begin{equation}
 T\coloneqq T_{\muunif}.
 \label{eq:uniform-operator-abbreviation}
\end{equation}

Since $\g$ has zero $\muunif$-mean, the interaction energy of every probability measure $\mu\ll\muunif$ admits the representation
\begin{equation}
 I_{\mu}
 =-\frac12\iint_{(\T^{\d})^2}\g(x-y)
 (\mu(x)-1)(\mu(y)-1)\,\dd\muunif(x)\dd\muunif(y)
 \in[-\infty,0],
 \label{eq:macroscopic-energy}
\end{equation}
which we use as the extended definition when the integral in \eqref{eq:hmu-imu-definition} is not finite.
For $\Hr(\mu\mid\muunif)<\infty$, the homogeneous macroscopic free energy is
\begin{equation}
 \Ec_{\be}(\mu)
 \coloneqq\Ec_{\be,0}(\mu)
 =\Hr(\mu\mid\muunif)+\be I_{\mu},
 \label{eq:macroscopic-free-energy}
\end{equation}
with the convention $\Ec_0(\mu)=\Hr(\mu\mid\muunif)$.  We introduce
\begin{align}
 \beta_{\coll}&\coloneqq2\d,
 \label{eq:collision-threshold-intro}\\
 \beta_{\spec}&\coloneqq\norm{T}_{\op}^{-1},
 \label{eq:spectral-threshold}
\end{align}
and define
\begin{equation}
 \beta_{\gm}
 \coloneqq\sup\left\{b\in[0,2\d]:
 \Ec_b(\mu)\geq0
 \text{ for every probability measure $\mu$ with $\Hr(\mu\mid\muunif)<\infty$}\right\},
 \label{eq:global-minimization-threshold}
\end{equation}
\begin{equation}
 \beta_{\pf}
 \coloneqq\sup\left\{b\geq0:
 \sup_{N\geq2}\sup_{0\leq\be\leq b_0}Z_{N,\be}<\infty
 \text{ for every }b_0<b\right\}.
 \label{eq:partition-threshold}
\end{equation}
Finally, define the macroscopic pressure by
\begin{equation}
 \Pc_{\d}(\be)
 \coloneqq\sup_{\substack{\mu\ll\muunif\\
 \Hr(\mu\mid\muunif)<\infty}}
 \left\{-\be I_{\mu}-\Hr(\mu\mid\muunif)\right\}
 =-\inf_{\substack{\mu\ll\muunif\\
 \Hr(\mu\mid\muunif)<\infty}}\Ec_{\be}(\mu).
 \label{eq:macroscopic-pressure}
\end{equation}

The collision, spectral, global-minimization, and partition thresholds concern, respectively, fixed-particle integrability, positivity of the Hessian at the uniform state, global minimality, and particle-uniform normalization. The collision and spectral thresholds are central in \cite{ChodrondeCourcelRosenzweigSerfaty2025}. The first theorem identifies the global-minimization and partition thresholds:
\begin{equation}
 \beta_{\pf}=\beta_{\gm}
 \leq\min\{\beta_{\coll},\beta_{\spec}\}.
 \label{eq:threshold-hierarchy-intro}
\end{equation}
\subsection{The principal theorem}
\label{subsec:principal-theorem}

\begin{thm}[Global minimization, partition thresholds, and determinant asymptotics]
\label{thm:principal}
Fix an integer $\d\geq2$ and use the preceding notation.

\smallskip
\noindent\textup{\textbf{Part A: the uniform reference.}}
The following assertions hold.
\begin{enumerate}[label=\textup{(\roman*)}]
\item For every $\be_0<\beta_{\gm}$, there is $C=C(\d,\be_0)<\infty$ such that
\begin{equation}
 \sup_{N\geq2}\sup_{0\leq\be\leq\be_0}Z_{N,\be}\leq C.
 \label{eq:principal-order-one}
\end{equation}
Consequently,
\begin{equation}
 \beta_{\pf}=\beta_{\gm}.
 \label{eq:principal-threshold-identity}
\end{equation}

\item For every $\be_0<\beta_{\gm}$, there is $C=C(\d,\be_0)<\infty$ such that, for every $N\geq2$,
\begin{equation}
 \sup_{0\leq\be\leq\be_0}
 \left|Z_{N,\be}-\det\nolimits_2(\Id-\be T)^{-1/2}\right|
 \leq \frac C N.
 \label{eq:principal-homogeneous-rate}
\end{equation}

\item If $\beta_{\gm}<\be<2\d$, then $\Pc_{\d}(\be)>0$ and
\begin{equation}
 \liminf_{N\to\infty}\frac1N\log Z_{N,\be}
 \geq\Pc_{\d}(\be)>0.
 \label{eq:principal-positive-pressure}
\end{equation}

\item For every fixed $N\geq2$,
\begin{equation}
 Z_{N,\be}<\infty
 \quad\Longleftrightarrow\quad
 \be<2\d.
 \label{eq:principal-fixed-particle-integrability}
\end{equation}
\end{enumerate}

\smallskip
\noindent\textup{\textbf{Part B: bi-Lipschitz transported references.}}
\label{thm:transported}
Let $\Phi:\T^{\d}\to\T^{\d}$ be a bi-Lipschitz bijection, set $\mu=\Phi_{\#}\muunif$, and define, off the diagonal,
\begin{equation}
 r_{\Phi}(y,z)\coloneqq\g(\Phi(y)-\Phi(z))-\g(y-z).
 \label{eq:pullback-remainder}
\end{equation}
Write $\tl r_{\Phi}$ for the Hoeffding centering of $r_{\Phi}$ with respect to $\muunif$, and set
\begin{equation}
 \beta_{\tr}(\Phi)
 \coloneqq\frac{\beta_{\gm}}
 {1+4\beta_{\gm}\norm{\tl r_{\Phi}}_{L^\infty(\muunif^{\otimes2})}}.
 \label{eq:transport-threshold}
\end{equation}
For every $\be_0<\beta_{\tr}(\Phi)$, there is a finite constant $C=C(\d,\Phi,\be_0)$ such that
\begin{equation}
 \sup_{N\geq2}\sup_{0\leq\be\leq\be_0}\K_{N,\be}(\mu)\leq C,
 \label{eq:principal-transport-order-one}
\end{equation}
and, for every $N\geq2$,
\begin{equation}
 \sup_{0\leq\be\leq\be_0}
 \left|\K_{N,\be}(\mu)
 -e^{\be I_{\mu}}\det\nolimits_2(\Id-\be T_{\mu})^{-1/2}\right|
 \leq \frac C N.
 \label{eq:principal-transport-rate}
\end{equation}

In particular, fix $\gamma\in(0,1)$ and $\be_0<\beta_{\gm}$.  There is $\vep=\vep(\d,\gamma,\be_0)>0$ such that, if $\mu\ll\muunif$ is a probability measure satisfying
\begin{equation}
 \mu\in C^\gamma(\T^{\d}),
 \qquad
 \inf_{\T^{\d}}\mu>0,
 \qquad
 \norm{\mu-1}_{C^\gamma}\leq\vep,
 \label{eq:near-uniform-hypothesis}
\end{equation}
then the conclusions \eqref{eq:principal-transport-order-one} and \eqref{eq:principal-transport-rate} hold on $0\leq\be\leq\be_0$.
\end{thm}

The error of order $N^{-1}$ is the order-zero consequence of the marked expansion in \Cref{thm:prescribed-marked-expansion}. The factor $e^{\be I_\mu}$ in \eqref{eq:principal-transport-rate} is the exact correction from diagonal exclusion; it equals one for the uniform reference. Both parts of the theorem concern strict compact temperature intervals. Neither the global endpoint nor the endpoint of the sufficient transport range is included.

When a thermal equilibrium $\mu_{\be,V}$ satisfies Part B, \eqref{eq:thermal-equilibrium-factorization} identifies the residual physical normalization and \eqref{eq:modulated-gibbs-law} identifies the canonical ensemble. This is a fixed-temperature interpretation. Part B varies $\be$ with $\mu$ fixed, whereas the equilibrium of a fixed one-body potential generally varies with temperature. The intrinsic nonuniform phase boundary may exceed the sufficient transport range.

\subsection{Periodic Beckner--Onofri input and dimension consequences}
\label{subsec:dimension-consequences}

The macroscopic threshold in \Cref{thm:principal} is determined by the sharp periodic Beckner--Onofri, equivalently logarithmic Hardy--Littlewood--Sobolev, inequality.  Concentration and the first Fourier shell give the elementary upper bound in \eqref{eq:threshold-hierarchy-intro}.  Which obstruction is sharp depends on the dimension, and a finite-amplitude transition can occur strictly before either one.

The two-dimensional statement below is due to Lin and Lucia \cite[Theorem~1.1 and (1.4)]{LinLucia2007}; the remaining statements are established by Lei and Rosenzweig \cite{LeiRosenzweig2026}.

\begin{thm}[Sharp periodic Beckner--Onofri inequality]
\label{thm:beckner-onofri-input}
The global-minimization threshold satisfies the following dimension-dependent conclusions.
\begin{enumerate}[label=\textup{(\roman*)}]
\item For $2\leq\d\leq10$,
\begin{equation}
 \beta_{\gm}(\d)=2\d.
 \label{eq:gm-low-dimensions}
\end{equation}
At the endpoint, the uniform measure is the unique minimizer of \eqref{eq:macroscopic-free-energy}.

\item In dimension eleven,
\begin{equation}
 \beta_{\gm}(11)<\beta_{\spec}(11)<22.
 \label{eq:dimension-eleven-strict}
\end{equation}

\item For every $\d\geq12$,
\begin{equation}
 \beta_{\gm}(\d)=\beta_{\spec}(\d).
 \label{eq:gm-high-dimensions}
\end{equation}
At the endpoint, the uniform measure is the unique minimizer.
\end{enumerate}
\end{thm}

Combining this input with \Cref{thm:principal} gives the full collision-subcritical partition conclusions in dimensions $2\leq\d\leq10$ and the spectral-subcritical conclusions in dimensions $\d\geq12$. In dimension eleven, a competing macroscopic profile is favorable for every $\be\in(\beta_{\gm}(11),\beta_{\spec}(11))$, although the Hessian at the uniform phase remains positive. The order-one regime ends at the presently unknown $\beta_{\gm}(11)$, and the positive-pressure lower bound holds on $(\beta_{\gm}(11),22)$.

\subsection{All-order marked expansions for prescribed references}
\label{subsec:prescribed-marked-introduction}

A partition bound at a strictly larger coupling, with the reference fixed, yields a strong expansion of each fixed labeled marginal density. It identifies the signed spatial coefficients and hence their Hoeffding projections and connected correlations. The marginal size is fixed in this all-order theorem; the growing-marginal and proportional-overlap results below address different regimes.

Set
\begin{equation}
 \psi_{\mu}\coloneqq h_{\mu}-2I_{\mu}
 \label{eq:prescribed-centered-source}
\end{equation}
and introduce the centered partition function
\begin{equation}
 \mathsf P_{N,\be}(\mu)
 \coloneqq e^{-\be I_{\mu}}\K_{N,\be}(\mu).
 \label{eq:prescribed-centered-partition}
\end{equation}
Thus, by \eqref{eq:modulated-energy-general-reference} and \eqref{eq:hoeffding-centered-kernel},
\begin{equation}
 \mathsf P_{N,\be}(\mu)
 =\int\exp\left\{
 \frac{\be}{N}\sum_{1\leq i<j\leq N}\tl\g_{\mu}(x_i,x_j)
 +\frac{\be}{N}\sum_{i=1}^N\psi_{\mu}(x_i)
 \right\}\,\dd\mu^{\otimes N}(\XN).
 \label{eq:prescribed-centered-partition-exponent}
\end{equation}
The source is centered in the sense that $\int\psi_{\mu}\,\dd\mu=0$. It is the one-body correction left by the modulated-energy splitting; its coefficient per particle in \eqref{eq:prescribed-centered-partition-exponent} is $\be/N$.

Let $\mathscr M$ be a family of probability measures whose densities with respect to $\muunif$ are bounded above by one common finite constant.  Fix $0\leq\be_0<b$ and assume the uniform partition bound
\begin{equation}
 \sup_{\mu\in\mathscr M}\sup_{N\geq1}
 \mathsf P_{N,b}(\mu)<\infty.
 \label{eq:prescribed-scalar-buffer}
\end{equation}
We call \eqref{eq:prescribed-scalar-buffer} a \emph{strict partition buffer}: the bound is available at the larger coupling $b$, while each reference $\mu$ is held fixed. The uniformity over $\mathscr M$ is part of the hypothesis.
No positive lower bound on the densities in $\mathscr M$ is required.  Under \eqref{eq:prescribed-scalar-buffer}, the operator $\Id-\be T_{\mu}$ is uniformly positive for $0\leq\be\leq\be_0$.  Define
\begin{equation}
 R_{\be,\mu}\coloneqq(\Id-\be T_{\mu})^{-1},
 \qquad
 B_{\be,\mu}\coloneqq R_{\be,\mu}-\Id,
 \qquad
 D_{\be,\mu}\coloneqq\det\nolimits_2(\Id-\be T_{\mu})^{-1/2}.
 \label{eq:prescribed-response-operators}
\end{equation}
Let $b_{\be,\mu}$ denote the integral kernel of $B_{\be,\mu}$.  The dressed diagonal is defined by the finite row formula
\begin{equation}
 \omega_{\be,\mu}(x)
 \coloneqq\be^2\ip*{\tl\g_{\mu}(x,\cdot),
 R_{\be,\mu}\tl\g_{\mu}(x,\cdot)}_{L^2(\mu)}.
 \label{eq:prescribed-omega}
\end{equation}
In particular, \eqref{eq:prescribed-omega} is not interpreted as the difference of two generally divergent diagonal values.  Write
\begin{equation}
 d_{\be,\mu}
 \coloneqq\be\psi_{\mu}
 +\frac12\left(\omega_{\be,\mu}
 -\int\omega_{\be,\mu}\,\dd\mu\right),
 \qquad
 a_{\be,\mu}\coloneqq R_{\be,\mu}d_{\be,\mu}.
 \label{eq:prescribed-singleton-response}
\end{equation}
Both $a_{\be,\mu}$ and $b_{\be,\mu}$ are centered in each of their variables.  For $k\geq1$, set
\begin{equation}
 G_{k,\be,\mu}(x_1,\ldots,x_k)
 \coloneqq\sum_{i=1}^ka_{\be,\mu}(x_i)
 +\sum_{1\leq i<j\leq k}b_{\be,\mu}(x_i,x_j).
 \label{eq:prescribed-leading-profile}
\end{equation}
Finally, let
\begin{equation}
 F_{N,\be,\mu}^{(k)}
 \coloneqq\frac{\dd\Q_{N,\be}^{(k)}(\mu)}{\dd\mu^{\otimes k}}
 \label{eq:prescribed-marginal-density}
\end{equation}
be the labeled $k$-particle marginal density.

\begin{thm}[All-order partition and strong marked expansions]
\label{thm:prescribed-marked-expansion}
Assume \eqref{eq:prescribed-scalar-buffer}.  There are scalar coefficients $z_j(\be,\mu)$ and symmetric marginal coefficients $f_{j,k,\be,\mu}$, independent of the order to which the expansion is truncated, with the following properties.

For every fixed integer $M\geq0$,
\begin{equation}
 \K_{N,\be}(\mu)
 =e^{\be I_{\mu}}D_{\be,\mu}
 \left(\sum_{j=0}^{M}\frac{z_j(\be,\mu)}{N^j}
 +O(N^{-M-1})\right),
 \qquad z_0=1,
 \label{eq:prescribed-partition-expansion}
\end{equation}
uniformly for $\mu\in\mathscr M$ and $0\leq\be\leq\be_0$.  For every fixed integer $k\geq1$, fixed integer $M\geq0$, and $1\leq p<\infty$,
\begin{equation}
 \left\|F_{N,\be,\mu}^{(k)}-1
 -\sum_{j=1}^{M}\frac{f_{j,k,\be,\mu}}{N^j}\right\|_{L^p(\mu^{\otimes k})}
 \leq C_{k,M,p}N^{-M-1}
 \label{eq:prescribed-marginal-expansion}
\end{equation}
for all sufficiently large $N$, uniformly in the same parameters.  Every coefficient belongs to every finite $L^p$, has integral zero, and is consistent under marginalization.  The first coefficient is
\begin{equation}
 f_{1,k,\be,\mu}=G_{k,\be,\mu},
 \label{eq:prescribed-first-coefficient}
\end{equation}
and the first scalar correction is
\begin{equation}
 z_1(\be,\mu)
 =\frac12\ip*{d_{\be,\mu},R_{\be,\mu}d_{\be,\mu}}_{L^2(\mu)}
 -\frac14\Tr(B_{\be,\mu}^2)
 +\frac1{12}\iint b_{\be,\mu}(x,y)^3\,\dd\mu(x)\dd\mu(y).
 \label{eq:prescribed-first-partition-correction}
\end{equation}
All inner products in \eqref{eq:prescribed-first-partition-correction} are in $L^2(\mu)$.
\end{thm}

Here contractions are the kernel products and source inner products obtained after Gaussian pairings, with the spatial variables integrated as indicated. The pointwise cubic integral in \eqref{eq:prescribed-first-partition-correction} uses three kernel factors at the same pair of points, whereas the operator trace contracts the arguments cyclically.

The parameters $M,k,p$ are fixed before $N\to\infty$. Enlarging them changes the constants and the lower bound on $N$, but not the reference family or the larger coupling $b$ in the partition bound. Thus the theorem gives an asymptotic expansion to each fixed finite order; it does not assert convergence of the infinite series.

\begin{cor}[Near-uniform prescribed references]
\label{cor:prescribed-near-uniform}
Fix $0\leq\be_0<\beta_{\gm}$ and $0<\gamma<1$.  There is $\vep>0$, depending on $\d,\gamma,\be_0$ and a chosen homogeneous partition bound at a strictly larger coupling, such that all conclusions of \Cref{thm:prescribed-marked-expansion} hold uniformly for
\begin{equation}
 \mu=\rho\muunif,
 \qquad
 \int\rho\,\dd\muunif=1,
 \qquad
 \norm{\rho-1}_{C^\gamma(\T^\d)}\leq\vep,
 \qquad
 0\leq\be\leq\be_0.
 \label{eq:prescribed-near-uniform-class}
\end{equation}
The same $\vep$ works for every prescribed fixed expansion order, fixed marginal size, and finite spatial moment exponent.
\end{cor}

The proof uses the Dacorogna--Moser transport already employed in Part B of \Cref{thm:principal}.  One first chooses
\begin{equation}
 \be_0<b<B<\beta_{\gm},
 \label{eq:prescribed-three-temperature-buffer}
\end{equation}
and then takes the transport sufficiently close to the identity that the bounded centered pullback remainder can be absorbed between $b$ and $B$.  The resulting same-reference bound is exactly \eqref{eq:prescribed-scalar-buffer}; the choice of the neighborhood precedes all later values of $M,k,p$.

\subsubsection{Growing marginals and the first-order profile}

The next theorem allows the marginal size to grow. Its estimates hold throughout a strict proportional window, but the first-order approximation becomes asymptotically exact on its natural scale when $k/N\to0$. The replica theorem below retains the additional information needed when the fraction stays positive.

\begin{thm}[First-order profile for growing prescribed-reference marginals]
\label{thm:prescribed-growing-marginals}
Assume \eqref{eq:prescribed-scalar-buffer}, and choose $0<\alpha<1$ with
\begin{equation}
 \be_0(1+\alpha)<b.
 \label{eq:prescribed-growing-window}
\end{equation}
For sufficiently large $N$ and $1\leq k\leq\lfloor\alpha N\rfloor$,
\begin{equation}
 \left\|F_{N,\be,\mu}^{(k)}-1-\frac{G_{k,\be,\mu}}N\right\|_2^2
 \leq C\left(\frac{k^2}{N^4}+\frac{k^3}{N^3}\right).
 \label{eq:prescribed-growing-profile}
\end{equation}
Consequently,
\begin{align}
 \left|\frac{N^2}{k^2}\chi^2
 -\frac{\norm{a_{\be,\mu}}_2^2}{k}
 -\frac{k-1}{2k}\norm{B_{\be,\mu}}_{\HS}^2\right|
 &\leq C\left(\frac1N+\frac{k}{N}\right),
 \label{eq:prescribed-growing-chi}\\
 \left|\frac{N^2}{k^2}\Hr
 -\frac{\norm{a_{\be,\mu}}_2^2}{2k}
 -\frac{k-1}{4k}\norm{B_{\be,\mu}}_{\HS}^2\right|
 &\leq C\left(\frac1N+\frac{k}{N}\right),
 \label{eq:prescribed-growing-entropy}\\
 \left|\frac Nk\TV-\frac1{2k}\norm{G_{k,\be,\mu}}_1\right|
 &\leq C\left(\frac1N+\sqrt{\frac{k}{N}}\right).
 \label{eq:prescribed-growing-TV}
\end{align}
The divergences compare $\Q_{N,\be}^{(k)}(\mu)$ with $\mu^{\otimes k}$.  Let
\begin{equation}
 \mathcal Q_{\be,\mu}
 \coloneqq\sum_{j\geq1}\lambda_j(B_{\be,\mu})(G_j^2-1),
 \label{eq:prescribed-second-chaos}
\end{equation}
where the eigenvalues, with multiplicity, are those on the real mean-zero space and the $G_j$ are independent standard real Gaussians.  The series converges in $L^2$ and has variance $2\norm{B_{\be,\mu}}_{\HS}^2$.  For $k\geq2$,
\begin{equation}
 \left|\frac Nk\TV-\frac14\E[|\mathcal Q_{\be,\mu}|]\right|
 \leq C\left(\frac1N+\sqrt{\frac{k}{N}}+k^{-1/6}\right).
 \label{eq:prescribed-growing-TV-chaos}
\end{equation}
Thus, for each fixed $\mu,\be$ and every sequence $k_N\to\infty$ with $k_N=o(N)$, the limits of $N^2k_N^{-2}\chi^2$, $N^2k_N^{-2}\Hr$, and $Nk_N^{-1}\TV$ are, respectively, $\frac12\norm{B_{\be,\mu}}_{\HS}^2$, $\frac14\norm{B_{\be,\mu}}_{\HS}^2$, and $\frac14\E[|\mathcal Q_{\be,\mu}|]$.  All displayed estimates are uniform in the original reference and temperature family.
\end{thm}

This is a first-order theorem for growing marginals. The all-order expansion in \Cref{thm:prescribed-marked-expansion} continues to require fixed marginal size.

\subsubsection{Fixed-marginal divergences and structural consequences}

\begin{cor}[Fixed-marginal divergences]
\label{cor:prescribed-divergence-expansions}
Under \eqref{eq:prescribed-scalar-buffer}, for every fixed $k\geq1$ and $L\geq2$,
\begin{align}
 \chi^2\!\left(\Q_{N,\be}^{(k)}(\mu)\mid\mu^{\otimes k}\right)
 &=\sum_{n=2}^{L}\frac{\mathcal C_{n,k}(\be,\mu)}{N^n}
 +O_{k,L}(N^{-L-1}),
 \label{eq:prescribed-chi-all-orders}\\
 \Hr\!\left(\Q_{N,\be}^{(k)}(\mu)\mid\mu^{\otimes k}\right)
 &=\sum_{n=2}^{L}\frac{\mathcal H_{n,k}(\be,\mu)}{N^n}
 +O_{k,L}(N^{-L-1}).
 \label{eq:prescribed-entropy-all-orders}
\end{align}
With all inner products and integrals over $\mu^{\otimes k}$,
\begin{align}
 \mathcal C_{n,k}
 &=\sum_{i=1}^{n-1}\ip*{f_{i,k},f_{n-i,k}}_{L^2(\mu^{\otimes k})},
 \label{eq:prescribed-chi-coefficients}\\
 \mathcal H_{n,k}
 &=\sum_{r=2}^{n}\frac{(-1)^r}{r(r-1)}
 \sum_{\substack{j_1,\ldots,j_r\geq1\\j_1+\cdots+j_r=n}}
 \int\prod_{a=1}^rf_{j_a,k}\,\dd\mu^{\otimes k}.
 \label{eq:prescribed-entropy-coefficients}
\end{align}
In particular,
\begin{equation}
 \mathcal C_{2,k}=2\mathcal H_{2,k}
 =k\norm{a_{\be,\mu}}_2^2
 +\binom{k}{2}\norm{B_{\be,\mu}}_{\HS}^2.
 \label{eq:prescribed-leading-divergences}
\end{equation}
Moreover,
\begin{equation}
 \left|N\TV\!\left(\Q_{N,\be}^{(k)}(\mu),\mu^{\otimes k}\right)
 -\frac12\norm{G_{k,\be,\mu}}_{L^1(\mu^{\otimes k})}\right|
 \leq\frac{C_k}{N}.
 \label{eq:prescribed-leading-TV}
\end{equation}
\end{cor}

The singleton term in \eqref{eq:prescribed-leading-divergences} is absent at the uniform reference because $\psi_{\muunif}=0$ and the dressed diagonal is constant.  Thus, $a_{\be,\muunif}=0$ and the pair profile in \eqref{eq:prescribed-leading-profile} becomes \eqref{eq:homogeneous-screened-pair-profile}.

A bound on the \emph{mark degree} means that each nonzero component of the canonical Hoeffding decomposition depends on at most that many distinct marked coordinates. This is unrelated to polynomial degree in the particle positions.

\begin{prop}[Marked structure, connected correlations, and linear statistics]
\label{prop:prescribed-structural-synopsis}
For every coefficient order $m\geq1$, the function $f_{m,k}$ has mark degree at most $2m$ and admits a canonical Hoeffding representation
\begin{equation}
 f_{m,k}(x_{[k]})
 =\sum_{r=1}^{\min\{2m,k\}}
 \sum_{\substack{A\subseteq[k]\\|A|=r}}h_{m,r}(x_A),
 \qquad
 \norm{f_{m,k}}_{L^p(\mu^{\otimes k})}\leq C_{m,p}(1+k)^m.
 \label{eq:prescribed-mark-degree-synopsis}
\end{equation}
Every fixed Hoeffding projection has its own all-order expansion, and its first possible inverse-particle order is $N^{-\lceil r/2\rceil}$.  The cumulant formula \eqref{eq:prescribed-connected-density} for labeled marginals defines the connected $r$-particle correlation density, whose expansion starts at order $N^{-(r-1)}$.  If factorial point-process correlations are used instead of labeled marginals, the deterministic factor $(N)_r/N^r$ must be retained.

For every fixed family of bounded centered real test functions, the joint Laplace transform of their $N^{-1/2}$-scaled linear statistics has an expansion to every prescribed order in $N^{-1/2}$, locally uniformly for complex transform parameters.  Its leading factor is the Gaussian transform with covariance \eqref{eq:static-clt-covariance}; taking an analytic logarithm near the origin gives all fixed joint moment and cumulant expansions.
\end{prop}

The transform expansions do not by themselves give Edgeworth expansions for distribution functions or densities; such conclusions require additional smoothing or Cram\'er hypotheses.

Section~\ref{sec:prescribed-marked-expansions} gives the coefficient formulas and explicit bounds in $k$. The polynomial exponent in the all-order remainder is a sufficient bound, with no optimality assertion.

\subsection{Homogeneous profiles and proportional-marginal replica expansions}
\label{subsec:replica-results-introduction}

We now specialize the response operators to the uniform reference. The strong fixed-marginal expansion and the first-order growing-marginal profile follow from the preceding prescribed-reference theorems. The replica expansion retains the nonvanishing marginal fraction and is stated separately.  Let $\Q_{N,\be}^{(k)}(\muunif)$ denote the $k$-particle marginal of $\Q_{N,\be}(\muunif)$, and define its density by
\begin{equation}
 F_{N,\be}^{(k)}
 \coloneqq
 \frac{\dd\Q_{N,\be}^{(k)}(\muunif)}{\dd\muunif^{\otimes k}}.
 \label{eq:homogeneous-marginal-density}
\end{equation}
Translation invariance gives $F_{N,\be}^{(1)}=1$ exactly.  On the real mean-zero space $L_0^2(\muunif)$, set
\begin{equation}
 B_{\be}
 \coloneqq(\Id-\be T)^{-1}-\Id,
 \qquad
 \mathcal C_{\be}\coloneqq\Tr(B_{\be}^2)=\norm{B_{\be}}_{\HS}^2.
 \label{eq:homogeneous-response-operator}
\end{equation}
The convolution kernel $b_{\be}$ of $B_{\be}$ has Fourier coefficients
\begin{equation}
 \wh b_{\be}(n)
 =\frac{\be}{\beta_{\spec}|n|^{\d}-\be},
 \qquad n\in\Z^{\d}\setminus\{0\},
 \label{eq:homogeneous-response-multiplier}
\end{equation}
and hence
\begin{equation}
 \mathcal C_{\be}
 =\sum_{n\in\Z^{\d}\setminus\{0\}}
 \left(\frac{\be}{\beta_{\spec}|n|^{\d}-\be}\right)^2.
 \label{eq:homogeneous-response-norm}
\end{equation}
The sum is over the full nonzero lattice; equivalently, it counts the real cosine and sine modes with their correct multiplicities.  For $k\geq2$, write
\begin{equation}
 S_{\be,k}(x_1,\ldots,x_k)
 \coloneqq\sum_{1\leq i<j\leq k}b_{\be}(x_i-x_j).
 \label{eq:homogeneous-screened-pair-profile}
\end{equation}
Thus, $b_{\be}$ is the pair response kernel appearing in the leading profile; it includes the resolvent correction to the original interaction kernel $\be\g$.\footnote{For every nonzero Fourier mode and $0<\be<\beta_{\spec}$, $\widehat b_\be(n)=\be\widehat\g(n)/(1-\be\widehat\g(n))>\be\widehat\g(n)$. Thus the attractive resolvent amplifies each mode relative to the first-order pair response.}

For the higher-order expansion, choose
\begin{equation}
 0\leq\be_0<\be_*<\beta_{\gm}.
 \label{eq:higher-order-buffer}
\end{equation}
The order-one estimate below the global-minimization threshold in \Cref{prop:compact-subcritical-partition-bound} gives
\begin{equation}
 \sup_{N\geq1}\sup_{0\leq c\leq\be_*}Z_{N,c}<\infty,
 \label{eq:higher-order-scalar-buffer}
\end{equation}
where $Z_{1,c}=1$.  This strict partition buffer is the only many-particle input used in the all-order expansion.  Put
\begin{equation}
 D_{\be}\coloneqq\det\nolimits_2(\Id-\be T)^{-1/2}
 \label{eq:higher-order-leading-determinant}
\end{equation}
and
\begin{equation}
 z_1(\be)
 \coloneqq\frac1{12}\int b_{\be}^3\,\dd\muunif-\frac14\mathcal C_{\be}.
 \label{eq:higher-order-z1}
\end{equation}

\begin{thm}[Homogeneous scalar and proportional-marginal replica expansions]
\label{thm:all-order-replica-expansions}
Assume \eqref{eq:higher-order-buffer}.  For every fixed integer $M\geq0$, there are coefficients $z_j(\be)$ such that
\begin{equation}
 Z_{N,\be}
 =D_{\be}\left(1+\sum_{j=1}^{M}\frac{z_j(\be)}{N^j}\right)
 +O(N^{-M-1}),
 \qquad 0\leq\be\leq\be_0,
 \label{eq:higher-order-scalar-expansion}
\end{equation}
uniformly in $\be$, and the first coefficient is given by \eqref{eq:higher-order-z1}.

Fix an integer $q\geq2$ and $0<\alpha<1$ satisfying
\begin{equation}
 \bigl[1+(q-1)\alpha\bigr]\be_0<\be_*.
 \label{eq:higher-order-qmargin}
\end{equation}
For every fixed integer $M\geq0$, there are coefficient functions $\mathcal R_{q,j}(\be,\theta)$ such that
\begin{equation}
 \int\left(F_{N,\be}^{(k)}\right)^q\,\dd\muunif^{\otimes k}
 =\sum_{j=0}^{M}\frac{\mathcal R_{q,j}(\be,k/N)}{N^j}
 +O(N^{-M-1})
 \label{eq:higher-order-replica-expansion}
\end{equation}
uniformly for integers $0\leq k\leq\alpha N$ and $0\leq\be\leq\be_0$.  The assertion is for sufficiently large $N$ and includes finiteness of the displayed moments.  The coefficient functions are smooth on $[0,\alpha]$ and analytic in $\theta$ on a fixed neighborhood of zero, uniformly in $\be$.  Their leading term is
\begin{equation}
 \mathcal R_{q,0}(\be,\theta)
 =\frac{D_{\be(1+(q-1)\theta)}D_{\be(1-\theta)}^{q-1}}{D_{\be}^q}.
 \label{eq:higher-order-leading-replica}
\end{equation}
At every prescribed order, the coefficients are finite Gaussian expectations of explicitly specified polynomials in spatially integrated Wick powers; see \Cref{app:higher-order-replica-expansions}.
\end{thm}

The expansion order $M$ and the replica number $q$ are fixed before $N$ grows.  No convergence of the full series in $N^{-1}$, and no uniformity for $q$ or $M$ growing with $N$, is asserted.  At two replicas, the leading coefficient is
\begin{equation}
 \mathcal R_{2,0}(\be,\theta)
 =\det(\Id-\theta^2B_{\be}^2)^{-1/2},
 \label{eq:higher-order-two-replica-determinant}
\end{equation}
where the right side is again an ordinary Fredholm determinant.

\begin{cor}[Homogeneous first-order profiles and the overlap determinant]
\label{thm:homogeneous-local-chaos}
Fix $0<\be_0<\beta_{\gm}$ and $0<\alpha<1$ such that
\begin{equation}
 (1+\alpha)\be_0<\beta_{\gm}.
 \label{eq:homogeneous-replica-window}
\end{equation}
There are $C,N_0<\infty$, depending only on $\d$, $\be_0$, and $\alpha$, such that the following statements hold.

For every $N\geq N_0$, $0\leq\be\leq\be_0$, and integer $0\leq\ell\leq\lfloor\alpha N\rfloor$, the marginal density $F_{N,\be}^{(\ell)}$ belongs to $L^2(\muunif^{\otimes\ell})$, and
\begin{equation}
 \left|\int\left(F_{N,\be}^{(\ell)}\right)^2\,\dd\muunif^{\otimes\ell}
 -\det\!\left(\Id-\frac{\ell^2}{N^2}B_{\be}^2\right)^{-1/2}\right|
 \leq C/N.
 \label{eq:homogeneous-partial-overlap-determinant}
\end{equation}
For $\ell=0$, both terms equal one.  The determinant in \eqref{eq:homogeneous-partial-overlap-determinant} is the ordinary Fredholm determinant, since $B_{\be}^2$ is trace class.

For every $2\leq k\leq\lfloor\alpha N\rfloor$,
\begin{equation}
 \left\|F_{N,\be}^{(k)}-1-\frac1N S_{\be,k}\right\|_{L^2(\muunif^{\otimes k})}^2
 \leq C\left(\frac{k^2}{N^4}+\frac{k^3}{N^3}\right).
 \label{eq:local-chaos-profile}
\end{equation}
Consequently,
\begin{align}
 \left|\frac{N^2}{k^2}
 \chi^2\!\left(\Q_{N,\be}^{(k)}(\muunif)\mid\muunif^{\otimes k}\right)
 -\frac{k-1}{2k}\mathcal C_{\be}\right|
 &\leq C\left(\frac1N+\frac{k}{N}\right),
 \label{eq:local-chaos-chi-rate}\\
 \left|\frac{N^2}{k^2}
 \Hr\!\left(\Q_{N,\be}^{(k)}(\muunif)\mid\muunif^{\otimes k}\right)
 -\frac{k-1}{4k}\mathcal C_{\be}\right|
 &\leq C\left(\frac1N+\frac{k}{N}\right),
 \label{eq:local-chaos-entropy-rate}\\
 \left|\frac Nk
 \TV\!\left(\Q_{N,\be}^{(k)}(\muunif),\muunif^{\otimes k}\right)
 -\frac1{2k}\int|S_{\be,k}|\,\dd\muunif^{\otimes k}\right|
 &\leq C\left(\frac1N+\sqrt{\frac{k}{N}}\right).
 \label{eq:local-chaos-TV-rate}
\end{align}
At the pair level,
\begin{align}
 \left|N^2\left\|F_{N,\be}^{(2)}-1\right\|_2^2-\mathcal C_{\be}\right|
 &\leq C/N,
 \label{eq:local-chaos-pair-norm}\\
 \left\|N\left(F_{N,\be}^{(2)}-1\right)-b_{\be}(x_1-x_2)\right\|_2^2
 &\leq C/N^2.
 \label{eq:local-chaos-pair-strong}
\end{align}
For $k=1$, all three divergences vanish identically.
\end{cor}


Let $(b_{\be,j})_{j\geq1}$ denote the real eigenvalues of $B_{\be}$, counted with multiplicity, and let $(G_j)_{j\geq1}$ be independent standard real Gaussian variables.  The centered second Gaussian chaos
\begin{equation}
 \Xi_{\be}\coloneqq\sum_{j\geq1}b_{\be,j}(G_j^2-1)
 \label{eq:homogeneous-second-chaos}
\end{equation}
converges in $L^2$ and satisfies
\begin{equation}
 \E[\Xi_{\be}^2]=2\mathcal C_{\be}.
 \label{eq:homogeneous-second-chaos-variance}
\end{equation}
The centering of each summand is essential because $B_{\be}$ need not be trace class.

\begin{cor}[Fixed- and growing-marginal constants]
\label{cor:homogeneous-local-chaos-limits}
Under the hypotheses of \Cref{thm:homogeneous-local-chaos}, for every fixed integer $k\geq2$, uniformly in $0\leq\be\leq\be_0$,
\begin{align}
 \lim_{N\to\infty}\frac{N^2}{k^2}
 \chi^2\!\left(\Q_{N,\be}^{(k)}(\muunif)\mid\muunif^{\otimes k}\right)
 &=\frac{k-1}{2k}\mathcal C_{\be},
 \label{eq:local-chaos-fixed-chi}\\
 \lim_{N\to\infty}\frac{N^2}{k^2}
 \Hr\!\left(\Q_{N,\be}^{(k)}(\muunif)\mid\muunif^{\otimes k}\right)
 &=\frac{k-1}{4k}\mathcal C_{\be},
 \label{eq:local-chaos-fixed-entropy}\\
 \lim_{N\to\infty}\frac Nk
 \TV\!\left(\Q_{N,\be}^{(k)}(\muunif),\muunif^{\otimes k}\right)
 &=\frac1{2k}\int|S_{\be,k}|\,\dd\muunif^{\otimes k}.
 \label{eq:local-chaos-fixed-TV}
\end{align}
If $k=k_N\to\infty$ and $k_N/N\to0$, then the three scaled quantities above converge, respectively, to
\begin{equation}
 \frac12\mathcal C_{\be},
 \qquad
 \frac14\mathcal C_{\be},
 \qquad
 \frac14\E[|\Xi_{\be}|].
 \label{eq:local-chaos-growing-constants}
\end{equation}
More precisely, for $2\leq k\leq\lfloor\alpha N\rfloor$,
\begin{equation}
 \left|\frac Nk
 \TV\!\left(\Q_{N,\be}^{(k)}(\muunif),\muunif^{\otimes k}\right)
 -\frac14\E[|\Xi_{\be}|]\right|
 \leq C\left[\frac1N+\sqrt{\frac{k}{N}}+k^{-1/6}\right].
 \label{eq:local-chaos-growing-TV-rate}
\end{equation}
To compare \eqref{eq:local-chaos-chi-rate} and \eqref{eq:local-chaos-entropy-rate} with the first two constants in \eqref{eq:local-chaos-growing-constants}, one adds $C/k$ to their right-hand sides.
\end{cor}

For proportional overlap $k/N\to\theta>0$ inside the admissible window, \eqref{eq:homogeneous-partial-overlap-determinant} yields the nonlinear chi-square limit
\begin{equation}
 \det(\Id-\theta^2B_{\be}^2)^{-1/2}-1.
 \label{eq:local-chaos-proportional-limit}
\end{equation}
The local quadratic profile is asserted only when $k=o(N)$.

At a fixed $\be>0$ and along a sequence $k_N/N\to\theta\in(0,\alpha]$, the quantity in \eqref{eq:local-chaos-proportional-limit} is finite and strictly positive. Thus, a positive fraction of the particles retains collective dependence detectable by chi-square divergence. The $O(N^{-1})$ error in \eqref{eq:homogeneous-partial-overlap-determinant} is relative to the determinant at the exact fraction $k/N$; replacing that fraction by its limit requires a separate estimate on its convergence.

The relation with the local-chaos constant is made explicit by the trace expansion
\begin{equation}
 \log\det(\Id-\theta^2B_\be^2)^{-1/2}
 =\frac12\sum_{r=1}^{\infty}\frac{\theta^{2r}}r\Tr(B_\be^{2r}),
 \label{eq:main-overlap-log}
\end{equation}
which converges under the strict gap. In particular,
\begin{equation}
 \det(\Id-\theta^2B_\be^2)^{-1/2}-1
 =\frac{\mathcal C_\be}{2}\theta^2
 +\left[\frac{\mathcal C_\be^2}{8}+\frac{\Tr(B_\be^4)}4\right]\theta^4
 +O(\theta^6).
 \label{eq:main-overlap-expansion}
\end{equation}
The quadratic term agrees with the growing local-chaos constant, while the higher terms contain information not resolved by the first-order profile at a positive fraction. These are expansions of the limiting determinant. The absolute error in the replica approximation may exceed the small divergence in a vanishing-fraction regime, so it does not replace \Cref{thm:prescribed-growing-marginals}. Integer replica moments alone give neither exact proportional-marginal entropy or total-variation limits nor a marginal density expansion in that regime.

Define
\begin{align}
 \mathcal L_{\be}
 &\coloneqq
 \int b_{\be}^3\,\dd\muunif
 +2\int b_{\be}^2(b_{\be}*b_{\be})\,\dd\muunif
 +\int b_{\be}^2(b_{\be}*b_{\be}*b_{\be})\,\dd\muunif
 \notag\\
 &\hspace{12mm}
 -4\Tr(B_{\be}^3)-2\Tr(B_{\be}^4).
 \label{eq:higher-order-L}
\end{align}
The pointwise cubic integral $\int b_{\be}^3\,\dd\muunif$ and the operator trace $\Tr(B_{\be}^3)$ are distinct quantities.

\begin{cor}[First finite-size corrections to chi-square divergence and relative entropy]
\label{cor:higher-order-divergences}
Under \eqref{eq:higher-order-buffer}, for every fixed integer $k\geq2$, uniformly in $0\leq\be\leq\be_0$,
\begin{equation}
 \chi^2\!\left(\Q_{N,\be}^{(k)}(\muunif)\mid\muunif^{\otimes k}\right)
 =\binom{k}{2}\frac{\mathcal C_{\be}}{N^2}
 +\binom{k}{2}\frac{\mathcal L_{\be}}{N^3}
 +O_k(N^{-4}),
 \label{eq:higher-order-chi}
\end{equation}
and
\begin{align}
 \Hr\!\left(\Q_{N,\be}^{(k)}(\muunif)\mid\muunif^{\otimes k}\right)
 &=\binom{k}{2}\frac{\mathcal C_{\be}}{2N^2}
 \notag\\
 &\quad+\frac1{N^3}\left[
 \binom{k}{2}\left(\frac{\mathcal L_{\be}}2-\frac16\int b_{\be}^3\,\dd\muunif\right)
 -\binom{k}{3}\Tr(B_{\be}^3)
 \right]
 +O_k(N^{-4}).
 \label{eq:higher-order-entropy}
\end{align}
Both expressions vanish identically for $k=1$.  Chi-square divergence admits an expansion to every fixed inverse-particle-number order by \eqref{eq:higher-order-replica-expansion}.  The strong marked theorem in \Cref{thm:prescribed-marked-expansion} also yields arbitrary finite-order relative-entropy expansions and every signed marginal coefficient.  Arbitrary-order total-variation expansions are not asserted.
\end{cor}

The homogeneous replica expansion and the prescribed-reference density expansion have complementary uniformity. The former keeps the exact marginal fraction in its coefficient functions; the latter identifies signed spatial coefficients for fixed marginal size. In both, the temperature interval is strictly below the relevant threshold, and the expansion order and replica number are fixed.

\subsection{Prescribed-reference local chaos and modulated log-HLS}
\label{subsec:prescribed-local-chaos-introduction}

The next result assumes an order-one partition bound around a prescribed reference. For some $B_\mu>0$,
\begin{equation}
 P_{\mu,b}\coloneqq
 \sup_{N\geq2}\sup_{0\leq\eta\leq b}\K_{N,\eta}(\mu)<\infty
 \qquad\text{for every }0\leq b<B_\mu.
 \label{eq:nonuniform-order-one-input}
\end{equation}
The number $B_\mu$ is a sufficient range supplied as a hypothesis, not a claimed intrinsic or optimal threshold.

\begin{thm}[Prescribed-reference local chaos]
\label{thm:prescribed-local-chaos}
Let $\mu\ll\muunif$ be a probability measure with density in $L^\infty(\muunif)$, and suppose \eqref{eq:nonuniform-order-one-input} holds with $0<B_\mu\leq2\d$. For every $0\leq\be_0<B_\mu$, there is a finite constant $C$, independent of $N$ and $\be$, such that
\begin{equation}
 \sup_{0\leq\be\leq\be_0}
 \chi^2\!\left(\Q_{N,\be}^{(1)}(\mu)\mid\mu\right)
 \leq\frac{C}{N^2},\qquad N\geq2.
 \label{eq:prescribed-one-particle-chi}
\end{equation}
If, in addition, $\al_0\in(0,1]$ and
\begin{equation}
 \be_0(1+\al_0)<B_\mu,
 \label{eq:prescribed-local-window}
\end{equation}
then there is $C=C(\d,\mu,B_\mu,\be_0,\al_0)<\infty$ such that
\begin{equation}
 \sup_{0\leq\be\leq\be_0}
 \chi^2\!\left(\Q_{N,\be}^{(k)}(\mu)\mid\mu^{\otimes k}\right)
 \leq C\frac{k^2}{N^2},
 \qquad N\geq2,\quad 1\leq k\leq\lfloor\al_0N\rfloor.
 \label{eq:prescribed-local-chi}
\end{equation}
For every $0\leq\be_0<B_\mu$, without the additional choice of $\al_0$, one has
\begin{equation}
 \sup_{0\leq\be\leq\be_0}
 \Hr\!\left(\Q_{N,\be}^{(k)}(\mu)\mid\mu^{\otimes k}\right)
 \leq C\frac{k^2}{N^2},
 \qquad N\geq2,\quad 1\leq k\leq N,
 \label{eq:prescribed-all-entropy}
\end{equation}
where $C=C(\d,\mu,B_\mu,\be_0)<\infty$. No positive lower bound on the density of $\mu$ is required.
\end{thm}

\begin{cor}[Transported local chaos]
\label{cor:transported-local-chaos}
For $\mu=\Phi_\#\muunif$ as in \Cref{thm:transported}, all conclusions of \Cref{thm:prescribed-local-chaos} hold with $B_\mu=\beta_{\tr}(\Phi)$ from \eqref{eq:transport-threshold}. In particular, the chi-square window requires $\be_0(1+\al_0)<\beta_{\tr}(\Phi)$, while the one-particle and all-marginal entropy bounds hold on every compact interval below $\beta_{\tr}(\Phi)$.
\end{cor}

The reference may be nonuniform, but it must satisfy the partition hypothesis \eqref{eq:nonuniform-order-one-input}. The transported corollary verifies that hypothesis in the range determined by the chosen map and its centered remainder. Only the order-one bound in \Cref{thm:transported} is used.

The fixed-temperature interpretation as a physical canonical ensemble is given by \Cref{rem:thermal-equilibrium-factorization}. The local-chaos estimates apply whenever the selected equilibrium satisfies the partition hypothesis.

\begin{cor}[Uniform homogeneous marginal upper bounds]
\label{cor:homogeneous-local-chaos-upper-bounds}
For every $N\geq2$ and $0\leq\be<2\d$,
\begin{equation}
 \Q_{N,\be}^{(1)}(\muunif)=\muunif.
 \label{eq:homogeneous-exact-first-marginal}
\end{equation}
If $0\leq\be_0<\beta_{\gm}$ and $\al_0\in(0,1]$ satisfy
\begin{equation}
 \be_0(1+\al_0)<\beta_{\gm},
 \label{eq:homogeneous-local-window}
\end{equation}
then there is $C=C(\d,\be_0,\al_0)<\infty$ such that
\begin{equation}
 \sup_{0\leq\be\leq\be_0}
 \chi^2\!\left(\Q_{N,\be}^{(k)}(\muunif)\mid\muunif^{\otimes k}\right)
 \leq C\frac{k(k-1)}{N^2},
 \qquad 2\leq k\leq\lfloor\al_0N\rfloor,
 \label{eq:homogeneous-local-chi}
\end{equation}
for every $N\geq2$. For every $0\leq\be_0<\beta_{\gm}$, there is $C=C(\d,\be_0)<\infty$ such that
\begin{equation}
 \sup_{0\leq\be\leq\be_0}
 \Hr\!\left(\Q_{N,\be}^{(k)}(\muunif)\mid\muunif^{\otimes k}\right)
 \leq C\frac{k^2}{N^2},
 \qquad N\geq2,\quad 1\leq k\leq N.
 \label{eq:homogeneous-all-entropy}
\end{equation}
\end{cor}
We write
\begin{equation}
 H_N(f_N\mid\mu^{\otimes N})\coloneqq\frac1N\Hr(f_N\mid\mu^{\otimes N})
 \label{eq:normalized-particle-entropy}
\end{equation}
for normalized particle entropy. The physical modulated free energy $\Er_N$ is the quantity defined in \eqref{eq:modulated-free-energy}; here $\be>0$ is implicit in $\Er_N$, and both terms in that definition are required to be finite.

Fix an integer $\ell>\d/2+2$. For the following static result, $\vep_{\d,\ell}>0$ is chosen once, depending only on $\d,\ell$, independently of $\be_0$. For $0<\be_0<\beta_{\gm}$, the common reference hypothesis is
\begin{equation}
 \mu\text{ is a probability density},\qquad
 \norm{\mu-1}_{H^\ell}
 \leq\vep_{\d,\ell}
 \frac{\beta_{\gm}-\be_0}{\beta_{\gm}(\beta_{\gm}+\be_0)}.
 \label{eq:dyn-smallness}
\end{equation}
The Sobolev norm is specified in \eqref{eq:sobolev-norm}. This choice also ensures $1/2\leq\mu\leq3/2$. Set
\begin{equation}
 \be_1=\frac{\be_0+\beta_{\gm}}2.
 \label{eq:dyn-beta1}
\end{equation}

\begin{cor}[Near-uniform modulated log-HLS inequality]
\label{cor:modulated-log-hls}
There is $C_{\mathrm{pf}}=C_{\mathrm{pf}}(\d,\ell,\be_0)\geq1$ such that, for every reference satisfying \eqref{eq:dyn-smallness}, every $N\geq2$, every $0<\be\leq\be_0$, and every probability law $f_N$ with finite $H_N(f_N\mid\mu^{\otimes N})$ and $\E_{f_N}[|\Fr_N(X_N,\mu)|]$,
\begin{equation}
 \Er_N(f_N,\mu)
 \geq\left(\frac1\be-\frac1{\be_1}\right)H_N(f_N\mid\mu^{\otimes N})
 -\frac{\log C_{\mathrm{pf}}}{\be_1N}.
 \label{eq:dyn-coercivity}
\end{equation}
Consequently,
\begin{equation}
 H_N(f_N\mid\mu^{\otimes N})
 \leq\frac{\be\be_1}{\be_1-\be_0}
 \left(\Er_N(f_N,\mu)+\frac{\log C_{\mathrm{pf}}}{\be_1N}\right).
 \label{eq:dyn-entropy-beta}
\end{equation}
The same constant works over the entire reference neighborhood. It is the partition constant in \Cref{prop:family-uniform-partition}.
\end{cor}



\subsection{Static joint linear statistics}
\label{subsec:static-clt-introduction}

The static central limit theorem requires only a strict partition buffer, integrability of the one-body term, and square integrability of the centered kernel. It therefore applies to reference measures beyond the transported classes in \Cref{thm:principal}.

Given a configuration $\XN=(x_1,\ldots,x_N)$, let
\begin{equation}
 \mu_N\coloneqq\frac1N\sum_{i=1}^N\delta_{x_i}.
 \label{eq:empirical-measure}
\end{equation}
For $\varphi\in L^2(\mu)$, set
\begin{equation}
 S_N(\varphi)
 \coloneqq\sqrt N\left(\int_{\T^{\d}}\varphi\,\dd\mu_N
 -\int_{\T^{\d}}\varphi\,\dd\mu\right).
 \label{eq:linear-statistic}
\end{equation}

\begin{cor}[Static linear statistics under a strict partition buffer]
\label{cor:static-clt}
Let $\mu$ be nonatomic, assume $h_{\mu}\in L^1(\mu)$ and $\tl{\g}_{\mu}\in L^2(\mu^{\otimes2})$, and fix $\be\geq0$.  Suppose that there is $b>\be$ such that
\begin{equation}
 \sup_{N\geq2}\K_{N,b}(\mu)<\infty.
 \label{eq:clt-buffer}
\end{equation}
Then
\begin{equation}
 b\norm{T_{\mu}}_{\op}<1.
 \label{eq:clt-spectral-gap}
\end{equation}
For every finite family of real-valued functions $\varphi_1,\ldots,\varphi_k\in L^2(\mu)$,
\begin{equation}
 \bigl(S_N(\varphi_1),\ldots,S_N(\varphi_k)\bigr)
 \Longrightarrow\mathcal N(0,\Sigma_{\be,\mu})
 \label{eq:static-vector-clt}
\end{equation}
under $\Q_{N,\be}(\mu)$, where
\begin{equation}
 \Sigma_{\be,\mu}(\varphi,\psi)
 \coloneqq\ip*{\varphi-\int\varphi\,\dd\mu,
 (\Id-\be T_{\mu})^{-1}
 \left(\psi-\int\psi\,\dd\mu\right)}_{L^2(\mu)}.
 \label{eq:static-clt-covariance}
\end{equation}
If the test functions are bounded, then their joint moment-generating functions converge on all of $\R^k$.

In particular, for $\mu=\muunif$, the conclusion holds for every $\be<\beta_{\gm}$.  It also holds for a transported reference throughout the strict range in \eqref{eq:transport-threshold}, and hence for the near-uniform densities covered by the final clause of \Cref{thm:principal}.  The scalar factor in \eqref{eq:principal-transport-rate} cancels from the normalized Gibbs law \eqref{eq:modulated-gibbs-law}.
\end{cor}


\section{Collision integrability and preliminary estimates}

\label{sec:collisions-preliminaries}

We first fix the functional domains and prove the collision criterion. A subset-marginal argument then gives a linear-in-$N$ bound for the logarithm of the partition function. The remaining estimates concern canonical arrays and marginal projections; they will be used in the covariance decomposition and the local-chaos arguments.

\subsection{Fourier positivity and functional domains}
\label{subsec:functional-domains}

We use the Fourier convention
\begin{equation}
 \wh f(k)=\int_{\T^{\d}}f(x)e^{-2\pi i k\cdot x}\,\dd x,
 \qquad
 \wh{\lvert\nabla\rvert^s f}(k)=(2\pi|k|)^s\wh f(k).
 \label{eq:torus-fourier-normalization}
\end{equation}
The mean-zero periodic logarithmic kernel is normalized by
\begin{equation}
 c_{\d}=2^{\d-1}\pi^{\d/2}\Gamma(\d/2),
 \qquad
 |\nabla|^{\d}\g=c_{\d}(\delta_0-1),
 \qquad
 \wh\g(k)=\frac{c_{\d}}{(2\pi|k|)^{\d}}
 \quad(k\neq0).
 \label{eq:kernel-normalization}
\end{equation}
Thus, $\g$ is real, even, bounded from below, and satisfies \eqref{eq:kernel-local-form}.  Its largest Fourier coefficient occurs on $|k|=1$, giving the spectral threshold in \eqref{eq:spectral-threshold}.
The ordered- and unordered-pair conventions agree exactly:
\begin{equation}
 \frac1{2N}\sum_{i\neq j}\g(x_i-x_j)
 =\frac1N\sum_{i<j}\g(x_i-x_j).
 \label{eq:normalization-pair-sum-crosswalk}
\end{equation}

For a finite signed measure $\nu$ on $\T^{\d}$ with $\nu(\T^{\d})=0$, extend the interaction-energy notation by
\begin{equation}
 I_{\nu}
 \coloneqq-\frac12\sum_{k\in\Z^{\d}\setminus\{0\}}
 \frac{c_{\d}}{(2\pi|k|)^{\d}}\abs{\wh\nu(k)}^2
 \in[-\infty,0].
 \label{eq:signed-measure-energy}
\end{equation}
If $\nu=f\muunif$ with $f\in L_0^2(\muunif)$, then
\begin{equation}
 -2I_{f\muunif}=\ip*{f,Tf}_{L^2(\muunif)}.
 \label{eq:energy-operator-form}
\end{equation}
For a probability measure $\mu$, the zero-mean normalization of $\g$ gives $I_{\mu}=I_{\mu-\muunif}$, and \eqref{eq:macroscopic-energy} is the density form of \eqref{eq:signed-measure-energy}.  Thus, $-I_{\mu}$ is nonnegative and lower semicontinuous.

For every $a\in(0,\d)$, \eqref{eq:kernel-local-form} gives $e^{a\g}\in L^1(\muunif^{\otimes2})$.  Applying the entropy inequality to $\mu^{\otimes2}$ therefore yields
\begin{equation}
 0\leq-2I_{\mu}
 \leq\frac{2}{a}\Hr(\mu\mid\muunif)
 +\frac1a\log\iint_{(\T^{\d})^2}e^{a\g(x-y)}
 \,\dd\muunif(x)\dd\muunif(y)<\infty
 \label{eq:finite-entropy-energy-bound}
\end{equation}
whenever $\Hr(\mu\mid\muunif)<\infty$.  In particular, the free energy \eqref{eq:macroscopic-free-energy} and pressure \eqref{eq:macroscopic-pressure} are well defined on their stated finite-entropy domains.

For every $b<\beta_{\gm}$, the definition \eqref{eq:global-minimization-threshold} gives
\begin{equation}
 \Hr(\mu\mid\muunif)\geq-bI_{\mu}
 \qquad\text{for every probability measure $\mu$ with $\Hr(\mu\mid\muunif)<\infty$}.
 \label{eq:strict-subcritical-macroscopic-inequality}
\end{equation}
The particle bound below uses this inequality only with $b<\beta_{\gm}$.

The concentration and Hessian tests give the upper bounds in \eqref{eq:threshold-hierarchy-intro}.  The Hessian computation is
\begin{align}
 \Hr(1+\vep f\mid\muunif)
 &=\frac{\vep^2}{2}\norm{f}_{L^2(\muunif)}^2+O_f(\vep^3),
 \label{eq:entropy-second-variation}\\
 I_{1+\vep f}
 &=-\frac{\vep^2}{2}\ip*{f,Tf}_{L^2(\muunif)},
 \label{eq:energy-second-variation}
\end{align}
for bounded mean-zero $f$ and sufficiently small $\vep$.  Taking the supremum of the Rayleigh quotient yields $\beta_{\gm}\leq\beta_{\spec}$.  A smooth probability density $\mu$ concentrated in a ball of radius $\vep$ satisfies
\begin{equation}
 \Hr(\mu\mid\muunif)=\d\log(1/\vep)+O(1),
 \qquad
 I_{\mu}=-\frac12\log(1/\vep)+O(1),
 \label{eq:concentration-threshold-test}
\end{equation}
which gives $\beta_{\gm}\leq2\d$.


We use the inhomogeneous Sobolev norm
\begin{equation}
 \norm f_{H^\ell}^2\coloneqq\sum_{k\in\Z^{\d}}
 (1+4\pi^2|k|^2)^\ell|\wh f(k)|^2.
 \label{eq:sobolev-norm}
\end{equation}
The local logarithmic form and smoothness away from zero give
\begin{equation}
 \sup_{x\neq0}\operatorname{dist}_{\T^{\d}}(x,0)|\nabla\g(x)|<\infty.
 \label{eq:logarithmic-gradient-bound}
\end{equation}


For bounded reference densities, the exact splitting \eqref{eq:intro-modulated-splitting} is equivalently
\begin{equation}
 \Fr_N(\XN,\mu)=-\frac12\iint_{x\neq y}\g(x-y)
 \,\dd(\mu_N-\mu)^{\otimes2}(x,y).
 \label{eq:empirical-modulated-energy}
\end{equation}
Only the empirical self-interactions are removed. The equivalent Hoeffding-centered formula is \eqref{eq:modulated-energy-general-reference}, and the modulated Gibbs law is defined in \eqref{eq:modulated-gibbs-law}. On the finite-entropy, finite-energy domain, the free-energy rewriting is \eqref{eq:modulated-free-energy-rewriting}.


\begin{remark}[Thermal equilibria and the residual physical normalization]
\label{rem:thermal-equilibrium-factorization}
Let $V$ be a one-body potential and consider
\begin{equation}
 Z_{N,\be}^{V}
 \coloneqq\int_{(\T^{\d})^N}
 \exp\left\{\frac{\be}{2N}\sum_{1\leq i\neq j\leq N}\g(x_i-x_j)
 -\be\sum_{i=1}^NV(x_i)\right\}
 \,\dd\muunif^{\otimes N}(\XN).
 \label{eq:external-potential-partition}
\end{equation}
Thus, $Z_{N,\be}^{0}=Z_{N,\be}$. The external-field free energy is
\begin{equation}
 \Ec_{\be,V}(\nu)
 \coloneqq\Hr(\nu\mid\muunif)+\be I_{\nu}+\be\int_{\T^{\d}}V\,\dd\nu.
 \label{eq:external-macroscopic-free-energy}
\end{equation}
A chosen thermal equilibrium $\mu_{\be,V}$ with positive density and for which these expressions are finite satisfies the Euler--Lagrange relation
\begin{equation}
 \frac{\dd\mu_{\be,V}}{\dd\muunif}(x)
 =\frac{e^{-\be(V(x)+h_{\mu_{\be,V}}(x))}}
 {\displaystyle\int_{\T^{\d}}e^{-\be(V(y)+h_{\mu_{\be,V}}(y))}\,\dd\muunif(y)}.
 \label{eq:thermal-equilibrium-euler-lagrange}
\end{equation}
Inserting \eqref{eq:thermal-equilibrium-euler-lagrange} into the splitting \eqref{eq:intro-modulated-splitting} gives the exact factorization
\begin{equation}
 Z_{N,\be}^{V}
 =e^{-N\Ec_{\be,V}(\mu_{\be,V})}\K_{N,\be}(\mu_{\be,V}).
 \label{eq:thermal-equilibrium-factorization}
\end{equation}
Moreover, normalizing the integrand in \eqref{eq:external-potential-partition} gives precisely the modulated Gibbs law $\Q_{N,\be}(\mu_{\be,V})$.  Thus, $\K_{N,\be}(\mu_{\be,V})$ is the residual normalization left after the equilibrium mean-field free energy has been removed.  At a fixed inverse temperature, the prescribed-reference formulation includes any thermal equilibrium to which the regularity and strict-range hypotheses of \Cref{thm:transported} apply.

\end{remark}

\subsection{The exact fixed-particle collision threshold}
\label{subsec:fixed-particle-collisions}

We begin with a Euclidean power-counting lemma that treats every partial diagonal.

\begin{lemma}[Graphical collision power counting]
\label{lem:nested-cluster-power-count}
Let $v\geq2$ and let $a_{ij}\geq0$ for $1\leq i<j\leq v$. Then
\begin{equation}
 \prod_{i<j}\operatorname{dist}_{\T^\d}(x_i,x_j)^{-a_{ij}}
 \in L^1((\T^\d)^v)
 \quad\Longleftrightarrow\quad
 \sum_{\substack{i<j\\ i,j\in S}}a_{ij}<\d(|S|-1)
 \quad\text{for every }S\subset\{1,\ldots,v\},\ |S|\geq2.
 \label{eq:graph-collision-criterion}
\end{equation}
The same criterion gives local integrability on $(\R^\d)^v$. In particular, for a common exponent $a\geq0$, the product
\begin{equation}
 \prod_{1\leq i<j\leq v}|x_i-x_j|^{-a}
 \label{eq:euclidean-collision-product}
\end{equation}
is locally integrable if and only if
\begin{equation}
 a\binom{s}{2}<\d(s-1),\qquad 2\leq s\leq v.
 \label{eq:cluster-integrability-criterion}
\end{equation}
\end{lemma}
\begin{proof}
For necessity, fix a set $S$ of $s\geq2$ vertices. Keep all other particles in disjoint neighborhoods away from a small ball containing the particles of $S$. In relative coordinates for $S$, restrict the angular variables to a compact set of pairwise separated shapes of positive surface measure. The radial part of the integral contains
\begin{equation}
 \int_0^{r_0}r^{\d(s-1)-1-\sum_{i<j,\ i,j\in S}a_{ij}}\,\dd r.
 \label{eq:graph-necessary-radial}
\end{equation}
It diverges when the corresponding strict inequality in \eqref{eq:graph-collision-criterion} fails, logarithmically when its exponent equals $-1$.

For sufficiency, ignore the null set of coincident particles and choose a minimum spanning tree of the complete graph with edge lengths equal to the particle distances. Order its edges so that $0<r_1\leq\cdots\leq r_{v-1}$, resolving ties measurably. There are only finitely many possible trees and orders. For any pair of vertices, its distance is at least the largest edge length on the tree path joining them; otherwise replacing that edge by the pair would shorten the tree. Assign $a_{ij}$ to the last edge on this path in the chosen order, and let $A_h$ be the total assigned to edge $h$. The singular product is bounded above by $\prod_h r_h^{-A_h}$.

The first $j$ edges form a forest. A pair has been assigned to one of these edges exactly when its endpoints belong to the same forest component. Consequently,
\begin{equation}
 \sum_{h=1}^j A_h
 =\sum_{C\text{ a component}}\ \sum_{\substack{i<l\\i,l\in C}}a_{il}
 <\d\sum_C(|C|-1)=\d j,
 \qquad 1\leq j\leq v-1.
 \label{eq:graph-prefix-weights}
\end{equation}
Singleton components contribute zero, and at least one component is nonsingleton when $j\geq1$.

Root the tree and use the root position and its oriented edge differences as coordinates. This is a measure-preserving change of variables on the product torus. Represent each difference in a fundamental cube, enlarge that cube to a fixed ball of radius $D$, and integrate the angular variables. Each tree/order sector is bounded by a constant times
\begin{equation}
 \int_{0<r_1<\cdots<r_{v-1}<D}
 \prod_{h=1}^{v-1}r_h^{\d-1-A_h}\,\dd r_1\cdots\dd r_{v-1}.
 \label{eq:graph-tree-sector}
\end{equation}
Substitute $r_h=D\prod_{j=h}^{v-1}u_j$, with $0<u_j<1$. Apart from a finite power of $D$, the integral becomes
\begin{equation}
 \prod_{j=1}^{v-1}\int_0^1
 u_j^{\d j-1-\sum_{h=1}^j A_h}\,\dd u_j<\infty
 \label{eq:graph-sector-product}
\end{equation}
by \eqref{eq:graph-prefix-weights}. On a bounded Euclidean configuration domain the same tree coordinates have Jacobian one and bounded edge lengths, so the argument is unchanged. Substitution of the common exponent gives \eqref{eq:cluster-integrability-criterion}.
\end{proof}

\begin{prop}[Fixed-particle integrability]
\label{prop:fixed-particle-integrability}
For every $N\geq2$,
\begin{equation}
 Z_{N,\be}<\infty
 \quad\Longleftrightarrow\quad
 \be<2\d.
 \label{eq:fixed-particle-threshold}
\end{equation}
At $\be=2\d$, the divergence at the full diagonal is logarithmic.
\end{prop}

\begin{proof}
Away from the pair diagonals, $\g$ is bounded.  Near a collision stratum, \eqref{eq:kernel-local-form} shows that the density in \eqref{eq:homogeneous-partition-function} is bounded above and below by positive smooth factors times
\begin{equation}
 \prod_{1\leq i<j\leq N}|x_i-x_j|^{-\be/N}.
 \label{eq:partition-collision-product}
\end{equation}
By \Cref{lem:nested-cluster-power-count}, a cluster of $r$ particles is integrable precisely when
\begin{equation}
 \frac{\be}{N}\binom r2<\d(r-1),
 \qquad\text{equivalently,}\qquad
 \be<\frac{2\d N}{r}.
 \label{eq:cluster-beta-condition}
\end{equation}
The right side of \eqref{eq:cluster-beta-condition} decreases with $r$, so the full cluster $r=N$ gives the strongest condition $\be<2\d$.  At equality, the exponent in the full-diagonal radial integral is $-1$, and the divergence is logarithmic.  For $\be>2\d$, it is a power divergence.
\end{proof}

\subsection{Subset marginals and a linear bound on the log partition function}
\label{subsec:linear-safety}

The bad cell-count event in \Cref{subsec:order-one-assembly} will be estimated by H\"older's inequality. For this step, a bound linear in $N$ for $\log Z_{N,\be}$ is sufficient. We obtain it directly from fixed-particle integrability, independently of the order-one theorem.

\begin{lemma}[Relative-entropy subset inequality]
\label{lem:relative-entropy-subset}
Let $P_N$ be a symmetric probability measure on $E^N$, let $\mu$ be a probability measure on $E$, and let $P_N^{(k)}$ denote the $k$-particle marginal.  Then
\begin{equation}
 \frac Nk\Hr(P_N^{(k)}\mid\mu^{\otimes k})
 \leq\Hr(P_N\mid\mu^{\otimes N})
 \qquad(1\leq k\leq N).
 \label{eq:relative-entropy-subset}
\end{equation}
\end{lemma}

\begin{proof}
Let $\mathcal S_k$ be the family of $k$-subsets of $\{1,\ldots,N\}$ and put $r=\binom{N-1}{k-1}$.  Every coordinate belongs to exactly $r$ members of $\mathcal S_k$.  Generalized H\"older gives, for bounded functions $\varphi_S$ on $E^S$,
\begin{equation}
 \int\exp\left\{\frac1r\sum_{S\in\mathcal S_k}\varphi_S(x_S)\right\}
 \dd\mu^{\otimes N}
 \leq\prod_{S\in\mathcal S_k}
 \left(\int e^{\varphi_S}\,\dd\mu^{\otimes S}\right)^{1/r}.
 \label{eq:holder-subset-product}
\end{equation}
Apply the entropy variational formula with the test function in \eqref{eq:holder-subset-product} and take the supremum independently in each $\varphi_S$.  This yields
\begin{equation}
 \sum_{S\in\mathcal S_k}\Hr(P_{N,S}\mid\mu^{\otimes S})
 \leq r\Hr(P_N\mid\mu^{\otimes N}).
 \label{eq:entropy-shearer}
\end{equation}
By symmetry, every term on the left side of \eqref{eq:entropy-shearer} equals $\Hr(P_N^{(k)}\mid\mu^{\otimes k})$, and $\binom Nk/r=N/k$.
\end{proof}

\begin{lemma}[Subset-marginal partition inequality]
\label{lem:subset-marginal-partition}
For $2\leq k\leq N$ and $\be\geq0$,
\begin{equation}
 \log Z_{N,\be}
 \leq\frac Nk\log Z_{k,\be k(N-1)/(N(k-1))}.
 \label{eq:subset-marginal-partition}
\end{equation}
\end{lemma}

\begin{proof}
First replace $\g$ by $\min\{\g,L\}$.  The bounded-kernel Gibbs variational formula and symmetrization reduce the supremum to symmetric probability measures $P_N$.  Exchangeability gives
\begin{equation}
 \E_{P_N}\left[\sum_{1\leq i<j\leq N}\g(x_i-x_j)\right]
 =\frac{N(N-1)}{k(k-1)}
 \E_{P_N^{(k)}}\left[\sum_{1\leq i<j\leq k}\g(x_i-x_j)\right].
 \label{eq:exchangeable-pair-reduction}
\end{equation}
Combining \eqref{eq:exchangeable-pair-reduction} with \Cref{lem:relative-entropy-subset} in the variational formula gives \eqref{eq:subset-marginal-partition} for the truncated kernel.  Since $\g$ is bounded from below, monotone convergence as $L\to\infty$ proves the result.
\end{proof}

\begin{cor}[Linear bound on the log partition function]
\label{cor:linear-partition-bound}
For every $b_0<2\d$, there is $C=C(\d,b_0)<\infty$ such that
\begin{equation}
 \sup_{0\leq\be\leq b_0}\log Z_{N,\be}
 \leq CN
 \qquad(N\geq2).
 \label{eq:linear-partition-bound}
\end{equation}
\end{cor}

\begin{proof}
Choose $k\geq2$ so that $b_0k/(k-1)<2\d$.  For $N\geq k$, \Cref{lem:subset-marginal-partition} and \Cref{prop:fixed-particle-integrability} give
\begin{equation}
 \log Z_{N,\be}
 \leq\frac Nk
 \sup_{0\leq c\leq b_0k/(k-1)}\log Z_{k,c}.
 \label{eq:linear-bound-fixed-k}
\end{equation}
The supremum in \eqref{eq:linear-bound-fixed-k} is finite.  Indeed, if
\begin{equation}
 Y_k\coloneqq\frac1k\sum_{1\leq i<j\leq k}\g(x_i-x_j),
 \label{eq:fixed-k-energy}
\end{equation}
then, for $0\leq c\leq c_*<2\d$,
\begin{equation}
 e^{cY_k}\leq1+e^{c_*Y_k}.
 \label{eq:fixed-k-dominating-function}
\end{equation}
The right side of \eqref{eq:fixed-k-dominating-function} is integrable by \Cref{prop:fixed-particle-integrability}, so dominated convergence gives continuity in $c$.  The finitely many values $2\leq N<k$ are absorbed by enlarging the constant.
\end{proof}

\subsection{Canonical arrays}
\label{subsec:canonical-arrays}

After conditioning on cell labels, the particles are independent but their laws depend on the cells. We therefore need a canonical-array estimate without an identical-distribution hypothesis.

\begin{lemma}[Canonical-array exponential estimate]
\label{lem:canonical-array-exponential}
There is a universal constant $C_{\mathrm u}<\infty$ with the following property.  Let $X_1,\ldots,X_m$ be independent random variables, and let $h_{ij}(X_i,X_j)$, $1\leq i<j\leq m$, be real kernels satisfying
\begin{align}
 \E[h_{ij}(X_i,y)]&=0
 &&\text{for every admissible }y,
 \label{eq:canonical-array-first}\\
 \E[h_{ij}(x,X_j)]&=0
 &&\text{for every admissible }x.
 \label{eq:canonical-array-second}
\end{align}
Assume that, for every integer $p\geq1$,
\begin{equation}
 \sup_y\E[|h_{ij}(X_i,y)|^p]
 +\sup_x\E[|h_{ij}(x,X_j)|^p]
 \leq2L^p p!
 \label{eq:canonical-array-kernel-moments}
\end{equation}
for every $i<j$.  Then, for every $k\geq1$,
\begin{equation}
 \E\left[\left|\sum_{i<j}h_{ij}(X_i,X_j)\right|^{2k}\right]
 \leq(C_{\mathrm u}L)^{2k}m^{2k}(2k)!.
 \label{eq:canonical-array-moment}
\end{equation}
Consequently, if $|\theta|C_{\mathrm u}L<1$, then
\begin{equation}
 \E\left[\exp\left\{\frac\theta m\sum_{i<j}h_{ij}(X_i,X_j)\right\}\right]
 \leq\frac{2}{1-(|\theta|C_{\mathrm u}L)^2}.
 \label{eq:canonical-array-exponential}
\end{equation}
The same conclusion holds when some kernels vanish identically.
\end{lemma}

\begin{proof}
Extend the array symmetrically by $h_{ji}(y,x)=h_{ij}(x,y)$ and $h_{ii}=0$, and let $X_1',\ldots,X_m'$ be an independent copy.  The degree-two decoupling inequality of de la Pe\~na and Montgomery-Smith \cite[Theorem~1]{delaPenaMontgomerySmith1995}, integrated against $2kt^{2k-1}\,\dd t$, yields a universal constant $C_{\mathrm{dec}}$ such that
\begin{equation}
 \E\left[\left|\sum_{i<j}h_{ij}(X_i,X_j)\right|^{2k}\right]
 \leq C_{\mathrm{dec}}^{2k}
 \E\left[\left|\sum_{i\neq j}h_{ij}(X_i,X_j')\right|^{2k}\right].
 \label{eq:canonical-array-decoupling}
\end{equation}
Expand the right side of \eqref{eq:canonical-array-decoupling}.  For words $\mathbf i=(i_1,\ldots,i_{2k})$ and $\mathbf j=(j_1,\ldots,j_{2k})$, set
\begin{equation}
 a_r\coloneqq\#\{\ell:i_\ell=r\},
 \qquad
 b_r\coloneqq\#\{\ell:j_\ell=r\}.
 \label{eq:word-multiplicities}
\end{equation}
By \eqref{eq:canonical-array-first}--\eqref{eq:canonical-array-second}, the expectation vanishes if any positive $a_r$ or $b_r$ equals one.  For every surviving pair of words, condition on all primed variables.  Independence, H\"older's inequality with exponent $a_r$ for the factors carrying the label $r$, and \eqref{eq:canonical-array-kernel-moments} give
\begin{equation}
 \left|\E\left[\prod_{\ell=1}^{2k}
 h_{i_\ell j_\ell}(X_{i_\ell},X_{j_\ell}')
 \mathrel{\Big|}X_1',\ldots,X_m'\right]\right|
 \leq2^kL^{2k}\prod_{r=1}^m a_r!.
 \label{eq:surviving-word-holder}
\end{equation}
A surviving word has at most $k$ active labels, which explains the factor $2^k$ in \eqref{eq:surviving-word-holder}.

For a fixed multiplicity vector $a$, the number of $\mathbf i$-words is $(2k)!/\prod_r a_r!$.  Hence the factorials in \eqref{eq:surviving-word-holder} cancel, and, after dropping the restrictions $i_\ell\neq j_\ell$, the remaining count is at most
\begin{equation}
 (2k)!\,
 \#\{\text{surviving multiplicity vectors }a\}\,
 \#\{\text{surviving }\mathbf j\text{-words}\}.
 \label{eq:canonical-array-count-factorization}
\end{equation}
If $m<4k$, the first count in \eqref{eq:canonical-array-count-factorization} is at most $\binom{2k+m-1}{m-1}\leq2^{6k}$ and the second is at most $m^{2k}$.  If $m\geq4k$, the number of surviving words with exactly $r$ active labels and the number of corresponding multiplicity vectors are bounded by
\begin{equation}
 \binom mr r^{2k}
 \qquad\text{and}\qquad
 \binom mr\binom{2k-r-1}{r-1},
 \quad 1\leq r\leq k,
 \label{eq:canonical-array-two-counts}
\end{equation}
respectively.  The elementary estimates $\binom mr\leq(em/r)^r$ and $m/k\geq4$ imply
\begin{align}
 \sum_{r=1}^k\binom mr r^{2k}
 &\leq C^k m^k k^k,
 \label{eq:canonical-array-word-bound}\\
 \sum_{r=1}^k\binom mr\binom{2k-r-1}{r-1}
 &\leq C^k\frac{m^k}{k^k}.
 \label{eq:canonical-array-multiplicity-bound}
\end{align}
Combining \eqref{eq:canonical-array-decoupling}--\eqref{eq:canonical-array-multiplicity-bound} and absorbing $2^k$ into the universal constant proves \eqref{eq:canonical-array-moment}.

Finally, $e^x\leq2\cosh x$ and \eqref{eq:canonical-array-moment} give
\begin{equation}
 \E\left[e^{\theta\sum_{i<j}h_{ij}/m}\right]
 \leq2\sum_{k=0}^{\infty}
 (|\theta|C_{\mathrm u}L)^{2k},
 \label{eq:canonical-array-cosh-series}
\end{equation}
which is \eqref{eq:canonical-array-exponential}.
\end{proof}


For the bounded pullback remainder in the transport argument, the next estimate gives an explicit admissible coupling in terms of the supremum norm.


\begin{lemma}[Bounded canonical kernel]
\label{lem:bounded-canonical-kernel}
Let $X_1,\ldots,X_N$ be independent with common law $\mu$, and let $h$ be a bounded real symmetric kernel satisfying
\begin{equation}
 \int h(x,y)\,\dd\mu(y)=0
 \qquad\text{for }\mu\text{-almost every }x.
 \label{eq:bounded-kernel-canonicality}
\end{equation}
Put $H=\norm{h}_{L^\infty(\mu^{\otimes2})}$.  If $4|t|H<1$, then
\begin{equation}
 \E\left[\exp\left\{\frac tN\sum_{i<j}h(X_i,X_j)\right\}\right]
 \leq\exp\left\{\frac{8t^2H^2}{1-16t^2H^2}\right\}.
 \label{eq:bounded-canonical-exponential}
\end{equation}
\end{lemma}

\begin{proof}
Let $I\sqcup J=\{1,\ldots,N\}$ be a random bipartition obtained from independent fair signs.  Every unordered pair is separated with probability $1/2$, so Jensen's inequality gives
\begin{equation}
 \E\left[e^{t\sum_{i<j}h(X_i,X_j)/N}\right]
 \leq\E_{I,J}\E\left[
 \exp\left\{\frac{2t}{N}\sum_{i\in I,\,j\in J}h(X_i,X_j)\right\}\right].
 \label{eq:random-bipartition-jensen}
\end{equation}
Fix $I,J$.  Apply conditional Jensen first in the $I$ variables and then in the $J$ variables, using independent copies $X_i'$ and $X_j'$.  The canonicality \eqref{eq:bounded-kernel-canonicality} and the symmetry obtained by swapping each pair with its copy allow independent Rademacher signs $\sigma_i$ and $\tau_j$ to be inserted.  The inner expectation in \eqref{eq:random-bipartition-jensen} is at most
\begin{equation}
 \E\left[\exp\left\{\frac{2t}{N}
 \sum_{i\in I,\,j\in J}\sigma_i\tau_j M_{ij}\right\}\right],
 \label{eq:double-symmetrized-bilinear-form}
\end{equation}
where
\begin{equation}
 M_{ij}\coloneqq h(X_i,X_j)-h(X_i',X_j)-h(X_i,X_j')+h(X_i',X_j').
 \label{eq:symmetrized-matrix}
\end{equation}
Thus, $|M_{ij}|\leq4H$ and
\begin{equation}
 \norm{M}_{\HS}\leq2HN,
 \qquad
 \norm{M}_{\op}\leq2HN.
 \label{eq:symmetrized-matrix-norms}
\end{equation}

For a deterministic matrix $M$ and independent Rademacher vectors $\sigma,\tau$, two applications of the scalar sub-Gaussian inequality, with a Hubbard--Stratonovich transformation between them, give
\begin{equation}
 \E_{\sigma,\tau}\left[e^{u\sigma^{\mathsf T}M\tau}\right]
 \leq\det(\Id-u^2MM^{\mathsf T})^{-1/2}
 \leq\exp\left\{\frac{u^2\norm{M}_{\HS}^2}
 {2(1-u^2\norm{M}_{\op}^2)}\right\}
 \label{eq:bilinear-rademacher-bound}
\end{equation}
whenever $|u|\norm{M}_{\op}<1$.  Taking $u=2t/N$ and using \eqref{eq:symmetrized-matrix-norms} gives \eqref{eq:bounded-canonical-exponential}.  The bound is uniform in the bipartition, so \eqref{eq:random-bipartition-jensen} completes the proof.
\end{proof}

\subsection{Product measures and marginal identities}
\label{subsec:shared-marginal-preliminaries}

For a bounded probability density $\mu$ and $\be>0$, direct evaluation of \eqref{eq:intro-modulated-splitting} under independent coordinates gives
\begin{equation}
 \Er_N(\mu^{\otimes N},\mu)=\E_{\mu^{\otimes N}}[\Fr_N(X_N,\mu)]= -\frac{I_\mu}{N}.
 \label{eq:dyn-product-exact}
\end{equation}
By the Fourier coefficients in \eqref{eq:kernel-normalization},
\begin{equation}
 0\leq-I_\mu=\frac1{2\beta_{\spec}}\sum_{k\ne0}|k|^{-\d}|\wh{\mu-1}(k)|^2
 \leq\frac1{2\beta_{\spec}}\norm{\mu-1}_{L^2}^2.
 \label{eq:dyn-product-I}
\end{equation}

The Hoeffding decomposition \cite{Hoeffding1948} gives a common orthogonality identity for the nonasymptotic bounds and the growing-marginal profile.

\begin{lemma}[Hoeffding decomposition and marginal projection inequality]
\label{lem:local-Hoeffding}
Let $F_M$ be a symmetric probability density in $L^2(\mu^{\otimes M})$. Write $F^{(k)}=F_M^{(k)}$ for its marginal density, with $F^{(0)}=1$. For $1\leq r\leq M$, define
\begin{equation}
 h_r(x_1,\ldots,x_r)
 =\sum_{A\subseteq\{1,\ldots,r\}}(-1)^{r-|A|}F^{(|A|)}(x_A).
 \label{eq:local-Hoeffding-projection}
\end{equation}
The kernels $h_r$ are canonical in every variable, and
\begin{equation}
 F^{(k)}=1+\sum_{r=1}^k\sum_{\substack{A\subseteq\{1,\ldots,k\}\\|A|=r}}h_r(x_A),
 \label{eq:prescribed-growing-Hoeffding-expansion}
\end{equation}
with the isometry
\begin{equation}
 \norm{F^{(k)}-1}_2^2=\sum_{r=1}^k\binom kr\norm{h_r}_2^2,
 \qquad 1\leq k\leq M.
 \label{eq:local-Hoeffding-orthogonality}
\end{equation}
Consequently,
\begin{equation}
 \norm{F_M^{(k)}-1}_2^2
 \leq k\norm{F_M^{(1)}-1}_2^2
 +\frac{k(k-1)}{M(M-1)}\norm{F_M-1}_2^2,
 \label{eq:local-Hoeffding-inequality}
\end{equation}
where the second term is omitted when $M=1$.
\end{lemma}
\begin{proof}
Marginal consistency cancels the terms in pairs after integration in any one variable, so each $h_r$ is canonical. Inclusion--exclusion gives the displayed decomposition. Summands indexed by different subsets are orthogonal: integrate in a coordinate belonging to one subset and not the other. Symmetry identifies the common kernel at each cardinality, which proves the isometry for every $k$, including $k=M$. The first kernel is $h_1=F^{(1)}-1$. For $2\leq r\leq k$, the product formula gives
\begin{equation}
 \frac{\binom kr}{\binom Mr}
 =\prod_{i=0}^{r-1}\frac{k-i}{M-i}
 \leq\frac{k(k-1)}{M(M-1)}.
 \label{eq:local-Hoeffding-binomial-ratio}
\end{equation}
Separate the first level in the isometry and use this inequality for all remaining levels.
\end{proof}

\subsection{A global entropy remainder}
\label{subsec:shared-entropy-remainder}

For $u\geq-1$, set $\phi(u)=(1+u)\log(1+u)-u$, with the continuous convention at $u=-1$.  We first establish the global scalar estimate
\begin{equation}\label{eq:entropy-scalar}
 \left|\phi(u)-\frac{u^2}{2}\right|
 \leq C\min\{|u|^3,u^2\}.
\end{equation}
For $|u|\leq1/2$, Taylor's theorem and $\phi'''(u)=-(1+u)^{-2}$ give the cubic bound.  On the remaining portion of $[-1,1]$, continuity and the fact that $|u|$ stays away from zero give the same bound.  Globally, convexity gives $\phi(u)\geq0$, and $\log(1+u)\leq u$ gives $\phi(u)\leq u^2$.  This proves the quadratic bound and hence \eqref{eq:entropy-scalar}.

If $u=w+r\geq-1$, then
\begin{equation}\label{eq:entropy-split}
 \left|\phi(u)-\frac{u^2}{2}\right|\leq C(|w|^3+r^2).
\end{equation}
Indeed, on $|r|>|w|$ use the quadratic bound and $|u|\leq2|r|$; on its complement use the cubic bound and $|u|\leq2|w|$.

\section{The exact coarse/canonical decomposition}
\label{sec:coarse-decomposition}

A fixed spatial partition defines both the covariance decomposition and the conditional particle law. We construct the regression covariance, derive the particle identity with its diagonal correction, and estimate the coarse partition function. The conditional product law is used for the residual estimates in \Cref{sec:canonical-residual}. The grid is fixed throughout this section and chosen in \Cref{subsec:order-one-assembly} before $N$ varies.

\subsection{A fixed spatial partition and complete label conditioning}
\label{subsec:cell-labels}

We work with real-valued functions and extend $T$ to $L^2(\muunif)$ by setting $T\one=0$, consistently with its convolution kernel $\g$. Fix an integer $J\geq2$, set $M=J^{\d}$, and partition $\T^{\d}$ into the half-open cubes
\begin{equation}
 A_a\coloneqq J^{-1}\bigl(a+[0,1)^{\d}\bigr)
 \qquad(a\in(\Z/J\Z)^{\d}).
 \label{eq:grid-cells}
\end{equation}
Their boundaries have Haar measure zero.  Let
\begin{equation}
 \nu_a\coloneqq M\one_{A_a}\muunif
 \label{eq:cell-uniform-law}
\end{equation}
be the normalized restriction of $\muunif$ to $A_a$, and let $Q$ be conditional expectation onto the sigma algebra generated by the cells:
\begin{equation}
 Qf\coloneqq\sum_a\left(\int f\,\dd\nu_a\right)\one_{A_a}.
 \label{eq:cell-conditional-expectation}
\end{equation}
Then $Q$ is the orthogonal projection in $L^2(\muunif)$ onto the cell-constant functions and preserves constants. Set $P=\Id-Q$; its range consists of the functions with zero average on each cell.

For independent $X_1,\ldots,X_N\sim\muunif$, let $L_i=a$ when $X_i\in A_a$.  The complete label vector $L=(L_1,\ldots,L_N)$ satisfies
\begin{equation}
 \Pb(L=\ell)=M^{-N},
 \qquad
 \mathcal L(X_1,\ldots,X_N\mid L=\ell)
 =\bigotimes_{i=1}^N\nu_{\ell_i}.
 \label{eq:complete-label-disintegration}
\end{equation}
This follows by restricting the product measure to the product of the assigned cells. In particular, conditioning on the complete label vector preserves independence. The occupancy numbers
\begin{equation}
 n_a(\ell)\coloneqq\#\{i:\ell_i=a\}
 \label{eq:occupancy-numbers}
\end{equation}
will enter through the local interaction coefficients and the large-deviation estimate in \Cref{sec:canonical-residual}. Conditioning only on these counts averages over the compatible label assignments and generally destroys the product structure in the labeled coordinates.

\subsection{Regression covariance and conditional covariance}
\label{subsec:regression-covariance}

Let
\begin{equation}
 \mathcal H\coloneqq L_0^2(\muunif),
 \qquad
 V\coloneqq Q\mathcal H,
 \qquad
 \mathcal W\coloneqq P\mathcal H.
 \label{eq:coarse-fine-spaces}
\end{equation}
The operator
\begin{equation}
 B\coloneqq QTQ\big|_V
 \label{eq:coarse-block}
\end{equation}
is strictly positive on the finite-dimensional space $V$.  Indeed, if $v\in V\setminus\{0\}$, then
\begin{equation}
 \ip*{v,Bv}_{L^2(\muunif)}=\ip*{v,Tv}_{L^2(\muunif)}>0
 \label{eq:coarse-block-positive}
\end{equation}
because every nonzero Fourier coefficient of $\g$ is positive.

To include the cross-covariance between the cell-constant and zero-cell-mean components, define on $\mathcal H$
\begin{align}
 A&\coloneqq TQB^{-1}QT,
 \label{eq:coarse-covariance-definition}\\
 R&\coloneqq PTP,
 \label{eq:raw-fine-covariance-definition}\\
 D&\coloneqq PTQB^{-1}QTP.
 \label{eq:predicted-fine-covariance-definition}
\end{align}

\begin{prop}[Exact regression covariance identity]
\label{prop:regression-covariance-identity}
The operators in \eqref{eq:coarse-covariance-definition}--\eqref{eq:predicted-fine-covariance-definition} satisfy
\begin{align}
 T&=A+R-D,
 \label{eq:exact-regression-identity}\\
 0&\leq A\leq T,
 \label{eq:coarse-operator-order}\\
 D&\geq0,
 \label{eq:predicted-fine-positive}\\
 QR&=RQ=0.
 \label{eq:raw-fine-canonical-operator}
\end{align}
The operators $A$ and $D$ have bounded finite-rank kernels.
\end{prop}

\begin{proof}
Relative to $\mathcal H=V\oplus\mathcal W$, write
\begin{equation}
 T=\begin{pmatrix}B&C^*\\ C&F\end{pmatrix},
 \qquad
 C\coloneqq PTQ\big|_V,
 \qquad
 F\coloneqq PTP\big|_{\mathcal W}.
 \label{eq:block-matrix-T}
\end{equation}
The definitions give
\begin{equation}
 A=\begin{pmatrix}B&C^*\\ C&CB^{-1}C^*\end{pmatrix},
 \qquad
 R=\begin{pmatrix}0&0\\0&F\end{pmatrix},
 \qquad
 D=\begin{pmatrix}0&0\\0&CB^{-1}C^*\end{pmatrix}.
 \label{eq:block-matrices-ARD}
\end{equation}
This proves \eqref{eq:exact-regression-identity} and \eqref{eq:raw-fine-canonical-operator}. The factorizations in \eqref{eq:coarse-covariance-definition} and \eqref{eq:predicted-fine-covariance-definition}, together with $B^{-1}>0$, give $A\geq0$ and $D\geq0$. The Schur complement of $B$ is nonnegative because $T\geq0$, so
\begin{equation}
 T-A
 =\begin{pmatrix}0&0\\0&F-CB^{-1}C^*\end{pmatrix}\geq0,
 \label{eq:schur-complement-positive}
\end{equation}
which proves \eqref{eq:coarse-operator-order}.  Equivalently, for $f\in\mathcal H$ and $v=B^{-1}QTf\in V$,
\begin{equation}
 \ip*{f,(T-A)f}_{L^2(\muunif)}
 =\ip*{f-v,T(f-v)}_{L^2(\muunif)}\geq0.
 \label{eq:schur-square-completion}
\end{equation}

The range of $TQ$ is generated by the finitely many functions $T\one_{A_a}$.  Since $\g\in L^1(\muunif)$,
\begin{equation}
 \norm{T\one_{A_a}}_{L^\infty(\muunif)}
 \leq\norm{\g}_{L^1(\muunif)}.
 \label{eq:cell-potential-bounded}
\end{equation}
Translation by $h$ changes $T\one_{A_a}$ in the uniform norm by at most $\norm{\g(\cdot+h)-\g}_{L^1(\muunif)}$, which tends to zero as $h\to0$. Thus, the range of $A$ is contained in a finite-dimensional space of bounded continuous functions. The range of $D$ is contained in the span of $PT\one_{A_a}=T\one_{A_a}-QT\one_{A_a}$. Since $Q$ is an $L^\infty$ contraction, these functions are bounded by $2\norm{\g}_{L^1(\muunif)}$; they need not be continuous across cell boundaries. The finite-dimensional spectral theorem now gives a bounded continuous kernel for $A$ and a bounded kernel for $D$.
\end{proof}

The covariance interpretation follows from Gaussian regression. Let $\mathsf X(f)$ be a centered Gaussian family indexed by $\mathcal H$ with covariance
\begin{equation}
 \E[\mathsf X(f)\mathsf X(h)]
 =\ip*{f,Th}_{L^2(\muunif)}.
 \label{eq:generalized-gaussian-field}
\end{equation}
Such a family is obtained by applying an isonormal Gaussian process to $T^{1/2}f$; it does not require an $L^2(\muunif)$-valued Gaussian random variable. For $f\in\mathcal H$ and $v\in V$,
\begin{equation}
 \E\left[(\mathsf X(f)-\mathsf X(B^{-1}QTf))\mathsf X(v)\right]
 =\ip*{QTf-BB^{-1}QTf,v}_{L^2(\muunif)}=0.
 \label{eq:coarse-predictor-orthogonality}
\end{equation}
The residual is therefore independent of the Gaussian variables indexed by a basis of $V$. Consequently, $\mathsf X(B^{-1}QTf)$ is the conditional expectation of $\mathsf X(f)$ given the cell-average variables. Its covariance operator is $A$, and the conditional covariance is
\begin{equation}
 T-A=R-D.
 \label{eq:conditional-covariance}
\end{equation}
Similarly, $R$ is the covariance of the fine component $\mathsf X(Pf)$, and $D$ is the covariance of its predictor. This Gaussian conditioning is distinct from conditioning the particle positions on their cell labels in \eqref{eq:complete-label-disintegration}.

The kernels of both $R$ and $R-D$ have zero cell averages in each variable, since $QR=RQ=QD=DQ=0$. We estimate $R$ through the explicit double-centering formula below. The subtraction of $D$ then has the favorable particle bound \eqref{eq:diagonal-correction-bound}.

\subsection{Kernels and the particle identity}
\label{subsec:kernels-particle-identity}

Extend $A,R,D$ to vanish on constants, and let $a,r,d$ be their symmetric kernels. For $a$ and $d$, choose the bounded finite-rank representatives furnished by the spectral theorem in the proof of \Cref{prop:regression-covariance-identity}. In particular, their diagonal values below refer to these representatives. For $x\in A_a$ and $y\in A_b$, the kernel of $R=PTP$ is the double centering
\begin{align}
 r(x,y)
 \coloneqq{}&\g(x-y)
 -\int_{A_a}\g(u-y)\,\dd\nu_a(u)
 -\int_{A_b}\g(x-v)\,\dd\nu_b(v)
 \notag\\
 &+\iint_{A_a\times A_b}\g(u-v)\,\dd\nu_a(u)\dd\nu_b(v).
 \label{eq:blockwise-centered-kernel}
\end{align}
Each cell average is finite because $\nu_a\leq M\muunif$ and $\g\in L^1(\muunif)$. Integrating the four terms in either variable gives
\begin{equation}
 \int_{A_a}r(x,y)\,\dd\nu_a(x)=0,
 \qquad
 \int_{A_b}r(x,y)\,\dd\nu_b(y)=0.
 \label{eq:blockwise-canonicality}
\end{equation}
The operator identity \eqref{eq:exact-regression-identity} is an equality of Hilbert--Schmidt kernels, and hence
\begin{equation}
 \g(x-y)=a(x,y)+r(x,y)-d(x,y)
 \quad\text{for }\muunif^{\otimes2}\text{-almost every }(x,y).
 \label{eq:kernel-identity}
\end{equation}
For a symmetric kernel $q$, set
\begin{equation}
 \mathsf W_N(q)
 \coloneqq\frac1N\sum_{1\leq i\neq j\leq N}q(x_i,x_j)
 =\frac2N\sum_{1\leq i<j\leq N}q(x_i,x_j).
 \label{eq:particle-kernel-sum}
\end{equation}
For each fixed $N$, the exceptional set in \eqref{eq:kernel-identity} has zero probability at every pair of independent coordinates. Taking the finite union over pairs gives
\begin{equation}
 \mathsf W_N(\g)=\mathsf W_N(a)+\mathsf W_N(r)-\mathsf W_N(d)
 \qquad \muunif^{\otimes N}\text{-almost surely}.
 \label{eq:particle-regression-identity}
\end{equation}

Since \(D\geq0\) has finite rank, there is a bounded map \(\mathbf d:\T^{\d}\to\R^q\) such that \(d(x,y)=\mathbf d(x)\cdot\mathbf d(y)\).  Therefore,
\begin{align}
 -\frac\be2\mathsf W_N(d)
 &=-\frac\be{2N}\left|\sum_{i=1}^N\mathbf d(x_i)\right|^2
 +\frac\be{2N}\sum_{i=1}^N|\mathbf d(x_i)|^2
 \notag\\
 &\leq\frac\be2\sup_{x\in\T^{\d}}d(x,x).
 \label{eq:diagonal-correction-bound}
\end{align}
Thus, for a fixed grid and a bounded inverse-temperature interval, the subtraction of $D$ costs a constant independent of $N$. This calculation uses the finite-rank diagonal of $D$ only; the singular diagonal of $\g$ never enters.

\subsection{Nonperturbative control of the coarse covariance}
\label{subsec:coarse-control}

The domination $A\leq T$ transfers the macroscopic entropy--energy inequality to the coarse feature moments. The resulting sub-Gaussian estimate controls the coarse partition function under the unconditioned product measure, as required by the H\"older estimate in \Cref{subsec:order-one-assembly}.

\begin{prop}[Coarse partition estimate]
\label{prop:coarse-partition-estimate}
Fix \(b<\beta_{\gm}\) and construct \(A\) from any fixed grid as above.  If \(0\leq\be<b\), then
\begin{equation}
 \sup_{N\geq1}
 \E_{\muunif^{\otimes N}}\left[
 \exp\left\{\frac\be2\mathsf W_N(a)\right\}\right]
 \leq\left(1-\frac\be b\right)^{-\rank(A)/2}.
 \label{eq:coarse-partition-estimate}
\end{equation}
\end{prop}

\begin{proof}
Set $q=\rank(A)$. The spectral representation of $A$ gives a bounded continuous mean-zero feature map $\mathbf a:\T^{\d}\to\R^q$ such that
\begin{equation}
 a(x,y)=\mathbf a(x)\cdot\mathbf a(y).
 \label{eq:coarse-feature-map}
\end{equation}
Let $\mu\ll\muunif$ be a probability measure with bounded density. Since $\mathbf a$ has zero $\muunif$-mean, the squared norm of its $\mu$-average equals $\ip*{\mu-1,A(\mu-1)}_{L^2(\muunif)}$. Hence, \eqref{eq:coarse-operator-order} and \eqref{eq:strict-subcritical-macroscopic-inequality} imply
\begin{equation}
 \Hr(\mu\mid\muunif)
 \geq-bI_{\mu}
 \geq\frac b2\left|\int_{\T^{\d}}\mathbf a(x)\,\dd\mu(x)\right|^2.
 \label{eq:coarse-feature-entropy}
\end{equation}
For a general $\mu$ with finite relative entropy, set $\mu_t=e^{t\Delta}\mu$, using the ordinary torus heat semigroup. Jensen's inequality gives $\Hr(\mu_t\mid\muunif)\leq\Hr(\mu\mid\muunif)$. Since $\mu_t\to\mu$ in $L^1(\muunif)$ and $\mathbf a$ is bounded, its moments converge. Applying \eqref{eq:coarse-feature-entropy} to $\mu_t$ and letting $t\downarrow0$ proves the same inequality for $\mu$.

For $z\in\R^q$, the Donsker--Varadhan lemma in variational form gives
\begin{align}
 \log\int_{\T^{\d}}e^{z\cdot\mathbf a(x)}\,\dd\muunif(x)
 &=\sup_{\Hr(\mu\mid\muunif)<\infty}
 \left\{z\cdot\int_{\T^{\d}}\mathbf a\,\dd\mu-\Hr(\mu\mid\muunif)\right\}
 \notag\\
 &\leq\sup_{w\in\R^q}\left\{z\cdot w-\frac b2|w|^2\right\}
 =\frac{|z|^2}{2b}.
 \label{eq:coarse-feature-subgaussian}
\end{align}
Here, the first supremum is over probability measures; equality follows by taking the normalized density proportional to $e^{z\cdot\mathbf a}$, which is bounded and has finite entropy.

Let \(Z\) be a standard Gaussian vector in \(\R^q\).  By \eqref{eq:coarse-feature-map},
\begin{equation}
 \frac\be2\mathsf W_N(a)
 =\frac\be{2N}\left|\sum_{i=1}^N\mathbf a(x_i)\right|^2
 -\frac\be{2N}\sum_{i=1}^N|\mathbf a(x_i)|^2
 \leq\frac\be{2N}\left|\sum_{i=1}^N\mathbf a(x_i)\right|^2.
 \label{eq:coarse-square-completion}
\end{equation}
The Hubbard--Stratonovich transformation (cf.\ \cite[(6.18)--(6.19)]{DelgadinoGvalaniRosenzweig2026}), Tonelli's theorem, and \eqref{eq:coarse-feature-subgaussian} yield
\begin{align}
 \E_{\muunif^{\otimes N}}\left[e^{\be\mathsf W_N(a)/2}\right]
 &\leq\E_Z\left[
 \left(\int e^{\sqrt{\be/N}\,Z\cdot\mathbf a(x)}\,\dd\muunif(x)\right)^N\right]
 \notag\\
 &\leq\E_Z\left[e^{\be|Z|^2/(2b)}\right]
 =\left(1-\frac\be b\right)^{-q/2}.
 \label{eq:coarse-hubbard-stratonovich}
\end{align}
Since $q=\rank(A)$, this is \eqref{eq:coarse-partition-estimate}.
\end{proof}

\section{Control of the canonical residual below the global-minimization threshold}
\label{sec:canonical-residual}

The explicit residual from \Cref{sec:coarse-decomposition} is estimated separately on near and far cell pairs. Logarithmic moment bounds control the near pairs, and mixed second differences control the far pairs. We then condition on the cell labels and choose the occupancy thresholds so that the canonical-array estimate applies. The complementary event is controlled by the linear bound on the log partition function.

\subsection{Near and far cell pairs}
\label{subsec:near-far-cell-pairs}

For $a,b\in(\Z/J\Z)^{\d}$, define the periodic cell-index distance
\begin{equation}
 \delta_J(a,b)\coloneqq\min_{n\in J\Z^{\d}}|a-b+n|_{\infty}.
 \label{eq:cell-index-distance}
\end{equation}
Fix an integer $R\geq2$ and assume $J\geq8(R+1)$.  We call $(a,b)$ near if $\delta_J(a,b)\leq R$ and far otherwise.  Accordingly, define
\begin{equation}
 r_{\mathrm{near}}(x,y)
 \coloneqq r(x,y)\one_{\{\delta_J(a,b)\leq R\}},
 \qquad
 r_{\mathrm{far}}\coloneqq r-r_{\mathrm{near}},
 \quad x\in A_a,\ y\in A_b.
 \label{eq:near-far-kernels}
\end{equation}
Both kernels retain the blockwise canonicality \eqref{eq:blockwise-canonicality}.

For near cells, rescaling introduces an additive logarithmic constant. Double centering cancels this constant, leaving moment bounds independent of the mesh size.

\begin{lemma}[Uniform near-cell logarithmic moments]
\label{lem:near-cell-logarithmic-moments}
For every fixed $R\geq2$, there is $L_R<\infty$, depending only on $\d$ and $R$, such that, for every $J\geq8(R+1)$, every near pair $(a,b)$, and every integer $p\geq1$,
\begin{align}
 \sup_{y\in A_b}\int_{A_a}|r(x,y)|^p\,\dd\nu_a(x)
 &\leq L_R^p p!,
 \label{eq:near-moment-first}\\
 \sup_{x\in A_a}\int_{A_b}|r(x,y)|^p\,\dd\nu_b(y)
 &\leq L_R^p p!.
 \label{eq:near-moment-second}
\end{align}
\end{lemma}

\begin{proof}
Choose compatible lifts of $A_a$ and $A_b$ to cubes of side length $J^{-1}$ in $\R^{\d}$.  Since the cell-index distance is at most $R$, every difference $x-y$ in these lifts has size at most $C_RJ^{-1}$.  By \eqref{eq:kernel-local-form},
\begin{equation}
 \g(x-y)-\log J
 =-\log|J(x-y)|+\psi(x-y),
 \label{eq:near-logarithmic-rescaling}
\end{equation}
where $\psi$ is bounded on the relevant neighborhood and the bound depends only on $\d$.  For every bounded set $B\subset\R^{\d}$, there is $C_B<\infty$ such that
\begin{equation}
 \sup_{z\in B}\int_B|\log|u-z||^p\,\dd u
 \leq C_B^p p!
 \qquad(p\geq1).
 \label{eq:logarithmic-factorial-moment}
\end{equation}
Indeed, the singular part follows from
\begin{equation}
 \int_0^1(-\log s)^p s^{\d-1}\,\dd s
 =\frac{p!}{\d^{p+1}},
 \label{eq:radial-logarithmic-moment}
\end{equation}
and the complement is bounded.  Rescaling each cell to a unit cube in \eqref{eq:near-logarithmic-rescaling} therefore gives
\begin{equation}
 \sup_{y\in A_b}\int_{A_a}
 |\g(x-y)-\log J|^p\,\dd\nu_a(x)
 \leq C_R^p p!,
 \label{eq:uncentered-near-moment}
\end{equation}
with the analogous bound after interchanging the variables.

The four terms in the double centering \eqref{eq:blockwise-centered-kernel} contain the same additive constant $\log J$, which cancels exactly.  Jensen's inequality applied to each cell average in \eqref{eq:blockwise-centered-kernel}, followed by \eqref{eq:uncentered-near-moment}, proves \eqref{eq:near-moment-first}--\eqref{eq:near-moment-second} after increasing $L_R$.
\end{proof}

For far pairs, double centering converts the interaction into a mixed second difference.

\begin{lemma}[Far-cell mixed-difference estimate]
\label{lem:far-cell-mixed-difference}
There is $C=C(\d)<\infty$ such that, for every $R\geq2$, $J\geq8(R+1)$, and every far pair $(a,b)$,
\begin{equation}
 \sup_{x\in A_a,\,y\in A_b}|r(x,y)|
 \leq\frac{C}{(R-1)^2}.
 \label{eq:far-cell-bound}
\end{equation}
Consequently, $r_{\mathrm{far}}$ is a globally canonical kernel and satisfies the hypotheses of \Cref{lem:canonical-array-exponential} with
\begin{equation}
 L_{\mathrm{far}}=\frac{C}{(R-1)^2}.
 \label{eq:far-moment-constant}
\end{equation}
\end{lemma}

\begin{proof}
The local form \eqref{eq:kernel-local-form} and smoothness away from the origin imply
\begin{equation}
 |\nabla^2\g(z)|
 \leq C\operatorname{dist}_{\T^{\d}}(z,0)^{-2}
 \qquad(z\neq0).
 \label{eq:logarithmic-hessian-bound}
\end{equation}
Using \eqref{eq:blockwise-centered-kernel}, write $r(x,y)$ as the average over $u\in A_a$ and $v\in A_b$ of
\begin{equation}
 \g(x-y)-\g(u-y)-\g(x-v)+\g(u-v).
 \label{eq:mixed-second-difference}
\end{equation}
Every segment appearing in the two applications of the fundamental theorem of calculus stays at distance at least $c(R-1)J^{-1}$ from the origin, while $|x-u|+|y-v|\leq CJ^{-1}$, for constants $c,C>0$ depending only on $\d$.  Equations \eqref{eq:logarithmic-hessian-bound} and \eqref{eq:mixed-second-difference} therefore give \eqref{eq:far-cell-bound}.  Blockwise canonicality and summation over the far cells imply global canonicality.  Since a bounded kernel satisfies \eqref{eq:canonical-array-kernel-moments} with its supremum norm as $L$, \eqref{eq:far-moment-constant} follows.
\end{proof}

\subsection{Good cell counts and the local coupling}
\label{subsec:good-occupancies}

Let $\mathcal E_R$ be the finite graph whose vertices are the cells and whose edges are the unordered near pairs, including loops.  Its maximal degree is bounded by a constant depending only on $\d$ and $R$.  A greedy coloring of its line graph gives a number
\begin{equation}
 \chi_R=\chi_R(\d,R)<\infty
 \label{eq:near-edge-color-number}
\end{equation}
of colors, independent of $J$, such that edges of the same color have disjoint vertex sets.

For a label vector $\ell$, let $S_e$ be the contribution of the edge $e=\{a,b\}$ to $\frac N2\mathsf W_N(r_{\mathrm{near}})$, and let $m_e$ be the number of labels belonging to the cells incident to $e$.  For a loop, $m_e=n_a$; otherwise, $m_e=n_a+n_b$.  Conditional on $L=\ell$, $S_e$ is a canonical quadratic statistic of the independent cell-uniform variables indexed by those labels.  Pairs that do not have the cell pattern specified by $e$ are represented by identically zero kernels, as allowed in \Cref{lem:canonical-array-exponential}.

For $\delta\in(0,1)$, define the good cell-count event
\begin{equation}
 \mathcal G_{N,\delta}
 \coloneqq\left\{\max_a n_a\leq\delta N\right\}.
 \label{eq:good-occupancy-event}
\end{equation}

\begin{prop}[Near residual on good cell counts]
\label{prop:near-residual-good}
For every $t_0<\infty$ and every fixed $R\geq2$, there is $\delta=\delta(\d,R,t_0)>0$ such that, for every $J\geq8(R+1)$,
\begin{equation}
 \sup_{N\geq1}\sup_{0\leq t\leq t_0}
 \E_{\muunif^{\otimes N}}\left[
 \one_{\mathcal G_{N,\delta}}
 \exp\left\{\frac t2\mathsf W_N(r_{\mathrm{near}})\right\}\right]
 <\infty.
 \label{eq:near-residual-good-bound}
\end{equation}
\end{prop}

\begin{proof}
Condition on the complete label vector.  H\"older's inequality over the $\chi_R$ edge colors gives
\begin{equation}
 \E\left[e^{t\sum_{e\in\mathcal E_R}S_e/N}\mid L\right]
 \leq\prod_{c=1}^{\chi_R}
 \left(\E\left[e^{t\chi_R\sum_{e\in c}S_e/N}\mid L\right]\right)^{1/\chi_R}.
 \label{eq:edge-color-holder}
\end{equation}
Within one color, the incident cell sets are disjoint, so the corresponding factors are conditionally independent.  For each edge with $m_e\geq2$,
\begin{equation}
 \frac{t\chi_R}{N}S_e
 =\frac{\theta_e}{m_e}S_e,
 \qquad
 \theta_e\coloneqq t\chi_R\frac{m_e}{N}.
 \label{eq:local-temperature-identity}
\end{equation}
On $\mathcal G_{N,\delta}$, one has $m_e\leq2\delta N$.  Choose $\delta$ so that
\begin{equation}
 2t_0\chi_R C_{\mathrm u}L_R\delta\leq\frac12.
 \label{eq:delta-choice-near}
\end{equation}
Edges with $m_e<2$ contribute zero. For all other edges, \Cref{lem:near-cell-logarithmic-moments} and \Cref{lem:canonical-array-exponential} bound every edge factor in \eqref{eq:edge-color-holder} by a constant independent of $N$ and the label vector.  The number of edges is finite once $J$ is fixed, which proves \eqref{eq:near-residual-good-bound} after averaging over the labels.
\end{proof}

\begin{prop}[Far residual]
\label{prop:far-residual}
For every $t_0<\infty$, there is $R=R(\d,t_0)\geq2$ such that, for every $J\geq8(R+1)$,
\begin{equation}
 \sup_{N\geq1}\sup_{0\leq t\leq t_0}
 \E_{\muunif^{\otimes N}}\left[
 \exp\left\{\frac t2\mathsf W_N(r_{\mathrm{far}})\right\}\right]
 <\infty.
 \label{eq:far-residual-bound}
\end{equation}
\end{prop}

\begin{proof}
By \Cref{lem:far-cell-mixed-difference}, the global far kernel is canonical and satisfies \eqref{eq:canonical-array-kernel-moments} with $L=C(R-1)^{-2}$.  Apply \Cref{lem:canonical-array-exponential} with $m=N$ and $\theta=t$.  Choose $R$ so that
\begin{equation}
 t_0C_{\mathrm u}C(R-1)^{-2}\leq\frac12.
 \label{eq:R-choice-far}
\end{equation}
\end{proof}

\begin{cor}[Full residual on good cell counts]
\label{cor:full-residual-good}
For every $t_0<\infty$, one can choose $R$, then $\delta$, and then $J\geq8(R+1)$ so that
\begin{equation}
 \sup_{N\geq1}\sup_{0\leq t\leq t_0}
 \E_{\muunif^{\otimes N}}\left[
 \one_{\mathcal G_{N,\delta}}
 \exp\left\{\frac t2\mathsf W_N(r)\right\}\right]
 <\infty.
 \label{eq:full-residual-good-bound}
\end{equation}
\end{cor}

\begin{proof}
Use $r=r_{\mathrm{near}}+r_{\mathrm{far}}$ and Cauchy--Schwarz:
\begin{align}
 &\E\left[\one_{\mathcal G_{N,\delta}}
 e^{t(\mathsf W_N(r_{\mathrm{near}})+\mathsf W_N(r_{\mathrm{far}}))/2}\right]
 \notag\\
 &\qquad\leq
 \left(\E\left[\one_{\mathcal G_{N,\delta}}
 e^{t\mathsf W_N(r_{\mathrm{near}})}\right]\right)^{1/2}
 \left(\E\left[e^{t\mathsf W_N(r_{\mathrm{far}})}\right]\right)^{1/2}.
 \label{eq:near-far-cauchy-schwarz}
\end{align}
Apply \Cref{prop:near-residual-good} and \Cref{prop:far-residual} with $2t_0$.
\end{proof}

\subsection{Overoccupied cells}
\label{subsec:bad-occupancies}

Under $\muunif^{\otimes N}$, the labels are independent and uniform on the $M$ cells.  If $M^{-1}<\delta$, the binomial Chernoff bound and a union bound give
\begin{equation}
 \Pb(\mathcal G_{N,\delta}^{\mathsf c})
 \leq M\exp\{-N\mathsf D(\delta\Vert M^{-1})\},
 \label{eq:bad-occupancy-probability}
\end{equation}
where
\begin{equation}
 \mathsf D(\delta\Vert q)
 \coloneqq\delta\log\frac\delta q+(1-\delta)\log\frac{1-\delta}{1-q}.
 \label{eq:binary-relative-entropy}
\end{equation}
For fixed $\delta>0$,
\begin{equation}
 \mathsf D(\delta\Vert M^{-1})\longrightarrow\infty
 \qquad\text{as }M\longrightarrow\infty.
 \label{eq:occupancy-rate-divergence}
\end{equation}
Thus, for a fixed occupancy fraction, refining the grid makes the exponential probability bound stronger than any prescribed linear-in-$N$ bound on the log partition function.

\subsection{Assembly of the bound below the global-minimization threshold}
\label{subsec:order-one-assembly}

\begin{prop}[Partition bound below the global-minimization threshold]
\label{prop:compact-subcritical-partition-bound}
For every $\be_0<\beta_{\gm}$,
\begin{equation}
 \sup_{N\geq2}\sup_{0\leq\be\leq\be_0}
 Z_{N,\be}<\infty.
 \label{eq:compact-subcritical-partition-bound}
\end{equation}
\end{prop}

\begin{proof}
Choose $p>1$ so that
\begin{equation}
 p\be_0<\beta_{\gm},
 \qquad
 \frac1p+\frac1q=1,
 \label{eq:good-holder-exponents}
\end{equation}
and then choose $b$ with $p\be_0<b<\beta_{\gm}$.  Apply \Cref{cor:full-residual-good} with $t_0=q\be_0$, choosing $R$, then $\delta$, and leaving the grid size $J$ to be fixed below.

On the good event, \eqref{eq:particle-regression-identity}, \eqref{eq:diagonal-correction-bound}, and H\"older's inequality give
\begin{align}
 &\E\left[\one_{\mathcal G_{N,\delta}}
 \exp\left\{\frac\be2\mathsf W_N(\g)\right\}\right]
 \notag\\
 &\qquad\leq e^{C}
 \left(\E\left[e^{p\be\mathsf W_N(a)/2}\right]\right)^{1/p}
 \left(\E\left[\one_{\mathcal G_{N,\delta}}
 e^{q\be\mathsf W_N(r)/2}\right]\right)^{1/q}.
 \label{eq:good-event-assembly}
\end{align}
For the fixed grid, the first factor on the right side of \eqref{eq:good-event-assembly} is uniformly bounded by \Cref{prop:coarse-partition-estimate}, applied with the coefficient $p\be\leq p\be_0<b$, and the second is bounded by \Cref{cor:full-residual-good}.

For the bad event, choose $s>1$ with $s\be_0<2\d$ and let $s'$ be its conjugate exponent.  By H\"older's inequality, \Cref{cor:linear-partition-bound}, and \eqref{eq:bad-occupancy-probability},
\begin{align}
 &\E\left[\one_{\mathcal G_{N,\delta}^{\mathsf c}}
 \exp\left\{\frac\be2\mathsf W_N(\g)\right\}\right]
 \notag\\
 &\qquad\leq
 \Pb(\mathcal G_{N,\delta}^{\mathsf c})^{1/s'}
 \left(Z_{N,s\be}\right)^{1/s}
 \notag\\
 &\qquad\leq
 M^{1/s'}\exp\left\{-\frac{N}{s'}\mathsf D(\delta\Vert M^{-1})
 +\frac{C^{\mathrm{lin}}_{s\be_0}}{s}N\right\}.
 \label{eq:bad-event-assembly}
\end{align}
After $R$ and $\delta$ have been fixed, choose $J$, equivalently $M=J^{\d}$, sufficiently large that
\begin{equation}
 \frac1{s'}\mathsf D(\delta\Vert M^{-1})
 >\frac{C^{\mathrm{lin}}_{s\be_0}}s.
 \label{eq:grid-choice-bad-event}
\end{equation}
The right side of \eqref{eq:bad-event-assembly} is then bounded uniformly in $N$.  Combining the good and bad contributions proves \eqref{eq:compact-subcritical-partition-bound}.  All parameters are fixed before $N$ varies, in the order
\begin{equation}
 p,q,b,\ R,\ \delta,\ s,s',\ J.
 \label{eq:parameter-choice-order}
\end{equation}
\end{proof}

\section{Partition thresholds, reference bounds, and Gaussian limits}
\label{sec:thresholds-asymptotics}

Product trials complement the order-one bound and identify the partition threshold. We next establish the reference estimates and Gaussian limits. The direct cutoff comparison applies to positive canonical kernels under its own moment hypotheses. The sharper error of order $N^{-1}$ follows from the marked theorem at order zero; that theorem uses only the scalar partition bounds established here.

\subsection{Product trials, positive pressure, and the threshold identity}
\label{subsec:pressure-threshold-identity}

The upper bound in \Cref{prop:compact-subcritical-partition-bound} is complemented by a product-measure lower bound.

\begin{prop}[Product trial]
\label{prop:product-trial}
Let $\mu$ be a probability measure satisfying $\Hr(\mu\mid\muunif)<\infty$.  Then, for every $N\geq2$ and $\be\geq0$,
\begin{equation}
 \frac1N\log Z_{N,\be}
 \geq-\left(1-\frac1N\right)\be I_{\mu}
 -\Hr(\mu\mid\muunif).
 \label{eq:product-trial-lower-bound}
\end{equation}
\end{prop}

\begin{proof}
Apply the Gibbs variational formula first with $\g$ replaced by $\g\wedge L$.  Since $\g$ is bounded from below, monotone convergence as $L\to\infty$ gives the extended identity
\begin{equation}
 \log Z_{N,\be}
 =\sup_{P_N}\left\{
 \frac\be2\E_{P_N}[\mathsf W_N(\g)]
 -\Hr(P_N\mid\muunif^{\otimes N})\right\}.
 \label{eq:gibbs-variational-partition}
\end{equation}
Assume first that $\mu$ has a smooth density.  Take $P_N=\mu^{\otimes N}$.  The mean-zero normalization of $\g$ and the Fourier representation of the energy give
\begin{equation}
 \E_{\mu^{\otimes N}}[\mathsf W_N(\g)]
 =-2(N-1)I_{\mu},
 \qquad
 \Hr(\mu^{\otimes N}\mid \muunif^{\otimes N})
 =N\Hr(\mu\mid\muunif).
 \label{eq:product-trial-energy-entropy}
\end{equation}
Substitution into \eqref{eq:gibbs-variational-partition} proves \eqref{eq:product-trial-lower-bound} for probability measures with smooth density.

For a general $\mu$ satisfying the hypotheses, let $\mu_t$ be its convolution with the torus heat kernel.  Then $\mu_t$ is smooth,
\begin{equation}
 \Hr(\mu_t\mid\muunif)\longrightarrow\Hr(\mu\mid\muunif),
 \qquad
 I_{\mu_t}\downarrow I_{\mu}
 \quad\text{as }t\downarrow0.
 \label{eq:product-trial-heat-approximation}
\end{equation}
The first convergence follows from entropy contraction, $L^1$ convergence, and lower semicontinuity; the second follows termwise from \eqref{eq:signed-measure-energy}.  Applying the smooth-measure estimate to $\mu_t$ and letting $t\downarrow0$ proves \eqref{eq:product-trial-lower-bound}.
\end{proof}

\begin{proof}[Proof of the threshold and integrability conclusions of \Cref{thm:principal}]
The integrability assertion is \Cref{prop:fixed-particle-integrability}, and \eqref{eq:principal-order-one} is \Cref{prop:compact-subcritical-partition-bound}.  Hence $\beta_{\pf}\geq\beta_{\gm}$.

If $\beta_{\gm}<\be<2\d$, the definition of $\beta_{\gm}$ provides a probability measure $\mu$ with finite entropy such that
\begin{equation}
 \Ec_{\be}(\mu)=\Hr(\mu\mid\muunif)+\be I_{\mu}<0.
 \label{eq:strict-free-energy-competitor}
\end{equation}
The finite-entropy estimate \eqref{eq:finite-entropy-energy-bound} shows that $I_{\mu}>-\infty$, so \Cref{prop:product-trial} applies.  Taking the lower limit in \eqref{eq:product-trial-lower-bound} and then the supremum over $\mu$ proves \eqref{eq:principal-positive-pressure}.  In particular, the partition functions cannot remain uniformly bounded at any such $\be$, so $\beta_{\pf}\leq\beta_{\gm}$ whenever $\beta_{\gm}<2\d$.  If $\beta_{\gm}=2\d$, fixed-particle divergence at $\be=2\d$ gives the same upper bound.  This proves \eqref{eq:principal-threshold-identity}.
\end{proof}

\subsection{Quantitative Carleman--Fredholm asymptotics}
\label{subsec:homogeneous-determinant-rate}

For a probability space $(X,\mu)$, let $\xi_1,\ldots,\xi_N$ be independent with common law $\mu$.  Given a real symmetric canonical kernel $q\in L^2(\mu^{\otimes2})$, let $\mathsf W_N(q)$ denote the statistic in \eqref{eq:particle-kernel-sum}, evaluated at $(\xi_1,\ldots,\xi_N)$.
Suppose that $q$ is nonnegative definite and has a centered feature representation
\begin{equation}
 q(x,y)=\ip*{\mathsf k(x),\mathsf k(y)}_{\mathcal H},
 \qquad
 \int_X\mathsf k\,\dd\mu=0,
 \qquad
 D(x)\coloneqq\norm{\mathsf k(x)}_{\mathcal H}^2,
 \label{eq:centered-feature-representation}
\end{equation}
for a separable real Hilbert space $\mathcal H$.  Let $T_q$ be the associated integral operator.

The first proposition treats a fixed feature kernel with an exponential diagonal moment. The second allows the kernel to vary with $N$, provided its diagonal grows at most logarithmically. Their proofs are in \Cref{app:feature-laplace}.

\begin{prop}[Hubbard--Stratonovich estimate under a higher-moment bound]
\label{prop:renyi-feature-space}
Fix $\be_0\geq0$, $r>1$, $b>r\be_0$, $a>0$, $\delta\in(0,1]$, and $H,P<\infty$.  Assume
\begin{align}
 r\be_0\norm{T_q}_{\op}&\leq1-\delta,
 \qquad
 \norm{T_q}_{\HS}\leq H,
 \label{eq:renyi-feature-operator-assumptions}\\
 \sup_{M\geq1}\E\left[e^{b\mathsf W_M(q)/2}\right]&\leq P,
 \label{eq:renyi-feature-particle-assumption}\\
 M_a(q)\coloneqq\int_Xe^{aD(x)}\,\dd\mu(x)&<\infty.
 \label{eq:renyi-feature-diagonal-assumption}
\end{align}
Then there is
\begin{equation}
 C=C(\be_0,r,b-r\be_0,a,\delta,H,P)<\infty
 \label{eq:renyi-feature-constant}
\end{equation}
such that, for every $N\geq1$,
\begin{equation}
 \sup_{0\leq\be\leq\be_0}
 \left|\E\left[e^{\be\mathsf W_N(q)/2}\right]
 -\det\nolimits_2(\Id-\be T_q)^{-1/2}\right|
 \leq C(1+M_a(q))^2N^{-1/2}.
 \label{eq:renyi-feature-rate}
\end{equation}
\end{prop}

\begin{prop}[Logarithmic-diagonal triangular array]
\label{prop:triangular-feature-space}
Fix $\be_0\geq0$, $r>1$, $b>r\be_0$, $\delta\in(0,1]$, and $H,P,L_*<\infty$.  For each $N\geq1$, let $q_N$ satisfy \eqref{eq:centered-feature-representation} and assume
\begin{align}
 r\be_0\norm{T_{q_N}}_{\op}&\leq1-\delta,
 \qquad
 \norm{T_{q_N}}_{\HS}\leq H,
 \label{eq:triangular-feature-operator-assumptions}\\
 \sup_{M\geq1}\E\left[e^{b\mathsf W_M(q_N)/2}\right]&\leq P,
 \label{eq:triangular-feature-particle-assumption}\\
 0\leq D_N(x)&\leq L_*(1+\log N)
 \quad\text{for }\mu\text{-almost every }x.
 \label{eq:triangular-feature-diagonal-assumption}
\end{align}
Then there is
\begin{equation}
 C=C(\be_0,r,b-r\be_0,\delta,H,P,L_*)<\infty
 \label{eq:triangular-feature-constant}
\end{equation}
such that
\begin{equation}
 \sup_{0\leq\be\leq\be_0}
 \left|\E\left[e^{\be\mathsf W_N(q_N)/2}\right]
 -\det\nolimits_2(\Id-\be T_{q_N})^{-1/2}\right|
 \leq C\frac{(1+\log N)^2}{\sqrt N}.
 \label{eq:triangular-feature-rate}
\end{equation}
\end{prop}

We regularize $\g$ by convolution with a nonnegative mollifier, preserving nonnegative definiteness.  Choose a smooth even probability density $\varphi$ supported in a sufficiently small Euclidean ball in the standard lift of the torus.  Let $\varphi_\eta(x)=\eta^{-\d}\varphi(x/\eta)$, put $\nu_\eta=\varphi_\eta*\check\varphi_\eta$, and define
\begin{equation}
 \g_{\eta}\coloneqq\g*\nu_\eta,
 \qquad
 \f_{\eta}\coloneqq\g-\g_{\eta}.
 \label{eq:positive-logarithmic-regularization}
\end{equation}
The properties needed below are proved in \Cref{app:homogeneous-regularization}.

\begin{prop}[Logarithmic regularization preserving nonnegative definiteness]
\label{prop:positive-logarithmic-regularization}
There is $\eta_0=\eta_0(\d)>0$ such that, for every $0<\eta\leq\eta_0$, the kernel $\g_{\eta}$ is smooth, real, even, mean zero, and nonnegative definite.  Its operator $T_{\eta}$ satisfies
\begin{equation}
 0\leq T_{\eta}\leq T.
 \label{eq:regularized-operator-order}
\end{equation}
For every $N\geq1$ and $\be\geq0$,
\begin{equation}
 Z_{N,\be}^{(\eta)}
 \coloneqq\E_{\muunif^{\otimes N}}\left[e^{\be\mathsf W_N(\g_{\eta})/2}\right]
 \leq Z_{N,\be}.
 \label{eq:regularized-partition-domination}
\end{equation}
Its feature diagonal is constant and, for a constant $C=C(\d)<\infty$, satisfies
\begin{equation}
 0\leq D_{\eta}=\g_{\eta}(0)
 \leq C(1+\log(1/\eta)).
 \label{eq:regularized-feature-diagonal}
\end{equation}
Finally, for every finite $u\geq2$ with $2u>\d$, there is $C=C(\d,u)<\infty$ such that
\begin{equation}
 \norm{\f_{\eta}}_{L^u(\muunif)}
 \leq C\eta^{\d/u}.
 \label{eq:regularization-Lu-rate}
\end{equation}
\end{prop}

The general triangular-array estimate gives a direct logarithmic comparison without using the all-order coefficient construction:
\begin{equation}
 \sup_{0\leq\be\leq\be_0}
 \left|Z_{N,\be}-\det\nolimits_2(\Id-\be T)^{-1/2}\right|
 \leq C\frac{(1+\log N)^2}{\sqrt N}.
 \label{eq:homogeneous-direct-rate}
\end{equation}
This argument shows how regularization preserving nonnegative definiteness and the feature-space estimate control the singular limit without computing higher coefficients. The marked expansion gives the sharper logarithmic rate afterward.

\begin{proof}[Proof of \eqref{eq:homogeneous-direct-rate}]
Fix $\be_0<\beta_{\gm}$.  Choose an even integer $u\geq2$ such that
\begin{equation}
 u>\frac\d2,
 \qquad
 r\be_0<\beta_{\gm},
 \qquad
 r\coloneqq\frac{u}{u-1},
 \label{eq:homogeneous-rate-u-choice}
\end{equation}
and choose $b$ with $r\be_0<b<\beta_{\gm}$.  By \Cref{prop:compact-subcritical-partition-bound},
\begin{equation}
 P_b\coloneqq\sup_{M\geq2}\sup_{0\leq\al\leq b}
 Z_{M,\al}<\infty.
 \label{eq:homogeneous-buffered-partition-constant}
\end{equation}

The elementary inequality $|e^x-e^y|\leq|x-y|(e^x+e^y)$, H\"older's inequality with exponents $u$ and $r$, the canonical even-moment bound in \Cref{lem:canonical-even-moment}, and \eqref{eq:regularized-partition-domination} imply
\begin{equation}
 \sup_{0\leq\be\leq\be_0}
 \left|Z_{N,\be}-Z_{N,\be}^{(\eta)}\right|
 \leq C\eta^{\d/u}.
 \label{eq:homogeneous-particle-cutoff-error}
\end{equation}
Since $\muunif$ is a probability measure and $u\geq2$, \eqref{eq:regularization-Lu-rate} also gives
\begin{equation}
 \norm{T-T_{\eta}}_{\HS}
 =\norm{\f_{\eta}}_{L^2(\muunif)}
 \leq C\eta^{\d/u}.
 \label{eq:homogeneous-HS-cutoff-error}
\end{equation}
The common spectral gap $\be_0\norm{T}_{\op}<1$ and \Cref{lem:determinant-continuity} therefore yield
\begin{equation}
 \sup_{0\leq\be\leq\be_0}
 \left|\det\nolimits_2(\Id-\be T)^{-1/2}
 -\det\nolimits_2(\Id-\be T_{\eta})^{-1/2}\right|
 \leq C\eta^{\d/u}.
 \label{eq:homogeneous-determinant-cutoff-error}
\end{equation}

Choose
\begin{equation}
 \eta_N\coloneqq\eta_0N^{-u/(2\d)},
 \label{eq:homogeneous-optimized-cutoff}
\end{equation}
where $\eta_0$ is small enough for \Cref{prop:positive-logarithmic-regularization}.  Then $\eta_N^{\d/u}=\eta_0^{\d/u}N^{-1/2}$ and \eqref{eq:regularized-feature-diagonal} gives
\begin{equation}
 D_{\eta_N}\leq L_*(1+\log N).
 \label{eq:homogeneous-logarithmic-diagonal}
\end{equation}
Moreover, \eqref{eq:regularized-operator-order}, \eqref{eq:regularized-partition-domination}, and \eqref{eq:homogeneous-buffered-partition-constant} verify every hypothesis of \Cref{prop:triangular-feature-space} for $q_N=\g_{\eta_N}$.  Hence
\begin{equation}
 \sup_{0\leq\be\leq\be_0}
 \left|Z_{N,\be}^{(\eta_N)}
 -\det\nolimits_2(\Id-\be T_{\eta_N})^{-1/2}\right|
 \leq C\frac{(1+\log N)^2}{\sqrt N}.
 \label{eq:homogeneous-regularized-rate}
\end{equation}
Combining \eqref{eq:homogeneous-particle-cutoff-error}, \eqref{eq:homogeneous-determinant-cutoff-error}, and \eqref{eq:homogeneous-regularized-rate} proves \eqref{eq:homogeneous-direct-rate}.
\end{proof}

\begin{proof}[Proof of \eqref{eq:principal-homogeneous-rate}]
Choose $\be_0<b<\beta_{\gm}$. The scalar bound \eqref{eq:principal-order-one} gives \eqref{eq:prescribed-scalar-buffer} for the singleton family $\{\muunif\}$, since $I_{\muunif}=0$ and $Z_{1,b}=1$. Apply \Cref{thm:prescribed-marked-expansion} with $M=0$. Its proof uses only that scalar bound and the kernel row estimates. The determinant is uniformly bounded by the strict spectral gap, so \eqref{eq:prescribed-partition-expansion} gives the error $C/N$ for sufficiently large $N$. The finitely many remaining particle numbers are absorbed by increasing $C$, using the same partition bound and determinant gap.
\end{proof}

\subsection{Reference partition bounds and modulated log-HLS}
\label{subsec:transported-references}

Let $\Phi:\T^{\d}\to\T^{\d}$ be bi-Lipschitz and $\mu=\Phi_{\#}\muunif$.  We write $Y_N=(Y_1,\ldots,Y_N)$ for the coordinate variables under $\muunif^{\otimes N}$.  The pullback remainder \eqref{eq:pullback-remainder} is bounded.  Indeed, the logarithmic singularities cancel because a bi-Lipschitz map distorts pair distances by bounded factors, while the smooth remainder in \eqref{eq:kernel-local-form} is bounded.  Define
\begin{equation}
 R_{\Phi}(y)\coloneqq\int_{\T^{\d}}r_{\Phi}(y,z)\,\dd\muunif(z),
 \qquad
 \overline r_{\Phi}\coloneqq\int_{\T^{\d}}R_{\Phi}\,\dd\muunif,
 \label{eq:transport-one-body-data}
\end{equation}
and
\begin{equation}
 \tl r_{\Phi}(y,z)
 \coloneqq r_{\Phi}(y,z)-R_{\Phi}(y)-R_{\Phi}(z)+\overline r_{\Phi}.
 \label{eq:transport-hoeffding-centering}
\end{equation}
The centered pullback kernel is
\begin{equation}
 q_{\Phi}(y,z)
 \coloneqq\tl{\g}_{\mu}(\Phi(y),\Phi(z))
 =\g(y-z)+\tl r_{\Phi}(y,z).
 \label{eq:centered-pullback-kernel}
\end{equation}
It is a real symmetric canonical nonnegative-definite kernel in $L^2(\muunif^{\otimes2})$.  The corresponding operator $T_{q_{\Phi}}$ is unitarily equivalent to $T_{\mu}$.
The one-body terms in the splitting satisfy
\begin{equation}
 h_\mu(\Phi(y))=-R_\Phi(y),
 \qquad
 I_{\mu}=-\frac12\overline r_\Phi.
 \label{eq:transport-potential-identities}
\end{equation}

\begin{prop}[Order-one transport bound]
\label{prop:order-one-transport}
Let $\Phi$ and $\mu$ be as above, and let $p,q>1$ be conjugate. If $\be\geq0$ and $4q\be\norm{\tl r_{\Phi}}_{L^\infty}<1$, then, for every $N\geq2$,
\begin{equation}
 \K_{N,\be}(\mu)\leq Z_{N,p\be}^{1/p}
 \exp\left\{\frac{3\be}{2}\norm{r_{\Phi}}_{L^\infty}
 +\frac{8q\be^2\norm{\tl r_{\Phi}}_{L^\infty}^2}
 {1-16q^2\be^2\norm{\tl r_{\Phi}}_{L^\infty}^2}\right\}.
 \label{eq:explicit-order-one-transport}
\end{equation}
More generally, let $B>0$ and suppose that
\begin{equation}
 \sup_{N\geq2}\sup_{0\leq\be\leq b} Z_{N,\be}<\infty
 \qquad\text{for every }b<B.
 \label{eq:abstract-homogeneous-order-one-range}
\end{equation}
If
\begin{equation}
 p\be_0<B,
 \qquad
 q\be_0\norm{\tl r_{\Phi}}_{L^\infty(\muunif^{\otimes2})}<\frac14,
 \label{eq:transport-holder-conditions}
\end{equation}
then
\begin{equation}
 \sup_{N\geq2}\sup_{0\leq\be\leq\be_0}
 \K_{N,\be}(\mu)<\infty.
 \label{eq:transport-order-one-bound}
\end{equation}
Such conjugate exponents exist whenever
\begin{equation}
 \be_0<\frac{B}{1+4B\norm{\tl r_{\Phi}}_{L^\infty(\muunif^{\otimes2})}}.
 \label{eq:optimized-transport-range}
\end{equation}
In particular, taking $B=\beta_{\gm}$ in \eqref{eq:abstract-homogeneous-order-one-range} gives \eqref{eq:transport-order-one-bound} for every $0\leq\be_0<\beta_{\tr}(\Phi)$.
\end{prop}

\begin{proof}
Pulling the modulated Gibbs weight back by $\Phi^{\otimes N}$ and using \eqref{eq:transport-potential-identities} gives the exact identity
\begin{align}
 \K_{N,\be}(\mu)
 ={}&\E_{\muunif^{\otimes N}}\left[
 \exp\left\{\frac\be2\mathsf W_N(\g)
 +\frac\be2\mathsf W_N(\tl r_{\Phi})
 -\frac\be N\sum_{i=1}^NR_{\Phi}(Y_i)
 +\frac\be2\overline r_{\Phi}\right\}\right].
 \label{eq:transport-pullback-identity}
\end{align}
H\"older's inequality therefore gives
\begin{align}
 \K_{N,\be}(\mu)
 \leq{}&\left(Z_{N,p\be}\right)^{1/p}
 \notag\\
 &\times\left(\E_{\muunif^{\otimes N}}\left[
 \exp\left\{\frac{q\be}{2}\mathsf W_N(\tl r_{\Phi})
 -\frac{q\be}{N}\sum_{i=1}^NR_{\Phi}(Y_i)
 +\frac{q\be}{2}\overline r_{\Phi}\right\}\right]\right)^{1/q}.
 \label{eq:transport-holder-splitting}
\end{align}
The last two terms in the second exponential have absolute value at most
\begin{equation}
 \frac{3q\be}{2}\norm{r_{\Phi}}_{L^\infty}.
 \label{eq:transport-one-body-absolute-bound}
\end{equation}
Apply \Cref{lem:bounded-canonical-kernel} to $\tl r_{\Phi}$ with $t=q\be$.  When $4q\be\norm{\tl r_{\Phi}}_{L^\infty}<1$,
\begin{align}
 &\left(\E_{\muunif^{\otimes N}}\left[
 \exp\left\{\frac{q\be}{2}\mathsf W_N(\tl r_{\Phi})
 -\frac{q\be}{N}\sum_{i=1}^NR_{\Phi}(Y_i)
 +\frac{q\be}{2}\overline r_{\Phi}\right\}\right]\right)^{1/q}
 \notag\\
 &\qquad\leq\exp\left\{
 \frac{3\be}{2}\norm{r_{\Phi}}_{L^\infty}
 +\frac{8q\be^2\norm{\tl r_{\Phi}}_{L^\infty}^2}
 {1-16q^2\be^2\norm{\tl r_{\Phi}}_{L^\infty}^2}
 \right\}.
 \label{eq:transport-remainder-bound}
\end{align}
Together with \eqref{eq:transport-holder-splitting}, this proves \eqref{eq:explicit-order-one-transport}. Under \eqref{eq:abstract-homogeneous-order-one-range} and \eqref{eq:transport-holder-conditions}, the homogeneous factor is uniformly bounded, proving \eqref{eq:transport-order-one-bound}.

For $\be_0=0$ there is nothing to prove. If $\tl r_{\Phi}=0$ and $0<\be_0<B$, any $1<p<B/\be_0$ works. Otherwise, exponents satisfying \eqref{eq:transport-holder-conditions} exist precisely when one can choose $p<B/\be_0$ and
\begin{equation}
 \frac{p}{p-1}
 <\frac{1}{4\be_0\norm{\tl r_{\Phi}}_{L^\infty}}.
 \label{eq:transport-optimization-conjugate}
\end{equation}
Letting $p\uparrow B/\be_0$ shows that this is equivalent to
\begin{equation}
 \frac{B\be_0\norm{\tl r_{\Phi}}_{L^\infty}}
 {B-\be_0}<\frac14,
 \label{eq:transport-optimization-algebra}
\end{equation}
which is \eqref{eq:optimized-transport-range}.  The final specialization follows from \eqref{eq:principal-order-one} and the definition \eqref{eq:transport-threshold}.
\end{proof}

\begin{lemma}[Logarithmic remainder near the identity]
\label{lem:near-identity-logarithmic-remainder}
For a $C^1$ diffeomorphism $\Phi$ in a sufficiently small $C^1$ neighborhood of the identity,
\begin{equation}
 \norm{r_{\Phi}}_{L^\infty((\T^{\d})^2\setminus\diag)}
 \leq C\norm{\Phi-\Id}_{C^1},
 \qquad
 \norm{\tl r_{\Phi}}_{L^\infty}
 \leq4\norm{r_{\Phi}}_{L^\infty}.
 \label{eq:near-identity-remainder-bound}
\end{equation}
where the difference from the identity is measured for a homotopic lift.
\end{lemma}
\begin{proof}
Write the lift as $\Phi(y)=y+\xi(y)$. For nearby points, choose the closest lifts; then
\begin{equation}
 |\xi(y)-\xi(z)|\leq\norm{D\xi}_{L^\infty}|y-z|.
 \label{eq:relative-displacement}
\end{equation}
The ratio $|\Phi(y)-\Phi(z)|/|y-z|$ consequently differs from one by at most $\norm{D\xi}_{L^\infty}$. The logarithmic part of \eqref{eq:kernel-local-form} changes by at most $C\norm{D\xi}_{L^\infty}$, and the smooth local remainder changes by at most $C\norm\xi_{C^1}$. Fix the local coordinate radius. On its complement, a sufficiently small $C^1$ perturbation keeps the difference variables away from zero, so the mean-value theorem applies to the smooth kernel. This proves the first estimate. The second follows from the four terms of Hoeffding centering in \eqref{eq:hoeffding-centering}.
\end{proof}

\begin{lemma}
\label{lem:sobolev-transport}
For $\ell>\d/2+2$, there are constants $\delta_{\d,\ell}>0$ and $C_{\d,\ell}<\infty$ such that every probability density $\mu$ with $\norm{\mu-1}_{H^\ell}\leq\delta_{\d,\ell}$ is $\Phi_\#\muunif$ for a $C^1$ diffeomorphism $\Phi$ satisfying
\begin{equation}
 \norm{\Phi-\Id}_{C^1}\leq C_{\d,\ell}\norm{\mu-1}_{H^\ell}.
 \label{eq:dyn-transport-C1}
\end{equation}
Here the difference is taken for a lift homotopic to the identity. The remainder \eqref{eq:pullback-remainder} and its Hoeffding centering with respect to $\muunif$ satisfy
\begin{equation}
 \norm{r_\Phi}_\infty+\norm{\tl r_\Phi}_\infty
 \leq C_{\d,\ell}\norm{\mu-1}_{H^\ell}.
 \label{eq:dyn-r-bound}
\end{equation}
\end{lemma}
\begin{proof}
We give the volume-form interpolation explicitly, following the construction of Moser \cite{Moser1965}. Set $h=\mu-1$, solve $-\Delta\psi=h$ with zero mean, and for $0\leq s\leq1$ set
\begin{equation}
 \rho_s=1+sh,\qquad v_s=\frac{\nabla\psi}{\rho_s}.
 \label{eq:dyn-moser}
\end{equation}
Choose $\delta_{\d,\ell}$ so that $\norm h_{C^1}\leq1/2$. Fourier inversion for $-\Delta$ and Sobolev embedding give
\begin{equation}
 \sup_{0\leq s\leq1}\norm{v_s}_{C^1}\leq C_{\d,\ell}\norm h_{H^\ell}.
 \label{eq:dyn-moser-v}
\end{equation}
Let $\Phi_s$ be the time-dependent flow of $v_s$ with $\Phi_0=\Id$. The flow exists on $[0,1]$ and is a $C^1$ diffeomorphism, by the spatial Lipschitz bound. Since
\begin{equation}
 \partial_s\rho_s+\diver(\rho_s v_s)=h+\Delta\psi=0,
 \label{eq:dyn-moser-continuity}
\end{equation}
differentiation of $\rho_s(\Phi_s(y))\det D\Phi_s(y)$ shows that it is identically one. Thus, $\Phi_\#\muunif=\mu$ for $\Phi=\Phi_1$. The flow equation and its differentiated equation, followed by Gr\"onwall, prove \eqref{eq:dyn-transport-C1}.

The remainder estimate follows from \Cref{lem:near-identity-logarithmic-remainder} and \eqref{eq:dyn-transport-C1}.
\end{proof}

\begin{prop}[Uniform partition bound on the near-uniform Sobolev neighborhood]
\label{prop:family-uniform-partition}
For the common choice of $\vep_{\d,\ell}$ in \eqref{eq:dyn-smallness} and the inverse temperature \eqref{eq:dyn-beta1}, the constant in \Cref{cor:modulated-log-hls} may be chosen so that every reference satisfying \eqref{eq:dyn-smallness} obeys
\begin{equation}
 \sup_{N\geq2}\K_{N,\be_1}(\mu)\leq C_{\mathrm{pf}},
 \label{eq:dyn-uniform-K}
\end{equation}
The bound holds uniformly over that entire neighborhood, with $C_{\mathrm{pf}}\geq1$ depending only on $\d,\ell,\be_0$.
\end{prop}
\begin{proof}
Use \Cref{lem:sobolev-transport}. Since the quotient in \eqref{eq:dyn-smallness} is at most $1/\beta_{\gm}$, decreasing $\vep_{\d,\ell}$ makes the transport construction applicable for every $\be_0$. Decrease it further so that
\begin{equation}
 \norm{\tl r_\Phi}_\infty
 \leq\frac{\beta_{\gm}-\be_0}{8\beta_{\gm}(\beta_{\gm}+\be_0)}.
 \label{eq:dyn-r-tube}
\end{equation}
The same construction gives a common bound on $\norm{r_\Phi}_\infty$.

The exact pullback identity \eqref{eq:transport-pullback-identity} includes both finite-particle one-body corrections. To apply \eqref{eq:explicit-order-one-transport} uniformly over the reference neighborhood, choose conjugate exponents
\begin{equation}
 p=\frac{3\beta_{\gm}+\be_0}{2(\beta_{\gm}+\be_0)},\qquad
 q=\frac{3\beta_{\gm}+\be_0}{\beta_{\gm}-\be_0}.
 \label{eq:dyn-pq}
\end{equation}
They are greater than one and satisfy
\begin{equation}
 p\be_1=\frac{3\beta_{\gm}+\be_0}4<\beta_{\gm},\qquad
 q\be_1\norm{\tl r_\Phi}_\infty
 \leq\frac{3\beta_{\gm}+\be_0}{16\beta_{\gm}}<\frac14.
 \label{eq:dyn-common-margin}
\end{equation}
For $0\leq b\leq\be_1$, the denominator in \eqref{eq:explicit-order-one-transport} has the common lower bound
\begin{equation}
 1-16q^2b^2\norm{\tl r_{\Phi}}_{L^\infty}^2
 \geq\frac{(\beta_{\gm}-\be_0)(7\beta_{\gm}+\be_0)}{16\beta_{\gm}^2}>0.
 \label{eq:family-common-denominator}
\end{equation}
Apply that bound at $b=\be_1$ and use \eqref{eq:principal-order-one} up to $p\be_1<\beta_{\gm}$. The common bounds on both remainders and \eqref{eq:family-common-denominator} give \eqref{eq:dyn-uniform-K}, with no reference-dependent loss. The required decreases of $\vep_{\d,\ell}$ depend only on $\d,\ell$, since $\beta_{\gm}$ is fixed by the dimension.
\end{proof}

\begin{proof}[Proof of \Cref{cor:modulated-log-hls}]
For $P=\mu^{\otimes N}$, Jensen applied under $f_N$ gives
\begin{equation}
 \int f_N A\leq\Hr(f_N\mid P)+\log\int e^A\dd P.
 \label{eq:dyn-entropy-duality}
\end{equation}
For bounded $A$, this follows by writing $\int e^A\dd P\geq\int_{\{\dd f_N/\dd P>0\}}e^A(\dd P/\dd f_N)\dd f_N$ on the support of $f_N$, applying Jensen to the logarithm, and discarding the nonnegative contribution from its complement. Truncation proves the inequality for $A=-\be_1N\Fr_N$, since its exponential is integrable by \eqref{eq:dyn-uniform-K}; on the finite-energy domain no indeterminate expression occurs. Substitute into \eqref{eq:modulated-free-energy} to get \eqref{eq:dyn-coercivity}. Invert its positive coefficient and use $\be_1-\be\geq\be_1-\be_0$ for \eqref{eq:dyn-entropy-beta}.
\end{proof}

\subsection{Transported determinant asymptotics}
\label{subsec:transported-determinants}

For a fixed bi-Lipschitz reference, the order-one partition bound is sufficient for the quantitative determinant asymptotic.

\begin{prop}[Quantitative determinant asymptotic under a strict partition buffer]
\label{prop:abstract-nonuniform-rate}
Let $\Phi$ be bi-Lipschitz, $\mu=\Phi_{\#}\muunif$, and suppose that \eqref{eq:nonuniform-order-one-input} holds for some $B_\mu>0$.
Then, for every $\be_0<B_{\mu}$, there is $C=C(\d,\Phi,\be_0,B_{\mu})<\infty$ such that
\begin{equation}
 \sup_{0\leq\be\leq\be_0}
 \left|\K_{N,\be}(\mu)
 -e^{\be I_{\mu}}
 \det\nolimits_2(\Id-\be T_{\mu})^{-1/2}\right|
 \leq C/N
 \label{eq:abstract-nonuniform-rate}
\end{equation}
for every $N\geq2$.
\end{prop}

We first remove the empirical one-body correction and regularize the centered kernel by independent convolution in its two variables. This gives a direct quantitative comparison with the determinant. The final step applies the marked expansion to the same reference to obtain the stated error of order $N^{-1}$. Throughout, $\Phi$ is fixed, $\mu=\Phi_\#\muunif$, and $q_\Phi$ is given by \eqref{eq:centered-pullback-kernel}.



\subsubsection{Exact factorization and the strict partition buffer}

By \eqref{eq:transport-potential-identities}, define the centered one-body term and centered quadratic Laplace transform by
\begin{equation}
 \tl R_\Phi(y)\coloneqq R_\Phi(y)-\overline r_\Phi,
 \qquad
 \mathsf L_{N,\be}^{\Phi}
 \coloneqq\E_{\muunif^{\otimes N}}\left[e^{\be\mathsf W_N(q_\Phi)/2}\right].
 \label{eq:transport-centered-onebody-partition}
\end{equation}
The explicit formula \eqref{eq:modulated-energy-general-reference} gives the exact factorization
\begin{equation}
 \K_{N,\be}(\mu)
 =e^{\be I_{\mu}}
 \E_{\muunif^{\otimes N}}\left[
 e^{\be\mathsf W_N(q_\Phi)/2}
 e^{-\be N^{-1}\sum_{i=1}^N\tl R_\Phi(Y_i)}
 \right].
 \label{eq:transport-exact-centered-factorization}
\end{equation}
In particular, for every $b<B_\mu$, \eqref{eq:nonuniform-order-one-input} implies
\begin{equation}
 \sup_{N\geq2}\sup_{0\leq\be\leq b}\mathsf L_{N,\be}^{\Phi}<\infty.
 \label{eq:transport-centered-partition-buffer}
\end{equation}
Indeed, the empirical one-body factor in \eqref{eq:transport-exact-centered-factorization} lies between $e^{-b\norm{\tl R_\Phi}_{\infty}}$ and $e^{b\norm{\tl R_\Phi}_{\infty}}$, while the deterministic factor $e^{\be I_\mu}$ is bounded above and below on $0\leq\be\leq b$.

The partition bound also forces a strict spectral gap.

\begin{lemma}[Partition control implies spectral control]
\label{lem:transport-partition-implies-gap}
Let $q\in L^2(\pi^{\otimes2})$ be a real symmetric canonical kernel on a probability space, with nonnegative integral operator $T_q$. If
\begin{equation}
 \sup_{N\geq2}\E_{\pi^{\otimes N}}[e^{b\mathsf W_N(q)/2}]<\infty
 \quad\text{for some }b>0,
 \label{eq:general-centered-partition-buffer}
\end{equation}
then
\begin{equation}
 b\norm{T_q}_{\op}<1.
 \label{eq:transport-spectral-gap}
\end{equation}
In particular, \eqref{eq:transport-centered-partition-buffer} implies this bound for $q=q_\Phi$.
\end{lemma}
\begin{proof}
The canonical Hilbert--Schmidt $U$-statistic limit \cite[Section~5.5.2]{Serfling1980} gives
\begin{equation}
 \mathsf W_N(q)\Longrightarrow
 \sum_{j\geq1}\lambda_j(G_j^2-1),
 \label{eq:transport-Serfling-quadratic-limit}
\end{equation}
where $(\lambda_j)$ are the nonzero eigenvalues of $T_q$ and the $G_j$ are independent standard real Gaussians. The series converges in $L^2$, since $\sum_j\lambda_j^2<\infty$. Portmanteau's theorem and \eqref{eq:general-centered-partition-buffer} give a finite exponential moment at coefficient $b/2$ for this limit. If $T_q\ne0$, isolate its largest eigenvalue $\lambda_1=\norm{T_q}_{\op}$. The remainder is independent of $G_1$ and finite almost surely. Tonelli's theorem therefore implies finiteness of $\E[e^{b\lambda_1(G_1^2-1)/2}]$, which holds exactly when $b\lambda_1<1$. The zero operator is immediate.
\end{proof}

\subsubsection{Independent mollification in the two torus variables and a fractional approximation rate}

For the mollifier $\varphi_\eta$ used in \eqref{eq:positive-logarithmic-regularization}, define the independent mollification in the two torus variables
\begin{equation}
 q_{\Phi,\eta}(y,z)
 \coloneqq\iint_{(\T^{\d})^2}
 q_\Phi(y-u,z-v)\varphi_\eta(u)\varphi_\eta(v)
 \,\dd\muunif(u)\dd\muunif(v).
 \label{eq:transport-independent-convolution}
\end{equation}

\begin{lemma}[Transported regularization preserving nonnegative definiteness]
\label{lem:transported-positive-cutoff}
There is $\eta_0=\eta_0(\d,\Phi)>0$ such that, for every $0<\eta\leq\eta_0$, the kernel $q_{\Phi,\eta}$ is smooth, real, symmetric, canonical, and nonnegative definite.  Its operator satisfies
\begin{equation}
 \norm{T_{q_{\Phi,\eta}}}_{\op}
 \leq\norm{T_{q_\Phi}}_{\op},
 \qquad
 \norm{T_{q_{\Phi,\eta}}}_{\HS}
 \leq\norm{T_{q_\Phi}}_{\HS}.
 \label{eq:transport-cutoff-operator-contraction}
\end{equation}
For every $N\geq1$ and $\be\geq0$,
\begin{equation}
 \E\left[e^{\be\mathsf W_N(q_{\Phi,\eta})/2}\right]
 \leq \mathsf L_{N,\be}^{\Phi}.
 \label{eq:transport-cutoff-partition-domination}
\end{equation}
For a constant $C=C(\d,\Phi)<\infty$, the feature diagonal satisfies
\begin{equation}
 0\leq D_{\Phi,\eta}(y)
 \coloneqq q_{\Phi,\eta}(y,y)
 \leq C(1+\log(1/\eta))
 \qquad(y\in\T^{\d}).
 \label{eq:transport-cutoff-diagonal}
\end{equation}
Finally, for every finite even integer $u\geq2$ and every
\begin{equation}
 0<\sigma<\min\left\{1,\frac\d u\right\},
 \label{eq:transport-fractional-exponent}
\end{equation}
there is $C=C(\d,\Phi,u,\sigma)<\infty$ such that
\begin{equation}
 \norm{q_\Phi-q_{\Phi,\eta}}_{L^u(\muunif^{\otimes2})}
 \leq C\eta^\sigma.
 \label{eq:transport-cutoff-Lu-rate}
\end{equation}
\end{lemma}

\begin{proof}
Let $C_\eta$ denote convolution by $\varphi_\eta$ on $L^2(\muunif)$.  At the operator level,
\begin{equation}
 T_{q_{\Phi,\eta}}=C_\eta T_{q_\Phi}C_\eta^*.
 \label{eq:transport-cutoff-operator-identity}
\end{equation}
Since $C_\eta$ is a self-adjoint contraction preserving constants, \eqref{eq:transport-cutoff-operator-identity} proves canonicality, nonnegative definiteness, and \eqref{eq:transport-cutoff-operator-contraction}.  Smoothness and symmetry follow directly from \eqref{eq:transport-independent-convolution}.

Let $U_1,\ldots,U_N$ be independent with density $\varphi_\eta$.  Conditional Jensen's inequality gives
\begin{align}
 &\exp\left\{\frac\be{2N}\sum_{i\neq j}
 q_{\Phi,\eta}(y_i,y_j)\right\}
 \notag\\
 &\qquad\leq
 \E_U\left[
 \exp\left\{\frac\be{2N}\sum_{i\neq j}
 q_\Phi(y_i-U_i,y_j-U_j)\right\}\right].
 \label{eq:transport-cutoff-Jensen}
\end{align}
Integration in the $y_i$ and Haar invariance prove \eqref{eq:transport-cutoff-partition-domination}.  For the singular kernel, apply Jensen first after truncation from above and then pass to the limit by monotone convergence; the kernel is bounded below.

The diagonal is
\begin{equation}
 D_{\Phi,\eta}(y)
 =\iint q_\Phi(y-u,y-v)
 \varphi_\eta(u)\varphi_\eta(v)\,\dd\muunif(u)\dd\muunif(v).
 \label{eq:transport-cutoff-diagonal-integral}
\end{equation}
The bounded term $\tl r_\Phi$ contributes $O_\Phi(1)$, while the logarithmic term is the same local convolution appearing in \eqref{eq:regularized-diagonal-asymptotic}.  This proves the upper bound in \eqref{eq:transport-cutoff-diagonal}; the lower bound follows from nonnegative definiteness.

We prove \eqref{eq:transport-cutoff-Lu-rate}.  Let
\begin{equation}
 G_\Phi(y,z)\coloneqq\g(\Phi(y)-\Phi(z)).
 \label{eq:transport-raw-pulled-kernel}
\end{equation}
Hoeffding centering is a bounded operator on $L^u(\muunif^{\otimes2})$, with norm at most four, and commutes with translations in the two variables.  It is therefore enough to prove a translation estimate for $G_\Phi$.

Let $a,b\in\T^{\d}$ be represented by sufficiently small vectors in the standard lift and put
\begin{equation}
 \delta\coloneqq|a|+|b|.
 \label{eq:transport-translation-size}
\end{equation}
Bi-Lipschitz continuity and \eqref{eq:kernel-smooth-local-form} show that, on the union of the two diagonal tubes
\begin{equation}
 \mathcal N_\delta
 \coloneqq\left\{\operatorname{dist}(y,z)\leq C_\Phi\delta\right\}
 \cup\left\{\operatorname{dist}(y-a,z-b)\leq C_\Phi\delta\right\},
 \label{eq:transport-near-diagonal-tubes}
\end{equation}
one has the integral bound
\begin{equation}
 \int_{\mathcal N_\delta}
 \abs{G_\Phi(y-a,z-b)-G_\Phi(y,z)}^u
 \,\dd\muunif(y)\dd\muunif(z)
 \leq C\delta^{\d}(1+|\log\delta|^u).
 \label{eq:transport-near-translation-bound}
\end{equation}
Indeed, a diagonal tube of radius $r$ has volume $O(r^{\d})$, and
\begin{equation}
 \int_0^{C\delta}r^{\d-1}(1+|\log r|^u)\,\dd r
 \leq C\delta^{\d}(1+|\log\delta|^u).
 \label{eq:transport-logarithmic-tube-integral}
\end{equation}

On the complement of \eqref{eq:transport-near-diagonal-tubes}, the segment joining $(y,z)$ to $(y-a,z-b)$ remains a distance comparable to $\operatorname{dist}(y,z)$ from the diagonal.  Rademacher's theorem, the Lipschitz bound on $\Phi$, and the gradient estimate in \eqref{eq:kernel-derivative-bounds} yield
\begin{equation}
 \abs{G_\Phi(y-a,z-b)-G_\Phi(y,z)}
 \leq C_\Phi\frac{\delta}{\operatorname{dist}(y,z)}.
 \label{eq:transport-far-translation-pointwise}
\end{equation}
Consequently,
\begin{align}
 &\int_{(\T^{\d})^2\setminus\mathcal N_\delta}
 \abs{G_\Phi(y-a,z-b)-G_\Phi(y,z)}^u
 \,\dd\muunif(y)\dd\muunif(z)
 \notag\\
 &\qquad\leq C\delta^u
 \int_{C\delta}^1r^{\d-1-u}\,\dd r
 \leq C_\sigma\delta^{\sigma u}
 \label{eq:transport-far-translation-integral}
\end{align}
for every $\sigma$ satisfying \eqref{eq:transport-fractional-exponent}.  The logarithmic factor in \eqref{eq:transport-near-translation-bound} is also absorbed by $\delta^{\sigma u}$ because $\sigma u<\d$.  We have proved
\begin{equation}
 \norm{q_\Phi(\cdot-a,\cdot-b)-q_\Phi}_{L^u(\muunif^{\otimes2})}
 \leq C\delta^\sigma.
 \label{eq:transport-translation-modulus}
\end{equation}
Finally, Minkowski's inequality and \eqref{eq:transport-independent-convolution} give
\begin{align}
 \norm{q_\Phi-q_{\Phi,\eta}}_{L^u}
 &\leq\iint\varphi_\eta(a)\varphi_\eta(b)
 \norm{q_\Phi-q_\Phi(\cdot-a,\cdot-b)}_{L^u}
 \,\dd\muunif(a)\dd\muunif(b)
 \notag\\
 &\leq C\eta^\sigma,
 \label{eq:transport-convolution-approximation}
\end{align}
which is \eqref{eq:transport-cutoff-Lu-rate}.
\end{proof}

\subsubsection{Assembly of the quantitative nonuniform estimate}

\begin{proof}[Proof of \Cref{prop:abstract-nonuniform-rate}]
Fix $\be_0<B_\mu$.  Choose an even integer $u\geq2$ so large that, with
\begin{equation}
 r\coloneqq\frac{u}{u-1},
 \label{eq:transport-rate-holder-conjugate}
\end{equation}
one has $r\be_0<B_\mu$.  Choose
\begin{equation}
 r\be_0<b<B_\mu
 \label{eq:transport-rate-buffer-choice}
\end{equation}
and then choose $\sigma$ as in \eqref{eq:transport-fractional-exponent}.  By \eqref{eq:transport-centered-partition-buffer},
\begin{equation}
 P_\Phi\coloneqq\sup_{M\geq2}\sup_{0\leq\al\leq b}
 \mathsf L_{M,\al}^{\Phi}<\infty.
 \label{eq:transport-centered-buffer-constant}
\end{equation}
By \Cref{lem:transport-partition-implies-gap},
\begin{equation}
 r\be_0\norm{T_{q_\Phi}}_{\op}
 \leq\frac{r\be_0}{b}
 =1-\delta
 \label{eq:transport-rate-common-gap}
\end{equation}
for some $\delta>0$.

Set
\begin{equation}
 \mathsf L_{N,\be}^{\Phi,\eta}
 \coloneqq\E\left[e^{\be\mathsf W_N(q_{\Phi,\eta})/2}\right].
 \label{eq:transport-regularized-centered-partition}
\end{equation}
The elementary exponential-difference inequality, H\"older's inequality with exponents $u,r$, \Cref{lem:canonical-even-moment}, and \eqref{eq:transport-cutoff-partition-domination} give
\begin{equation}
 \sup_{0\leq\be\leq\be_0}
 \abs{\mathsf L_{N,\be}^{\Phi}-\mathsf L_{N,\be}^{\Phi,\eta}}
 \leq C\norm{q_\Phi-q_{\Phi,\eta}}_{L^u(\muunif^{\otimes2})}
 \leq C\eta^\sigma.
 \label{eq:transport-particle-cutoff-error}
\end{equation}
Because $u\geq2$, \eqref{eq:transport-cutoff-Lu-rate} also gives
\begin{equation}
 \norm{T_{q_\Phi}-T_{q_{\Phi,\eta}}}_{\HS}
 \leq C\eta^\sigma.
 \label{eq:transport-HS-cutoff-error}
\end{equation}
The operator contraction \eqref{eq:transport-cutoff-operator-contraction}, the common gap \eqref{eq:transport-rate-common-gap}, and \Cref{lem:determinant-continuity} imply
\begin{equation}
 \sup_{0\leq\be\leq\be_0}
 \left|\det\nolimits_2(\Id-\be T_{q_\Phi})^{-1/2}
 -\det\nolimits_2(\Id-\be T_{q_{\Phi,\eta}})^{-1/2}\right|
 \leq C\eta^\sigma.
 \label{eq:transport-determinant-cutoff-error}
\end{equation}

Choose
\begin{equation}
 \eta_N\coloneqq\eta_0N^{-1/(2\sigma)}.
 \label{eq:transport-optimized-cutoff}
\end{equation}
Then \eqref{eq:transport-cutoff-diagonal} gives
\begin{equation}
 D_{\Phi,\eta_N}(y)\leq L_*(1+\log N),
 \label{eq:transport-logarithmic-diagonal}
\end{equation}
while \eqref{eq:transport-cutoff-operator-contraction}, \eqref{eq:transport-cutoff-partition-domination}, and \eqref{eq:transport-centered-buffer-constant} verify the remaining assumptions of \Cref{prop:triangular-feature-space}.  Therefore,
\begin{equation}
 \sup_{0\leq\be\leq\be_0}
 \left|\mathsf L_{N,\be}^{\Phi,\eta_N}
 -\det\nolimits_2(\Id-\be T_{q_{\Phi,\eta_N}})^{-1/2}\right|
 \leq C\frac{(1+\log N)^2}{\sqrt N}.
 \label{eq:transport-regularized-determinant-rate}
\end{equation}
Combining \eqref{eq:transport-particle-cutoff-error}, \eqref{eq:transport-determinant-cutoff-error}, and \eqref{eq:transport-regularized-determinant-rate} yields
\begin{equation}
 \sup_{0\leq\be\leq\be_0}
 \left|\mathsf L_{N,\be}^{\Phi}
 -\det\nolimits_2(\Id-\be T_{q_\Phi})^{-1/2}\right|
 \leq C\frac{(1+\log N)^2}{\sqrt N}.
 \label{eq:transport-centered-final-rate}
\end{equation}

It remains to restore the one-body factor in \eqref{eq:transport-exact-centered-factorization}.  Since $\tl R_\Phi$ is bounded and centered, the Marcinkiewicz--Zygmund inequality gives
\begin{equation}
 \norm*{\frac1N\sum_{i=1}^N\tl R_\Phi(Y_i)}_{L^u(\muunif^{\otimes N})}
 \leq C_u\frac{\norm{\tl R_\Phi}_{L^u(\muunif)}}{\sqrt N}.
 \label{eq:transport-onebody-sample-mean}
\end{equation}
Using
\begin{equation}
 \abs{e^{-\be x}-1}
 \leq\be e^{\be\norm{\tl R_\Phi}_\infty}\abs x
 \qquad(\abs x\leq\norm{\tl R_\Phi}_\infty),
 \label{eq:transport-onebody-exponential-difference}
\end{equation}
H\"older's inequality with exponents $u,r$, \eqref{eq:transport-centered-buffer-constant}, and \eqref{eq:transport-onebody-sample-mean} imply
\begin{align}
 &\sup_{0\leq\be\leq\be_0}
 \left|\K_{N,\be}(\mu)
 -e^{\be I_{\mu}}\mathsf L_{N,\be}^{\Phi}\right|
 \notag\\
 &\qquad\leq C N^{-1/2}.
 \label{eq:transport-onebody-removal-rate}
\end{align}
Finally, the unitary map
\begin{equation}
 \mathcal U:L^2(\mu)\longrightarrow L^2(\muunif),
 \qquad
 \mathcal U f=f\circ\Phi,
 \label{eq:transport-unitary-map}
\end{equation}
intertwines the centered operators:
\begin{equation}
 \mathcal U T_{\mu}\mathcal U^{-1}=T_{q_\Phi}.
 \label{eq:transport-unitary-equivalence}
\end{equation}
The Carleman--Fredholm determinant is invariant under unitary conjugation. Combining \eqref{eq:transport-centered-final-rate}, \eqref{eq:transport-onebody-removal-rate}, and \eqref{eq:transport-unitary-equivalence} proves the direct comparison with error $C(1+\log N)^2N^{-1/2}$.

For the stronger rate in \eqref{eq:abstract-nonuniform-rate}, use \Cref{thm:prescribed-marked-expansion} at $M=0$ for the fixed reference $\mu$. The bi-Lipschitz pushforward has bounded density, so $h_\mu$ and $\psi_\mu$ are bounded. For $\be_0<b<B_\mu$, hypothesis \eqref{eq:nonuniform-order-one-input} and the exact factor $\mathsf P_{N,b}=e^{-bI_\mu}\K_{N,b}(\mu)$ give the strict partition buffer for $N\geq2$; boundedness of $\psi_\mu$ supplies the $N=1$ term. Thus, \eqref{eq:prescribed-partition-expansion} gives the claimed $C/N$ error for large $N$. The common determinant gap and the assumed partition bound absorb the finitely many smaller $N$. The constant may depend on the fixed transport $\Phi$, as stated in the proposition.
\end{proof}


\begin{proof}[Proof of \Cref{thm:transported}]
Apply \Cref{prop:order-one-transport}, using the homogeneous bound \eqref{eq:principal-order-one}.  This proves \eqref{eq:principal-transport-order-one} for every compact interval below \eqref{eq:transport-threshold}.  Since the same statement holds on every smaller compact interval, \Cref{prop:abstract-nonuniform-rate} applies with $B_{\mu}=\beta_{\tr}(\Phi)$ and proves \eqref{eq:principal-transport-rate}.
\end{proof}

\begin{proof}[Proof of the H\"older-density clause in \Cref{thm:transported}]
The Dacorogna--Moser refinement of Moser's volume-form theorem \cite{Moser1965,DacorognaMoser1990} gives, for $\mu$ satisfying \eqref{eq:near-uniform-hypothesis} with sufficiently small norm, a $C^{1,\gamma}$ diffeomorphism $\Phi$ and a constant $C=C(\d,\gamma)$ such that
\begin{equation}
 \Phi_{\#}\muunif=\mu,
 \qquad
 \norm{\Phi-\Id}_{C^{1,\gamma}}
 \leq C\norm{\mu-1}_{C^\gamma}.
 \label{eq:dacorogna-moser-stability}
\end{equation}
Applying \Cref{lem:near-identity-logarithmic-remainder} gives the required bound on the centered remainder. Choose $\vep$ so that the right side of \eqref{eq:near-identity-remainder-bound} makes $\be_0<\beta_{\tr}(\Phi)$.  \Cref{thm:transported} completes the proof.
\end{proof}

\subsection{Static joint linear statistics}
\label{subsec:static-linear-statistics}

We prove \Cref{cor:static-clt} by jointly approximating the linear statistic and the canonical pair statistic, and then passing to the Gibbs tilt using the strict partition buffer.

\begin{proof}[Proof of \Cref{cor:static-clt}]
We use the pair statistic $\mathsf W_N(\tl{\g}_{\mu})$ defined in \eqref{eq:particle-kernel-sum}.
Since $T_{\mu}$ is nonnegative, self-adjoint, and Hilbert--Schmidt, there are nonnegative eigenvalues $(\lambda_\ell)_{\ell\geq1}$, repeated according to multiplicity, and an orthonormal family $(e_\ell)_{\ell\geq1}$ in $L_0^2(\mu)$ such that
\begin{equation}
 T_{\mu}f
 =\sum_{\ell\geq1}\lambda_\ell
 \ip*{f,e_\ell}_{L^2(\mu)}e_\ell,
 \qquad
 \sum_{\ell\geq1}\lambda_\ell^2<\infty.
 \label{eq:clt-spectral-resolution}
\end{equation}
Let $\mathsf Z$ be an isonormal Gaussian process over $L_0^2(\mu)$, write $G_\ell=\mathsf Z(e_\ell)$, and define
\begin{equation}
 \mathsf W_\infty
 \coloneqq\sum_{\ell\geq1}\lambda_\ell(G_\ell^2-1).
 \label{eq:clt-second-chaos}
\end{equation}
The series in \eqref{eq:clt-second-chaos} converges in $L^2$.

Serfling's Hilbert--Schmidt theorem for canonical degree-two $U$-statistics \cite[Section~5.5.2]{Serfling1980}, together with the ordinary multivariate central limit theorem, gives the joint convergence
\begin{equation}
 \bigl(S_N(\varphi_1),\ldots,S_N(\varphi_k),\mathsf W_N(\tl{\g}_{\mu})\bigr)
 \Longrightarrow
 \bigl(\mathsf Z(\overline\varphi_1),\ldots,
 \mathsf Z(\overline\varphi_k),\mathsf W_\infty\bigr),
 \label{eq:joint-serfling-limit}
\end{equation}
where
\begin{equation}
 \overline\varphi_j
 \coloneqq\varphi_j-\int_{\T^{\d}}\varphi_j\,\dd\mu.
 \label{eq:centered-test-function}
\end{equation}
For a finite spectral sum, the joint convergence follows from the multivariate central limit theorem applied to the test functions and eigenfunctions. The canonical variance identity controls the omitted modes uniformly in $N$:
\begin{equation}
 \sup_{N\geq2}\E_{\mu^{\otimes N}}\left[
 \left|\mathsf W_N(\tl{\g}_{\mu})
 -\frac1N\sum_{i\ne j}\sum_{\ell=1}^{M}
 \lambda_\ell e_\ell(X_i)e_\ell(X_j)\right|^2\right]
 \leq 2\sum_{\ell>M}\lambda_\ell^2
 \xrightarrow[M\to\infty]{}0.
 \label{eq:clt-spectral-tail-uniform}
\end{equation}
The same tail bound holds for the limiting Gaussian chaos. Finite-rank approximation therefore proves \eqref{eq:joint-serfling-limit} using only square-integrable test functions.

The explicit formula \eqref{eq:modulated-energy-general-reference} gives the exact identity
\begin{equation}
 -\be N\Fr_N(\XN,\mu)
 =\frac\be2\mathsf W_N(\tl{\g}_{\mu})
 +\frac\be N\sum_{i=1}^Nh_\mu(X_i)
 -\be I_\mu.
 \label{eq:clt-exact-weight-decomposition}
\end{equation}
Since \(\int h_\mu\,\dd\mu=2I_\mu\), the weak law of large numbers and \(h_\mu\in L^1(\mu)\) give
\begin{equation}
 \frac1N\sum_{i=1}^Nh_\mu(X_i)
 \longrightarrow2I_\mu
 \qquad\text{in probability under }\mu^{\otimes N}.
 \label{eq:clt-one-body-law-large-numbers}
\end{equation}
Consequently, \eqref{eq:joint-serfling-limit}--\eqref{eq:clt-exact-weight-decomposition} imply
\begin{equation}
 \left(S_N(\varphi_1),\ldots,S_N(\varphi_k),
 -\be N\Fr_N\right)
 \Longrightarrow
 \left(\mathsf Z(\overline\varphi_1),\ldots,
 \mathsf Z(\overline\varphi_k),
 \frac\be2\mathsf W_\infty+\be I_\mu\right).
 \label{eq:clt-joint-weight-limit}
\end{equation}

The assumed bound at $b$ controls every smaller coefficient.  More precisely, for $0\leq\al\leq b$, concavity of $x\mapsto x^{\al/b}$ gives
\begin{equation}
 \K_{N,\al}(\mu)
 \leq\left(\K_{N,b}(\mu)\right)^{\al/b}.
 \label{eq:clt-temperature-monotonicity}
\end{equation}
Choose $r>1$ such that $r\be<b$.  Then
\begin{equation}
 \sup_{N\geq2}
 \E_{\mu^{\otimes N}}
 \left[e^{-r\be N\Fr_N(\XN,\mu)}\right]
 =\sup_{N\geq2}\K_{N,r\be}(\mu)<\infty.
 \label{eq:clt-weight-uniform-integrability}
\end{equation}
Thus, the Gibbs weights at inverse temperature $\be$ are uniformly integrable.

Applying the Portmanteau theorem to the nonnegative lower-semicontinuous exponential at coefficient $b$, using \eqref{eq:clt-joint-weight-limit} with $\be=b$, yields
\begin{equation}
 \E\left[e^{b\mathsf W_\infty/2}\right]<\infty.
 \label{eq:clt-gaussian-buffer}
\end{equation}
The deterministic factor $e^{bI_{\mu}}$ is finite and nonzero, so it does not affect exponential integrability.  Since $T_{\mu}$ is compact and nonnegative, \eqref{eq:clt-gaussian-buffer} is equivalent to
\begin{equation}
 b\norm{T_{\mu}}_{\op}<1.
 \label{eq:clt-strict-spectral-gap-proof}
\end{equation}
To see this directly, if $\lambda_1=\norm{T_{\mu}}_{\op}>0$, then the expectation in \eqref{eq:clt-gaussian-buffer} contains the factor
\begin{equation}
 \E\left[e^{b\lambda_1(G_1^2-1)/2}\right],
 \label{eq:clt-top-mode-factor}
\end{equation}
which is finite exactly when $b\lambda_1<1$; the case $T_{\mu}=0$ is immediate.

Let $f\in L_0^2(\mu)$.  The strict gap \eqref{eq:clt-strict-spectral-gap-proof} permits evaluation of the Gaussian quadratic tilt.  Decomposing $f$ into the eigenspaces in \eqref{eq:clt-spectral-resolution} and the kernel of $T_{\mu}$, finite-dimensional Gaussian integration gives
\begin{equation}
 \frac{\E\left[
 e^{i\mathsf Z(f)+\be\mathsf W_\infty/2}\right]}
 {\E\left[e^{\be\mathsf W_\infty/2}\right]}
 =\exp\left\{-\frac12
 \ip*{f,(\Id-\be T_{\mu})^{-1}f}_{L^2(\mu)}\right\}.
 \label{eq:clt-gaussian-quadratic-tilt}
\end{equation}
The finite spectral sums converge in $L^2$, and the common exponential moment at $b$ gives uniform integrability of their tilts. Moreover,
\begin{equation}
 \left\|(\Id-\be T_\mu^{(M)})^{-1}f
       -(\Id-\be T_\mu)^{-1}f\right\|_2
 \leq\frac{\be\|T_\mu^{(M)}-T_\mu\|_{\op}}
 {(1-\be\|T_\mu\|_{\op})^2}\|f\|_2
 \xrightarrow[M\to\infty]{}0.
 \label{eq:clt-resolvent-tail}
\end{equation}
Here $T_\mu^{(M)}$ is the first $M$ terms of \eqref{eq:clt-spectral-resolution}. These limits justify passage to the full operator in \eqref{eq:clt-gaussian-quadratic-tilt}.

We now apply \eqref{eq:clt-joint-weight-limit}.  For $t=(t_1,\ldots,t_k)\in\R^k$, set
\begin{equation}
 f_t\coloneqq\sum_{j=1}^kt_j\overline\varphi_j.
 \label{eq:clt-linear-combination}
\end{equation}
Weak convergence and \eqref{eq:clt-weight-uniform-integrability} imply
\begin{equation}
 \lim_{N\to\infty}
 \E_{\Q_{N,\be}(\mu)}
 \left[e^{iS_N(f_t)}\right]
 =\frac{\E\left[e^{i\mathsf Z(f_t)+\be\mathsf W_\infty/2}\right]}
 {\E\left[e^{\be\mathsf W_\infty/2}\right]}.
 \label{eq:clt-tilted-characteristic-function}
\end{equation}
The deterministic factor involving $I_{\mu}$ cancels between numerator and denominator.  Combining \eqref{eq:clt-gaussian-quadratic-tilt} and \eqref{eq:clt-tilted-characteristic-function} gives
\begin{equation}
 \lim_{N\to\infty}
 \E_{\Q_{N,\be}(\mu)}
 \left[e^{iS_N(f_t)}\right]
 =\exp\left\{-\frac12
 \ip*{f_t,(\Id-\be T_{\mu})^{-1}f_t}_{L^2(\mu)}\right\}.
 \label{eq:clt-limiting-characteristic-function}
\end{equation}
This is the characteristic function of the centered Gaussian vector with covariance \eqref{eq:static-clt-covariance}, and proves \eqref{eq:static-vector-clt}.

It remains to prove moment-generating-function convergence for bounded tests.  Fix $t\in\R^k$ and suppose $f_t\in L^\infty(\mu)$.  Choose $s>1$ and conjugate exponents $p,q>1$ such that
\begin{equation}
 qs\be<b.
 \label{eq:clt-mgf-exponent-choice}
\end{equation}
Hoeffding's lemma gives
\begin{equation}
 \sup_{N\geq1}
 \E_{\mu^{\otimes N}}
 \left[e^{psS_N(f_t)}\right]
 \leq\exp\left\{\frac{p^2s^2}{2}\norm{f_t}_{L^\infty(\mu)}^2\right\}.
 \label{eq:clt-product-mgf-bound}
\end{equation}
H\"older's inequality, \eqref{eq:clt-temperature-monotonicity}, and \eqref{eq:clt-mgf-exponent-choice} therefore imply
\begin{align}
 &\sup_{N\geq2}
 \E_{\mu^{\otimes N}}
 \left[e^{sS_N(f_t)-s\be N\Fr_N(\XN,\mu)}\right]
 \notag\\
 &\qquad\leq
 \left(\sup_N\E\left[e^{psS_N(f_t)}\right]\right)^{1/p}
 \left(\sup_N \K_{N,qs\be}(\mu)\right)^{1/q}<\infty.
 \label{eq:clt-mgf-uniform-integrability}
\end{align}
Thus, the real exponential numerators are uniformly integrable.  Repeating the limiting calculation leading to \eqref{eq:clt-limiting-characteristic-function}, now with $i f_t$ replaced by $f_t$, proves convergence of the joint moment-generating functions on all of $\R^k$.
\end{proof}

\begin{remark}[Equilibrium terminology]
When $\mu=\muunif$, or when $\mu$ is a thermal equilibrium measure selected by an external field, \Cref{cor:static-clt} is naturally an equilibrium central limit theorem.  For a general reference measure, it is a central limit theorem under the modulated Gibbs law \eqref{eq:modulated-gibbs-law}.
\end{remark}

\section{Local chaos under a strict partition buffer}
\label{sec:local-chaos}

We prove \Cref{thm:prescribed-local-chaos} from the partition bound \eqref{eq:nonuniform-order-one-input}. Throughout this section, $\mu\ll\muunif$ has bounded density and $0<B_\mu\leq2\d$. The first Hoeffding projection need not vanish for this reference; an exact conditional-density identity gives its required bound. A two-replica estimate controls a proportional marginal, and Hoeffding orthogonality then controls the higher projections of smaller marginals. We first introduce finite-rank conditional expectations that preserve the partition bound, so that every Gaussian calculation can be performed in finite dimension.

\subsection{Reference-dependent conditional expectations and partition control}
\label{subsec:local-cutoffs}

For a two-variable kernel, its row at $x$ is the one-variable function obtained by holding $x$ fixed. The boundedness of the density gives the following uniform row estimates for the centered kernel in \eqref{eq:hoeffding-centered-kernel}:
\begin{equation}
 \norm{h_\mu}_\infty<\infty,\qquad
 \mathop{\operatorname{ess\,sup}}_x
 \int e^{a|\tl g_\mu(x,y)|}(1+|\tl g_\mu(x,y)|)^p\,\dd\mu(y)<\infty,
 \quad 0<a<\d,\quad 1\leq p<\infty.
 \label{eq:local-row-moments}
\end{equation}
Indeed, the negative part of $\tl g_\mu$ is bounded, and its positive singularity is $-\log\operatorname{dist}(x,y)+O(1)$. The resulting radial integral is bounded by a constant times $\int_0^{1/2}r^{\d-1-a}(1+|\log r|)^p\,\dd r$. The same argument gives every finite polynomial section moment. In particular, $T_\mu$ is Hilbert--Schmidt. It is nonnegative: for $f\in L^2(\mu)$, the signed density $(f-\int f\,\dd\mu)\mu$ belongs to $L^2(\muunif)$ and has mean zero, so the Fourier positivity of $\g$ applies. This also proves positivity when the reference density vanishes.

Jensen's inequality and the product-energy identity \eqref{eq:dyn-product-exact} give
\begin{equation}
 \K_{N,\be}(\mu)\geq e^{\be I_\mu},\qquad N\geq2,\quad \be\geq0.
 \label{eq:local-partition-lower}
\end{equation}
The bounded one-body term in \eqref{eq:clt-exact-weight-decomposition} transfers \eqref{eq:nonuniform-order-one-input} to the centered quadratic Laplace transform. Applying \Cref{lem:transport-partition-implies-gap} at each $0<b<B_\mu$ and then increasing $b$ to $B_\mu$ yields
\begin{equation}
 B_\mu\norm{T_\mu}_{\op}\leq1.
 \label{eq:local-spectral-bound}
\end{equation}
The spectral-control lemma uses the weak limit of the quadratic statistic and the partition bound. The joint central limit theorem under the Gibbs law is not needed for this step.

We use conditional expectations adapted to $\mu$ to preserve the partition bound at every approximation scale. Let $\mathscr P_n$ be the refining dyadic partitions of the torus, and let $Q_{\mu,n}$ be conditional expectation in $L^2(\mu)$ onto their cells. This is the construction used for the cell projection in Section~\ref{sec:coarse-decomposition}, with the reference measure and the scale now specified. Cells of zero $\mu$-mass are omitted. Set
\begin{equation}
 q_n=(Q_{\mu,n}\otimes Q_{\mu,n})\tl g_\mu,\qquad
 h_{\mu,n}=Q_{\mu,n}h_\mu,\qquad
 T_{\mu,n}=Q_{\mu,n}T_\mu Q_{\mu,n}.
 \label{eq:local-Markov-cutoff}
\end{equation}
Thus, $q_n$ is bounded, finite rank, symmetric, and canonical, $T_{\mu,n}\geq0$, and $\norm{T_{\mu,n}}_{\op}\leq\norm{T_\mu}_{\op}$. Define the cutoff exponent and its normalization by
\begin{equation}
 V_{N,\be,n}(X_N)
 =\frac{\be}{2}\mathsf W_N(q_n)
  +\frac{\be}{N}\sum_{i=1}^N(h_{\mu,n}(x_i)-I_\mu),
 \qquad
 \K_{N,\be,n}(\mu)=\E_{\mu^{\otimes N}}[e^{V_{N,\be,n}}].
 \label{eq:local-cutoff-exponent}
\end{equation}
Write $f_{N,\be,n}=e^{V_{N,\be,n}}/\K_{N,\be,n}(\mu)$ and let $f_{N,\be,n}^{(k)}$ be its $k$-coordinate marginal density relative to $\mu^{\otimes k}$. The uncut densities are denoted by $f_{N,\be}$ and $f_{N,\be}^{(k)}$.

\begin{prop}[Uniform partition control under reference-adapted cutoffs]
\label{prop:local-Markov-cutoff}
For $N\geq2$ and $\be\geq0$,
\begin{equation}
 e^{\be I_\mu}\leq\K_{N,\be,n}(\mu)\leq\K_{N,\be}(\mu).
 \label{eq:local-cutoff-domination}
\end{equation}
For fixed $N$ and $0\leq\be<B_\mu$, the cutoff densities converge to $f_{N,\be}$ in $L^1(\mu^{\otimes N})$, and their marginals converge in the corresponding $L^1$ spaces. The bounds \eqref{eq:local-row-moments} hold with $q_n$ in place of $\tl g_\mu$, uniformly in $n$.
\end{prop}
\begin{proof}
Let $X_1,\ldots,X_N$ be independent with law $\mu$. Conditional on their cells, resample independent $Z_i$ with the law of $\mu$ conditioned on the cell of $X_i$. This includes independent resampling of particles assigned to the same cell. The $Z_i$ are unconditionally independent with law $\mu$, and the conditional expectation of the full exponent $-\be N\Fr_N(Z_N,\mu)$ is precisely \eqref{eq:local-cutoff-exponent}. Here both the pair term and the one-body term must be conditioned. Conditional Jensen proves the upper bound in \eqref{eq:local-cutoff-domination}. The product mean of the cutoff exponent is $\be I_\mu$, since $q_n$ is canonical and $\int h_{\mu,n}\,\dd\mu=2I_\mu$. A second application of Jensen gives the lower bound.

Martingale convergence gives $q_n\to\tl g_\mu$ in $L^2(\mu^{\otimes2})$ and almost everywhere, and $h_{\mu,n}\to h_\mu$ in $L^2(\mu)$ and almost everywhere. Consequently, the unnormalized densities converge almost everywhere for each fixed $N$. Choose $r>1$ with $r\be<B_\mu$, when $\be>0$. Then
\begin{equation}
 \int e^{rV_{N,\be,n}}\,\dd\mu^{\otimes N}
 =\K_{N,r\be,n}(\mu)\leq\K_{N,r\be}(\mu)<\infty.
 \label{eq:local-cutoff-uniform-integrability}
\end{equation}
Uniform integrability implies $L^1$ convergence. The normalizations converge and have a common positive lower bound, which proves convergence of the normalized densities. Conditional expectation is an $L^1$ contraction, so the same is true of each marginal. At $\be=0$ every density is one.

For a positive-mass cell $A$ containing $x$ and a convex nonnegative function $\Psi$, conditional Jensen in the two variables gives
\begin{equation}
 \int\Psi(q_n(x,y))\,\dd\mu(y)
 \leq\frac1{\mu(A)}\int_A\int\Psi(\tl g_\mu(z,y))\,\dd\mu(y)\dd\mu(z).
 \label{eq:local-cutoff-row-Jensen}
\end{equation}
Use $\Psi(t)=|t|^p$, or $\Psi(t)=e^{a|t|}(1+|t|)^p$. The latter is convex because it is the composition of $|t|$ with a convex nondecreasing function on $[0,\infty)$. This proves the uniform section bounds. The final bound depends on the section moment of the original kernel, independently of the cell masses.
\end{proof}

For later use, fix $\be_0<B_\mu$ and choose $1<\theta\leq2$ with $\theta\be_0<B_\mu$. Raising the density to the power $\theta$ gives the exact identity
\begin{equation}
 \norm{f_{N,\be,n}}_{L^\theta(\mu^{\otimes N})}^\theta
 =\frac{\K_{N,\theta\be,n}(\mu)}{\K_{N,\be,n}(\mu)^\theta}
 \leq e^{-\theta\be I_\mu}P_{\mu,\theta\be_0},
 \qquad 0\leq\be\leq\be_0.
 \label{eq:local-full-density-buffer}
\end{equation}
The same bound holds without the cutoff and for every marginal, by conditional Jensen. In particular, it applies to the reservoir marginal of the $N$-particle law. That reservoir is not replaced by an $(N-1)$-particle Gibbs law.

\subsection{Finite Gaussian moments and a proportional marginal}
\label{subsec:local-replicas}

Let $Y_n^1,Y_n^2$ be independent centered finite-dimensional Gaussian fields with covariance $q_n$, and write $Y_n=(Y_n^1,Y_n^2)$. Such fields exist by the finite-rank positivity established above. Their cellwise constant versions define every diagonal covariance $q_n(x,x)$. For $a\in\R^2$, set
\begin{equation}
 M_{a,n}=\int\exp\left\{a\cdot Y_n(x)
 -\frac{|a|^2}{2}q_n(x,x)+|a|^2(h_{\mu,n}(x)-I_\mu)\right\}\,\dd\mu(x).
 \label{eq:local-Gaussian-mass}
\end{equation}
The diagonal covariance in the Wick factor cancels the self-interactions produced by Gaussian integration. Expanding an integer power and applying the Gaussian exponential identity gives
\begin{equation}
 \E[M_{a,n}^M]=\K_{M,M|a|^2,n}(\mu),\qquad M\geq2.
 \label{eq:local-integer-Wick-identity}
\end{equation}
Indeed, the diagonal variance cancels the Wick factors. Gaussian integration leaves coefficient $|a|^2$ on each unordered pair and on each one-body term. These are exactly the coefficients in \eqref{eq:local-cutoff-exponent} at inverse temperature $M|a|^2$.

Put $\lambda^2=\be/N$ and choose the Gaussian coefficient vectors for the shared and separate coordinates
\begin{equation}
 a_0=(\lambda,\lambda),\qquad a_1=(\lambda,0),\qquad a_2=(0,\lambda).
 \label{eq:local-replica-charges}
\end{equation}
Their Gram matrix is
\begin{equation}
 (a_i\cdot a_j)_{0\leq i,j\leq2}
 =\frac{\be}{N}
 \begin{pmatrix}2&1&1\\1&1&0\\1&0&1\end{pmatrix}.
 \label{eq:local-replica-Gram}
\end{equation}
Squaring the marginal density and introducing the independent second field yields
\begin{equation}
 \int (f_{N,\be,n}^{(k)})^2\,\dd\mu^{\otimes k}
 =\frac{\E[M_{a_0,n}^kM_{a_1,n}^{N-k}M_{a_2,n}^{N-k}]}
 {\K_{N,\be,n}(\mu)^2}.
 \label{eq:local-exact-replica}
\end{equation}
To check this identity directly, express the $k$-particle marginal density by integrating $N-k$ insertions from one Gaussian field and keeping the remaining $k$ insertions fixed. Upon squaring, the shared insertions have coefficient vector $a_0$, while the two separate groups have coefficient vectors $a_1,a_2$. Their one-body coefficients are the squared lengths of these vectors. All interchanges are justified by positivity and finite-dimensional Gaussian exponential moments.

The following finite-graph calculation handles particle numbers too small for the uniform moment argument.

\begin{lemma}[Collision condition for two overlapping replicas]
\label{lem:local-replica-collisions}
For fixed $N\geq2$ and $1\leq k\leq N$, the uncut numerator in \eqref{eq:local-exact-replica}, interpreted by squaring the original marginal integral, is finite whenever
\begin{equation}
 \be+\d\frac{k-1}{N-1}<2\d.
 \label{eq:local-replica-collision-condition}
\end{equation}
For the homogeneous reference this condition is also necessary. If it holds at $\be_0$, the numerator is bounded uniformly on $0\leq\be\leq\be_0$.
\end{lemma}
\begin{proof}
There are $k$ shared vertices and two separate groups of $r=N-k$ vertices. An induced cluster containing $a$ shared vertices, $b$ vertices from the first group, and $c$ from the second has total edge multiplicity
\begin{equation}
 E(a,b,c)=\binom{a+b}{2}+\binom{a+c}{2},
 \quad 0\leq a\leq k,\quad 0\leq b,c\leq r,
 \quad a+b+c\geq2.
 \label{eq:local-replica-edge-count}
\end{equation}
We claim
\begin{equation}
 \max\frac{E(a,b,c)}{a+b+c-1}
 =\frac{N(N-1)}{2N-k-1}.
 \label{eq:local-maximal-cluster-density}
\end{equation}
When $r=0$, the ratio is $a\leq N$ and the assertion follows. Otherwise, subtract the right-hand side times $a+b+c-1$ from $E(a,b,c)$. This is a convex quadratic function, since its quadratic part is $((a+b)^2+(a+c)^2)/2$. Its maximum over the box cut by $a+b+c\geq2$ occurs at a vertex. At uncut box vertices of size at least two, the possible ratios are
\begin{equation}
 k,\quad \frac r2,\quad\frac{r(r-1)}{2r-1},\quad
 \frac{N(N-1)+k(k-1)}{2(N-1)},\quad
 \frac{N(N-1)}{2N-k-1},
 \label{eq:local-cluster-corner-ratios}
\end{equation}
with entries omitted if their cluster has fewer than two vertices. The first and fourth entries are bounded by the last because the last minus the first, and the last minus the fourth, respectively equal
\begin{equation}
 \frac{r(r-1)}{2N-k-1},\qquad
 \frac{(k-1)r(r-1)}{2(N-1)(2N-k-1)}.
 \label{eq:local-cluster-corner-differences}
\end{equation}
The middle two are at most $N/2$, and the last is at least $N/2$. At a new vertex on the cutting plane the total size is two. If $k\geq2$, its edge multiplicity is at most two, and the last ratio in \eqref{eq:local-cluster-corner-ratios} is at least $k\geq2$. If $k=1$, maximize instead over the larger triangle cut by $a\leq1$ and $a+b+c=2$; its vertices are $(0,2,0),(0,0,2),(1,1,0),(1,0,1)$, all of multiplicity one. The last ratio then equals $N/2\geq1$. This proves the upper bound in \eqref{eq:local-maximal-cluster-density}; the full cluster attains it.

Each edge contributes exponent $\be/N$ times its multiplicity. The induction and polar-coordinate argument of \Cref{lem:nested-cluster-power-count} applies to this weighted graph: for each cluster of $r$ particles, the sum of its edge exponents must be less than $\d(r-1)$. In view of \eqref{eq:local-maximal-cluster-density}, the resulting product is integrable precisely when $(\be/N)N(N-1)/(2N-k-1)<\d$, which is \eqref{eq:local-replica-collision-condition}. The centered and uncentered interactions differ by bounded one-body terms, and the density of $\mu^{\otimes(2N-k)}$ is bounded. These facts prove sufficiency for $\mu$; for the homogeneous reference the comparison is two-sided and proves necessity. For the complete uncut exponent $\be V$, the inequality $e^{\be V}\leq1+e^{\be_0V}$ gives the asserted uniformity. Necessity is not asserted for a density which may vanish.
\end{proof}

\begin{prop}[A uniformly bounded proportional marginal]
\label{prop:local-replica-anchor}
If \eqref{eq:prescribed-local-window} holds, then
\begin{equation}
 \sup_{N\geq2}\sup_{0\leq\be\leq\be_0}
 \sup_{1\leq k\leq\lfloor\al_0N\rfloor}
 \chi^2\!\left(\Q_{N,\be}^{(k)}(\mu)\mid\mu^{\otimes k}\right)<\infty,
 \label{eq:local-replica-anchor}
\end{equation}
where an empty marginal range imposes no condition.
\end{prop}
\begin{proof}
The case $\be_0=0$ is immediate. Fix a number $b$ strictly between $\be_0(1+\al_0)$ and $B_\mu$, and first take $N$ large enough that $\be_0(1+\al_0+N^{-1})\leq b$. For $1\leq k<N$, put
\begin{equation}
 \omega=\frac{2k}{N+k},\qquad
 q_0=\frac{k}{\omega}=\frac{N+k}{2},\qquad
 q_1=\frac{N-k}{1-\omega}=N+k.
 \label{eq:local-replica-Holder-exponents}
\end{equation}
Apply H\"older to $M_{a_0,n}$ for the shared coordinates and to the product of $M_{a_1,n}$ and $M_{a_2,n}$ for the separate coordinates. The latter two Gaussian integrals are independent. Rotational invariance of the two Gaussian fields therefore gives
\begin{equation}
 \E[M_{a_0,n}^kM_{a_1,n}^{N-k}M_{a_2,n}^{N-k}]
 \leq \E[M_{\sqrt2\lambda,n}^{q_0}]^\omega
       \E[M_{\lambda,n}^{q_1}]^{2(1-\omega)},
 \label{eq:local-replica-Holder}
\end{equation}
where a scalar subscript denotes the coefficient of a single Gaussian field. The exponent $2(1-\omega)$ results from grouping the two independent Gaussian integrals for the separate coordinates before applying H\"older.

Set $m_0=\lceil(N+k)/2\rceil\geq2$. Monotonicity of probability-space $L^p$ norms, followed by \eqref{eq:local-integer-Wick-identity} and \eqref{eq:local-cutoff-domination}, gives
\begin{align}
 \E[M_{\sqrt2\lambda,n}^{q_0}]^\omega
 &\leq\K_{m_0,\,2\be m_0/N}(\mu)^{k/m_0},
 \label{eq:local-shared-moment-temperature}\\
 \E[M_{\lambda,n}^{q_1}]
 &\leq\K_{N+k,\,\be(1+k/N)}(\mu).
 \label{eq:local-private-moment-temperature}
\end{align}
Both inverse temperatures are at most $b$, since $2m_0\leq N+k+1$. Since $k/m_0\leq\omega$ and $P_{\mu,b}\geq1$, the right-hand side of \eqref{eq:local-replica-Holder} is at most $P_{\mu,b}^{2}$. The denominator in \eqref{eq:local-exact-replica} is at least $e^{2\be I_\mu}$. This bounds the cutoff marginal in $L^2$, uniformly in $n,N,k,\be$ in the displayed ranges. Its $L^1$ convergence from \Cref{prop:local-Markov-cutoff} and weak compactness in $L^2$ identify the uncut limit and transfer the bound. If $k=N$ is admissible, then $\al_0=1$ and $2\be_0<B_\mu$; the numerator is simply $\K_{N,2\be,n}(\mu)$, so the same conclusion follows directly.

Only finitely many particle numbers remain. For $\al=k/N\leq\al_0$,
\begin{equation}
 \be_0<\frac{2\d}{1+\al_0}
 \leq\frac{2\d}{1+\al}
 \leq\d(2-\al)
 \leq\d\left(2-\frac{k-1}{N-1}\right).
 \label{eq:local-window-implies-collision}
\end{equation}
The third inequality follows from $\al(1-\al)\geq0$, and the last from $(k-1)/(N-1)\leq k/N$. Thus \Cref{lem:local-replica-collisions} supplies a finite bound uniform in $\be$ for each remaining pair $(N,k)$. Taking their maximum proves \eqref{eq:local-replica-anchor}. This treats the remaining particle numbers directly with the original interaction.
\end{proof}

\subsection{The one-particle conditional density}
\label{subsec:local-cavity}

The first projection need not vanish for a prescribed reference. We control it by a conditional expansion averaged under the actual interacting reservoir. Fix $\be_0<B_\mu$, choose $\theta$ as in \eqref{eq:local-full-density-buffer}, and set $s=\theta/(\theta-1)\geq2$. For a reservoir configuration $Y=(y_2,\ldots,y_N)$, define
\begin{equation}
 B_Y(x)=\frac{\be}{N}
 \left(\sum_{j=2}^Nq_n(x,y_j)+h_{\mu,n}(x)-2I_\mu\right),
 \qquad \int B_Y\,\dd\mu=0.
 \label{eq:local-cavity-field}
\end{equation}
This notation denotes the field for fixed $N,n,\be$ throughout this subsection.

\begin{lemma}[Weighted moments under the product law]
\label{lem:local-product-moments}
For each fixed $p\geq2$ and $a\geq0$, there are $N_0$ and $C<\infty$ such that
\begin{equation}
 \E_{\mu^{\otimes N}}\!\left[e^{a(B_Y(X))_+}|B_Y(X)|^p\right]
 \leq C N^{-p/2},
 \quad N\geq N_0,\quad n\geq1,\quad 0\leq\be\leq\be_0.
 \label{eq:local-weighted-product-moments}
\end{equation}
Here $X,y_2,\ldots,y_N$ are independent with law $\mu$.
\end{lemma}
\begin{proof}
Condition on $X=x$. The variables $q_n(x,y_j)$ are independent and centered. To bound a moment of order $p$, choose an even integer $2r\geq p$ and expand the $2r$-th power of their sum. A summand with an index appearing once has expectation zero. Every other summand contains at most $r$ distinct indices, so there are at most $C_rN^r$ such summands. H\"older's inequality and the uniform section moment of order $2r$ bound each one by a constant. Taking the power $p/(2r)$ therefore gives
\begin{equation}
 \E\!\left[\left|\sum_{j=2}^Nq_n(x,y_j)\right|^p\right]
 \leq C_pN^{p/2}.
 \label{eq:local-centered-sum-moments}
\end{equation}
The bounded one-body term in \eqref{eq:local-cavity-field} implies $\E[|B_Y(X)|^p]\leq C_pN^{-p/2}$, uniformly also in $x$.

For a fixed $t\geq0$, independence gives
\begin{equation}
 \E[e^{tB_Y(x)}]
 =e^{t\be(h_{\mu,n}(x)-2I_\mu)/N}
 \left(\int e^{t\be q_n(x,y)/N}\,\dd\mu(y)\right)^{N-1}.
 \label{eq:local-product-exponential}
\end{equation}
The elementary remainder inequality
\begin{equation}
 0\leq e^z-1-z\leq\tfrac12 z^2e^{z_+}
 \label{eq:local-exponential-remainder}
\end{equation}
and row centering bound the integral in parentheses by $1+C/N^2$ for all sufficiently large $N$. To obtain a uniform constant, choose $a'\in(0,\d)$ and then $N_0$ so that $t\be_0/N_0\leq a'$, and use \eqref{eq:local-row-moments} at exponent $a'$. The prefactor in \eqref{eq:local-product-exponential} is bounded as well. Thus $\E[e^{tB_Y(X)}]\leq C_t$. Finally, $e^{2az_+}\leq1+e^{2az}$ and Cauchy--Schwarz yield
\begin{equation}
 \E[e^{a(B_Y(X))_+}|B_Y(X)|^p]
 \leq\E[e^{2a(B_Y(X))_+}]^{1/2}
      \E[|B_Y(X)|^{2p}]^{1/2}
 \leq C N^{-p/2}.
 \label{eq:local-weighted-product-closure}
\end{equation}
This is an averaged estimate. It requires no cutoff-uniform pointwise bound on the attractive field.
\end{proof}

Let $G_{N,\be,n}$ be the reservoir density relative to $\mu^{\otimes(N-1)}$. Its $L^\theta$ norm is uniformly bounded by \eqref{eq:local-full-density-buffer}. The marked conditional density is
\begin{equation}
 p_Y(x)=\frac{e^{B_Y(x)}}{\int e^{B_Y}\,\dd\mu}.
 \label{eq:local-marked-conditional-density}
\end{equation}
Indeed, the raw marked exponent in \eqref{eq:local-cutoff-exponent} has $h_{\mu,n}-I_\mu$ in place of $h_{\mu,n}-2I_\mu$; the difference is constant in $x$ and cancels upon normalization. Put $\epsilon(z)=e^z-1-z$ and $\delta_Y=\int\epsilon(B_Y)\,\dd\mu\geq0$. Then
\begin{equation}
 p_Y-1=B_Y+
 \frac{\epsilon(B_Y)-\delta_Y-B_Y\delta_Y}{1+\delta_Y}.
 \label{eq:local-normalized-exponential-identity}
\end{equation}

\begin{prop}[The averaged nonlinear remainder]
\label{prop:local-cavity-remainder}
For all sufficiently large $N$, the centered function
\begin{equation}
 r_{N,\be,n}(x)=\E_{G_{N,\be,n}}\!\left[
 \frac{\epsilon(B_Y(x))-\delta_Y-B_Y(x)\delta_Y}{1+\delta_Y}\right]
 \label{eq:local-cavity-remainder-definition}
\end{equation}
satisfies
\begin{equation}
 \sup_n\sup_{0\leq\be\leq\be_0}
 \norm{r_{N,\be,n}}_{L^2(\mu)}\leq\frac{C}{N}.
 \label{eq:local-cavity-remainder-bound}
\end{equation}
\end{prop}
\begin{proof}
Write $G=G_{N,\be,n}$ only within this proof. Positivity of $\epsilon$ and $\delta_Y$, followed by Minkowski for the signed mixed term, gives
\begin{equation}
 \norm{r_{N,\be,n}}_2
 \leq\norm{\E_G[\epsilon(B_Y)]}_2+\E_G[\delta_Y]
     +\E_G[\delta_Y\norm{B_Y}_2].
 \label{eq:local-remainder-three-terms}
\end{equation}
For the mixed term, this uses
$\norm{\E_G[B_Y\delta_Y/(1+\delta_Y)]}_2
\leq\E_G[\delta_Y\norm{B_Y}_2]$.
Minkowski is applied before bounding the denominator, since the integrand in this expectation has no fixed sign.

H\"older in the reservoir variables and then Jensen, with $2/s\leq1$, give
\begin{equation}
 \norm{\E_G[\epsilon(B_Y)]}_2
 \leq C\left(\E_{\mu^{\otimes N}}[\epsilon(B_Y(X))^s]\right)^{1/s}
 \leq\frac{C}{N}.
 \label{eq:local-remainder-first-term}
\end{equation}
The last inequality follows from \eqref{eq:local-exponential-remainder} and \Cref{lem:local-product-moments} with $p=2s$, $a=s$. Jensen in the marked variable gives the same bound for $\E_G[\delta_Y]$. For the remaining term, H\"older in the reservoir and Jensen in the marked variable yield
\begin{align}
 \E_G[\delta_Y^2]
 &\leq C\left(\E_{\mu^{\otimes N}}[\epsilon(B_Y(X))^{2s}]\right)^{1/s}
 \leq\frac{C}{N^2},
 \label{eq:local-remainder-delta-second}\\
 \E_G[\norm{B_Y}_2^2]
 &\leq C\left(\E_{\mu^{\otimes N}}[|B_Y(X)|^{2s}]\right)^{1/s}
 \leq\frac{C}{N}.
 \label{eq:local-remainder-field-second}
\end{align}
The first estimate uses $p=4s$, $a=2s$, and the second uses $p=2s$, $a=0$. Cauchy--Schwarz therefore bounds the last term of \eqref{eq:local-remainder-three-terms} by $CN^{-3/2}$. This proves \eqref{eq:local-cavity-remainder-bound}. Integrating \eqref{eq:local-normalized-exponential-identity} in $x$ shows that each remainder, and hence its average, has zero $\mu$-mean.
\end{proof}

\begin{prop}[Resolvent estimate for the first projection]
\label{prop:local-screened-first-projection}
Set $u_{N,\be,n}=f_{N,\be,n}^{(1)}-1$. Then
\begin{equation}
 \left(\Id-\be\frac{N-1}{N}T_{\mu,n}\right)u_{N,\be,n}
 =\frac{\be}{N}(h_{\mu,n}-2I_\mu)+r_{N,\be,n}.
 \label{eq:local-screened-equation}
\end{equation}
For all sufficiently large $N$,
\begin{equation}
 \sup_n\sup_{0\leq\be\leq\be_0}\norm{u_{N,\be,n}}_2\leq\frac{C}{N}.
 \label{eq:local-first-projection-bound}
\end{equation}
\end{prop}
\begin{proof}
Average \eqref{eq:local-normalized-exponential-identity} under the reservoir law. Exchangeability identifies each reservoir one-coordinate marginal with $f_{N,\be,n}^{(1)}$. Canonicality of $q_n$ then gives exactly \eqref{eq:local-screened-equation}. Both sides have zero mean. Since $\theta\be_0<B_\mu$, \eqref{eq:local-spectral-bound} and compression by $Q_{\mu,n}$ imply
\begin{equation}
 \left\|\left(\Id-\be\frac{N-1}{N}T_{\mu,n}\right)^{-1}\right\|_{L^2_0(\mu)\to L^2_0(\mu)}
 \leq\frac{\theta N}{(\theta-1)N+1}\leq\frac{\theta}{\theta-1}.
 \label{eq:local-screened-resolvent-bound}
\end{equation}
Apply this bound to \eqref{eq:local-screened-equation}, using \Cref{prop:local-cavity-remainder} and the $L^2$ contraction of $Q_{\mu,n}$ on the fixed centered one-body function $h_\mu-2I_\mu$. This proves \eqref{eq:local-first-projection-bound}.
\end{proof}

For each fixed $N$ and $\be\leq\be_0$, \Cref{prop:local-Markov-cutoff} gives convergence of the first marginal in $L^1(\mu)$. The uniform $L^2$ bound \eqref{eq:local-first-projection-bound} passes to this limit by weak compactness and lower semicontinuity. For the finitely many remaining $N$, \Cref{lem:local-replica-collisions} with $k=1$ applies because $\be_0<B_\mu\leq2\d$. Its bound uniform in $\be$, together with \eqref{eq:local-partition-lower}, absorbs these particle numbers into the constant. This proves \eqref{eq:prescribed-one-particle-chi} for every $N\geq2$.

\subsection{Hoeffding projections and completion of the static bounds}
\label{subsec:local-Hoeffding-entropy}

The proportional-marginal bound and the first-projection estimate now supply the two inputs to \Cref{lem:local-Hoeffding}. Orthogonality gives the chi-square estimate for smaller marginals. A separate full-entropy bound then extends the entropy conclusion to every marginal size.

\begin{proof}[Proof of the chi-square assertion in \Cref{thm:prescribed-local-chaos}]
The one-particle assertion was established above. Under \eqref{eq:prescribed-local-window}, take $M=\lfloor\al_0N\rfloor$. For sufficiently large $N$, one has $M\geq2$ and $M\geq\al_0N/2$. Apply \Cref{lem:local-Hoeffding} to the $M$-particle marginal of the original $N$-particle law. \Cref{prop:local-replica-anchor} bounds its full chi-square divergence, while \eqref{eq:prescribed-one-particle-chi} bounds its first projection. Thus
\begin{equation}
 \chi^2\!\left(\Q_{N,\be}^{(k)}(\mu)\mid\mu^{\otimes k}\right)
 \leq\frac{Ck}{N^2}+C\frac{k(k-1)}{M(M-1)}
 \leq C\frac{k^2}{N^2},\qquad 1\leq k\leq M.
 \label{eq:local-chi-square-closure}
\end{equation}
The finitely many smaller $N$ are absorbed into the constant using \Cref{prop:local-replica-anchor}. Empty ranges require no estimate. This proves \eqref{eq:prescribed-local-chi}.
\end{proof}

\begin{prop}[Full relative entropy under a strict partition buffer]
\label{prop:local-full-entropy}
For every $\be_0<B_\mu$,
\begin{equation}
 \sup_{N\geq2}\sup_{0\leq\be\leq\be_0}
 \Hr\!\left(\Q_{N,\be}(\mu)\mid\mu^{\otimes N}\right)<\infty.
 \label{eq:local-full-entropy-bound}
\end{equation}
\end{prop}
\begin{proof}
Fix $\be_0<b<B_\mu$ and set $L_N(\be)=\log\K_{N,\be}(\mu)$. For fixed $N$, $-N\Fr_N$ is bounded below, and the partition estimate at $b$ supplies an exponential moment above every $\be\leq\be_0$. It therefore justifies differentiation under the integral, including the right derivative at zero. In particular, entropy and energy are separately finite, and
\begin{equation}
 \Hr\!\left(\Q_{N,\be}(\mu)\mid\mu^{\otimes N}\right)
 =\be L_N'(\be)-L_N(\be).
 \label{eq:local-entropy-log-partition}
\end{equation}
For example, the positive part of $-N\Fr_N$ times $e^{-\be N\Fr_N}$ is dominated by a constant times $e^{-bN\Fr_N}$, and its negative part is bounded for fixed $N$. Convexity and \eqref{eq:local-partition-lower} give
\begin{equation}
 L_N'(\be)\leq\frac{L_N(b)-L_N(\be)}{b-\be}
 \leq\frac{\log P_{\mu,b}+\be_0|I_\mu|}{b-\be_0},
 \qquad -L_N(\be)\leq\be_0|I_\mu|.
 \label{eq:local-entropy-partition-bound}
\end{equation}
Substitution in \eqref{eq:local-entropy-log-partition} proves \eqref{eq:local-full-entropy-bound} uniformly in $N$ and $\be$.
\end{proof}

\begin{proof}[Proof of the entropy assertion in \Cref{thm:prescribed-local-chaos}]
Choose $\al_0\in(0,1]$ with $\be_0(1+\al_0)<B_\mu$. For $k\leq\lfloor\al_0N\rfloor$, the inequality $t\log t-t+1\leq(t-1)^2$ and \eqref{eq:prescribed-local-chi} give the desired entropy bound. For $k>\lfloor\al_0N\rfloor$, one has $k>\al_0N$. \Cref{lem:relative-entropy-subset,prop:local-full-entropy} give a constant bound on the marginal entropy, which is at most $C\al_0^{-2}k^2/N^2$. This argument includes all small particle numbers and all marginal sizes without a divisibility assumption. It proves \eqref{eq:prescribed-all-entropy}.
\end{proof}

\begin{proof}[Proof of \Cref{cor:homogeneous-local-chaos-upper-bounds,cor:transported-local-chaos}]
For the homogeneous law, simultaneous translation of all coordinates preserves its density and its product reference. Its one-particle marginal is consequently translation invariant and is the normalized Haar measure. This proves \eqref{eq:homogeneous-exact-first-marginal} for every integrable inverse temperature, by \Cref{prop:fixed-particle-integrability}. In the strict global-minimization range, \eqref{eq:principal-order-one} supplies \eqref{eq:nonuniform-order-one-input} with $B_\mu=\beta_{\gm}$. Apply \Cref{prop:local-replica-anchor,lem:local-Hoeffding} as above, now with the first projection exactly zero. This gives $Ck(k-1)/N^2$ for large $N$; the finite set of remaining admissible pairs is covered by \Cref{lem:local-replica-collisions}. At $k=1$ both divergences are zero. The entropy argument above proves \eqref{eq:homogeneous-all-entropy}. Translation invariance is used for the original law, not for an arbitrary finite conditional-expectation cutoff.

For a bi-Lipschitz $\Phi$, its pushforward of normalized Haar measure has bounded density. \Cref{prop:order-one-transport} supplies \eqref{eq:nonuniform-order-one-input} with the already defined $B_\mu=\beta_{\tr}(\Phi)\leq\beta_{\gm}\leq2\d$. Thus \Cref{thm:prescribed-local-chaos} applies and proves \Cref{cor:transported-local-chaos}. No additional transport estimate or determinant limit is needed.
\end{proof}


\subsection{Homogeneous shared-label identities and partition bounds}\label{subsec:replica-identities}


\subsubsection{The exact shared-label integral}
Fix $\be_0$ and $\alpha$ as in \eqref{eq:homogeneous-replica-window}. Let $\pi_{N,\ell}$ be the law of two $N$-tuples $(X^{(1)},X^{(2)})$ sharing their first $\ell$ coordinates, with the remaining $2N-2\ell$ coordinates and the $\ell$ common coordinates independent with law $\muunif$.  Each individual replica has law $\muunif^{\otimes N}$, and $\mathsf W_N^{(r)}$ denotes the statistic \eqref{eq:particle-kernel-sum} evaluated on replica $r\in\{1,2\}$.  Put
\begin{equation}\label{eq:J}
 J_{N,\ell}(\be)
 =\E_{\pi_{N,\ell}}\!\left[
 e^{\frac\be2(\mathsf W_N^{(1)}(\g)+\mathsf W_N^{(2)}(\g))}\right].
\end{equation}
Tonelli's theorem gives the extended identity
\begin{equation}\label{eq:replica-identity}
 R_{N,\be}(\ell):=\int\bigl(F_{N,\be}^{(\ell)}\bigr)^2\,\dd \muunif^{\otimes\ell}
 =\frac{J_{N,\ell}(\be)}{Z_{N,\be}^2}.
\end{equation}
Finiteness will follow from the estimates below; it is not assumed in this identity.

Write $J_{N,\ell}^{\eta}$ for the same integral with $\g_\eta$ in place of $\g$.  Let $\mathsf Z_1,\mathsf Z_2$ be independent centered real Gaussian fields with covariance $\g_\eta(x-y)$.  Smoothness and the rapidly decaying Fourier coefficients give jointly measurable smooth versions.  Their spatial integrals vanish almost surely.  With $D=D_\eta$, define
\begin{align}
 A_j&=\int\exp\!\left\{\sqrt{\frac\be N}\mathsf Z_j(x)-\frac{\be D}{2N}\right\}\,\dd \muunif(x),\quad j=1,2,\label{eq:Aj}\\
 A_{12}&=\int\exp\!\left\{\sqrt{\frac\be N}(\mathsf Z_1(x)+\mathsf Z_2(x))-
 \frac{\be D}{N}\right\}\,\dd \muunif(x),\label{eq:A12}\\
 P_{N,\ell}^{\eta}&=(A_1A_2)^{N-\ell}A_{12}^{\ell}.\label{eq:P}
\end{align}
The Hubbard--Stratonovich transformation gives
\begin{equation}\label{eq:HS-replica}
 J_{N,\ell}^{\eta}(\be)=\E_G[P_{N,\ell}^{\eta}].
\end{equation}
Indeed, each Gaussian field produces one copy of $\be\mathsf W_N(\g_\eta)/2$.  Each shared label contributes to both fields, and its Wick subtraction is $\be D/N$.  Thus, a pair of shared labels is counted twice, while the self-interactions are removed once in each replica.  This verifies all diagonal and factor-of-two terms in \eqref{eq:HS-replica}.

\subsubsection{The scalar partition bound controls the replica moment}
Choose an even integer $u\geq2$ sufficiently large that, with $r=u/(u-1)$,
\begin{equation}\label{eq:r-choice}
 2u>\d,\qquad r\be_0(1+\alpha)<\beta_{\gm}.
\end{equation}
Fix $b$ strictly between the last two quantities, and then choose $N_0$ so large that
\begin{equation}\label{eq:ceiling-buffer}
 r\be_0(1+\alpha)+\frac{2r\be_0}{N_0}<b.
\end{equation}
The exponents and the threshold $N_0$ depend only on the fixed temperature interval and overlap window. They are independent of $N$, $\ell$, $\eta$, and $\be$.

\begin{lemma}[Higher-moment bound for the replica exponential]\label{lem:replica-buffer}
Under \eqref{eq:r-choice}--\eqref{eq:ceiling-buffer},
\begin{equation}\label{eq:replica-Gaussian-buffer}
 \sup_{N\geq N_0}\sup_{0\leq\ell\leq\alpha N}\sup_{0\leq\be\leq\be_0}\sup_{0<\eta\leq\eta_0}
 \|P_{N,\ell}^{\eta}\|_{L^r(G)}<\infty.
\end{equation}
Moreover, both with and without the cutoff,
\begin{equation}\label{eq:replica-particle-buffer}
 \sup_{N\geq N_0}\sup_{0\leq\ell\leq\alpha N}\sup_{0\leq\be\leq\be_0}\sup_{0<\eta\leq\eta_0}
 J_{N,\ell}^{\eta}(r\be)<\infty,
 \qquad
 \sup_{N\geq N_0}\sup_{0\leq\ell\leq\alpha N}\sup_{0\leq\be\leq\be_0}
 J_{N,\ell}(r\be)<\infty.
\end{equation}
\end{lemma}
\begin{proof}
For an integer $M\geq1$, Gaussian integration gives
\begin{equation}\label{eq:scalar-Wick-moments}
 \E_G[A_j^M]=Z_{M,\be M/N}^{(\eta)},\qquad
 \E_G[A_{12}^M]=Z_{M,2\be M/N}^{(\eta)}.
\end{equation}
The second identity uses that $(\mathsf Z_1+\mathsf Z_2)/\sqrt2$ has covariance $\g_\eta$.  For any real $v>0$, Lyapunov's inequality gives
\begin{equation}\label{eq:rounding}
 \E_G[A^v]\leq\bigl(\E_G[A^{\lceil v\rceil}]\bigr)^{v/\lceil v\rceil}.
\end{equation}
For $0<\theta=\ell/N<1$, use the conjugate exponents
\begin{equation}\label{eq:replica-holder}
 p_\theta=\frac{1+\theta}{1-\theta},\qquad
 q_\theta=\frac{1+\theta}{2\theta},\qquad
 \frac1{p_\theta}+\frac1{q_\theta}=1.
\end{equation}
H\"older's inequality yields
\begin{equation}\label{eq:holder-moment}
 \|P_{N,\ell}^{\eta}\|_r
 \leq\Bigl(\E_G[(A_1A_2)^{p_\theta r(N-\ell)}]\Bigr)^{1/(p_\theta r)}
 \Bigl(\E_G[A_{12}^{q_\theta r\ell}]\Bigr)^{1/(q_\theta r)}.
\end{equation}
The two powers to be rounded are exactly
\begin{equation}\label{eq:balanced-powers}
 p_\theta r(N-\ell)=r(N+\ell),\qquad
 q_\theta r\ell=\frac{r(N+\ell)}2.
\end{equation}
Consequently, the scalar couplings in \eqref{eq:scalar-Wick-moments} are at most
\begin{equation}\label{eq:effective-coupling}
 r\be(1+\theta)+\frac{\be}{N},\qquad
 r\be(1+\theta)+\frac{2\be}{N},
\end{equation}
respectively.  Both are below $b$.  Independence of $A_1,A_2$, \eqref{eq:rounding}, \eqref{eq:regularized-partition-domination}, and \eqref{eq:principal-order-one} bound \eqref{eq:holder-moment} by a constant depending only on the partition supremum at $b<\beta_{\gm}$.  The exponents $p_\theta,q_\theta$ themselves need not be bounded.  For $\ell=0$, independence directly gives the same conclusion using the rounded power $\lceil rN\rceil$.

To bound $J_{N,\ell}^{\eta}(r\be)$, repeat the argument with Gaussian moment order $1$ and coefficient $r\be$.  The ceiling error is now at most $2r\be/N$, which is covered by \eqref{eq:ceiling-buffer}.  Finally, for $\pi_{N,\ell}$-almost every configuration, no two distinct coordinates within either replica coincide, and $\g_\eta$ converges to $\g$ at every relevant difference.  Fatou's lemma proves the unregularized bound in \eqref{eq:replica-particle-buffer}.
\end{proof}

The balanced powers in \eqref{eq:balanced-powers} yield the effective coupling \eqref{eq:effective-coupling}. Thus the larger-coupling scalar bound controls the replica integral even though the exponents in H\"older's inequality vary with the overlap. Appendix~\ref{app:higher-order-replica-expansions} uses the fixed-integer extension of the same calculation.

\section{All-order marked expansions for prescribed references and homogeneous replica asymptotics}
\label{sec:prescribed-marked-expansions}

We prove the prescribed-reference expansions in \Cref{subsec:prescribed-marked-introduction}, assuming only the strict partition buffer as a many-particle input. The main step is a strong expansion of the marked density with a remainder uniform in the cutoff. An exact Cameron--Martin identity moves the marked fields into a centered source; a shared-label moment bound then controls the exponential before the spatial $L^p$ norm is taken. We derive the local-chaos and correlation consequences from this expansion, and conclude with the complementary homogeneous replica asymptotics at nonvanishing overlap.

Throughout this section, $\mathscr M$, $\be_0$, and $b$ satisfy \eqref{eq:prescribed-scalar-buffer}.  Constants may depend on the common upper density bound, $b$, $\be_0$, and the fixed indices displayed in a statement, but not on $N$, $\mu\in\mathscr M$, or $0\leq\be\leq\be_0$.

\subsection{Cutoffs preserving nonnegative definiteness and the spectral gap}
\label{subsec:prescribed-cutoffs}

Let $Q_j$ be conditional expectation in $L^2(\mu)$ onto an increasing sequence of finite measurable partitions generating the Borel sigma algebra, and set
\begin{equation}
 \tl\g_{\mu,j}\coloneqq(Q_j\otimes Q_j)\tl\g_{\mu},
 \qquad
 \psi_{\mu,j}\coloneqq Q_j\psi_{\mu},
 \qquad
 T_{\mu,j}\coloneqq Q_jT_{\mu}Q_j.
 \label{eq:prescribed-positive-cutoffs}
\end{equation}
On a cell of zero $\mu$-mass, choose the zero representative.  The regularized centered partition function is
\begin{equation}
 \mathsf P^j_{N,\be}(\mu)
 \coloneqq\int\exp\left\{
 \frac{\be}{N}\sum_{i<l}\tl\g_{\mu,j}(x_i,x_l)
 +\frac{\be}{N}\sum_i\psi_{\mu,j}(x_i)
 \right\}\,\dd\mu^{\otimes N}.
 \label{eq:prescribed-cutoff-partition}
\end{equation}

\begin{lemma}[Uniform row bounds and a forced spectral gap]
\label{lem:prescribed-row-spectral}
There are $a_0>0$ and $C<\infty$ such that
\begin{equation}
 \mathop{\rm ess\,sup}_{x}
 \int e^{a_0|\tl\g_{\mu}(x,y)|}\,\dd\mu(y)\leq C,
 \qquad
 \tl\g_{\mu}(x,y)\geq-C,
 \qquad
 \norm{\psi_{\mu}}_{\infty}\leq C.
 \label{eq:prescribed-row-exponential}
\end{equation}
The same bounds hold for the cutoffs in \eqref{eq:prescribed-positive-cutoffs}.  Moreover,
\begin{equation}
 1\leq\mathsf P^j_{N,\be}(\mu)
 \leq\mathsf P_{N,\be}(\mu)\leq C_b,
 \qquad 0\leq\be\leq b,
 \label{eq:prescribed-cutoff-partition-order}
\end{equation}
and
\begin{equation}
 b\norm{T_{\mu}}_{\op}\leq1-\delta,
 \qquad
 \delta\coloneqq e^{-1}C_b^{-2}>0.
 \label{eq:prescribed-uniform-spectral-margin}
\end{equation}
Consequently, the resolvents and Carleman--Fredholm determinants in \eqref{eq:prescribed-response-operators} are uniformly bounded on $0\leq\be\leq\be_0$.
\end{lemma}

\begin{proof}
The one-body potentials are bounded by the common upper density bound times $\norm{\g}_{L^1(\muunif)}$.  The logarithmic singularity gives $\int e^{a|\g|}\,\dd\muunif<\infty$ for every $0<a<\d$.  Translation invariance, the upper density bound, and the bounded centering terms prove \eqref{eq:prescribed-row-exponential}.  The lower bound follows from the lower boundedness of $\g$.  For $f\in L_0^2(\mu)$, the density $\rho f$ belongs to $L^2(\muunif)$, and Fourier positivity gives
\begin{equation}
 \ip*{f,T_{\mu}f}_{L^2(\mu)}
 =\sum_{n\neq0}\wh\g(n)|\wh{\rho f}(n)|^2\geq0.
 \label{eq:prescribed-Fourier-positivity}
\end{equation}
The row bound with exponent two gives the Hilbert--Schmidt property.  Conditional Jensen preserves all stated bounds for the cutoffs.  Sampling each point conditionally on its partition cell identifies the cutoff exponent as the conditional expectation of the original exponent.  Jensen's inequality gives the upper bound in \eqref{eq:prescribed-cutoff-partition-order}; the lower bound follows because the exponent has product expectation zero.

For the spectral estimate, let $\lambda_1\geq\lambda_2\geq\cdots\geq0$ be the eigenvalues of $T_{\mu}$.  A finite-rank truncation of the centered quadratic statistic converges by the multivariate central limit theorem and the law of large numbers to
\begin{equation}
 \frac b2\sum_{i=1}^L\lambda_i(G_i^2-1).
 \label{eq:prescribed-finite-chaos-limit}
\end{equation}
The Hilbert--Schmidt tail tends to zero in $L^2$, uniformly in $N$, by canonical orthogonality. Thus the full quadratic statistic converges in law to the infinite second-chaos series. The Portmanteau theorem and \eqref{eq:prescribed-cutoff-partition-order} bound its exponential moment by $C_b$. Its finite partial sums are independent of their centered tails, whose exponential expectations are at least one by Jensen's inequality. The same bound therefore holds for every finite partial sum. In particular,
\begin{equation}
 C_b\geq e^{-b\lambda_1/2}(1-b\lambda_1)^{-1/2}
 \geq e^{-1/2}(1-b\lambda_1)^{-1/2}.
 \label{eq:prescribed-top-mode-gap}
\end{equation}
This is \eqref{eq:prescribed-uniform-spectral-margin}.  The determinant bounds follow from the Hilbert--Schmidt norm bound and the common spectral gap.
\end{proof}

The cutoffs converge in every finite $L^p(\mu^{\otimes2})$, and the sources converge in every finite $L^p(\mu)$.  All cutoff removals below are performed only after establishing uniform integrability at a slightly amplified inverse temperature.

\subsection{The exact marked identity and its exponential-moment bound}
\label{subsec:prescribed-marked-identity}

At a fixed cutoff, let $\phi_j$ be a centered finite-dimensional Gaussian field with covariance
\begin{equation}
 c_j(x,y)\coloneqq\be\tl\g_{\mu,j}(x,y),
 \qquad
 s_j(x)\coloneqq\be\psi_{\mu,j}(x).
 \label{eq:prescribed-covariance-source}
\end{equation}
For a centered source $u$, define
\begin{equation}
 A_{j,u}(t)
 \coloneqq\int\exp\left\{t\phi_j(y)
 -\frac{t^2}{2}c_j(y,y)+t^2u(y)\right\}\,\dd\mu(y).
 \label{eq:prescribed-one-site-Gaussian}
\end{equation}
At $t=N^{-1/2}$, Gaussian integration gives
\begin{equation}
 \mathsf P^j_{N,\be}(\mu)
 =\E_G\left[A_{j,s_j}(N^{-1/2})^N\right].
 \label{eq:prescribed-scalar-Gaussian-identity}
\end{equation}
For a marked configuration $x=(x_1,\ldots,x_k)$, set
\begin{equation}
 u_{j,x}(y)\coloneqq s_j(y)+\sum_{a=1}^kc_j(y,x_a),
 \qquad
 J_{k,j}(x)\coloneqq\sum_{a=1}^ks_j(x_a)
 +\sum_{1\leq a<c\leq k}c_j(x_a,x_c).
 \label{eq:prescribed-marked-source}
\end{equation}
Here $s_j$ is the one-body correction, whereas $u_{j,x}$ also contains the interaction of an integrated particle with the fixed marked coordinates. Both enter the one-site exponent multiplied by $t^2=N^{-1}$; a fixed number of marked coordinates does not assert boundedness of the source after cutoff removal.

\begin{lemma}[Marked Cameron--Martin identity]
\label{lem:prescribed-marked-CM}
Let $F_{N,\be,\mu}^{(k),j}$ be the cutoff marginal density.  For $0\leq k\leq N$ and $t=N^{-1/2}$,
\begin{equation}
 \mathsf P^j_{N,\be}(\mu)F_{N,\be,\mu}^{(k),j}(x)
 =e^{t^2J_{k,j}(x)}
 \E_G\left[A_{j,u_{j,x}}(t)^{N-k}\right].
 \label{eq:prescribed-marked-CM}
\end{equation}
\end{lemma}

All mark--mark, mark--bulk, source, and diagonal terms in this identity have the original finite-particle denominator $N$.

\begin{proof}
Choose a finite feature representation $\phi_j(y)=v_j(y)\cdot Z$ on the range of the covariance.  In the Gaussian representation of the marked marginal numerator, the mark fields contribute $t\sum_av_j(x_a)\cdot Z$.  The Cameron--Martin formula
\begin{equation}
 \E_G[e^{tq\cdot Z}H(Z)]
 =e^{t^2|q|^2/2}\E_G[H(Z+tq)]
 \label{eq:prescribed-finite-CM}
\end{equation}
shifts each bulk field by $t\sum_ac_j(\cdot,x_a)$.  The self-diagonal pieces of $|q|^2$ cancel the diagonal Wick factors attached to the marks.  The remaining terms are precisely the mark--mark part of $J_{k,j}$, while the source at the marks gives its singleton part.  The shifted one-site integral is $A_{j,u_{j,x}}$.  This proves \eqref{eq:prescribed-marked-CM}.  A degenerate covariance causes no difficulty because the calculation is performed on its range.
\end{proof}

The next estimate separates spatial integrability from the partition bound at a larger inverse temperature. A fixed spatial $L^p$ norm can be controlled using the strict partition buffer above $\be_0$, without assuming a partition bound at $p\be$.

\begin{lemma}[Moment bound with fixed marked coordinates]
\label{lem:prescribed-anchored-moment}
Fix $k\geq0$ and $1\leq p<\infty$.  There is $r>1$, depending only on the strict margin between the two inverse temperatures, such that, for all sufficiently large $N$,
\begin{equation}
 \left\|\E_G\left[A_{j,u_{j,x}}(N^{-1/2})^{r(N-k)}\right]
 \right\|_{L^p(\mu^{\otimes k})}\leq C_{p,k}.
 \label{eq:prescribed-anchored-moment}
\end{equation}
The estimate is uniform in the cutoff and the reference family.
\end{lemma}

\begin{proof}
Choose $r>1$ and $b'<b$ with $r\be_0<b'$.  Put $n=N-k$, $M=\lceil rn\rceil$, $L=M+k$, and $\eta=\be L/N$.  Since $A^{rn}\leq1+A^M$ for $A\geq0$, it suffices to control the integer moment.  Completing the interactions among the $k$ marks gives the exact identity
\begin{equation}
 \E_G[A_{j,u_{j,x}}(N^{-1/2})^M]
 =e^{-N^{-1}J_{k,j}(x)}
 \mathsf P^j_{L,\eta}(\mu)F_{L,\eta,\mu}^{(k),j}(x).
 \label{eq:prescribed-anchored-completion}
\end{equation}
The row lower bound gives $J_{k,j}\geq-C_k$.  Choose an integer $q\geq p$.  The exact shared-label replica identity, followed by H\"older's inequality, bounds the $L^q$ norm of the last marginal whenever its largest effective scalar coupling lies below $b$.  Here that coupling is
\begin{equation}
 \eta\left(1+(q-1)\frac{k}{L}\right)+O_q(N^{-1})
 =\frac{\be}{N}(M+qk)+O_q(N^{-1})
 =r\be+O_{k,q}(N^{-1})<b.
 \label{eq:prescribed-anchored-effective-temperature}
\end{equation}
The scalar factor in \eqref{eq:prescribed-anchored-completion} is bounded by \eqref{eq:prescribed-cutoff-partition-order}.  This proves \eqref{eq:prescribed-anchored-moment}.  The spatial moment introduces only a fixed number of shared particles; it does not multiply the bulk inverse temperature by $p$.
\end{proof}

For later use, the same shared-label argument gives the following uniform marginal moment statement.  Fix an integer $q\geq2$ and $0<\alpha<1$ with
\begin{equation}
 \be_0[1+(q-1)\alpha]<b.
 \label{eq:prescribed-replica-window}
\end{equation}
Then, for sufficiently large $N$ and $0\leq k\leq\alpha N$,
\begin{equation}
 \sup_{j,\mu,\be}
 \int\left(F_{N,\be,\mu}^{(k),j}\right)^q\,\dd\mu^{\otimes k}<\infty.
 \label{eq:prescribed-replica-Lq}
\end{equation}
Indeed, H\"older's inequality in the exact $q$-replica Gaussian identity uses conjugate exponents
\begin{equation}
 a=\frac{1+(q-1)k/N}{1-k/N},
 \qquad
 a'=\frac{1+(q-1)k/N}{qk/N},
 \label{eq:prescribed-replica-holder-exponents}
\end{equation}
and after integer rounding both scalar couplings are bounded by
\begin{equation}
 r\be[1+(q-1)k/N]+q\be/N<b
 \label{eq:prescribed-replica-effective-temperature}
\end{equation}
for a fixed $r>1$ sufficiently close to one.  This is the precise reason that no scalar bound at $q\be$ is required.

\subsection{Strong coefficient expansion and the first profiles}
\label{subsec:prescribed-coefficients}

The cutoff-uniform Wick construction and exact exponential remainder are stated precisely in \Cref{prop:prescribed-marked-remainder}.  They give coefficient functions $U_{m,k}$ before scalar normalization.  For completeness, the finite recursion is as follows.  Let
\begin{equation}
 \log\left(1+\sum_{n\geq2}t^nA_{n,j}(u)\right)
 =\sum_{n\geq2}t^n\ell_{n,j}(u)
 \label{eq:prescribed-log-coefficients}
\end{equation}
be interpreted coefficientwise at every prescribed finite order, and let
\begin{equation}
 \mathcal B_n(L_1,\ldots,L_n)
 \coloneqq
 \sum_{\substack{a_1,\ldots,a_n\geq0\\\sum_{l=1}^nla_l=n}}
 \prod_{l=1}^n\frac{L_l^{a_l}}{a_l!},
 \qquad \mathcal B_0=1.
 \label{eq:prescribed-Bell-polynomial}
\end{equation}
Define
\begin{equation}
 V_{n,k,j}(x)
 \coloneqq\ell_{n+2,j}(u_{j,x})-k\ell_{n,j}(u_{j,x})
 +\one_{\{n=2\}}J_{k,j}(x),
 \qquad \ell_{0,j}=\ell_{1,j}=0,
 \label{eq:prescribed-marked-V}
\end{equation}
and
\begin{equation}
 U_{m,k}(x)
 \coloneqq D_{\be,\mu}^{-1}
 \lim_{j\to\infty}\E_G\left[
 e^{\ell_{2,j}}\mathcal B_{2m}(V_{1,k,j},\ldots,V_{2m,k,j})
 \right].
 \label{eq:prescribed-U-coefficients}
\end{equation}
The limit is taken in every finite spatial $L^p$.  Then
\begin{equation}
 f_{0,k}=1,
 \qquad
 z_m=U_{m,0},
 \qquad
 f_{m,k}=U_{m,k}-\sum_{i=1}^mz_if_{m-i,k}.
 \label{eq:prescribed-density-recursion}
\end{equation}

\begin{proof}[Proof of \Cref{thm:prescribed-marked-expansion}]
At every cutoff, \Cref{prop:prescribed-marked-remainder} expands the exact right side of \eqref{eq:prescribed-marked-CM}, after Gaussian expectation, with an error $O_{k,M,p}(N^{-M-1})$ in $L^p(\mu^{\otimes k})$.  The proof first expands on the event $|A_{j,u_{j,x}}(N^{-1/2})-1|\leq1/2$ and uses \Cref{lem:prescribed-anchored-moment} to remove its complement by conditional H\"older.  Thus, the spatial $L^p$ norm is taken only after the Gaussian exact exponential has been controlled.  Simultaneous reversal of the Gaussian field kills every odd power of $N^{-1/2}$.  The external factor $e^{J_{k,j}/N}$ has all fixed spatial moments by \eqref{eq:prescribed-row-exponential} and is expanded separately.

The coefficient variables converge in every finite joint $L^p$ under cutoff removal.  The cutoff-free marginal itself converges in the required $L^p$ by \eqref{eq:prescribed-replica-Lq} at a larger fixed integer moment and interpolation.  Hence the remainder passes to the singular limit.  Scalar finite-series division is legitimate because $\mathsf P_{N,\be}\geq1$ and $D_{\be,\mu}$ is uniformly bounded above and below.  It produces precisely \eqref{eq:prescribed-density-recursion}.  Uniqueness of asymptotic coefficients, exact normalization of the marginal, and exact consistency under marginalization show successively that each $f_{m,k}$ has integral zero and that the family is consistent.  This proves \eqref{eq:prescribed-partition-expansion} and \eqref{eq:prescribed-marginal-expansion}.  It remains to identify the first coefficients.

Keep a finite cutoff until all diagonal cancellations have been made.  Write $C_j=\be T_{\mu,j}$, $R_j=(\Id-C_j)^{-1}$, and $B_j=R_j-\Id$.  If $c_j$ is the kernel of $C_j$ and $b_j$ the kernel of $B_j$, put
\begin{equation}
 \omega_j(x)\coloneqq b_j(x,x)-c_j(x,x)
 =\ip*{c_j(x,\cdot),R_jc_j(x,\cdot)}_{L^2(\mu)}.
 \label{eq:prescribed-finite-dressed-diagonal}
\end{equation}
The right side is uniformly bounded and converges in every finite $L^p$ to \eqref{eq:prescribed-omega}. The difference of diagonal values is therefore taken at finite cutoff, before passing to the limit.  The kernel identity $b_j=c_j+C_jR_jC_j$ also shows that the correction $b_j-c_j$ is uniformly bounded by Cauchy--Schwarz and the common spectral gap.  In particular, every contraction below is finite and has a cutoff limit.  Set
\begin{equation}
 s_j\coloneqq\be\psi_{\mu,j},
 \qquad
 d_j\coloneqq s_j+\frac12\omega_j^\circ.
 \label{eq:prescribed-finite-source-dressed-diagonal}
\end{equation}
Abbreviate the first integrated Wick coefficients by
\begin{align}
 A_{2,j}
 &\coloneqq\frac12\int:\!\phi_j^2\!:\!_{c_j}\,\dd\mu,
 \label{eq:prescribed-finite-A2}\\
 A_{3,j}(u)
 &\coloneqq\frac16\int:\!\phi_j^3\!:\!_{c_j}\,\dd\mu
 +\int u\phi_j\,\dd\mu,
 \label{eq:prescribed-finite-A3}\\
 A_{4,j}(u)
 &\coloneqq\frac1{24}\int:\!\phi_j^4\!:\!_{c_j}\,\dd\mu
 +\frac12\int u:\!\phi_j^2\!:\!_{c_j}\,\dd\mu
 +\frac12\int u^2\,\dd\mu.
 \label{eq:prescribed-finite-A4}
\end{align}
Let $\E_{B_j}$ denote expectation under the normalized quadratic tilt, whose centered covariance is $B_j$.  The order-$N^{-1}$ scalar coefficient with centered deterministic source $u$ is
\begin{equation}
 z_{1,j}(u)
 =\E_{B_j}\left[A_{4,j}(u)-\frac12A_{2,j}^2
 +\frac12A_{3,j}(u)^2\right].
 \label{eq:prescribed-finite-z1-expectation}
\end{equation}
Wick contraction gives
\begin{align}
 \E_{B_j}[A_{4,j}(u)]
 &=\frac18\int\omega_j^2\,\dd\mu
 +\frac12\int u\omega_j\,\dd\mu
 +\frac12\int u^2\,\dd\mu,
 \label{eq:prescribed-finite-A4-expectation}\\
 \E_{B_j}[A_{2,j}^2]
 &=\frac14\left[
 \left(\int\omega_j\,\dd\mu\right)^2
 +2\Tr(B_j^2)\right].
 \label{eq:prescribed-finite-A2-expectation}
\end{align}
For the cubic term, the change of Wick ordering is
\begin{equation}
 :\!\phi_j(x)^3\!:\!_{c_j}
 =:\!\phi_j(x)^3\!:\!_{b_j}+3\omega_j(x)\phi_j(x).
 \label{eq:prescribed-finite-change-Wick}
\end{equation}
The field has zero spatial average, so constants in $u+\omega_j/2$ do not contribute.  Orthogonality of the first and third Gaussian chaoses therefore yields
\begin{equation}
 \E_{B_j}[A_{3,j}(u)^2]
 =\frac16\iint b_j(x,y)^3\,\dd\mu(x)\dd\mu(y)
 +\ip*{u+\frac12\omega_j^\circ,
 B_j\left(u+\frac12\omega_j^\circ\right)}_{L^2(\mu)}.
 \label{eq:prescribed-finite-A3-expectation}
\end{equation}
Combining these three identities and using $\Id+B_j=R_j$ gives
\begin{equation}
 z_{1,j}(u)
 =\frac12\ip*{u+\frac12\omega_j^\circ,
 R_j\left(u+\frac12\omega_j^\circ\right)}_{L^2(\mu)}
 -\frac14\Tr(B_j^2)
 +\frac1{12}\iint b_j(x,y)^3\,\dd\mu(x)\dd\mu(y).
 \label{eq:prescribed-finite-z1-source}
\end{equation}
Taking $u=s_j$ and passing to the cutoff limit proves
\begin{equation}
 z_1
 =\frac12\ip*{d_{\be,\mu},R_{\be,\mu}d_{\be,\mu}}_{L^2(\mu)}
 -\frac14\Tr(B_{\be,\mu}^2)
 +\frac1{12}\iint b_{\be,\mu}(x,y)^3\,\dd\mu(x)\dd\mu(y).
 \label{eq:prescribed-z1-proof}
\end{equation}
Every term has a finite cutoff limit.  The first contraction is an $L^2(\mu)$ inner product, $\Tr(B_{\be,\mu}^2)$ is the squared Hilbert--Schmidt norm, and the cubic term is the pointwise spatial integral; no infinite original diagonal remains.

For the marked coefficient, the external term and the quadratic completion combine as
\begin{equation}
 s_j(x)+\frac12\omega_j^\circ(x)+(C_jR_jd_j)(x)
 =d_j(x)+(B_jd_j)(x)=R_jd_j(x),
 \label{eq:prescribed-singleton-calculation}
\end{equation}
while for two distinct marks $c_j+C_jR_jC_j$ is the kernel of $B_j$.  Passing to the limit yields exactly \eqref{eq:prescribed-first-coefficient} and \eqref{eq:prescribed-first-partition-correction}.
\end{proof}

Three specializations clarify the coefficient formula.  At $\be=0$, all positive-order coefficients vanish.  At $\mu=\muunif$, the source is zero and the dressed diagonal is constant, so $a_{\be,\muunif}=0$.  At small $\be$,
\begin{equation}
 a_{\be,\mu}=\be\psi_{\mu}+O(\be^2),
 \qquad
 b_{\be,\mu}=\be\tl\g_{\mu}+O(\be^2),
 \label{eq:prescribed-small-beta-check}
\end{equation}
in agreement with direct perturbation of \eqref{eq:prescribed-centered-partition-exponent}.

\begin{proof}[Proof of \Cref{cor:prescribed-near-uniform}]
Choose \eqref{eq:prescribed-three-temperature-buffer}.  The homogeneous estimate below the global-minimization threshold \eqref{eq:principal-order-one} supplies the scalar bound up to $B$.  The Dacorogna--Moser map in \eqref{eq:dacorogna-moser-stability} pulls the reference back to $\muunif$.  By \eqref{eq:near-identity-remainder-bound}, its centered kernel remainder tends to zero in $L^\infty$ with $\norm{\rho-1}_{C^\gamma}$.  Choose $\vep$ so that H\"older's inequality in \eqref{eq:transport-holder-splitting} absorbs this remainder while retaining the strict gap from $b$ to $B$.  This gives \eqref{eq:prescribed-scalar-buffer} uniformly on \eqref{eq:prescribed-near-uniform-class}.  The choice of $\vep$ involves only $\d,\gamma,\be_0,b,B$ and therefore precedes every fixed $M,k,p$.  Apply \Cref{thm:prescribed-marked-expansion}.
\end{proof}

\subsection{Divergences, total variation, and growing marginals}
\label{subsec:prescribed-divergences}

\begin{proof}[Proof of \Cref{cor:prescribed-divergence-expansions}]
Put $u_N=F_{N,\be,\mu}^{(k)}-1$.  The strong expansion in every finite $L^p$ permits multiplication of any fixed number of coefficient series, with an $L^1$ remainder of the expected inverse-particle order.  Since $\int u_N\,\dd\mu^{\otimes k}=0$,
\begin{equation}
 \chi^2=\int u_N^2\,\dd\mu^{\otimes k},
 \qquad
 \Hr=\int\bigl[(1+u_N)\log(1+u_N)-u_N\bigr]\,\dd\mu^{\otimes k}.
 \label{eq:prescribed-divergence-identities}
\end{equation}
Expanding the first expression gives \eqref{eq:prescribed-chi-coefficients}.  For the second, use the global finite Taylor remainder
\begin{equation}
 \left|(1+u)\log(1+u)-u
 -\sum_{r=2}^{L}\frac{(-1)^ru^r}{r(r-1)}\right|
 \leq C_L\bigl(|u|^{L+1}+|u|^{L+2}\bigr),
 \qquad u\geq-1,
 \label{eq:prescribed-entropy-global-remainder}
\end{equation}
where on a compact neighborhood of $-1$ the left side is continuous and away from it Taylor's theorem applies, while for large $u$ the right side dominates $u\log(1+u)$.  The bound $\norm{u_N}_{p}=O(N^{-1})$ for every fixed finite $p$ makes the integrated remainder $O(N^{-L-1})$.  Collecting the products gives \eqref{eq:prescribed-entropy-coefficients}.

The first profile in \eqref{eq:prescribed-leading-profile} is an orthogonal sum of a first Hoeffding level and a second Hoeffding level.  Hence
\begin{equation}
 \norm{G_{k,\be,\mu}}_2^2
 =k\norm{a_{\be,\mu}}_2^2
 +\binom{k}{2}\norm{b_{\be,\mu}}_2^2,
 \label{eq:prescribed-profile-orthogonality}
\end{equation}
and $\norm{b_{\be,\mu}}_2=\norm{B_{\be,\mu}}_{\HS}$.  This proves \eqref{eq:prescribed-leading-divergences}.  The reverse triangle inequality in $L^1$, applied to \eqref{eq:prescribed-marginal-expansion} with $M=1$, proves \eqref{eq:prescribed-leading-TV}.
\end{proof}

The zero set of the leading profile contributes to the second total-variation coefficient. The strong signed expansion gives  For fixed $\be,\mu,k$,
\begin{align}
 \TV\!\left(\Q_{N,\be}^{(k)}(\mu),\mu^{\otimes k}\right)
 &=\frac1{2N}\int|G_{k,\be,\mu}|\,\dd\mu^{\otimes k}
 \notag\\
 &\quad+\frac1{2N^2}\left[
 \int_{\{G_{k,\be,\mu}\neq0\}}\sgn(G_{k,\be,\mu})f_{2,k}\,\dd\mu^{\otimes k}
 +\int_{\{G_{k,\be,\mu}=0\}}|f_{2,k}|\,\dd\mu^{\otimes k}
 \right]
 +o(N^{-2}).
 \label{eq:prescribed-second-TV-coefficient}
\end{align}
Indeed, the one-sided difference quotient of $|G+\varepsilon f_2|$ is bounded by $|f_2|$ and converges to the displayed integrand; dominated convergence and the $O(N^{-3})$ signed $L^1$ remainder complete the argument.  No arbitrary-order expansion of the absolute value is asserted.

The growing-marginal result is stated in \Cref{thm:prescribed-growing-marginals}. We now prove its profile and divergence estimates.

\begin{proof}[Proof of \eqref{eq:prescribed-growing-profile}--\eqref{eq:prescribed-growing-TV}]
Let $L=\lfloor\alpha N\rfloor$.  Equation \eqref{eq:prescribed-replica-Lq} with $q=2$ gives $\norm{F_{N,\be,\mu}^{(L)}}_2\leq C$.  Apply \Cref{lem:local-Hoeffding} to $F_{N,\be,\mu}^{(L)}$, and write $h_r$ for its common canonical kernel at level $r$. The expansion and isometry are \eqref{eq:prescribed-growing-Hoeffding-expansion}--\eqref{eq:local-Hoeffding-orthogonality}. The fixed-marginal marked expansion at $k=1,2$ gives
\begin{equation}
 \norm{h_1-a_{\be,\mu}/N}_2\leq C N^{-2},
 \qquad
 \norm{h_2-b_{\be,\mu}/N}_2\leq C N^{-2}.
 \label{eq:prescribed-growing-first-Hoeffding-rates}
\end{equation}
For $r\geq3$,
\begin{equation}
 \frac{\binom kr}{\binom Lr}\leq\left(\frac{k}{L}\right)^r
 \leq\left(\frac{k}{L}\right)^3.
 \label{eq:prescribed-growing-binomial-comparison}
\end{equation}
Using \eqref{eq:local-Hoeffding-orthogonality} at $L$ and the uniform $L^2$ bound gives
\begin{equation}
 \sum_{r=3}^k\binom kr\norm{h_r}_2^2
 \leq\left(\frac{k}{L}\right)^3\norm{F^{(L)}-1}_2^2
 \leq C\left(\frac{k}{N}\right)^3.
 \label{eq:prescribed-growing-higher-Hoeffding}
\end{equation}
Orthogonality applied to the difference between the exact Hoeffding expansion and $G_{k,\be,\mu}/N$ proves \eqref{eq:prescribed-growing-profile}.  Moreover, the squared norms of $h_1$ and $h_2$ differ from $N^{-2}\norm{a_{\be,\mu}}_2^2$ and $N^{-2}\norm{b_{\be,\mu}}_2^2$ by $O(N^{-3})$.  Substitution into \eqref{eq:local-Hoeffding-orthogonality}, together with \eqref{eq:prescribed-growing-higher-Hoeffding}, proves \eqref{eq:prescribed-growing-chi}.

For entropy, apply \eqref{eq:entropy-split} with $u=F_{N,\be,\mu}^{(k)}-1$, $w=G_{k,\be,\mu}/N$, and $r=u-w$.  The leading profile satisfies $\norm{G_{k,\be,\mu}}_4\leq Ck$.  For the pair sum, expand its fourth moment.  Canonicality eliminates every term whose graph has a vertex of degree one; every surviving four-edge graph therefore has at most four vertices, so there are at most $Ck^4$ patterns, each controlled by the uniform $L^4$ kernel bound.  The centered singleton sum has fourth moment $O(k^2)$ by independence and boundedness.  Equation \eqref{eq:prescribed-growing-profile} consequently gives
\begin{equation}
 \left|\Hr-\frac12\chi^2\right|
 \leq C\left(\frac{k^3}{N^3}+\frac{k^2}{N^4}\right).
 \label{eq:prescribed-growing-entropy-remainder}
\end{equation}
Combining this with \eqref{eq:prescribed-growing-chi} proves \eqref{eq:prescribed-growing-entropy}.  The reverse triangle inequality, $L^1\leq L^2$, and \eqref{eq:prescribed-growing-profile} prove \eqref{eq:prescribed-growing-TV}.
\end{proof}

For the pair part of the profile, set
\begin{equation}
 S_k\coloneqq\sum_{1\leq i<l\leq k}b_{\be,\mu}(X_i,X_l),
 \qquad X_1,\ldots,X_k\sim\mu^{\otimes k}.
 \label{eq:prescribed-pair-profile-sum}
\end{equation}
The required rank-independent estimate is
\begin{equation}
 \left|\frac1k\E[|S_k|]-\frac12\E[|\mathcal Q_{\be,\mu}|]\right|
 \leq Ck^{-1/6}.
 \label{eq:prescribed-chaos-absolute}
\end{equation}

\begin{proof}[Proof of \eqref{eq:prescribed-chaos-absolute}]
Approximate the nonnegative canonical Hilbert--Schmidt operator $B_{\be,\mu}$ by $B^{(j)}\coloneqq Q_jB_{\be,\mu}Q_j$.  Conditional Jensen preserves the row-moment bounds, and $B^{(j)}\to B_{\be,\mu}$ in Hilbert--Schmidt norm.  Write the kernel of $B^{(j)}$ as
\begin{equation}
 b^{(j)}(x,y)=\zeta(x)\cdot\zeta(y),
 \qquad
 \Lambda\coloneqq\E_\mu[\zeta\zeta^{\mathsf T}],
 \label{eq:prescribed-finite-rank-feature}
\end{equation}
where $\zeta$ is centered.  The following quantities are bounded independently of the rank:
\begin{equation}
 \Tr(\Lambda^2),
 \qquad
 \int\!\left(\zeta(x)^{\mathsf T}\Lambda\zeta(x)\right)^2\,\dd\mu(x),
 \qquad
 \iint|\zeta(x)\cdot\zeta(y)|^4\,\dd\mu(x)\dd\mu(y).
 \label{eq:prescribed-feature-moment-bounds}
\end{equation}
The middle term can be written explicitly as
\begin{equation}
 \int\!\left(\int |b^{(j)}(x,y)|^2\,\dd\mu(y)\right)^2\,\dd\mu(x).
 \label{eq:prescribed-integrated-row-norm}
\end{equation}
The other two quantities are the squared Hilbert--Schmidt norm and the fourth kernel moment.  The uniform row bounds for $b_{\be,\mu}=\be\tl\g_\mu+(B_{\be,\mu}-\be T_\mu)$ follow from \eqref{eq:prescribed-row-exponential}, the common spectral gap, and Cauchy--Schwarz for the bounded correction kernel.

Replace the independent feature vectors one at a time by independent centered Gaussian vectors with covariance $\Lambda$.  The statistic divided by $k$ is affine in each vector, with increment
\begin{equation}
 \Delta_i
 =\frac1kV_i\cdot\sum_{l\neq i}V_l.
 \label{eq:prescribed-Lindeberg-increment}
\end{equation}
For every intermediate mixture of original and Gaussian feature laws,
\begin{equation}
 \E[|\Delta_i|^4]\leq Ck^{-2},
 \qquad
 \E[|\Delta_i|^3]\leq Ck^{-3/2}.
 \label{eq:prescribed-Lindeberg-increment-moments}
\end{equation}
To verify the fourth-moment estimate, condition on $V_i$ and expand the fourth moment of the sum of independent centered scalar variables.  The individual fourth moments and the pairwise products of variances are controlled by \eqref{eq:prescribed-feature-moment-bounds}.  If one or both vectors are Gaussian, the corresponding terms are bounded by $3\E[(\zeta^{\mathsf T}\Lambda\zeta)^2]$ or by a constant times $(\Tr\Lambda^2)^2$; the integrated square of the conditional variance is bounded by the middle quantity in \eqref{eq:prescribed-feature-moment-bounds}, or by $(\Tr\Lambda^2)^2+2\Tr\Lambda^4$ in the Gaussian case.  There are $O(k^2)$ terms before division by $k^4$.  The third-moment estimate follows by H\"older's inequality.

For every $C^3$ function $\vartheta$, Taylor expansion about the statistic with the $i$th vector removed has matching first and second conditional moments under the original and Gaussian laws.  Summing the third-order remainders gives
\begin{equation}
 \left|\E\left[\vartheta(S_k^{(j)}/k)\right]
 -\E\left[\vartheta(S_k^{G,(j)}/k)\right]\right|
 \leq Ck^{-1/2}\norm{\vartheta'''}_\infty.
 \label{eq:prescribed-smooth-replacement}
\end{equation}
For Gaussian features,
\begin{equation}
 \frac{S_k^{G,(j)}}k
 =\frac12\left(\left|k^{-1/2}\sum_iV_i\right|^2-\Tr\Lambda\right)
 -\frac1{2k}\sum_i\left(|V_i|^2-\Tr\Lambda\right).
 \label{eq:prescribed-Gaussian-diagonal}
\end{equation}
The first term has the law of half the finite-rank second chaos.  The second has $L^2$ norm $\sqrt{\Tr(\Lambda^2)/(2k)}$ and is centered at every finite rank.

Take $\vartheta_\delta(x)=\sqrt{x^2+\delta^2}$.  Its uniform distance from $|x|$ is at most $\delta$ and $\norm{\vartheta_\delta'''}_\infty\leq C\delta^{-2}$.  Equations \eqref{eq:prescribed-smooth-replacement} and \eqref{eq:prescribed-Gaussian-diagonal} give an error bounded by
\begin{equation}
 C\left(\delta+k^{-1/2}\delta^{-2}+k^{-1/2}\right).
 \label{eq:prescribed-smoothed-absolute-error}
\end{equation}
Choosing $\delta=k^{-1/6}$ proves the estimate uniformly in the finite rank.  Finally, the canonical isometry gives $S_k^{(j)}/k\to S_k/k$ in $L^2$, while the Gaussian second-chaos isometry gives variance $2\norm{B^{(j)}-B_{\be,\mu}}_{\HS}^2$ for the difference of the chaos variables.  Passing to the limit proves \eqref{eq:prescribed-chaos-absolute}.  No trace of the infinite-rank operator $B_{\be,\mu}$ is used.
\end{proof}

\begin{proof}[Completion of \Cref{thm:prescribed-growing-marginals}]
The centered singleton sum obeys
\begin{equation}
 \frac1k\norm{\sum_{i=1}^ka_{\be,\mu}(X_i)}_1
 \leq\frac{\norm{a_{\be,\mu}}_2}{\sqrt{k}}.
 \label{eq:prescribed-singleton-absolute-bound}
\end{equation}
Combine this estimate with \eqref{eq:prescribed-chaos-absolute} and \eqref{eq:prescribed-growing-TV} to obtain \eqref{eq:prescribed-growing-TV-chaos}.  The chi-square and entropy limits follow from \eqref{eq:prescribed-growing-chi}--\eqref{eq:prescribed-growing-entropy}, because their singleton contributions and the pair-count correction vanish as $k\to\infty$.  The real spectral series in \eqref{eq:prescribed-second-chaos} has the stated $L^2$ variance by independence and square summability.
\end{proof}

\subsection{Finite mark degree and explicit dependence on the marginal size}
\label{subsec:prescribed-mark-degree}

For $A\subseteq[k]$, write $x_A=(x_i)_{i\in A}$.

\begin{prop}[Finite mark degree]
\label{prop:prescribed-mark-degree}
For every $m\geq1$, there are symmetric kernels
\begin{equation}
 h_{m,r}\in\bigcap_{1\leq p<\infty}L^p(\mu^{\otimes r}),
 \qquad 1\leq r\leq2m,
 \label{eq:prescribed-hmr-domain}
\end{equation}
canonical in every variable, such that
\begin{equation}
 f_{m,k}(x_{[k]})
 =\sum_{r=1}^{\min\{2m,k\}}
 \sum_{\substack{A\subseteq[k]\\|A|=r}}h_{m,r}(x_A).
 \label{eq:prescribed-coefficient-Hoeffding}
\end{equation}
They are given by
\begin{equation}
 h_{m,r}(x_{[r]})
 =\sum_{A\subseteq[r]}(-1)^{r-|A|}f_{m,|A|}(x_A),
 \qquad f_{m,0}=0,
 \label{eq:prescribed-hmr-Mobius}
\end{equation}
and satisfy, for fixed $m,p$,
\begin{equation}
 \norm{f_{m,k}}_{L^p(\mu^{\otimes k})}
 \leq C_{m,p}(1+k)^m.
 \label{eq:prescribed-coefficient-k-bound}
\end{equation}
In particular, $h_{1,1}=a_{\be,\mu}$, $h_{1,2}=b_{\be,\mu}$, and $h_{1,r}=0$ for $r\geq3$.
\end{prop}

The alternating sum in \eqref{eq:prescribed-hmr-Mobius} is inclusion--exclusion over coordinate subsets. It extracts a canonical Hoeffding component; the cumulant formula \eqref{eq:prescribed-connected-density} instead sums over partitions of the coordinate set.

\begin{proof}
In \eqref{eq:prescribed-marked-V}, the source degree of $\ell_n(u_{j,x})$ is at most $\lfloor n/2\rfloor$, and $\ell_2$ is independent of the marks.  Every monomial in $V_{n,k,j}$ therefore involves at most $n$ distinct marks.  A Bell monomial of total weight $2m$ involves at most $2m$ marks.  Gaussian expectation, cutoff removal, and scalar series division do not increase this number.  Consistency of the coefficients identifies their fixed Hoeffding projections with the alternating sums in \eqref{eq:prescribed-hmr-Mobius}; one-coordinate integration proves canonicality and \eqref{eq:prescribed-coefficient-Hoeffding}.

To estimate the dependence on $k$,  Fix $r$ and an even integer $q\geq\max\{p,2\}$.  Expanding the $q$th moment of
\begin{equation}
 \sum_{\substack{A\subseteq[k]\\|A|=r}}h_{m,r}(X_A)
 \label{eq:prescribed-canonical-sum}
\end{equation}
produces zero whenever some coordinate occurs in exactly one selected subset.  Thus, every surviving coordinate appears at least twice, the union of all selected subsets has cardinality at most $rq/2$, and there are at most $C_{r,q}k^{rq/2}$ surviving patterns.  Generalized H\"older bounds every pattern uniformly.  Taking the $q$th root gives a bound $Ck^{r/2}$, and summing over $r\leq2m$ proves \eqref{eq:prescribed-coefficient-k-bound}.
\end{proof}

\begin{prop}[Conservative explicit remainder]
\label{prop:prescribed-k-remainder}
For every fixed $M\geq0$ and $1\leq p<\infty$, there are finite constants $C_{M,p},c_{M,p}$ such that, whenever
\begin{equation}
 N\geq c_{M,p}(1+k),
 \label{eq:prescribed-k-remainder-range}
\end{equation}
one has
\begin{equation}
 \left\|F_{N,\be,\mu}^{(k)}-1
 -\sum_{m=1}^{M}\frac{f_{m,k}}{N^m}\right\|_{L^p(\mu^{\otimes k})}
 \leq C_{M,p}\frac{(1+k)^{2M+2}}{N^{M+1}}.
 \label{eq:prescribed-explicit-k-remainder}
\end{equation}
\end{prop}

The exponent $2M+2$ follows by counting occurrences of marked variables in the remainder. No optimality of this exponent is asserted.

\begin{proof}
Write $\kappa=1+k$, $t=N^{-1/2}$, and $y=\kappa t$.  On $y\geq\epsilon$, the marginal moment bound \eqref{eq:prescribed-replica-Lq} and \eqref{eq:prescribed-coefficient-k-bound} make the left side uniformly bounded, while the right side is at least a fixed positive constant after enlarging $C_{M,p}$.  On $y\leq\epsilon$, normalize the marked source by $\kappa$.  Convexity of the exponential row bound in \eqref{eq:prescribed-row-exponential} gives cutoff-uniform weighted-Wick estimates with source degree $l$ bounded by $C^{a+2l}\kappa^l$.  Hence
\begin{equation}
 \norm{\ell_n(u_{j,x})}_{L^v(G)}\leq C_{n,v}\kappa^{n-2},
 \qquad n\geq2,
 \label{eq:prescribed-k-log-coefficients}
\end{equation}
and all one-site series tails are geometric in $\kappa t$.

With $K=2M+1$, the good-event logarithmic remainder is $O(y^{K+1})$.  In the exact moment identity \eqref{eq:prescribed-anchored-completion}, $J_{k,j}\geq-C\kappa^2$, so the prefactor is at most $e^{Cy^2}$; integer rounding changes the effective inverse temperature only by $O_{M,p}(k/N)$.  The moment bound with fixed marked coordinates therefore controls the exact exponential at every fixed higher spatial exponent required by conditional H\"older.  The bad event has arbitrarily high order in $t$, hence also order $y^{K+1}$ because $t\leq y$.  Finally, $\norm{J_{k,j}}_q\leq C_q\kappa^2$.  More importantly, H\"older's inequality over the $\binom{k}{2}$ pair factors and the $k$ bounded source factors, together with the logarithmic row-exponential bound \eqref{eq:prescribed-row-exponential}, gives
\begin{equation}
 \norm{e^{t^2|J_{k,j}|}}_{L^q(\mu^{\otimes k})}\leq C_q,
 \qquad y\leq\epsilon,
 \label{eq:prescribed-external-factor-exponential-moment}
\end{equation}
provided $\epsilon$ is chosen small for the finitely many exponents in use.  Indeed, after the H\"older split, the exponent attached to each pair is at most $C_qt^2\kappa^2\leq C_q\epsilon^2$.  The Taylor remainder
\begin{equation}
 \left|e^x-\sum_{r=0}^M\frac{x^r}{r!}\right|
 \leq C_M|x|^{M+1}e^{|x|}
 \label{eq:prescribed-external-factor-Taylor-remainder}
\end{equation}
with $x=t^2J_{k,j}$, followed by H\"older's inequality, is therefore bounded by $C(t^2\kappa^2)^{M+1}=Cy^{2M+2}$ in every required spatial norm.  Multiplying the two expansions and dividing by the scalar series gives \eqref{eq:prescribed-explicit-k-remainder}.
\end{proof}

Combining \eqref{eq:prescribed-coefficient-k-bound} and \eqref{eq:prescribed-explicit-k-remainder} with the product formulas in \eqref{eq:prescribed-chi-coefficients}--\eqref{eq:prescribed-entropy-coefficients} yields
\begin{equation}
 |\mathcal C_{n,k}|+|\mathcal H_{n,k}|
 \leq C_n(1+k)^n
 \label{eq:prescribed-divergence-coefficient-k-bound}
\end{equation}
and, for every fixed $L$, a remainder bounded by
\begin{equation}
 C_L\frac{(1+k)^{4L+4}}{N^{L+1}}
 \label{eq:prescribed-divergence-k-remainder}
\end{equation}
throughout $N\geq c_L(1+k)$.  In particular,
\begin{equation}
 \left|\TV-\frac1{2N}\norm{f_{1,k}}_1\right|
 \leq C\left(\frac{(1+k)^2}{N^2}
 +\frac{(1+k)^6}{N^3}\right).
 \label{eq:prescribed-TV-explicit-error}
\end{equation}

\subsection{Hoeffding and connected-correlation expansions}
\label{subsec:prescribed-hoeffding-clusters}

For $r\geq1$, define the exact $r$th Hoeffding kernel of the likelihood ratio by
\begin{equation}
 H_{N,r}(x_{[r]})
 \coloneqq\sum_{A\subseteq[r]}(-1)^{r-|A|}
 F_{N,\be,\mu}^{(|A|)}(x_A),
 \qquad F_{N,\be,\mu}^{(0)}\coloneqq1.
 \label{eq:prescribed-exact-Hoeffding-kernel}
\end{equation}

\begin{prop}[All fixed Hoeffding projections]
\label{prop:prescribed-Hoeffding-expansion}
For every fixed $r,M,p$,
\begin{equation}
 H_{N,r}
 =\sum_{m=\lceil r/2\rceil}^{M}\frac{h_{m,r}}{N^m}
 +O_{L^p(\mu^{\otimes r})}(N^{-M-1}).
 \label{eq:prescribed-Hoeffding-expansion}
\end{equation}
Consequently,
\begin{equation}
 NH_{N,1}\longrightarrow a_{\be,\mu},
 \qquad
 NH_{N,2}\longrightarrow b_{\be,\mu},
 \qquad
 NH_{N,r}\longrightarrow0\quad(r\geq3)
 \label{eq:prescribed-first-Hoeffding-limits}
\end{equation}
in every fixed finite $L^p$.
\end{prop}

\begin{proof}
Insert \eqref{eq:prescribed-marginal-expansion} into the finite alternating sum \eqref{eq:prescribed-exact-Hoeffding-kernel}.  The subset-lattice inversion \eqref{eq:prescribed-hmr-Mobius} extracts exactly $h_{m,r}$.  By \Cref{prop:prescribed-mark-degree}, it vanishes for $r>2m$, which gives the lower limit $m\geq\lceil r/2\rceil$.  The remainder is a finite sum of fixed-$r$ strong remainders.
\end{proof}

Define the connected correlation density of the labeled marginals by the cumulant formula
\begin{equation}
 U_{N,r}(x_{[r]})
 \coloneqq\sum_{\pi\in\Part_r}
 (-1)^{|\pi|-1}(|\pi|-1)!
 \prod_{B\in\pi}F_{N,\be,\mu}^{(|B|)}(x_B),
 \label{eq:prescribed-connected-density}
\end{equation}
where $\Part_r$ denotes the set of partitions of $[r]$. This is a cumulant combination relative to $\mu^{\otimes r}$, not a probability density or a factorial point-process correlation.

\begin{prop}[Normalized connected correlations]
\label{prop:prescribed-cluster-expansion}
For every fixed $r\geq2$, $M\geq r-1$, and finite $p$,
\begin{equation}
 U_{N,r}
 =\sum_{m=r-1}^{M}\frac{u_{m,r}}{N^m}
 +O_{L^p(\mu^{\otimes r})}(N^{-M-1}).
 \label{eq:prescribed-cluster-expansion}
\end{equation}
In particular,
\begin{equation}
 NU_{N,2}\longrightarrow b_{\be,\mu}
 \label{eq:prescribed-pair-connected-limit}
\end{equation}
in every finite $L^p$.  If $\rho_{N,r}^{\mathrm{fac}}$ denotes the factorial point-process correlation normalized by the macroscopic intensity $N\mu$, then
\begin{equation}
 \rho_{N,r}^{\mathrm{fac}}
 =\frac{(N)_r}{N^r}F_{N,\be,\mu}^{(r)}.
 \label{eq:prescribed-factorial-correction}
\end{equation}
Equivalently, the unnormalized factorial correlation density relative to $\mu^{\otimes r}$ is $(N)_rF_{N,\be,\mu}^{(r)}$.  The deterministic prefactor in \eqref{eq:prescribed-factorial-correction} must be expanded along with the labeled density.
\end{prop}

\begin{proof}
Existence of an expansion follows by substituting \eqref{eq:prescribed-marginal-expansion} into the finite polynomial \eqref{eq:prescribed-connected-density} and applying H\"older's inequality.  To identify the first nonzero order, we work with the normalized Gaussian expectation, including the marked one-site denominators.

At finite cutoff, write
\begin{equation}
 E_x(t)\coloneqq\exp\left\{t\phi_j(x)-\frac{t^2}{2}c_j(x,x)+t^2s_j(x)\right\},
 \quad
 A(t)\coloneqq\int E_x(t)\,\dd\mu(x),
 \quad
 H_x(t)\coloneqq\frac{E_x(t)}{A(t)}.
 \label{eq:prescribed-cluster-normalized-insertion}
\end{equation}
For $t=N^{-1/2}$, define a probability measure on the finite-dimensional Gaussian space by
\begin{equation}
 \dd\nu_t\coloneqq
 \frac{A(t)^N}{\E_G[A(t)^N]}\,\dd G.
 \label{eq:prescribed-cluster-Gaussian-tilt}
\end{equation}
The exact Gaussian marginal representation gives, simultaneously for every fixed $r\leq N$,
\begin{equation}
 F_{N,\be,\mu}^{(r)}(x_1,\ldots,x_r)
 =\E_{\nu_t}\left[\prod_{a=1}^rH_{x_a}(t)\right].
 \label{eq:prescribed-cluster-exact-moment}
\end{equation}
Consequently, $U_{N,r}$ is exactly the joint cumulant of $H_{x_1}(t),\ldots,H_{x_r}(t)$ under $\nu_t$.

Only finite coefficient identities are used from this point.  Tilt the original Gaussian measure by the quadratic factor $e^{A_2}$, where
\begin{equation}
 A_2=\frac12\int:\!\phi_j^2\!:\,\dd\mu.
 \label{eq:prescribed-cluster-quadratic-tilt}
\end{equation}
The resulting centered Gaussian has covariance $B_j=C_j(\Id-C_j)^{-1}$.  Put
\begin{equation}
 V(t)\coloneqq t^{-2}\log A(t)-A_2
 =\sum_{n\geq1}t^n\ell_{n+2,j}.
 \label{eq:prescribed-cluster-bulk-series}
\end{equation}
At each fixed order, the coefficient of $H_x(t)$ is a Gaussian polynomial whose field degree is at most its $t$-degree, while the coefficient of $t^n$ in $V(t)$ has field degree at most $n+2$.  These bounds follow from the weights one for a Gaussian field and two for a source or diagonal subtraction, and they are preserved by the finite products defining $A(t)^{-1}$ and $\log A(t)$.

The desired cumulant is the coefficient of $z_1\cdots z_r$ in the finite formal expansion of
\begin{equation}
 \log\E_{B_j}\exp\left\{V(t)+\sum_{a=1}^rz_aH_{x_a}(t)\right\}
 -\log\E_{B_j}[e^{V(t)}].
 \label{eq:prescribed-cluster-formal-cumulant}
\end{equation}
At every fixed power of $t$, this is a finite sum of joint Gaussian cumulants of $r$ marked polynomial blocks and some number $v$ of bulk polynomial blocks coming from $V$. In the diagram description, each polynomial block is a vertex and each covariance pairing of two Gaussian factors is an edge. Marked vertices correspond to the fixed coordinates; the blocks from $V$, called bulk blocks, contain integrated spatial variables. A loop pairs two factors belonging to the same block.  Wick's rule and the partition M\"obius identity retain exactly the pairings that connect all these blocks: a pairing factors over its connected-block partition, and the cumulant M\"obius sum vanishes unless that partition consists of one block.

For a surviving monomial, let $h_a$ be the Gaussian degree of marked block $a$, let $q_j$ be the degree of bulk block $j$, and let $e$ be the number of Wick edges, including loops.  Its total $t$-degree $D$ satisfies
\begin{equation}
 D\geq\sum_{a=1}^rh_a+\sum_{j=1}^v(q_j-2)=2e-2v.
 \label{eq:prescribed-cluster-degree-count}
\end{equation}
Connectivity requires at least $r+v-1$ edges between the $r+v$ blocks, and therefore $e\geq r+v-1$.  Hence $D\geq2(r-1)$.  Gaussian parity eliminates odd powers of $t$, so all coefficients before $t^{2(r-1)}=N^{-(r-1)}$ vanish identically at every finite cutoff.

Pass to the cutoff limit using the convergent marked coefficient functions and the finite polynomial \eqref{eq:prescribed-connected-density}; no pointwise Gaussian-field or diagram limit is taken.  For $r=2$,
\begin{equation}
 U_{N,2}=F_N^{(2)}-F_N^{(1)}(x_1)F_N^{(1)}(x_2),
 \label{eq:prescribed-pair-connected-formula}
\end{equation}
and the first-order singleton terms cancel, leaving $b_{\be,\mu}$.  This proves \eqref{eq:prescribed-cluster-expansion} and \eqref{eq:prescribed-pair-connected-limit}.
\end{proof}

\subsection{Joint transforms, moments, and cumulants}
\label{subsec:prescribed-linear-statistic-expansions}

Fix bounded centered real functions $\varphi_1,\ldots,\varphi_q$ and define
\begin{equation}
 L_{N,a}\coloneqq\frac1{\sqrt N}\sum_{i=1}^N\varphi_a(x_i),
 \qquad 1\leq a\leq q,
 \label{eq:prescribed-bounded-linear-statistics}
\end{equation}
and
\begin{equation}
 \Sigma_{ab}\coloneqq\ip*{\varphi_a,R_{\be,\mu}\varphi_b}_{L^2(\mu)}.
 \label{eq:prescribed-linear-covariance}
\end{equation}
Pairings in the complex transform below are extended bilinearly.

\begin{prop}[All-order joint transform and cumulant expansions]
\label{prop:prescribed-linear-statistic-expansions}
For every fixed $J\geq0$ and compact $K\subset\C^q$, there are polynomials $P_0,\ldots,P_J$, with $P_0=1$, such that
\begin{equation}
 \E_{\Q_{N,\be}(\mu)}\left[
 \exp\left\{\sum_{a=1}^qz_aL_{N,a}\right\}\right]
 =\exp\left\{\frac12z^{\mathsf T}\Sigma z\right\}
 \sum_{j=0}^{J}\frac{P_j(z)}{N^{j/2}}
 +O_K(N^{-(J+1)/2})
 \label{eq:prescribed-mgf-expansion}
\end{equation}
uniformly for $z\in K$.  On a sufficiently small fixed complex polydisc, the exact transform is zero-free for all large $N$, and its analytic logarithm has a corresponding expansion.  Hence every fixed joint moment and cumulant has an expansion to every prescribed order.  Moreover,
\begin{equation}
 \partial_{z_a}P_1(0)=\ip*{\varphi_a,a_{\be,\mu}}_{L^2(\mu)}.
 \label{eq:prescribed-first-mean-correction}
\end{equation}
\end{prop}

\begin{proof}
For complex functions, all quadratic pairings in this proof are extended bilinearly rather than sesquilinearly.  At finite cutoff, abbreviate
\begin{equation}
 R_{\be,\mu,j}\coloneqq(\Id-\be T_{\mu,j})^{-1},
 \qquad
 D_{\be,\mu,j}\coloneqq\det\nolimits_2(\Id-\be T_{\mu,j})^{-1/2},
 \label{eq:prescribed-linear-cutoff-response}
\end{equation}
and, for $t=N^{-1/2}$ and $\varphi_z=\sum_az_a\varphi_a$, define
\begin{equation}
 A_{j,z}(t)\coloneqq\int\exp\left\{t\phi_j(x)+t\varphi_z(x)
 -\frac{t^2}{2}c_j(x,x)+t^2s_j(x)\right\}\,\dd\mu(x).
 \label{eq:prescribed-linear-test-one-site}
\end{equation}
The numerator of the exact transform is $\E_G[A_{j,z}(t)^N]$.  The weighted-Wick estimates in \Cref{app:prescribed-marked-technical} remain uniform on compact complex $z$-sets because the tests are bounded.  Since both the Gaussian field and the tests are centered, $A_{j,z}(t)=1+O_{L^v(G)}(t^2)$ for every fixed $v<\infty$.  The leading logarithmic exponent is
\begin{equation}
 L_{0,j}(z)
 =\frac12\int:\!\phi_j^2\!:\,\dd\mu
 +\int\phi_j\varphi_z\,\dd\mu
 +\frac12\int\varphi_z^2\,\dd\mu,
 \label{eq:prescribed-linear-leading-exponent}
\end{equation}
and finite-dimensional Gaussian completion gives
\begin{equation}
 \E_G[e^{L_{0,j}(z)}]
 =D_{\be,\mu,j}
 \exp\left\{\frac12\ip*{\varphi_z,R_{\be,\mu,j}\varphi_z}_{L^2(\mu)}\right\}.
 \label{eq:prescribed-linear-leading-Gaussian}
\end{equation}

To control the exact exponential on the bad event, choose $\gamma>1$ and $h>1$ with $h\gamma\be_0<b$, set $L=\lceil\gamma N\rceil$, and write $w=\Re\varphi_z$.  Let $A_{j,w}(t)$ be defined by \eqref{eq:prescribed-linear-test-one-site} with $\varphi_z$ replaced by $w$.  Since $|A_{j,z}(t)|\leq A_{j,w}(t)$,
\begin{equation}
 \E_G|A_{j,z}(t)|^{\gamma N}
 \leq\left(\E_GA_{j,w}(t)^L\right)^{\gamma N/L}.
 \label{eq:prescribed-linear-higher-Gaussian-moment}
\end{equation}
Gaussian integration of the integer power and particle-space H\"older give
\begin{equation}
 \E_GA_{j,w}(t)^L
 \leq\left(\mathsf P^j_{L,h\be L/N}(\mu)\right)^{1/h}
 \left(\int e^{h'tw}\,\dd\mu\right)^{L/h'}\leq C_K.
 \label{eq:prescribed-linear-particle-holder}
\end{equation}
The coupling in the first factor lies within the partition-bound interval for all large $N$, and centering and boundedness of $w$ control the second.  This is a uniform higher Gaussian moment of the exact integer power.

On $|A_{j,z}-1|\leq1/2$, use the analytic logarithm in that disc and the complex exponential remainder
\begin{equation}
 \left|e^X-e^Y\sum_{l=0}^{J}\frac{(X-Y)^l}{l!}\right|
 \leq C_J|X-Y|^{J+1}(e^{\Re X}+e^{\Re Y}).
 \label{eq:prescribed-complex-exp-remainder}
\end{equation}
On the complement, keep the exact integer power $A_{j,z}^N$ and use \eqref{eq:prescribed-linear-higher-Gaussian-moment}; no global logarithm is introduced.  This yields a cutoff-uniform expansion through every fixed power of $t$.  Unlike the partition expansion without a test-function insertion, odd powers need not vanish.  Mixed-cutoff convergence and division by the scalar partition expansion prove \eqref{eq:prescribed-mgf-expansion}.

The leading Gaussian factor is bounded away from zero on a sufficiently small polydisc.  Uniform convergence makes the exact transform zero-free there.  Finite logarithmic series division and Cauchy estimates on a slightly larger polydisc justify termwise differentiation.  Finally,
\begin{equation}
 \E_{\Q_{N,\be}(\mu)}[L_{N,a}]
 =\sqrt N\int\varphi_a(F_{N,\be,\mu}^{(1)}-1)\,\dd\mu
 =N^{-1/2}\ip*{\varphi_a,a_{\be,\mu}}_{L^2(\mu)}+O(N^{-3/2}),
 \label{eq:prescribed-mean-from-onepoint}
\end{equation}
which proves \eqref{eq:prescribed-first-mean-correction}.
\end{proof}

\begin{remark}[Transform expansions versus Edgeworth expansions]
\label{rem:prescribed-no-Edgeworth}
The preceding proposition gives joint transform, moment, and cumulant corrections.  A distribution-function or density Edgeworth expansion requires additional nonlattice, Cram\'er, or smoothing hypotheses on the test vector.  No such conclusion is asserted for arbitrary bounded torus observables.
\end{remark}

\subsection{Relation with the homogeneous replica expansion}
\label{subsec:prescribed-versus-homogeneous}

For $\mu=\muunif$, one has $a_{\be,\muunif}=0$ and $b_{\be,\muunif}=b_\be$. The marked theorem therefore gives the homogeneous fixed-marginal density, Hoeffding, connected-correlation, and relative-entropy expansions. The replica theorem in Subsection~\ref{sec:homogeneous-replica-chaos} addresses a different scale: the fraction $k/N$ may converge to a positive limit. Its coefficient functions preserve that fraction, whereas the fixed-$k$ density remainder alone gives no proportional-marginal asymptotic.

\subsection{Homogeneous replica expansions and finite-size corrections}
\label{sec:homogeneous-replica-chaos}

For the uniform reference, the singleton coefficient vanishes and the general marked results give the local profiles and their quantitative limits. We first make this specialization. We then prove the fixed-integer replica expansion, using the shared-label bound from \Cref{subsec:replica-identities} and the coefficient calculation in \Cref{app:higher-order-replica-expansions}. The latter calculation keeps the overlap fraction in each coefficient.

\subsubsection{Homogeneous specialization of the first-order profile}

\begin{proof}[Proof of \Cref{thm:homogeneous-local-chaos}]
Choose $b$ and $\be_*$ strictly between $(1+\alpha)\be_0$ and $\beta_{\gm}$. The order-one estimate supplies \eqref{eq:prescribed-scalar-buffer} for the singleton family $\{\muunif\}$. The source vanishes and the dressed diagonal is constant, so $a_{\be,\muunif}=0$ and $G_{k,\be,\muunif}=S_{\be,k}$. Apply \Cref{thm:prescribed-growing-marginals} to obtain \eqref{eq:local-chaos-profile}--\eqref{eq:local-chaos-TV-rate}.

For the fixed pair, \Cref{thm:prescribed-marked-expansion} at $k=2$, $M=1$, and $p=2$ gives
\begin{equation}
 \left\|N(F_{N,\be}^{(2)}-1)-b_\be\right\|_2\leq C/N.
 \label{eq:homogeneous-pair-from-marked}
\end{equation}
Squaring gives \eqref{eq:local-chaos-pair-strong}. Expanding the squared norm and using the uniform $L^2$ bound for $b_\be$ gives \eqref{eq:local-chaos-pair-norm}. Translation invariance gives the exact one-particle statement, as in \eqref{eq:homogeneous-exact-first-marginal}.

Finally, \Cref{thm:all-order-replica-expansions} with $q=2$ and $M=0$, together with \eqref{eq:higher-order-two-replica-determinant}, gives \eqref{eq:homogeneous-partial-overlap-determinant}, including its $C/N$ error. That replica proof uses the strict partition buffer and the identities from \Cref{subsec:replica-identities}, not the present first-order corollary.
\end{proof}

\begin{proof}[Proof of \Cref{cor:homogeneous-local-chaos-limits}]
The fixed-$k$ limits follow from \eqref{eq:local-chaos-chi-rate}--\eqref{eq:local-chaos-TV-rate}. In the growing-marginal theorem, $\mathcal Q_{\be,\muunif}=\Xi_\be$ and the singleton term vanishes. Thus, \eqref{eq:prescribed-growing-TV-chaos} gives \eqref{eq:local-chaos-growing-TV-rate}; the chi-square and entropy constants follow by taking $k\to\infty$ and $k/N\to0$ in the same estimates.
\end{proof}

\subsubsection{The pair response kernel and the replica coefficients}
\label{subsec:screened-pair-response}

\paragraph{The response differs from the original kernel by a bounded kernel.}
The resolvent identity gives
\begin{equation}\label{eq:b-split}
 b_\be=\be \g+s_\be,\qquad
 \wh s_\be(n)=\frac{\be^2\wh\g(n)^2}{1-\be\wh\g(n)},\qquad
 s_\be=\be \g*b_\be.
\end{equation}
The Fourier series of $s_\be$ is absolutely convergent and nonnegative coefficientwise.  Hence it is a bounded continuous canonical kernel, with
\begin{equation}\label{eq:s-bound}
 \|s_\be\|_\infty\leq s_\be(0)
 \leq S:=\frac{\be_0^2\|T\|_{\HS}^2}{1-\be_0/\beta_{\spec}},
 \qquad
 \sup_{0\leq\be\leq\be_0}\|b_\be\|_{L^p(\muunif)}<\infty\quad(1\leq p<\infty).
\end{equation}

The row and contraction bounds just obtained are the input to \Cref{app:higher-order-replica-expansions}. There, \Cref{lem:qbuffer} supplies the exact exponential bound, and \Cref{lem:expansion} gives the arbitrary finite-order expansion. They prove \Cref{thm:all-order-replica-expansions}, while \Cref{prop:moments-next} and the subsequent contraction identities prove \Cref{cor:higher-order-divergences}.

\subsubsection{Proportional marginals and the local limit}

The determinant in \eqref{eq:main-overlap-log}--\eqref{eq:main-overlap-expansion} describes the full dependence on the marginal fraction within the strict amplified-coupling window. Its quadratic term agrees with the first-order local-chaos constant. At a fixed positive fraction, the higher even powers give additional contributions. Appendix~\ref{app:higher-order-replica-expansions} separates these overlap powers from the inverse-particle corrections; their relative ordering changes near $k$ of order $\sqrt N$. In all these formulas the coefficients are evaluated at the exact fraction $k/N$ and use the exact pair count.


\section{Planar confined equilibria and partition bounds}
\label{sec:planar-confined-static-package}

We now consider the planar logarithmic gas with quadratic confinement. The two estimates needed near its thermal equilibrium are a strict global entropy--energy margin and an order-one centered partition bound at a larger coupling. The next theorem gives both estimates uniformly on a fixed interval below the critical inverse temperature. These are also the whole-space static inputs used in Part~II.

Throughout,
\begin{equation}\label{pcsp:eq:kernel}
 \g(x)=-\log|x|,\qquad x\in\R^2\setminus\{0\}.
\end{equation}
For a probability measure $\mu$ for which the expressions are finite, set
\begin{equation}\label{pcsp:eq:potential-energy}
 U_\mu(x)=\int_{\R^2}\g(x-y)\,\mu(\dd y),
 \qquad
 \Gcal(\mu)=\iint_{\R^2\times\R^2}\g(x-y)\,\mu(\dd x)\mu(\dd y).
\end{equation}
The canonical centered kernel is
\begin{equation}\label{pcsp:eq:centered-kernel}
 \widetilde\g_\mu(x,y)
 =\g(x-y)-U_\mu(x)-U_\mu(y)+\Gcal(\mu).
\end{equation}
It satisfies
\begin{equation}\label{pcsp:eq:canonicality}
 \int_{\R^2}\widetilde\g_\mu(x,y)\,\mu(\dd y)=0
 \quad\text{for $\mu$-almost every $x$}.
\end{equation}
For $N\geq2$ and $s\geq0$, define the centered pair Laplace transform
\begin{equation}\label{pcsp:eq:centered-laplace}
 \K_N(s;\mu)
 =\E_{\mu^{\otimes N}}\left[
   \exp\left\{\frac{s}{N}\sum_{1\leq i<j\leq N}
   \widetilde\g_\mu(X_i,X_j)\right\}\right].
\end{equation}
The sum in \eqref{pcsp:eq:centered-laplace} is over unordered pairs, and the coefficient is $s/N$.  Part~II uses the same unordered-pair convention.

Let
\begin{equation}\label{pcsp:eq:gaussian-reference}
 \gamma(x)=\pi^{-1}e^{-|x|^2}.
\end{equation}
For $0\leq\be<4$, the dimensionless mean-field functional is
\begin{equation}\label{pcsp:eq:mean-field-functional}
 \mathcal I_\be(\rho)
 =\Ent(\rho\mid\gamma)-\frac\be2\Gcal(\rho).
\end{equation}
Its equilibrium will be denoted by $m_\be$.

\begin{thm}[Planar equilibrium margin and partition bound]\label{thm:planar-confined-static-package}
Fix $0<\be_0<4$.  There are constants
\begin{equation}\label{pcsp:eq:package-constants}
 R\geq1,
 \quad
 \eta_{\rm eq}\in(0,4-\be_0),
 \quad
 \delta\in(0,4-\be_0),
 \quad
 C_{\rm eq}\geq1,
 \quad
 C_{\rm buf}\geq1,
\end{equation}
with the following properties.

For every $0\leq\be\leq\be_0$, there is a unique smooth positive probability density $m_\be$ satisfying
\begin{equation}\label{pcsp:eq:euler-lagrange}
 \log m_\be(x)+|x|^2-\be U_{m_\be}(x)=c_\be
\end{equation}
for a constant $c_\be$.  If
\begin{equation}\label{pcsp:eq:comparison-density}
 n_\be(x)=Z_\be^{-1}e^{-|x|^2}(R^2+|x|^2)^{-\be/2},
\end{equation}
then
\begin{equation}\label{pcsp:eq:equilibrium-envelope}
 C_{\rm eq}^{-1}n_\be(x)\leq m_\be(x)\leq C_{\rm eq}n_\be(x),
\end{equation}
and, for $j=1,2,3$,
\begin{equation}\label{pcsp:eq:equilibrium-derivatives}
 \abs{D^jm_\be(x)}\leq C_j(1+|x|)^jm_\be(x),
\end{equation}
with constants common to $\be\in[0,\be_0]$.

For every probability density $\rho$ for which both terms are finite,
\begin{equation}\label{pcsp:eq:F-EQ}
 \Ent(\rho\mid m_\be)
 -\frac{\be+\eta_{\rm eq}}2
 \iint\g(x-y)(\rho-m_\be)(x)(\rho-m_\be)(y)\,\dd x\dd y
 \geq0.
\end{equation}
Set
\begin{equation}\label{pcsp:eq:frozen-radius}
 r_*=\min\left\{\frac14,
 \frac{\eta_{\rm eq}}{4(\eta_{\rm eq}+\be_0)}\right\},
 \qquad
 M_*=(1-r_*)^{-1}.
\end{equation}
For every probability density $\nu$ satisfying
\begin{equation}\label{pcsp:eq:tube}
 \norm{\nu/m_\be-1}_{L^\infty(m_\be)}\leq r_*,
\end{equation}
one has
\begin{equation}\label{pcsp:eq:F-BUF}
 \sup_{N\geq2}\K_N(\be+\delta;\nu)\leq C_{\rm buf}.
\end{equation}
Every constant in \eqref{pcsp:eq:package-constants}--\eqref{pcsp:eq:F-BUF} is independent of $\be\in[0,\be_0]$ and of $\nu$ in the tube \eqref{pcsp:eq:tube}.
\end{thm}

\begin{remark}[Scope]\label{pcsp:rem:scope}
The theorem concerns quadratic confinement in dimension two and reference measures in the bounded-ratio neighborhood \eqref{pcsp:eq:tube}. Its constants are uniform on the stated compact interval below $\be=4$. Extensions to the critical endpoint, higher dimensions, other confinements, or larger reference classes require additional estimates.
\end{remark}



The proof is given in \Cref{app:planar-confined-static-package}.

\section{Perspectives and open problems}
\label{sec:perspectives}

\subsection{The endpoint and intrinsic nonuniform phase boundaries}

At a collision endpoint, the unregularized partition function is infinite for every $N\geq2$, by \Cref{prop:fixed-particle-integrability}. Thus, in dimensions $2\leq\d\leq10$, a critical normalization must specify a regularization or a joint limiting procedure. For $\d\geq11$, the global endpoint lies strictly below collision, so the fixed-particle integrals remain finite there; the question is their growth with $N$.

At the global-minimization endpoint, \eqref{eq:finite-entropy-energy-bound} makes $I_\mu$ finite on the stated finite-entropy domain. Passing to the endpoint in \eqref{eq:strict-subcritical-macroscopic-inequality} therefore gives
\begin{equation}
 \Ec_{\beta_{\gm}}(\mu)
 =\lim_{b\uparrow\beta_{\gm}}\bigl(\Hr(\mu\mid\muunif)+bI_\mu\bigr)\geq0,
 \qquad \Pc_\d(\beta_{\gm})=0.
 \label{eq:global-endpoint-macroscopic-pressure}
\end{equation}
The uniform measure attains zero. Thus, a strictly lower-free-energy competing phase is not the macroscopic mechanism at this endpoint. This conclusion does not prove a microscopic pressure upper bound or classify the endpoint fluctuations. For $\d\geq12$, \Cref{thm:beckner-onofri-input} identifies the global endpoint with the spectral threshold and gives the uniform measure as its unique minimizer. The finite-particle normalization and the degenerate fluctuations at this endpoint remain to be determined; a natural question is whether the first Fourier shell governs those fluctuations. In dimension eleven, the strict inequality $\beta_{\gm}(11)<\beta_{\spec}(11)$ places the spectral point inside the exponentially growing phase.

For a general nonatomic reference measure $\mu$, one may define the intrinsic macroscopic functional
\begin{equation}
 \Ec_{\mu,\be}(\nu)
 \coloneqq\Hr(\nu\mid\mu)
 -\frac\be2\iint\tl{\g}_{\mu}(x,y)
 (\nu(x)-1)(\nu(y)-1)\,\dd\mu(x)\dd\mu(y),
 \label{eq:intrinsic-nonuniform-free-energy}
\end{equation}
where the same symbol $\nu$ denotes the measure and its density with respect to $\mu$.  One may then define the corresponding global-minimization and order-one partition thresholds.  The transported estimate in Part B gives an explicit lower bound on the latter threshold.  It does not prove equality between the two intrinsic nonuniform thresholds.  Establishing that equality would require a version of the coarse/canonical argument adapted directly to the geometry and cell masses of $\mu$, together with a sharp macroscopic inequality for \eqref{eq:intrinsic-nonuniform-free-energy}.


\subsection{Local-chaos endpoints and finer asymptotics}
\label{subsec:local-chaos-perspectives}


The leading prescribed-reference profiles and divergence constants are identified in \Cref{thm:prescribed-marked-expansion,cor:prescribed-divergence-expansions,thm:prescribed-growing-marginals}. When the leading coefficient is nonzero, these asymptotics also give matching lower bounds. All fixed-marginal density and relative-entropy orders are covered by the marked theorem. The remaining questions concern the optimal overlap window and extensions to its boundary, critical parameters, or symmetry-broken phases. The finite-particle replica collision condition in \eqref{eq:local-replica-collision-condition} is distinct from a particle-uniform proportional bound. Static sharp marginal estimates do not improve the time-dependent rates addressed in Part II without a new evolution argument.



\subsection{Higher total-variation orders and growing marked expansions}

The strong marked theorem identifies the second signed marginal profile and yields the second total-variation coefficient in \eqref{eq:prescribed-second-TV-coefficient}.  This does not make total variation an ordinary analytic functional of the signed coefficient series.  Already at second order, the coefficient separates the regions where the leading profile is nonzero and zero.  At higher orders, additional strata on which successive signed profiles vanish can enter.  An arbitrary-order total-variation expansion therefore requires geometric or transversality information on those zero sets; it does not follow from the $L^p$ density expansion alone.

The explicit estimate \eqref{eq:prescribed-explicit-k-remainder} is uniform for $k$ bounded by a fixed multiple of $N$, but its polynomial dependence on $k$ need not be optimal.  The all-order theorem still fixes $k$ and the expansion order before $N$ tends to infinity.  An all-order expansion for growing $k$ requires a comparison of the coefficient scale $(1+k)^m$, the remainder scale, and the powers of the overlap fraction $k/N$ with constants controlled as $k$ grows.  The homogeneous calculation after \eqref{eq:overlap-series} shows that even the ordering of the first finite-size and nonlinear-overlap corrections changes near $k\asymp\sqrt N$.

\subsection{Attractive interactions with phase transitions}

The identity $\beta_{\pf}=\beta_{\gm}$ isolates a general mechanism: the strict order-one regime and the strict global-minimization regime have the same boundary.  The proof uses additional logarithmic structure, especially the cancellation of the additive scale constant, the collision threshold $2\d$, and the mixed-second-difference bound.  A natural target is to extend the argument to a logarithmic core plus a bounded sufficiently regular perturbation, assuming the corresponding macroscopic global-minimization threshold and local residual estimates can be established.  For sign-indefinite or multi-well interactions, the same question becomes phase dependent: a determinant describes fluctuations around a selected locally minimizing phase, while the full partition function also receives exponentially large contributions from any phase with lower macroscopic free energy.

In a symmetry-broken regime, one expects the macroscopic minimizers to determine the exponential contribution to the partition function. A nondegenerate minimizer should have a determinant correction, while phase coexistence, continuous symmetry orbits, or a degenerate Hessian can change the subexponential normalization. Determining these corrections lies beyond the results proved here strictly below the global-minimization threshold.

\subsection{Whole-space thermal equilibria under confinement}

For the attractive logarithmic interaction on $\R^{\d}$ with confinement, the natural reference measure is a nonuniform thermal equilibrium.  Formally, whenever the Hoeffding-centered interaction kernel defines a nonnegative Hilbert--Schmidt operator, the corresponding candidate normalization is
\begin{equation}
 e^{\be I_{\mu}}
 \det\nolimits_2(\Id-\be T_{\mu})^{-1/2}.
 \label{eq:whole-space-candidate-determinant}
\end{equation}
In the planar quadratic setting, \Cref{thm:planar-confined-static-package} gives the strict global equilibrium margin and the prescribed-reference order-one moment bound. The finite-partition argument in \Cref{app:planar-confined-static-package} includes the tail cells, but it does not identify the determinant limit in \eqref{eq:whole-space-candidate-determinant}. A determinant asymptotic requires approximation estimates for the centered interaction and its covariance operator that remain uniform under the same partition buffer. Higher dimensions, nonquadratic confinements, the endpoint $\be=4$, and references outside the stated bounded-ratio neighborhood are not covered by the planar theorem.

\appendix

\section{Hubbard--Stratonovich estimate under a higher-moment bound}
\label{app:feature-laplace}

We prove \Cref{prop:renyi-feature-space,prop:triangular-feature-space} by comparing the exact Gaussian single-site integral with its quadratic approximation. A particle partition bound at a coefficient strictly larger than $r\be$, for some $r>1$, controls the exact exponential in Gaussian $L^r$. The spectral gap controls the quadratic exponential in the same space. These two bounds justify Taylor expansion on a good event and control its complement; a bound at $2\be$ is unnecessary.

Let $\mathsf Z$ be an isonormal Gaussian process over $\mathcal H$, and write
\begin{equation}
 \mathsf Z_x\coloneqq\ip*{\mathsf Z,\mathsf k(x)}_{\mathcal H}.
 \label{eq:feature-gaussian-field}
\end{equation}
For $\tau\geq0$, define
\begin{equation}
 f_{\be}(\tau)
 \coloneqq\int_X\exp\left\{
 \tau\sqrt\be\,\mathsf Z_x-\frac{\tau^2\be}{2}D(x)
 \right\}\,\dd\mu(x).
 \label{eq:feature-single-site-transform}
\end{equation}
The Hubbard--Stratonovich transformation gives
\begin{equation}
 \E_{\mu^{\otimes N}}
 \left[e^{\be\mathsf W_N(q)/2}\right]
 =\E_{\mathsf Z}\left[f_{\be}(N^{-1/2})^N\right].
 \label{eq:feature-hubbard-stratonovich}
\end{equation}
Indeed, conditional on $(\xi_1,\ldots,\xi_N)$, the Gaussian expectation produces
\begin{equation}
 \exp\left\{\frac\be{2N}
 \norm*{\sum_{i=1}^N\mathsf k(\xi_i)}_{\mathcal H}^2
 -\frac\be{2N}\sum_{i=1}^ND(\xi_i)\right\}
 =e^{\be\mathsf W_N(q)/2}.
 \label{eq:feature-diagonal-cancellation}
\end{equation}

We first estimate the error in the logarithm of the single-site integral. These estimates will be paired with the Gaussian $L^r$ bounds for the two exponentials. Put
\begin{equation}
 p\coloneqq\frac r{r-1},
 \qquad
 s\coloneqq2p.
 \label{eq:feature-conjugate-exponents}
\end{equation}
For fixed $\be$, define the second- and third-chaos terms
\begin{equation}
 B_{\be}
 \coloneqq\frac\be2\int_X(\mathsf Z_x^2-D(x))\,\dd\mu(x),
 \label{eq:feature-second-chaos}
\end{equation}
\begin{equation}
 R_{\be}
 \coloneqq\frac{\be^{3/2}}6\int_X
 (\mathsf Z_x^3-3D(x)\mathsf Z_x)\,\dd\mu(x).
 \label{eq:feature-third-chaos}
\end{equation}

\begin{lemma}[Feature-space Taylor bounds]
\label{lem:feature-taylor-bounds}
Assume the hypotheses of \Cref{prop:renyi-feature-space}.  There are $N_0<\infty$ and $C<\infty$, depending only on the quantities in \eqref{eq:renyi-feature-constant}, such that, for $N\geq N_0$, $0\leq\be\leq\be_0$, and
\begin{equation}
 A_N\coloneqq f_{\be}(N^{-1/2}),
 \qquad
 Q_a\coloneqq1+M_a(q),
 \label{eq:feature-AN-Qa}
\end{equation}
one has
\begin{equation}
 A_N
 =1+\frac{B_{\be}}N+\frac{R_{\be}}{N^{3/2}}
 +\frac{E_{N,\be}}{N^2},
 \label{eq:feature-Taylor-expansion}
\end{equation}
with
\begin{equation}
 \norm{B_{\be}}_{L^s(\mathsf Z)}
 +\norm{R_{\be}}_{L^s(\mathsf Z)}
 +\norm{E_{N,\be}}_{L^s(\mathsf Z)}
 \leq C Q_a.
 \label{eq:feature-chaos-remainder-bound}
\end{equation}
Consequently,
\begin{equation}
 \norm{A_N-1}_{L^s(\mathsf Z)}
 \leq\frac{CQ_a}{N}.
 \label{eq:feature-AN-minus-one}
\end{equation}
If
\begin{equation}
 G_N\coloneqq\{\abs{A_N-1}\leq1/2\},
 \label{eq:feature-good-event}
\end{equation}
then
\begin{equation}
 \norm{(N\log A_N-B_{\be})\one_{G_N}}_{L^p(\mathsf Z)}
 \leq C\frac{Q_a^2}{\sqrt N},
 \label{eq:feature-logarithmic-error}
\end{equation}
and
\begin{equation}
 \Pb_{\mathsf Z}(G_N^c)^{1/p}
 \leq C\frac{Q_a^2}{N^2}.
 \label{eq:feature-bad-event-probability}
\end{equation}
\end{lemma}

\begin{proof}
For fixed $x$, put
\begin{equation}
 a_x\coloneqq\sqrt\be\,\mathsf Z_x,
 \qquad
 c_x\coloneqq\be D(x),
 \qquad
 y_x(\tau)\coloneqq a_x-c_x\tau.
 \label{eq:feature-derivative-data}
\end{equation}
Four differentiations give
\begin{equation}
 \frac{\dd^4}{\dd\tau^4}
 e^{a_x\tau-c_x\tau^2/2}
 =\left(y_x(\tau)^4-6c_xy_x(\tau)^2+3c_x^2\right)
 e^{a_x\tau-c_x\tau^2/2}.
 \label{eq:feature-fourth-derivative}
\end{equation}
Also, canonicality in \eqref{eq:centered-feature-representation} gives $f_{\be}'(0)=0$, while direct differentiation gives
\begin{equation}
 \frac12f_{\be}''(0)=B_{\be},
 \qquad
 \frac16f_{\be}'''(0)=R_{\be}.
 \label{eq:feature-second-third-derivatives}
\end{equation}
Taylor's formula with integral remainder therefore yields \eqref{eq:feature-Taylor-expansion}, where
\begin{equation}
 E_{N,\be}
 \coloneqq\frac{N^2}{6}
 \int_0^{N^{-1/2}}(N^{-1/2}-\tau)^3
 f_{\be}^{(4)}(\tau)\,\dd\tau.
 \label{eq:feature-remainder-definition}
\end{equation}

To bound the fourth derivative in Gaussian $L^s$, absorb its exponential factor into a scalar Gaussian shift. Under the probability measure with density proportional to
\begin{equation}
 \exp\left\{s\tau\sqrt\be\,\mathsf Z_x
 -\frac{s\tau^2\be}{2}D(x)\right\},
 \label{eq:feature-gaussian-shift-density}
\end{equation}
the scalar $\sqrt\be\,\mathsf Z_x$ is Gaussian with variance $\be D(x)$ and mean $s\tau\be D(x)$.  Hence $y_x(\tau)$ has variance $\be D(x)$ and mean $(s-1)\tau\be D(x)$.  Taking the $L^s(\mathsf Z)$ norm in \eqref{eq:feature-fourth-derivative} and using the elementary Gaussian moment bounds gives
\begin{align}
 &\norm*{\frac{\dd^4}{\dd\tau^4}
 e^{a_x\tau-c_x\tau^2/2}}_{L^s(\mathsf Z)}
 \notag\\
 &\qquad\leq C_{s,\be_0}
 e^{(s-1)\tau^2\be D(x)/2}
 \left(D(x)^2+\tau^2D(x)^3+\tau^4D(x)^4\right).
 \label{eq:feature-fourth-derivative-bound}
\end{align}
Choose $N_0$ so large that
\begin{equation}
 \frac{(s-1)\be_0}{2N}\leq\frac a2
 \qquad(N\geq N_0).
 \label{eq:feature-diagonal-coefficient-choice}
\end{equation}
For $0\leq\tau\leq N^{-1/2}$, \eqref{eq:feature-diagonal-coefficient-choice} and the inequalities
\begin{equation}
 z^j e^{az/2}\leq C_{a,j}e^{az}
 \qquad(z\geq0,
 \ j=2,3,4)
 \label{eq:feature-polynomial-absorption}
\end{equation}
imply, by Minkowski's integral inequality,
\begin{equation}
 \sup_{0\leq\tau\leq N^{-1/2}}
 \norm{f_{\be}^{(4)}(\tau)}_{L^s(\mathsf Z)}
 \leq C Q_a.
 \label{eq:feature-fourth-derivative-Ls}
\end{equation}
Equation \eqref{eq:feature-remainder-definition} now gives
\begin{equation}
 \norm{E_{N,\be}}_{L^s(\mathsf Z)}\leq C Q_a.
 \label{eq:feature-remainder-Ls}
\end{equation}

The second chaos has the spectral expansion
\begin{equation}
 B_{\be}
 =\frac\be2\sum_{\ell\geq1}\lambda_\ell(G_\ell^2-1).
 \label{eq:feature-B-spectral-expansion}
\end{equation}
Gaussian hypercontractivity and $\sum_\ell\lambda_\ell^2\leq H^2$ give
\begin{equation}
 \norm{B_{\be}}_{L^s(\mathsf Z)}\leq C_{s,\be_0}H.
 \label{eq:feature-B-Ls}
\end{equation}
For the third chaos, Minkowski's inequality and
\begin{equation}
 \norm{\mathsf Z_x^3-3D(x)\mathsf Z_x}_{L^s(\mathsf Z)}
 \leq C_sD(x)^{3/2}
 \label{eq:feature-cubic-gaussian-moment}
\end{equation}
give
\begin{equation}
 \norm{R_{\be}}_{L^s(\mathsf Z)}
 \leq C_{s,\be_0}\int_XD(x)^{3/2}\,\dd\mu(x)
 \leq C_{s,\be_0,a}Q_a.
 \label{eq:feature-R-Ls}
\end{equation}
Equations \eqref{eq:feature-remainder-Ls}, \eqref{eq:feature-B-Ls}, and \eqref{eq:feature-R-Ls} prove \eqref{eq:feature-chaos-remainder-bound}, after enlarging the constant to absorb $H$.  Equation \eqref{eq:feature-AN-minus-one} follows from \eqref{eq:feature-Taylor-expansion}.

On $G_N$, the elementary inequality
\begin{equation}
 \abs{\log(1+z)-z}\leq2z^2
 \qquad(\abs z\leq1/2)
 \label{eq:feature-logarithm-remainder}
\end{equation}
gives
\begin{align}
 \abs{N\log A_N-B_{\be}}
 &\leq\frac{\abs{R_{\be}}}{\sqrt N}
 +\frac{\abs{E_{N,\be}}}{N}
 +2N\abs{A_N-1}^2.
 \label{eq:feature-logarithm-pointwise}
\end{align}
Since $s=2p$, taking $L^p$ norms and using \eqref{eq:feature-chaos-remainder-bound}--\eqref{eq:feature-AN-minus-one} proves \eqref{eq:feature-logarithmic-error}.  Finally, Markov's inequality and $s/p=2$ give
\begin{equation}
 \Pb(G_N^c)^{1/p}
 \leq2^{s/p}\norm{A_N-1}_{L^s}^{s/p}
 \leq C\frac{Q_a^2}{N^2},
 \label{eq:feature-bad-event-Markov}
\end{equation}
which is \eqref{eq:feature-bad-event-probability}.
\end{proof}

\begin{proof}[Proof of \Cref{prop:renyi-feature-space}]
Let $A_N$ be as in \eqref{eq:feature-AN-Qa}.  For an integer $M\geq1$, Fubini's theorem and Gaussian integration give
\begin{equation}
 \E_{\mathsf Z}[A_N^M]
 =\E_{\mu^{\otimes M}}
 \left[e^{\be M\mathsf W_M(q)/(2N)}\right].
 \label{eq:feature-integer-power-identity}
\end{equation}
Let
\begin{equation}
 M_N\coloneqq\lceil rN\rceil.
 \label{eq:feature-ceiling}
\end{equation}
Since $b-r\be_0>0$, there is $N_1<\infty$ such that
\begin{equation}
 \be\frac{M_N}{N}\leq b
 \qquad(N\geq N_1,
 \ 0\leq\be\leq\be_0).
 \label{eq:feature-ceiling-buffer}
\end{equation}
Monotonicity of $L^u$ norms, \eqref{eq:feature-integer-power-identity}, and \eqref{eq:renyi-feature-particle-assumption} yield
\begin{align}
 \norm{A_N^N}_{L^r(\mathsf Z)}
 &=\left(\E\left[A_N^{rN}\right]\right)^{1/r}
 \notag\\
 &\leq\left(\E\left[A_N^{M_N}\right]\right)^{N/M_N}
 \leq P
 \qquad(N\geq N_1).
 \label{eq:feature-AN-power-Lr}
\end{align}
Here, the particle assumption at coefficient $b$ controls all smaller coefficients by concavity, and $P\geq1$ because the $M=1$ statistic vanishes.

The Gaussian second chaos satisfies
\begin{equation}
 \E_{\mathsf Z}[e^{B_{\be}}]
 =\prod_{\ell\geq1}
 e^{-\be\lambda_\ell/2}(1-\be\lambda_\ell)^{-1/2}
 =\det\nolimits_2(\Id-\be T_q)^{-1/2}.
 \label{eq:feature-B-determinant}
\end{equation}
Moreover, the spectral gap in \eqref{eq:renyi-feature-operator-assumptions} implies
\begin{equation}
 \sup_{0\leq\be\leq\be_0}
 \norm{e^{B_{\be}}}_{L^r(\mathsf Z)}<\infty.
 \label{eq:feature-eB-Lr}
\end{equation}
Indeed, for $0\leq x\leq1-\delta$,
\begin{equation}
 -x-\log(1-x)\leq\frac{x^2}{2\delta},
 \label{eq:feature-scalar-determinant-bound}
\end{equation}
so the logarithm of the $r$th moment in \eqref{eq:feature-eB-Lr} is bounded by $C_{r,\delta}\be_0^2\norm{T_q}_{\HS}^2$.

On the good event \eqref{eq:feature-good-event},
\begin{equation}
 \abs{A_N^N-e^{B_{\be}}}
 \leq\abs{N\log A_N-B_{\be}}
 \left(A_N^N+e^{B_{\be}}\right).
 \label{eq:feature-exponential-difference}
\end{equation}
H\"older's inequality with exponents $p,r$, together with \eqref{eq:feature-logarithmic-error}, \eqref{eq:feature-AN-power-Lr}, and \eqref{eq:feature-eB-Lr}, gives
\begin{equation}
 \E\left[\abs{A_N^N-e^{B_{\be}}}\one_{G_N}\right]
 \leq C\frac{Q_a^2}{\sqrt N}.
 \label{eq:feature-good-event-comparison}
\end{equation}
On $G_N^c$, the same H\"older pair and \eqref{eq:feature-bad-event-probability} give
\begin{equation}
 \E\left[(A_N^N+e^{B_{\be}})\one_{G_N^c}\right]
 \leq C\frac{Q_a^2}{N^2}.
 \label{eq:feature-bad-event-comparison}
\end{equation}
Combining \eqref{eq:feature-hubbard-stratonovich}, \eqref{eq:feature-B-determinant}, \eqref{eq:feature-good-event-comparison}, and \eqref{eq:feature-bad-event-comparison} proves \eqref{eq:renyi-feature-rate} for $N\geq\max\{N_0,N_1\}$.  The remaining finite set of particle numbers is controlled by \eqref{eq:renyi-feature-particle-assumption}, \eqref{eq:feature-scalar-determinant-bound}, and $Q_a\geq1$, after enlarging $C$.
\end{proof}

\begin{proof}[Proof of \Cref{prop:triangular-feature-space}]
For each $N$, repeat the preceding construction with $q=q_N$ and write
\begin{equation}
 L_N\coloneqq L_*(1+\log N).
 \label{eq:triangular-LN}
\end{equation}
The integer-power estimate \eqref{eq:feature-AN-power-Lr} and the Gaussian bound \eqref{eq:feature-eB-Lr} are uniform by \eqref{eq:triangular-feature-operator-assumptions}--\eqref{eq:triangular-feature-particle-assumption}.  It remains to replace the exponential diagonal moment in \Cref{lem:feature-taylor-bounds}.

The pointwise fourth-derivative bound \eqref{eq:feature-fourth-derivative-bound}, now with $D_N\leq L_N$, gives, for $0\leq\tau\leq N^{-1/2}$,
\begin{align}
 \norm{f_{N,\be}^{(4)}(\tau)}_{L^s(\mathsf Z)}
 &\leq C
 e^{C L_N/N}
 \left(L_N^2+N^{-1}L_N^3+N^{-2}L_N^4\right)
 \notag\\
 &\leq C(1+L_N)^2.
 \label{eq:triangular-fourth-derivative}
\end{align}
The second inequality uses
\begin{equation}
 \sup_{N\geq1}\frac{1+\log N}{N}<\infty.
 \label{eq:triangular-log-over-N}
\end{equation}
Likewise,
\begin{equation}
 \norm{R_{N,\be}}_{L^s(\mathsf Z)}
 \leq C L_N^{3/2}.
 \label{eq:triangular-third-chaos}
\end{equation}
Taylor's formula therefore yields
\begin{equation}
 \norm{A_N-1}_{L^s(\mathsf Z)}
 \leq C\frac{1+L_N}{N}.
 \label{eq:triangular-AN-minus-one}
\end{equation}
The same logarithmic argument as in \eqref{eq:feature-logarithm-pointwise}, using $s=2p$, gives
\begin{equation}
 \norm{(N\log A_N-B_{N,\be})\one_{G_N}}_{L^p(\mathsf Z)}
 \leq C\frac{(1+L_N)^2}{\sqrt N},
 \label{eq:triangular-logarithmic-error}
\end{equation}
and Markov's inequality gives
\begin{equation}
 \Pb(G_N^c)^{1/p}
 \leq C\frac{(1+L_N)^2}{N^2}.
 \label{eq:triangular-bad-event}
\end{equation}
Applying the good- and bad-event comparison \eqref{eq:feature-exponential-difference}--\eqref{eq:feature-bad-event-comparison} proves
\begin{equation}
 \sup_{0\leq\be\leq\be_0}
 \left|\E\left[e^{\be\mathsf W_N(q_N)/2}\right]
 -\det\nolimits_2(\Id-\be T_{q_N})^{-1/2}\right|
 \leq C\frac{(1+L_N)^2}{\sqrt N}.
 \label{eq:triangular-prelog-rate}
\end{equation}
Since $1+L_N\leq C_{L_*}(1+\log N)$, \eqref{eq:triangular-prelog-rate} proves \eqref{eq:triangular-feature-rate} for all sufficiently large $N$.  The finitely many smaller particle numbers are controlled uniformly by \eqref{eq:triangular-feature-particle-assumption} and the common determinant gap in \eqref{eq:triangular-feature-operator-assumptions}; enlarging $C$ proves the estimate for every $N\geq1$.
\end{proof}

\section{Logarithmic regularization and determinant continuity}
\label{app:homogeneous-regularization}

We prove the estimates for regularization preserving nonnegative definiteness used in Section~\ref{sec:thresholds-asymptotics}. Two preliminary bounds control the effect of changing the kernel: a canonical even-moment estimate for the particle statistic and Hilbert--Schmidt continuity of the Carleman--Fredholm determinant under a common spectral gap. We then estimate the logarithmic convolution error and show that convolution with the prescribed nonnegative mollifier preserves the scalar partition bound.

\begin{lemma}[Canonical even moment]
\label{lem:canonical-even-moment}
Let $(X,\mu)$ be a probability space, let $u\geq2$ be an even integer, and let $h\in L^u(\mu^{\otimes2})$ be real, symmetric, and canonical:
\begin{equation}
 \int_Xh(x,y)\,\dd\mu(y)=0
 \qquad\text{for }\mu\text{-almost every }x.
 \label{eq:canonical-even-moment-hypothesis}
\end{equation}
For independent $\xi_1,\ldots,\xi_N\sim\mu$, define
\begin{equation}
 \mathsf W_N(h)\coloneqq\frac1N\sum_{i\neq j}h(\xi_i,\xi_j).
 \label{eq:canonical-even-moment-statistic}
\end{equation}
Then
\begin{equation}
 \norm{\mathsf W_N(h)}_{L^u(\mu^{\otimes N})}
 \leq C_u\norm{h}_{L^u(\mu^{\otimes2})}
 \qquad(N\geq1).
 \label{eq:canonical-even-moment-conclusion}
\end{equation}
\end{lemma}

\begin{proof}
Expanding the even power gives
\begin{equation}
 \E\left[\abs{\mathsf W_N(h)}^u\right]
 =N^{-u}\sum_{(i_\ell,j_\ell)_{\ell=1}^u}
 \int\prod_{\ell=1}^uh(x_{i_\ell},x_{j_\ell})
 \,\dd\mu^{\otimes N}(\XN),
 \label{eq:canonical-even-moment-expansion}
\end{equation}
where $i_\ell\neq j_\ell$ for every $\ell$.  If a particle label occurs exactly once among the $2u$ endpoints of a summand, integration in that variable and \eqref{eq:canonical-even-moment-hypothesis} make the summand vanish.  Every surviving summand therefore uses at most $u$ distinct particle labels.

For each $1\leq v\leq u$, the number of endpoint patterns on $v$ abstract labels is at most $v^{2u}$, and the number of injections of those labels into $\{1,\ldots,N\}$ is at most $N^v$.  Hence the number of potentially nonzero summands is bounded by
\begin{equation}
 \sum_{v=1}^uv^{2u}N^v\leq C_uN^u.
 \label{eq:canonical-even-moment-count}
\end{equation}
For each summand, H\"older's inequality on the full product space, with exponent $u$ for each of the $u$ factors, gives
\begin{equation}
 \left|\int\prod_{\ell=1}^uh(x_{i_\ell},x_{j_\ell})
 \,\dd\mu^{\otimes N}\right|
 \leq\norm{h}_{L^u(\mu^{\otimes2})}^u.
 \label{eq:canonical-even-moment-holder}
\end{equation}
Combining \eqref{eq:canonical-even-moment-expansion}--\eqref{eq:canonical-even-moment-holder} proves \eqref{eq:canonical-even-moment-conclusion}.
\end{proof}

For a Hilbert--Schmidt operator $A$ with $\norm A_{\op}<1$, the Carleman--Fredholm determinant is characterized by
\begin{equation}
 \log\det\nolimits_2(\Id-A)
 =-\sum_{n=2}^\infty\frac1n\Tr(A^n).
 \label{eq:det2-series-definition}
\end{equation}
We use the standard properties of this determinant as developed, for example, in Simon \cite[Chapter~9]{Simon2005}.

\begin{lemma}[Hilbert--Schmidt continuity under a common gap]
\label{lem:determinant-continuity}
Let $A$ and $B$ be nonnegative self-adjoint Hilbert--Schmidt operators satisfying
\begin{equation}
 \norm A_{\op},\norm B_{\op}\leq L,
 \qquad
 \norm A_{\HS},\norm B_{\HS}\leq H,
 \qquad
 \be_0L<1.
 \label{eq:determinant-continuity-hypotheses}
\end{equation}
Then there is $C=C(\be_0,L,H)<\infty$ such that
\begin{equation}
 \sup_{0\leq\be\leq\be_0}
 \left|\det\nolimits_2(\Id-\be A)^{-1/2}
 -\det\nolimits_2(\Id-\be B)^{-1/2}\right|
 \leq C\norm{A-B}_{\HS}.
 \label{eq:determinant-continuity-conclusion}
\end{equation}
\end{lemma}

\begin{proof}
Define
\begin{equation}
 \Phi_{\be}(A)
 \coloneqq-\frac12\log\det\nolimits_2(\Id-\be A)
 =\frac12\sum_{n=2}^\infty\frac{\be^n}{n}\Tr(A^n).
 \label{eq:determinant-Phi}
\end{equation}
The noncommutative telescoping identity
\begin{equation}
 A^n-B^n
 =\sum_{j=0}^{n-1}A^j(A-B)B^{n-1-j}
 \label{eq:determinant-power-telescoping}
\end{equation}
and the Hilbert--Schmidt trace inequality imply
\begin{equation}
 \abs{\Tr(A^n-B^n)}
 \leq nHL^{n-2}\norm{A-B}_{\HS}
 \qquad(n\geq2).
 \label{eq:determinant-trace-power-bound}
\end{equation}
Indeed, after cyclically moving $A-B$ to the first position, the remaining product has Hilbert--Schmidt norm at most $HL^{n-2}$.

Summing \eqref{eq:determinant-trace-power-bound} gives
\begin{equation}
 \abs{\Phi_{\be}(A)-\Phi_{\be}(B)}
 \leq\frac{\be_0^2H}{2(1-\be_0L)}
 \norm{A-B}_{\HS}.
 \label{eq:determinant-Phi-Lipschitz}
\end{equation}
Also,
\begin{equation}
 0\leq\Phi_{\be}(A),\Phi_{\be}(B)
 \leq\frac{\be_0^2H^2}{2(1-\be_0L)}.
 \label{eq:determinant-Phi-bound}
\end{equation}
The inequality $|e^x-e^y|\leq e^{\max\{x,y\}}|x-y|$ now proves \eqref{eq:determinant-continuity-conclusion}.
\end{proof}

We turn to the periodic logarithmic cutoff.  The normalization \eqref{eq:kernel-normalization} has the local refinement
\begin{equation}
 \g(x)=-\log|x|+\psi(x)
 \qquad(0<|x|<r_0),
 \label{eq:kernel-smooth-local-form}
\end{equation}
where $r_0>0$ and $\psi\in C^\infty(B(0,r_0))$ depend only on $\d$.  This is the standard local parametrix for the logarithmic Green kernel and is consistent with the normalization in Chodron de Courcel, Rosenzweig, and Serfaty \cite{ChodrondeCourcelRosenzweigSerfaty2025}.  In particular,
\begin{equation}
 \norm{\nabla\g(x)}\leq C|x|^{-1},
 \qquad
 \norm{\nabla^2\g(x)}\leq C|x|^{-2}
 \qquad(0<|x|<r_0).
 \label{eq:kernel-derivative-bounds}
\end{equation}

\begin{proof}[Proof of \Cref{prop:positive-logarithmic-regularization}]
The Fourier multiplier of $\nu_\eta$ is
\begin{equation}
 \wh\nu_\eta(k)=\abs{\wh\varphi_\eta(k)}^2\in[0,1].
 \label{eq:positive-mollifier-Fourier}
\end{equation}
Consequently,
\begin{equation}
 \wh\g_{\eta}(k)
 =\wh\g(k)\abs{\wh\varphi_\eta(k)}^2
 \qquad(k\neq0),
 \label{eq:regularized-kernel-Fourier}
\end{equation}
which proves smoothness, reality, evenness, mean zero, nonnegative definiteness, and the operator order \eqref{eq:regularized-operator-order}.

We next prove \eqref{eq:regularized-partition-domination}.  Let $U_1,\ldots,U_N$ be independent with density $\varphi_\eta$.  For fixed $(x_1,\ldots,x_N)$,
\begin{equation}
 \g_{\eta}(x_i-x_j)
 =\E\left[\g(x_i+U_i-x_j-U_j)\right].
 \label{eq:regularized-pair-shift}
\end{equation}
Apply conditional Jensen's inequality to the exponential of the sum in \eqref{eq:regularized-pair-shift}, then integrate in the $x_i$.  Haar invariance shows that $(x_i+U_i)_{i=1}^N$ are again independent with law $\muunif$, and hence
\begin{equation}
 \E_{\muunif^{\otimes N}}\left[e^{\be\mathsf W_N(\g_{\eta})/2}\right]
 \leq\E_{\muunif^{\otimes N}}\left[e^{\be\mathsf W_N(\g)/2}\right].
 \label{eq:regularized-Jensen-domination}
\end{equation}
To justify conditional Jensen at the singularity, first replace $\g$ by $\g\wedge L$.  Both sides then involve bounded kernels.  Since $\g$ is bounded below, monotone convergence as $L\to\infty$ gives \eqref{eq:regularized-Jensen-domination} for the original kernels.

A centered feature representation of the smooth kernel \eqref{eq:regularized-kernel-Fourier} has constant diagonal
\begin{equation}
 D_{\eta}
 =\sum_{k\neq0}\wh\g(k)
 \abs{\wh\varphi_\eta(k)}^2
 =\g_{\eta}(0)\geq0.
 \label{eq:regularized-diagonal-Fourier}
\end{equation}
Since $\nu_\eta$ is the rescaling of a fixed smooth probability density supported in a ball of radius $O(\eta)$, \eqref{eq:kernel-smooth-local-form} gives
\begin{equation}
 \g_{\eta}(0)
 =-\log\eta-
 \int_{\R^{\d}}\log|z|\,\nu_1(z)\,\dd z+O_{\d}(1),
 \label{eq:regularized-diagonal-asymptotic}
\end{equation}
which proves \eqref{eq:regularized-feature-diagonal}.

It remains to establish \eqref{eq:regularization-Lu-rate}.  Write
\begin{equation}
 r(x)\coloneqq\operatorname{dist}_{\T^{\d}}(x,0).
 \label{eq:regularization-torus-distance}
\end{equation}
On $r(x)\leq\eta/2$, rescale $x=\eta z$ in the standard lift.  The logarithmic parts of $\g(x)$ and $\g_{\eta}(x)$ differ by a fixed rescaled logarithmic potential, while the smooth remainder is uniformly bounded.  Therefore,
\begin{equation}
 \abs{\f_{\eta}(x)}
 \leq C\left(1+
 \abs*{\log\frac{r(x)}\eta}\right)
 \qquad(0<r(x)\leq\eta/2).
 \label{eq:regularization-near-pointwise}
\end{equation}
It follows that
\begin{equation}
 \int_{r(x)\leq\eta/2}\abs{\f_{\eta}(x)}^u\,\dd\muunif(x)
 \leq C_u\eta^{\d}.
 \label{eq:regularization-near-integral}
\end{equation}

On $\eta/2\leq r(x)\leq r_0/2$, evenness of $\nu_\eta$ gives zero first moment.  Taylor's formula and \eqref{eq:kernel-derivative-bounds} yield
\begin{equation}
 \abs{\f_{\eta}(x)}
 \leq C\eta^2r(x)^{-2}.
 \label{eq:regularization-far-pointwise}
\end{equation}
Consequently, since $2u>\d$,
\begin{align}
 \int_{\eta/2\leq r(x)\leq r_0/2}
 \abs{\f_{\eta}(x)}^u\,\dd\muunif(x)
 &\leq C_u\eta^{2u}
 \int_{\eta/2}^{r_0/2}r^{\d-1-2u}\,\dd r
 \notag\\
 &\leq C_u\eta^{\d}.
 \label{eq:regularization-far-integral}
\end{align}
Away from the origin, smoothness gives $\abs{\f_{\eta}}\leq C\eta^2$, whose $u$th power is $O(\eta^{2u})=O(\eta^{\d})$.  Combining this with \eqref{eq:regularization-near-integral} and \eqref{eq:regularization-far-integral}, and taking the $u$th root, proves \eqref{eq:regularization-Lu-rate}.
\end{proof}

\section{Cutoff-uniform integrated Wick expansions}
\label{app:higher-order-replica-expansions}

This appendix proves \Cref{thm:all-order-replica-expansions,cor:higher-order-divergences}.  Fix throughout $0\leq\be_0<\be_*<\beta_{\gm}$ as in \eqref{eq:higher-order-buffer}.  The sole many-particle input is the strict partition buffer \eqref{eq:higher-order-scalar-buffer}, supplied by \Cref{prop:compact-subcritical-partition-bound}.  We reuse the regularization preserving nonnegative definiteness in \Cref{prop:positive-logarithmic-regularization} and its domination \eqref{eq:regularized-partition-domination}.

\subsection{Kernel integrability and the pair response}

\begin{lemma}[Logarithmic moments and higher-order coefficients]\label{lem:kernel}
There are constants $C_0,C_1<\infty$ such that, for every integer $j\geq1$,
\begin{equation}\label{eq:factorial}
 \int |\g|^j\,\dd\muunif\leq C_0 C_1^j j!.
\end{equation}
Moreover, $b_\be\in L^p(\muunif)$ for every finite $p$, uniformly for $0\leq\be\leq\be_0$, and every term in \eqref{eq:higher-order-L} is finite.
\end{lemma}
\begin{proof}
Choose $0<a<\d$.  The logarithmic singularity and boundedness away from it give $\int e^{a|\g|}\,\dd\muunif<\infty$.  The pointwise bound $a^j|\g|^j/j!\leq e^{a|\g|}$ proves \eqref{eq:factorial}.  The decomposition \eqref{eq:b-split} and bound \eqref{eq:s-bound} imply that $b_\be$ is the sum of $\be\g$ and a bounded continuous function, uniformly on the temperature window.  Thus, $b_\be\in L^p$ for every finite $p$.  Cauchy--Schwarz gives $\|b_\be*b_\be\|_\infty\leq\|b_\be\|_2^2$, and further convolution preserves boundedness.  Finally, $B_\be^j$ is trace class for every $j\geq2$.  This proves the claim.
\end{proof}

\subsection{Integrated Wick powers and cutoff-uniform power series}

Let $\phi_\eta$ be a centered real Gaussian field of covariance $c \g_\eta$, where $c$ ranges over a fixed bounded nonnegative interval. Its spatial integral vanishes almost surely. Write $v_\eta=c\g_\eta(0)$, and let
\begin{equation}\label{eq:Wick}
 H_{j,\eta}=\int :\!\phi_\eta(x)^j\!:\,\dd \muunif(x),\qquad
 A_\eta(t)=\int e^{t\phi_\eta(x)-t^2v_\eta/2}\,\dd \muunif(x).
\end{equation}
Wick ordering is with respect to the original covariance $c\g_\eta$. In particular, $H_{1,\eta}=0$.

\begin{lemma}[Uniform Gaussian power series]\label{lem:Wickseries}
For every fixed $p\geq2$, there is $K_p<\infty$, independent of $\eta,c,j$, such that
\begin{equation}\label{eq:Wickbound}
 \left\|\frac{H_{j,\eta}}{j!}\right\|_{L^p(G)}\leq K_p^j,
 \qquad j\geq2.
\end{equation}
For $|t|\leq(2K_p)^{-1}$,
\begin{equation}\label{eq:Wickseries}
 A_\eta(t)=1+\sum_{j\geq2}\frac{t^j}{j!}H_{j,\eta}
 \quad\hbox{in }L^p(G).
\end{equation}
For every fixed integer $J\geq2$, its tail after degree $J$ is $O_{p,J}(|t|^{J+1})$ in $L^p$, uniformly in the cutoff. In particular, $\|A_\eta(t)-1\|_p=O_p(t^2)$.
\end{lemma}
\begin{proof}
The covariance identity for Wick powers gives
\begin{equation}\label{eq:Wickvariance}
 \E[H_{j,\eta}^2]
 =j!\iint (c \g_\eta(x-y))^j\,\dd \muunif(x)\dd \muunif(y).
\end{equation}
The right side is nonnegative; to estimate it, take the absolute value inside the integral. Gaussian hypercontractivity on the $j$th homogeneous Wiener chaos gives
\begin{equation}\label{eq:hyper}
 \|H_{j,\eta}\|_p\leq(p-1)^{j/2}\|H_{j,\eta}\|_2.
\end{equation}
The inequality is the standard hypercontractivity estimate; an explicit statement is \cite[Proposition~2.6]{NourdinPeccatiReinert2010}. Combining \eqref{eq:factorial}, convolution contraction, and \eqref{eq:Wickvariance} yields
\begin{equation}\label{eq:geometric}
 \left\|\frac{H_{j,\eta}}{j!}\right\|_p
 \leq (p-1)^{j/2}
       \left(\frac{\int|c \g_\eta|^j\,\dd \muunif}{j!}\right)^{1/2}
 \leq K_p^j.
\end{equation}
For fixed $\eta$, the Hermite generating series is the Wick exponential in \eqref{eq:Wick}. It can first be integrated using its $L^p$ convergence at that fixed cutoff. The geometric estimate then proves the cutoff-uniform series and tail bound for the indicated range of $t$. This argument does not bound the field at a point uniformly in $\eta$.
\end{proof}

\step{Cutoff convergence of the coefficients} Use the even mollifier $\varphi$ fixed in \eqref{eq:positive-logarithmic-regularization}. Couple the regularized fields using the same standard real Fourier Gaussian variables, with mode amplitudes $\sqrt{c\widehat{\g}(n)}\,\widehat{\varphi_\eta}(n)$ for $n\ne0$; the mollifier multipliers are real and retain their signs. Each field then has covariance $c\g_\eta$, and the mixed covariance for two cutoffs is $c\g*\varphi_\eta*\varphi_\delta$. Formula \eqref{eq:Wickvariance}, applied to the mixed covariance, shows that every fixed $H_{j,\eta}$ is Cauchy in $L^2(G)$ as $\eta\to0$. Indeed, the mixed covariances converge in $L^j$ to $c\g$, and powers converge in $L^1$ by H\"older. Hypercontractivity upgrades this convergence to every finite $L^p$. The statement also applies to sums of any fixed number of independent fields and to finite products of the resulting integrated Wick powers. Only finite products of integrated Wick powers are passed to the limit in this coefficient construction.

\subsection{Fixed-integer replicas under a strict partition buffer}

Fix $q\geq2$. Let $\phi_1,\ldots,\phi_q$ be independent fields of covariance $\be \g_\eta$. Define $A_i(t)$ as in \eqref{eq:Wick} for $\phi_i$, and define $A_\Sigma(t)$ using $\phi_1+\cdots+\phi_q$, with Wick variance $q\be \g_\eta(0)$. In this section $t=N^{-1/2}$.

The unnormalized integral for $q$ replicas sharing exactly $k$ labels is
\begin{equation}\label{eq:qHS}
 J_{q,N,k}^{\eta}
 =\E_G\left[\left(\prod_{i=1}^q A_i(t)\right)^{N-k}A_\Sigma(t)^k\right],
 \qquad
 \int(F_{N,\be}^{(k),\eta})^q\,\dd \muunif^{\otimes k}
 =\frac{J_{q,N,k}^{\eta}}{(Z_{N,\be}^{(\eta)})^q}.
\end{equation}
To verify the identity, integrate each Gaussian field. Each pair in each replica contributes $\be \g_\eta/N$; every shared--shared pair is counted $q$ times. The subtractions $\be \g_\eta(0)/(2N)$ in each single-field factor and $q\be \g_\eta(0)/(2N)$ in the shared factor remove exactly the self-interactions of all replicas.

\begin{lemma}[Uniform exponential bounds]\label{lem:qbuffer}
Under \eqref{eq:higher-order-qmargin}, there is $r>1$ such that both
\begin{equation}\label{eq:qbuffer}
 \sup_{N\geq N_0}\sup_{k\leq\alpha N}\sup_{\be\leq\be_0}\sup_\eta
 \E_G\left[\left(\left(\prod_{i=1}^q A_i(t)\right)^{N-k}
 A_\Sigma(t)^k\right)^r\right]<\infty
\end{equation}
and the corresponding unnormalized particle overlap moment with coupling $r\be$ are finite uniformly in the same parameters.
\end{lemma}
\begin{proof}
Choose $r>1$ so that $r[1+(q-1)\alpha]\be_0<\be_*$. For $\theta=k/N>0$, apply H\"older with conjugate exponents
\begin{equation}\label{eq:Holder}
 a=\frac{1+(q-1)\theta}{1-\theta},\qquad
 a'=\frac{1+(q-1)\theta}{q\theta}.
\end{equation}
The unshared powers become $ra(N-k)=r[N+(q-1)k]$, while the shared power becomes $ra'k=r[N+(q-1)k]/q$. For every integer $L\geq1$, Gaussian integration gives
\begin{equation}\label{eq:integer-moments}
 \E_G[A_i(t)^L]=Z_{L,\be L/N}^{(\eta)},\qquad
 \E_G[A_\Sigma(t)^L]=Z_{L,q\be L/N}^{(\eta)}.
\end{equation}
Round the required moment exponents upward to integers and use Lyapunov's inequality. Both scalar couplings are at most
\begin{equation}\label{eq:effective}
 r\be[1+(q-1)\theta]+\frac{q\be}{N}.
\end{equation}
For sufficiently large $N$, this is below $\be_*$. Independence of the $A_i$ and \eqref{eq:regularized-partition-domination} prove \eqref{eq:qbuffer}. For $k=0$, independence and \eqref{eq:integer-moments} suffice. Repeat the same argument with Gaussian moment exponent $1$ and coupling $r\be$ to get the amplified particle integral. The rounding error is again absorbed by the strict margin.
\end{proof}

For fixed $q$ and $k$, the effective coupling in \eqref{eq:effective} tends to $\be$ as $N\to\infty$. For a positive limiting fraction $k/N$, it tends to $[1+(q-1)\theta]\be$. Thus increasing the replica order restricts the admissible overlap window, rather than requiring a scalar bound at $q\be$ for the original system.

\subsection{The cutoff-uniform expansion}

For a field indexed by $s$ (one of the individual fields or their sum), define polynomial coefficients by the formal identity
\begin{equation}\label{eq:logcoeff}
 \log\left(1+\sum_{j\geq2}\frac{t^jH_{j,s}}{j!}\right)
 =\sum_{n\geq0}t^{n+2}K_{n,s}.
\end{equation}
At each fixed order this is finite algebra. More explicitly,
\begin{equation}\label{eq:logcoeff-explicit}
 K_{n,s}=\sum_{a=1}^{\lfloor(n+2)/2\rfloor}\frac{(-1)^{a+1}}a
 \sum_{\substack{j_1,\ldots,j_a\geq2\\j_1+\cdots+j_a=n+2}}
 \prod_{i=1}^a\frac{H_{j_i,s}}{j_i!}.
\end{equation}
Set
\begin{equation}\label{eq:Lcoeff}
 L_n(\theta)=(1-\theta)\sum_{i=1}^qK_{n,i}+\theta K_{n,\Sigma},
\end{equation}
and define $P_n(\theta)$ by
\begin{equation}\label{eq:Bell}
 P_n(\theta)=
 \sum_{\substack{a_1,\ldots,a_n\geq0\\a_1+2a_2+\cdots+na_n=n}}
 \prod_{j=1}^n\frac{L_j(\theta)^{a_j}}{a_j!},\qquad P_0=1.
\end{equation}
All variables here are at the finite cutoff until the limit is specified.

\begin{lemma}[Expansion under a higher-moment bound for the exponential]\label{lem:expansion}
For every fixed integer $K\geq0$,
\begin{equation}\label{eq:expansion-half}
 J_{q,N,k}^{\eta}
 =\sum_{n=0}^{K} t^n\E_G[e^{L_0(\theta)}P_n(\theta)]
   +O_K(t^{K+1}),\qquad t=N^{-1/2},\quad\theta=k/N,
\end{equation}
uniformly in the cutoff and the parameters of \Cref{lem:qbuffer}. Every coefficient with $n$ odd is zero.
\end{lemma}
\begin{proof}
The case $K=0$ follows from the proof for $K=1$, so take $K\geq1$. Put $z_s=A_s(t)-1$ and let $E_t$ be the event $\max_s|z_s|\leq1/2$. By \Cref{lem:Wickseries}, for every fixed $p<\infty$,
\begin{equation}\label{eq:good-event}
 \|z_s\|_p\leq C_pt^2,\qquad
 \Pb(E_t^c)\leq C_pt^{2p}.
\end{equation}
We may choose $p$ as large as needed; its value is fixed before $N$ grows.

On $E_t$, Taylor's formula for $\log(1+z_s)$ gives
\begin{equation}\label{eq:log-error}
 \left|\log(1+z_s)-\sum_{a=1}^{A}\frac{(-1)^{a+1}}a z_s^a\right|
 \leq C_A|z_s|^{A+1}.
\end{equation}
Choose $A$ so that $2A\geq K+1$. After division by $t^2$, \eqref{eq:good-event} makes this error $O(t^{K+1})$ in every prescribed finite $L^p$, with the indicator of $E_t$. Replace $z_s$ in the finite polynomial by its series through degree $K+2$. Its omitted tail is $O(t^{K+3})$ in every finite $L^p$ by \Cref{lem:Wickseries}. H\"older's inequality applied to the finitely many products shows that their resulting error, after division by $t^2$, is $O(t^{K+1})$. The remaining terms of degree at most $K$ are exactly \eqref{eq:logcoeff}.

Let
\begin{equation}\label{eq:exactX}
 X=\frac1{t^2}\left[(1-\theta)\sum_{i=1}^q\log A_i(t)+\theta\log A_\Sigma(t)\right].
\end{equation}
Thus, $e^X$ is the exact Gaussian overlap integrand. On $E_t$,
\begin{equation}\label{eq:Delta}
 \left\|\mathbf1_{E_t}\left(X-\sum_{n=0}^{K}t^nL_n(\theta)\right)\right\|_p
 \leq C_pt^{K+1},\qquad
 \|\mathbf1_{E_t}(X-L_0)\|_p\leq C_pt.
\end{equation}
The constants are uniform in $\theta\in[0,\alpha]$, because all weights in \eqref{eq:exactX} are bounded and there are only $q+1$ fields.

We also need an exponential moment of $L_0$. Its quadratic form has replica eigenvalues $1+(q-1)\theta$ and $1-\theta$ (the latter with multiplicity $q-1$). Gaussian integration therefore gives, for $r>1$ sufficiently close to one,
\begin{equation}\label{eq:quadratic-moment}
 \E_G[e^{rL_0(\theta)}]
 =D_{r\be(1+(q-1)\theta),\eta}
  D_{r\be(1-\theta),\eta}^{q-1}\leq C.
\end{equation}
Here, $D_{c,\eta}=\det\nolimits_2(\Id-cT_\eta)^{-1/2}$. The bound follows from the strict spectral margin and a uniform Hilbert--Schmidt norm: the trace series for the logarithm starts at degree two. By decreasing $r$ if necessary, both \eqref{eq:quadratic-moment} and \eqref{eq:qbuffer} hold.

For real $x,y$, the integral remainder for the exponential gives
\begin{equation}\label{eq:expTaylor}
 \left|e^x-e^y\sum_{j=0}^{K}\frac{(x-y)^j}{j!}\right|
 \leq C_K|x-y|^{K+1}(e^x+e^y).
\end{equation}
Apply this on $E_t$ with $x=X$, $y=L_0$, and use H\"older with \eqref{eq:Delta}, \eqref{eq:qbuffer}, and \eqref{eq:quadratic-moment}. The expectation of the remainder is $O(t^{K+1})$. Replacing $X-L_0$ by its truncated polynomial in \eqref{eq:Delta}, and then discarding powers above $K$, produces another $O(t^{K+1})$ error. To justify this explicitly, use
\begin{equation}\label{eq:power-difference}
 |u^j-v^j|\leq j|u-v|(|u|+|v|)^{j-1}
\end{equation}
and H\"older, with all moments chosen beforehand. Every coefficient polynomial has moments of every fixed order uniformly in the cutoff. The polynomial of degree at most $K$ is \eqref{eq:Bell}.

It remains to remove the indicator of $E_t$. H\"older and \eqref{eq:qbuffer} bound $\E[\mathbf1_{E_t^c}e^X]$ by a fixed constant times $\Pb(E_t^c)^{1-1/r}$. The same argument, with a slightly smaller exponential exponent and a further H\"older factor for the polynomial, bounds $\E[\mathbf1_{E_t^c}e^{L_0}|P_n|]$. Choosing $p$ in \eqref{eq:good-event} large enough makes all these errors $O(t^{K+1})$.

Finally, under the simultaneous sign reversal of all Gaussian fields, $H_{j,s}$ acquires the sign $(-1)^j$. Thus, $K_{n,s}$ and $L_n$ acquire $(-1)^n$, $P_n$ acquires $(-1)^n$, and $L_0$ is unchanged. Gaussian symmetry makes each odd coefficient vanish exactly.
\end{proof}

\begin{proof}[Proof of \Cref{thm:all-order-replica-expansions}]
Take $K=2M+1$ in \Cref{lem:expansion}. The odd terms vanish and the remainder is $O(N^{-M-1})$. The same proof with one field gives the scalar expansion.

All polynomial coefficient variables converge in every finite $L^p$ as the cutoff is removed, by the discussion after \Cref{lem:Wickseries}. The quadratic variables $L_{0,\eta}$ converge in every finite $L^p$. The inequality
\begin{equation}\label{eq:exponential-continuity}
 |e^x-e^y|\leq |x-y|(e^x+e^y)
\end{equation}
and the strict exponential bound imply convergence of $e^{L_{0,\eta}}$ in some $L^{1+\delta}$. H\"older then proves convergence of each coefficient expectation. The same estimates are uniform in $\be,\theta$ on the stated compact set, because, after writing each field as $\sqrt{\be}$ times a fixed Gaussian field, the polynomial coefficients are polynomial in $\sqrt{\be}$ and $\theta$, and the quadratic forms have a uniform spectral margin.

For each fixed $N,k$, convergence $\g_\eta\to \g$ in finite $L^p$ implies convergence in measure of the finite sum of particle interactions, and hence of its exponential. Its $L^r$ norm on the underlying product space is uniformly bounded by the amplified particle estimate in \Cref{lem:qbuffer}. Uniform integrability therefore gives $J_{q,N,k}^{\eta}\to J_{q,N,k}$ for each sufficiently large $N,k$. The scalar convergence follows in the same way from \eqref{eq:regularized-partition-domination}. The cutoff-independent remainder estimate consequently survives unchanged. Jensen gives $Z_{N,\be}\geq1$, so division by its fixed-$q$ power preserves the asymptotic expansion and its uniform remainder.

For completeness, write each field as $\sqrt{\be}$ times a field at unit coupling. Then $H_j=\be^{j/2}H_j^{(1)}$ and $K_n=\be^{(n+2)/2}K_n^{(1)}$. Every monomial $\prod_{n\geq1}L_n^{a_n}$ in $P_{2j}$ satisfies $\sum_n n a_n=2j$ and has the integer $\be$-degree $j+\sum_n a_n$. Thus these coefficient polynomials have no square-root singularity at $\be=0$. Since $L_0$ is linear in $\be$, the same spectral margin controls all fixed $\be$ derivatives of their expectations.

Gaussian integration of $e^{L_0}$ gives the numerator of \eqref{eq:higher-order-leading-replica}. Smoothness of the coefficient functions follows by differentiating $e^{L_0(\theta)}P_n(\theta)$: every derivative inserts only a fixed polynomial of quadratic Gaussian variables, and the same exponential margin applies. For analyticity near zero, allow $\theta$ to be complex in a sufficiently small disk. The modulus of the exponential is $e^{L_0(\Re\theta)}$; choose the disk so that the quadratic eigenvalues for its real parts still have a strict spectral margin. Polynomial factors and their derivatives remain integrable by H\"older. Dominated contour integration gives the asserted holomorphy, uniformly in $\be$.
\end{proof}

\subsection{Explicit corrections and fixed-marginal moments}

\subsubsection{The scalar coefficient}
At a fixed cutoff, tilt the Gaussian measure by $e^{H_2/2}/D_{\be,\eta}$. Its covariance is $b_{\be,\eta}$, where $B_{\be,\eta}=\be T_\eta(\Id-\be T_\eta)^{-1}$. If $M_j=\int\phi^j\,\dd \muunif$, the first relevant logarithmic coefficients are
\begin{equation}\label{eq:first-polys}
 K_1=\frac{M_3}{6},\qquad
 K_2=\frac{M_4}{24}-\frac{M_2^2}{8}.
\end{equation}
The Wick variance cancels from these formulas: $H_3=M_3$ because $\int\phi=0$, and substitution of $H_4=M_4-6vM_2+3v^2$ and $H_2=M_2-v$ cancels the other terms. Consequently,
\begin{equation}\label{eq:first-expect}
 z_{1,\eta}=\E_{b_{\be,\eta}}\left[\frac{M_4}{24}-\frac{M_2^2}{8}+\frac{M_3^2}{72}\right].
\end{equation}
Write $w=b_{\be,\eta}(0)$. Wick's formula gives
\begin{align}
 \E[M_4]&=3w^2,\label{eq:M4}\\
 \E[M_2^2]&=w^2+2\int b_{\be,\eta}^2\,\dd \muunif,\label{eq:M22}\\
 \E[M_3^2]&=\iint\left[9w^2 b_{\be,\eta}(x-y)+6b_{\be,\eta}(x-y)^3\right]\dd \muunif(x)\dd \muunif(y)
 =6\int b_{\be,\eta}^3\,\dd \muunif.\label{eq:M32}
\end{align}
The zero Fourier mode vanishes in the last line. Substitution cancels $w^2$ exactly and gives \eqref{eq:higher-order-z1} after cutoff removal. Convergence of $b_{\be,\eta}$ in every finite $L^p$ follows from the Fourier multiplier formula for $B_{\be,\eta}$, the decomposition \eqref{eq:b-split}, and dominated convergence of the absolutely summable bounded correction.

\subsubsection{The first replica correction}
For this calculation only, write $b=b_\be$. Under the quadratic tilted law of $q$ replicas, let $u_\theta$ be the same-replica covariance kernel, $v_\theta$ the covariance between different replicas, and $s_\theta$ the covariance between one replica and the sum of the fields. Spectral diagonalization gives the convolution operators
\begin{align}
 U_\theta&=\frac1q\left[B(\Id-(q-1)\theta B)^{-1}
            +(q-1)B(\Id+\theta B)^{-1}\right],\label{eq:U}\\
 V_\theta&=\frac1q\left[B(\Id-(q-1)\theta B)^{-1}
            -B(\Id+\theta B)^{-1}\right],\label{eq:V}\\
 S_\theta&=B(\Id-(q-1)\theta B)^{-1}=U_\theta+(q-1)V_\theta.\label{eq:Sop}
\end{align}
The field sum has covariance $q s_\theta$. These kernels are mean-zero. The difference of each of $u_\theta,s_\theta$ from $b$ is bounded near $\theta=0$, and $v_\theta$ is bounded there.

The first relative correction to the unnormalized replica determinant is
\begin{align}\label{eq:cq}
 c_q(\theta)
 &=-\frac14\left[q(1-\theta)\int u_\theta^2+q^2\theta\int s_\theta^2\right]\notag\\
 &\quad+\frac1{12}\left[
 q(1-\theta)^2\int u_\theta^3
 +q(q-1)(1-\theta)^2\int v_\theta^3
 +(2q\theta(1-\theta)+q^3\theta^2)\int s_\theta^3\right].
\end{align}
Here and below in this calculation, $\dd \muunif$ is suppressed in spatial integrals. To verify \eqref{eq:cq}, apply \eqref{eq:first-polys} to the $q$ individual fields and their sum, with weights $1-\theta$ and $\theta$. For the quartic logarithmic coefficient, the single-field identity is $-\frac14\int C^2$ for covariance $C$. For the square of the cubic coefficient, two jointly Gaussian mean-zero fields with cross covariance $C_{ab}$ contribute $\frac1{12}\int C_{ab}^3$. The apparent additional term proportional to the product of their pointwise variances integrates to zero, because $\int C_{ab}=0$. Summing the diagonal, distinct-replica, replica--sum, and sum--sum contributions gives the four terms in \eqref{eq:cq}.

Normalization by $Z_{N,\be}^q$ gives
\begin{equation}\label{eq:Rq1}
 \mathcal R_{q,1}(\be,\theta)
 =\mathcal R_{q,0}(\be,\theta)[c_q(\theta)-qz_1(\be)].
\end{equation}
For $q=1$, this correction is identically zero, as required.

To expand near zero, put $b_2=b*b$ and $b_3=b*b*b$ within this calculation. Resolvent expansion gives
\begin{align}
 u_\theta&=b+(q-1)\theta^2b_3+O(\theta^3),\label{eq:u-expand}\\
 v_\theta&=\theta b_2+(q-2)\theta^2b_3+O(\theta^3),\label{eq:v-expand}\\
 s_\theta&=b+(q-1)\theta b_2+(q-1)^2\theta^2b_3+O(\theta^3).\label{eq:s-expand}
\end{align}
The remainders hold in every fixed $L^p$. In particular,
\begin{align}
 \int u_\theta^2&=\mathcal C_\be+2(q-1)\theta^2\Tr(B^4)+O(\theta^3),\label{eq:u2}\\
 \int s_\theta^2&=\mathcal C_\be+2(q-1)\theta\Tr(B^3)+O(\theta^2),\label{eq:s2}\\
 \int u_\theta^3&=\int b^3+3(q-1)\theta^2\int b^2b_3+O(\theta^3),\label{eq:u3}\\
 \int v_\theta^3&=O(\theta^3),\label{eq:v3}\\
 \int s_\theta^3&=\int b^3+3(q-1)\theta\int b^2b_2+O(\theta^2).\label{eq:s3}
\end{align}
Substitution into \eqref{eq:cq} yields
\begin{align}\label{eq:cq-expand}
 c_q(\theta)-qz_1
 &=-\frac{q(q-1)}4\mathcal C_\be\theta\notag\\
 &\quad+\left[\binom q2\frac{\mathcal L_\be}{2}
 +\binom q3\left(\frac12\int b_\be^3-3\Tr(B_\be^3)\right)\right]\theta^2
 +O(\theta^3).
\end{align}
All coefficients in this calculation are uniform in $\be\leq\be_0$.

\subsubsection{The exact one-particle marginal determines the remaining terms}
The Gaussian determinant has the trace expansion
\begin{equation}\label{eq:Rq-log}
 \log\mathcal R_{q,0}(\be,\theta)
 =\sum_{j\geq2}\frac{(q-1)^j+(q-1)(-1)^j}{2j}\,
 \theta^j\Tr(B_\be^j).
\end{equation}
Thus,
\begin{equation}\label{eq:Rq-leading-expand}
 \mathcal R_{q,0}
 =1+\frac12\binom q2\mathcal C_\be\theta^2
    +\binom q3\Tr(B_\be^3)\theta^3+O(\theta^4).
\end{equation}
For every $j\geq1$, $\mathcal R_{q,j}(\be,0)=0$. This follows successively from the exact identity at $k=0$ and uniqueness of coefficients in an asymptotic expansion. The identity at $k=1$, where $F^{(1)}=1$ exactly, gives
\begin{align}
 \mathcal R_{q,1}'(\be,0)&=-\frac12\binom q2\mathcal C_\be,\label{eq:k1-second}\\
 \mathcal R_{q,2}'(\be,0)&=-\binom q3\Tr(B_\be^3)-\frac12\mathcal R_{q,1}''(\be,0).\label{eq:k1-third}
\end{align}
These identities are derived at each fixed $\be$ by multiplying the expansion at $k=1$ by the relevant power of $N$ and taking the limit. No analytic interpolation in the finite-particle label $k$ is used.

\begin{prop}[Integer moments through order three]\label{prop:moments-next}
For fixed integers $q,k\geq2$, under \eqref{eq:higher-order-buffer},
\begin{align}\label{eq:Mq-fixed}
 \int(F_{N,\be}^{(k)})^q\,\dd \muunif^{\otimes k}
 &=1+\binom q2\binom k2\frac{\mathcal C_\be}{N^2}\notag\\
 &\quad+\frac1{N^3}\left\{
 \binom q2\binom k2\mathcal L_\be
 +\binom q3\left[\binom k2\int b_\be^3\,\dd \muunif
              +6\binom k3\Tr(B_\be^3)\right]\right\}
 +O_{q,k}(N^{-4}).
\end{align}
The error is uniform for $\be\leq\be_0$.
\end{prop}
\begin{proof}
Choose $\alpha>0$ satisfying \eqref{eq:higher-order-qmargin}. Apply \Cref{thm:all-order-replica-expansions} through order three, and expand its smooth coefficients at $\theta=k/N$. Use \eqref{eq:Rq1}, \eqref{eq:cq-expand}, \eqref{eq:Rq-leading-expand}, \eqref{eq:k1-second}, and \eqref{eq:k1-third}. If $A$ denotes the coefficient of $\theta^2$ in \eqref{eq:Rq-leading-expand}, $D$ the coefficient of $\theta^3$, and $E$ the coefficient of $\theta^2$ in \eqref{eq:cq-expand}, the result is
\begin{equation}\label{eq:fixed-intermediate}
 1+\frac{A(k^2-k)}{N^2}+\frac{D(k^3-k)+E(k^2-k)}{N^3}+O_{q,k}(N^{-4}).
\end{equation}
Insert the displayed values and use $k^3-3k^2+2k=6\binom k3$. This gives \eqref{eq:Mq-fixed}.
\end{proof}

\begin{proof}[Proof of \Cref{cor:higher-order-divergences}]
Chi-square is the $q=2$ moment minus one, which proves \eqref{eq:higher-order-chi}. Arbitrary fixed-order expansions follow by the same argument, using the arbitrary fixed-order expansion in \Cref{thm:all-order-replica-expansions}; at inverse order $j$, the coefficient is a polynomial in fixed $k$ of degree at most $j$.

For entropy, put $u=F_{N,\be}^{(k)}-1$. Use \eqref{eq:Mq-fixed} for $q=2,3,4$, together with the trivial moments for $q=0,1$. The binomial formulas for $(F-1)^3$ and $(F-1)^4$ give
\begin{align}
 \int u^3\,\dd \muunif^{\otimes k}
 &=\frac1{N^3}\left[\binom k2\int b_\be^3\,\dd \muunif
              +6\binom k3\Tr(B_\be^3)\right]+O_k(N^{-4}),\label{eq:central3}\\
 \int u^4\,\dd \muunif^{\otimes k}&=O_k(N^{-4}).\label{eq:central4}
\end{align}
The fourth finite difference annihilates the polynomials $\binom q2$ and $\binom q3$ in \eqref{eq:Mq-fixed}. This proves \eqref{eq:central4}, including a nonnegative upper bound, because the left side is nonnegative.

For every real $u\geq-1$,
\begin{equation}\label{eq:entropy-fourth}
 \left|(1+u)\log(1+u)-u-\frac{u^2}{2}+\frac{u^3}{6}\right|\leq C u^4.
\end{equation}
On $[-1/2,1/2]$, this is Taylor's theorem with bounded fourth derivative. On $[-1,-1/2]\cup[1/2,1]$, divide by $u^4$ and use continuity, interpreting $0\log0=0$. For $u\geq1$, $\log(1+u)\leq u$ bounds each term by a constant times $u^4$. Thus, the estimate is global. Since $\int u=0$, integration gives
\begin{equation}\label{eq:entropy-identity}
 \Hr=\frac12\chi^2-\frac16\int u^3\,\dd \muunif^{\otimes k}+O_k(N^{-4}).
\end{equation}
Equations \eqref{eq:higher-order-chi} and \eqref{eq:central3} prove \eqref{eq:higher-order-entropy}. No Taylor expansion uniform in the particle configuration is used.
\end{proof}

The cubic coefficient has a direct combinatorial interpretation. For $S_{\be,k}=\sum_{i<j}b_\be(x_i-x_j)$, canonicality eliminates every three-edge graph with a leaf. The surviving graphs are a triple edge and a triangle, so
\begin{equation}\label{eq:S-cube}
 \int S_{\be,k}^3\,\dd \muunif^{\otimes k}
 =\binom k2\int b_\be^3\,\dd \muunif+6\binom k3\Tr(B_\be^3).
\end{equation}
The coefficient $6$ counts the orderings of the three distinct triangle edges. This independently explains \eqref{eq:central3}.

\subsection{The marginal fraction as a separate expansion parameter}

At two replicas, \eqref{eq:Rq-log} gives
\begin{align}\label{eq:overlap-series}
 \mathcal R_{2,0}(\be,\theta)-1
 &=\frac{\mathcal C_\be}{2}\theta^2
 +\left[\frac{\mathcal C_\be^2}{8}+\frac{\Tr(B_\be^4)}4\right]\theta^4\notag\\
 &\quad+\left[\frac{\mathcal C_\be^3}{48}
             +\frac{\mathcal C_\be\Tr(B_\be^4)}8
             +\frac{\Tr(B_\be^6)}6\right]\theta^6
 +O(\theta^8).
\end{align}
This is an expansion of the limiting overlap determinant. It does not by itself supply the finite-$N$ coefficients in \eqref{eq:higher-order-chi}. The absence of a cubic power of $\theta$ in \eqref{eq:overlap-series} is compatible with a nonzero order-$N^{-3}$ coefficient at fixed $k$, because the latter comes from $\mathcal R_{2,1}$ and $\mathcal R_{2,2}$.

For growing $k$, the coefficient functions in \eqref{eq:higher-order-replica-expansion} preserve the dependence on $k/N$. The quartic overlap term has size $k^4/N^4$, whereas the first finite-size correction displayed in \eqref{eq:higher-order-chi} has size $k^2/N^3$. Their ratio is $k^2/N$, so their ordering changes around $k\asymp\sqrt N$. The exact pair-count factor is $(k^2-k)/(2N^2)$; replacing it by $k^2/(2N^2)$ introduces a separate error $k/N^2$. Each of these contributions must be compared with the remainder for the specified sequence $k=k_N$.

\section{Cutoff-uniform marked Wick expansions under a strict partition buffer}
\label{app:prescribed-marked-technical}

This appendix supplies the analytic expansion used in \Cref{thm:prescribed-marked-expansion}.  The notation is that of \Cref{subsec:prescribed-cutoffs,subsec:prescribed-marked-identity}.  Every estimate is uniform in the finite cutoff, the reference family, and $0\leq\be\leq\be_0$.  The order of expansion, number of marks, and spatial moment exponent are fixed unless an explicit dependence on $k$ is stated.

\subsection{Weighted integrated Wick powers}

For integers $r,l\geq0$, define
\begin{equation}
 W_{r,l,j}(u)
 \coloneqq\frac1{r!l!}\int u(y)^l:\!\phi_j(y)^r\!:\,\dd\mu(y),
 \label{eq:prescribed-weighted-Wick}
\end{equation}
where Wick ordering uses the covariance $c_j$ and $:\!\phi^0\!:=1$.  Put
\begin{equation}
 A_{n,j}(u)\coloneqq\sum_{r+2l=n}W_{r,l,j}(u).
 \label{eq:prescribed-A-coefficients}
\end{equation}

\begin{lemma}[Weighted power series]
\label{lem:prescribed-weighted-Wick}
For every fixed $p\geq2$ and $k$, there is $K_{p,k}<\infty$ such that
\begin{equation}
 \mathop{\rm ess\,sup}_{x}
 \norm{W_{r,l,j}(u_{j,x})}_{L^p(G)}
 \leq K_{p,k}^{r+2l},
 \qquad r+2l\geq1.
 \label{eq:prescribed-weighted-Wick-bound}
\end{equation}
For sufficiently small $|t|$, depending only on the fixed $p,k$,
\begin{equation}
 A_{j,u}(t)=1+\sum_{n\geq2}t^nA_{n,j}(u)
 \quad\text{in }L^p(G),
 \label{eq:prescribed-A-series}
\end{equation}
with the tail after degree $J$ bounded by $C_{p,k,J}|t|^{J+1}$, uniformly in the essential supremum over the marks.  Moreover,
\begin{equation}
 A_{1,j}(u)=0,
 \qquad
 A_{2,j}(u)=\frac12\int:\!\phi_j(y)^2\!:\,\dd\mu(y),
 \label{eq:prescribed-first-Wick-coefficients}
\end{equation}
so the second coefficient is independent of the centered source.
\end{lemma}

\begin{proof}
The covariance identity for Wick powers and the hypercontractivity estimate \eqref{eq:hyper} give
\begin{equation}
 \norm{W_{r,l,j}(u)}_p
 \leq\frac{(p-1)^{r/2}}{l!\sqrt{r!}}
 \left(\iint|c_j(y,z)|^r|u(y)|^l|u(z)|^l
 \,\dd\mu(y)\dd\mu(z)\right)^{1/2}.
 \label{eq:prescribed-weighted-hypercontractivity}
\end{equation}
For $r=0$, the same right side is understood as a deterministic integral estimate.  By \eqref{eq:prescribed-marked-source} and H\"older's inequality over the finitely many kernel rows, there is $a_k>0$ such that
\begin{equation}
 \mathop{\rm ess\,sup}_{x}\iint
 e^{a_k(|c_j(y,z)|+|u_{j,x}(y)|+|u_{j,x}(z)|)}
 \,\dd\mu(y)\dd\mu(z)\leq C_k.
 \label{eq:prescribed-weighted-source-exponential}
\end{equation}
Since $v^n\leq n!a^{-n}e^{av}$ for $v\geq0$, the integral in \eqref{eq:prescribed-weighted-hypercontractivity} is at most
\begin{equation}
 C_k r!(l!)^2a_k^{-r-2l}.
 \label{eq:prescribed-factorial-cancellation}
\end{equation}
The factorials cancel.  After enlarging the constant, this proves \eqref{eq:prescribed-weighted-Wick-bound}.  There are at most $n+1$ terms in \eqref{eq:prescribed-A-coefficients}, so the coefficient norms are geometric in $n$.  Expanding the Hermite and source exponentials at finite cutoff, multiplying the two absolutely convergent series in $L^p$, and using the geometric tail proves \eqref{eq:prescribed-A-series}.  Centering of the field and of $u$ gives \eqref{eq:prescribed-first-Wick-coefficients}.
\end{proof}

\subsection{Convergence of the coefficient variables}

Couple the cutoff fields by one isonormal Gaussian process on $L^2(\mu)$, using the operators $Q_j(\be T_\mu)^{1/2}$.  The mixed covariance between cutoffs $j$ and $l$ is the kernel of $\be Q_jT_\mu Q_l$.  It converges to $\be\tl\g_\mu$ in every finite $L^r$, and the mixed kernels satisfy the same row bounds as in \eqref{eq:prescribed-row-exponential}.

For fixed $r,l$, the $L^2(G\times\mu^{\otimes k})$ inner product of two variables in \eqref{eq:prescribed-weighted-Wick} is a finite integral of the two source weights multiplied by the $r$th power of the mixed covariance.  Each factor converges in arbitrarily high finite $L^p$ and has a uniform bound there.  H\"older's inequality therefore proves convergence of the inner products and hence the Cauchy property in joint $L^2$.  The higher moments from \Cref{lem:prescribed-weighted-Wick} upgrade this to every finite joint $L^p$ by interpolation.  Finite sums and products of the coefficient variables have the same convergence.  This argument is applied separately to each reference measure; the uniform bounds suffice to pass the uniform inequalities to the limit, and no uniform rate of conditional-expectation approximation over the full reference family is needed.

\subsection{Finite logarithmic polynomials and exact exponential control}

The finite logarithmic coefficients in \eqref{eq:prescribed-log-coefficients} are explicitly
\begin{equation}
 \ell_{n,j}(u)
 =\sum_{a=1}^{\lfloor n/2\rfloor}\frac{(-1)^{a+1}}a
 \sum_{\substack{n_1,\ldots,n_a\geq2\\n_1+\cdots+n_a=n}}
 \prod_{i=1}^aA_{n_i,j}(u).
 \label{eq:prescribed-ell-recursion}
\end{equation}
This is a finite polynomial identity at every prescribed order; it is not a claim that the random logarithmic series converges on the entire Gaussian space.

For scalar or fixed-integer replica expansions, write the logarithmic exponent as
\begin{equation}
 X=L_0+\sum_{n=1}^{K}t^nL_n+\mathcal R_{K,t},
 \qquad
 L_0=\ell_{2,j},
 \label{eq:prescribed-generic-log-exponent}
\end{equation}
with $t=N^{-1/2}$.  The precise $L_n$ are finite linear combinations of the $\ell_{n+2,j}$ associated with individual replica fields and their sum.  The common spectral gap gives an $r>1$ such that
\begin{equation}
 \sup_j\E_G[e^{rL_0}]<\infty.
 \label{eq:prescribed-quadratic-buffer}
\end{equation}
The shared-label replica estimate \eqref{eq:prescribed-replica-Lq} supplies the corresponding bound for the exact exponential.

Let
\begin{equation}
 E_t\coloneqq\{|A_{j,u}(t)-1|\leq1/2\}.
 \label{eq:prescribed-good-event}
\end{equation}
For every fixed $K,v$,
\begin{equation}
 \norm{\one_{E_t}\mathcal R_{K,t}}_{L^v(G)}\leq C_{K,v}t^{K+1},
 \qquad
 \norm{\one_{E_t}(X-L_0)}_{L^v(G)}\leq C_vt,
 \label{eq:prescribed-log-remainder}
\end{equation}
and, for every fixed $v$,
\begin{equation}
 \Pb_G(E_t^c)\leq C_vt^{2v}.
 \label{eq:prescribed-good-event-probability}
\end{equation}
These follow from \Cref{lem:prescribed-weighted-Wick}, the finite Taylor formula for the logarithm on the disc $|z-1|\leq1/2$, and Markov's inequality.

For real $x,y$ and integer $K\geq0$, the integral remainder for the exponential gives
\begin{equation}
 \left|e^x-e^y\sum_{j=0}^K\frac{(x-y)^j}{j!}\right|
 \leq C_K|x-y|^{K+1}(e^x+e^y).
 \label{eq:prescribed-global-exponential-remainder}
\end{equation}
Apply this on $E_t$ with $x=X$ and $y=L_0$.  Equations \eqref{eq:prescribed-log-remainder}, \eqref{eq:prescribed-quadratic-buffer}, and the exact-exponential bound, combined by H\"older, give expectation error $O(t^{K+1})$.  Replacing $X-L_0$ by its finite polynomial and discarding terms of degree above $K$ gives another error of the same order.  On $E_t^c$, conditional H\"older and \eqref{eq:prescribed-good-event-probability} give
\begin{equation}
 \E_G[\one_{E_t^c}e^X]
 \leq \Pb_G(E_t^c)^{1-1/r}(\E_G[e^{rX}])^{1/r},
 \label{eq:prescribed-bad-event-bound}
\end{equation}
and the probability power can be made arbitrarily large by choosing the fixed moment $v$.  A quadratic exponential multiplied by any fixed coefficient polynomial is treated similarly after decreasing the exponential margin and taking one higher polynomial moment.

Under simultaneous reversal of all Gaussian fields, $A_{n,j}$ and $\ell_{n,j}$ acquire the sign $(-1)^n$, while $L_0$ remains unchanged.  Consequently, every odd coefficient expectation in the scalar and marginal expansions without a test-function insertion vanishes.  Taking $K=2M+1$ yields an error $O(t^{2M+2})=O(N^{-M-1})$.  The argument uses only the exact exponential and the quadratic exponential; it never exponentiates a truncated cubic or higher Gaussian polynomial.

\subsection{The marked strong expansion}

\begin{prop}[Cutoff-uniform unnormalized marked remainder]
\label{prop:prescribed-marked-remainder}
Assume \eqref{eq:prescribed-scalar-buffer} and use the finite-cutoff objects in \eqref{eq:prescribed-marked-CM} and the coefficient polynomials in \eqref{eq:prescribed-marked-V}. For every pair of fixed integers $k,M\geq0$ and $1\leq p<\infty$, for sufficiently large $N$,
\begin{align}
 &\left\|e^{J_{k,j}/N}\E_G\!\left[A_{j,u_{j,x}}(N^{-1/2})^{N-k}\right]
 -\sum_{m=0}^{M}\frac{1}{N^m}
 \E_G\!\left[e^{\ell_{2,j}}\mathcal B_{2m}
 (V_{1,k,j},\ldots,V_{2m,k,j})\right]\right\|_{L^p(\mu^{\otimes k})}
 \notag\\
 &\hspace{35mm}\leq C_{k,M,p}N^{-M-1}.
 \label{eq:prescribed-marked-unnormalized-remainder}
\end{align}
Here, $k=0$ is the scalar case and $\mathcal B_0=1$. The constant is uniform in the finite cutoff, the reference family, and $0\leq\be\leq\be_0$. Each coefficient has a cutoff limit in every finite spatial $L^p$, equal to $D_{\be,\mu}U_{m,k}$ in \eqref{eq:prescribed-U-coefficients}. \end{prop}

This proposition estimates the unnormalized numerator. Scalar normalization in Section~\ref{sec:prescribed-marked-expansions} then gives the marginal expansion.

\begin{proof}
For the marked exponent, put
\begin{equation}
 X_x=(t^{-2}-k)\log A_{j,u_{j,x}}(t),
 \qquad
 L_0=\ell_{2,j}.
 \label{eq:prescribed-marked-X}
\end{equation}
On the good event $E_{t,x}=\{|A_{j,u_{j,x}}(t)-1|\leq1/2\}$, the preceding argument gives
\begin{equation}
 \left\|\one_{E_{t,x}}\left(
 X_x-L_0-\sum_{n=1}^Kt^n[\ell_{n+2,j}(u_{j,x})-k\ell_{n,j}(u_{j,x})]
 \right)\right\|_{L^v(G)}
 \leq C_{K,v,k}t^{K+1},
 \label{eq:prescribed-marked-log-error}
\end{equation}
uniformly in the essential supremum over $x$.  Conditional H\"older bounds the Gaussian expectation error by
\begin{equation}
 C_Kt^{K+1}\left[1+(\E_G[e^{rX_x}])^{1/r}\right].
 \label{eq:prescribed-conditional-error}
\end{equation}
The bad-event contribution has the same form with an arbitrarily high factor from its probability.  Taking the spatial $L^p$ norm after these Gaussian estimates and applying \Cref{lem:prescribed-anchored-moment} proves an $O_{K,k,p}(t^{K+1})$ remainder.  In particular, the proof does not take the spatial $p$th power inside the Gaussian expectation of the exact exponential.

The factor $e^{t^2J_{k,j}}$ in \eqref{eq:prescribed-marked-CM} is handled separately.  For fixed $k$, the row bound gives $a_k>0$ such that
\begin{equation}
 \int e^{a_k|J_{k,j}(x)|}\,\dd\mu^{\otimes k}(x)\leq C_k.
 \label{eq:prescribed-external-moments}
\end{equation}
Its Taylor remainder therefore has the required order in every fixed $L^p$.  Multiplying it with the Gaussian expansion uses H\"older in the marked variables and the already established estimate at a larger fixed spatial exponent.  Field reversal eliminates odd total powers.  The coefficient of $t^{2m}$ is exactly the expression in \eqref{eq:prescribed-U-coefficients}, and cutoff removal follows from the joint coefficient convergence above.  This proves \eqref{eq:prescribed-marked-unnormalized-remainder} and the asserted coefficient convergence.
\end{proof}

\section{Repulsive marked expansions by the oscillatory Hubbard--Stratonovich method}
\label{app:repulsive-marked-oscillatory}

We prove the one-component repulsive counterpart of the marked expansion on the same volume-one torus. The logarithmic kernel $\g$ and its Fourier normalization are unchanged. Repulsion is encoded by the physical interaction $W^-=\g$, in contrast with $W=-\g$ in the main text. Define
\begin{equation}
 h^-_{\mu}(x)\coloneqq\int_{\T^{\d}}\g(x-y)\,\dd\mu(y),
 \qquad
 I^-_{\mu}\coloneqq\frac12\int h^-_{\mu}\,\dd\mu,
 \qquad
 \psi^-_{\mu}\coloneqq h^-_{\mu}-2I^-_{\mu}.
 \label{rep:eq:normalization-crosswalk}
\end{equation}
Thus, relative to the attractive quantities from \eqref{eq:hmu-imu-definition} and \eqref{eq:prescribed-centered-source},
\begin{equation}
 h^-_{\mu}=-h_{\mu},
 \qquad
 I^-_{\mu}=-I_{\mu},
 \qquad
 \psi^-_{\mu}=-\psi_{\mu},
 \qquad
 q^-_{\mu}=\tl\g_{\mu},
 \qquad
 T^-_{\mu}=T_{\mu}.
 \label{rep:eq:sign-crosswalk}
\end{equation}
The repulsive resolvent is $(\Id+\be T_\mu)^{-1}$, and the weak one-body source is $+\be\psi^-_\mu$. The finite logarithmic and Bell-polynomial recursions have the same form as in \Cref{app:prescribed-marked-technical}. The new analytic step is control of the oscillatory exponential as the cutoff is removed. We obtain that control from a moving heat cutoff and uniform repulsive partition estimates.

\subsection{Repulsive normalization and all-order theorem}
\label{rep:subsec:one-component-theorem}

Throughout both appendices, the reference measure satisfies
\begin{equation}
 \mu=\rho\muunif,
 \qquad 0\leq\rho\leq M_0,
 \qquad \int_{\T^{\d}}\rho\,\dd\muunif=1.
 \label{rep:eq:mu-class}
\end{equation}
No positive lower bound for $\rho$ is assumed.


For the repulsive physical interaction \(W^-(x,y)=\g(x-y)\), let
\begin{equation}
 q^-_{\mu}(x,y)
 =\g(x-y)-h^-_{\mu}(x)-h^-_{\mu}(y)+2I^-_{\mu}
 =\tl\g_{\mu}(x,y).
 \label{rep:eq:qmu}
\end{equation}
Its nonnegative self-adjoint Hilbert--Schmidt operator is exactly the operator \(T_{\mu}\) from \eqref{eq:positive-covariance-operator}; we extend it by zero on constants.

We denote by \(\Fr_N^-(X_N,\mu)\) the repulsive modulated energy determined by the diagonal-free DGR splitting below; this notation is distinct from the attractive \(\Fr_N\) used in the main text.

The exact DGR splitting is
\begin{equation}\label{rep:eq:repulsive-splitting}
 -\be N\Fr_N^-(X_N,\mu)
 =-\frac\be N\sum_{1\le i<j\le N}q^-_\mu(x_i,x_j)
  +\frac\be N\sum_{i=1}^N\psi^-_\mu(x_i)+\be I^-_\mu.
\end{equation}
Thus, with
\begin{equation}\label{rep:eq:centered-P}
 P^-_{N,\be}(\mu)=\int\exp\left\{
 -\frac\be N\sum_{i<j}q^-_\mu(x_i,x_j)
 +\frac\be N\sum_i\psi^-_\mu(x_i)
 \right\}\,\dd\mu^{\otimes N},
\end{equation}
we have
\begin{equation}\label{rep:eq:K-P}
 K^-_{N,\be}(\mu)=e^{\be I^-_\mu}P^-_{N,\be}(\mu).
\end{equation}
Denote by $F_{N,\be,\mu}^{-,(k)}$ the density of the $k$-particle marginal of the corresponding Gibbs law relative to $\mu^{\otimes k}$.

Define the repulsive resolvent and response operator by
\begin{equation}\label{rep:eq:Rminus-Bminus}
 R^-_{\be,\mu}=(\Id+\be T_\mu)^{-1},\qquad
 B^-_{\be,\mu}=R^-_{\be,\mu}-\Id
 =-\be T_\mu R^-_{\be,\mu}.
\end{equation}
Let $b^-_{\be,\mu}(x,y)$ be the kernel of $B^-_{\be,\mu}$. Set
\begin{equation}\label{rep:eq:omega-minus}
 \omega^-_{\be,\mu}(x)
 =\be^2\left\langle q^-_\mu(x,\cdot),
 R^-_{\be,\mu}q^-_\mu(x,\cdot)\right\rangle_{L^2(\mu)},
\end{equation}
write $(\omega^-)^\circ=\omega^- -\int\omega^-\,\dd\mu$, and define
\begin{equation}\label{rep:eq:a-minus}
 d^-_{\be,\mu}=\be\psi^-_\mu+\frac12(\omega^-_{\be,\mu})^\circ,
 \qquad
 a^-_{\be,\mu}=R^-_{\be,\mu}d^-_{\be,\mu}.
\end{equation}

\begin{remark}[Sign of the one-body source]\label{rep:rem:sign-correction}
The positive sign of the weak one-body insertion in \eqref{rep:eq:repulsive-splitting} is essential: it is $+\be N^{-1}\sum_i\psi^-_\mu(x_i)$, and the corresponding singleton is \eqref{rep:eq:a-minus}. The direct small-temperature check is
\begin{equation}\label{rep:eq:small-beta-a}
 a^-_{\be,\mu}
 =\be\psi^-_\mu-\be^2T_\mu\psi^-_\mu
 +\frac{\be^2}{2}\left(
 x\mapsto\|q^-_\mu(x,\cdot)\|_{L^2(\mu)}^2
 \right)^\circ+O_{L^2}(\be^3),
\end{equation}
which agrees with the conditional expansion of \eqref{rep:eq:centered-P}.
\end{remark}

\begin{thm}[Repulsive prescribed-reference all-order theorem]\label{rep:thm:main}
Fix $\d\ge2$, $M_0<\infty$, and $\be_0<\infty$. Uniformly over all references satisfying \eqref{rep:eq:mu-class} and all $0\le\be\le\be_0$, the following conclusions hold.

For every fixed integer $M\ge0$, there are scalar coefficients $z^-_j(\be,\mu)$, independent of $M$, such that
\begin{equation}\label{rep:eq:partition-all-order}
 K^-_{N,\be}(\mu)
 =e^{\be I^-_\mu}\det{}_2(\Id+\be T_\mu)^{-1/2}
 \left(\sum_{j=0}^{M}\frac{z^-_j(\be,\mu)}{N^j}
 +O_{M}(N^{-M-1})\right),\qquad z^-_0=1.
\end{equation}

For every fixed $k\ge1$, $M\ge0$, and $1\le p<\infty$, there are symmetric coefficient functions $f^-_{j,k,\be,\mu}$, independent of $M$, such that
\begin{equation}\label{rep:eq:marginal-all-order}
 \left\|F_{N,\be,\mu}^{-,(k)}-1-
 \sum_{j=1}^{M}\frac{f^-_{j,k,\be,\mu}}{N^j}
 \right\|_{L^p(\mu^{\otimes k})}
 \le C_{M,k,p}N^{-M-1}.
\end{equation}
Every coefficient belongs to every finite $L^p$, has integral zero, and is consistent under marginalization. The first coefficient is
\begin{equation}\label{rep:eq:first-profile-minus}
 f^-_{1,k,\be,\mu}(x_1,\ldots,x_k)
 =\sum_{i=1}^ka^-_{\be,\mu}(x_i)
 +\sum_{1\le i<j\le k}b^-_{\be,\mu}(x_i,x_j).
\end{equation}
The first scalar correction is
\begin{equation}\label{rep:eq:z1-minus}
 z^-_1(\be,\mu)
 =\frac12\langle d^-_{\be,\mu},R^-_{\be,\mu}d^-_{\be,\mu}\rangle
 -\frac14\Tr\big((B^-_{\be,\mu})^2\big)
 +\frac1{12}\iint (b^-_{\be,\mu}(x,y))^3\,\dd\mu(x)\dd\mu(y).
\end{equation}

For each fixed coefficient order $j$ and finite $p$,
\begin{equation}\label{rep:eq:k-polynomial}
 \|f^-_{j,k,\be,\mu}\|_{L^p(\mu^{\otimes k})}
 \le C_{j,p}(1+k)^j.
\end{equation}
The fixed-$k$ remainder constant can be bounded as follows:
\begin{equation}\label{rep:eq:k-remainder}
 \left\|F_{N,\be,\mu}^{-,(k)}-1-
 \sum_{j=1}^{M}\frac{f^-_{j,k,\be,\mu}}{N^j}
 \right\|_{L^p}
 \le C_{M,p}^{\,k+1}(1+k)^{C_{M,p}}N^{-M-1},
 \qquad N\ge N_0(M,k,p,\be_0).
\end{equation}
\end{thm}

The exponential dependence on $k$ in \eqref{rep:eq:k-remainder} arises from the exact marked exponential on the bad event. The estimate is not claimed to be sharp and does not by itself give an all-order expansion when $k$ grows with $N$.

\begin{cor}[Divergences, Hoeffding projections, clusters, and linear statistics]\label{rep:cor:consequences}
Under the hypotheses of \Cref{rep:thm:main}, every fixed marginal has arbitrary finite-order chi-square and relative-entropy expansions. In particular,
\begin{align}
 \chi^2\!\left(\Q_{N,\be}^{-,(k)}(\mu)\mid\mu^{\otimes k}\right)
 &=\frac1{N^2}\left[k\|a^-_{\be,\mu}\|_2^2
 +\binom{k}{2}\|B^-_{\be,\mu}\|_{\HS}^2\right]+O_k(N^{-3}),\label{rep:eq:chi-leading}\\
 \Hr\!\left(\Q_{N,\be}^{-,(k)}(\mu)\mid\mu^{\otimes k}\right)
 &=\frac1{2N^2}\left[k\|a^-_{\be,\mu}\|_2^2
 +\binom{k}{2}\|B^-_{\be,\mu}\|_{\HS}^2\right]+O_k(N^{-3}),\label{rep:eq:kl-leading}\\
 \TV\!\left(\Q_{N,\be}^{-,(k)}(\mu),\mu^{\otimes k}\right)
 &=\frac1{2N}\left\|\sum_i a^-_{\be,\mu}(x_i)
 +\sum_{i<j}b^-_{\be,\mu}(x_i,x_j)\right\|_{L^1}
 +O_k(N^{-2}).\label{rep:eq:tv-leading}
\end{align}
For every fixed Hoeffding level and every fixed normalized connected correlation, substitution into the corresponding finite M\"obius inversion gives an expansion to every prescribed order in every finite $L^p$. For bounded centered test functions, the joint Laplace transform of the $N^{-1/2}$-scaled linear statistics has an expansion to every prescribed order in $N^{-1/2}$, uniformly on compact subsets of a fixed complex neighborhood of the origin. Its leading covariance is
\begin{equation}\label{rep:eq:linear-covariance}
 \Sigma^-_{ab}=\left\langle\varphi_a,
 R^-_{\be,\mu}\varphi_b\right\rangle_{L^2(\mu)}.
\end{equation}
\end{cor}

\subsection{Kernel estimates and approximation preserving nonnegative definiteness}

We record the analytic properties used by both proofs.

\begin{lemma}[Uniform logarithmic row estimates]\label{rep:lem:rows}
Under \eqref{rep:eq:mu-class}, $q^-_\mu$ is canonical, nonnegative definite, and Hilbert--Schmidt. There are constants $a_0,C>0$, depending only on $\d$ and $M_0$, such that
\begin{equation}\label{rep:eq:row-exp}
 \operatorname*{ess\,sup}_x\int e^{a_0|q^-_\mu(x,y)|}\,\dd\mu(y)\le C,
 \qquad q^-_\mu(x,y)\ge-C,
 \qquad \|\psi^-_\mu\|_\infty\le C.
\end{equation}
Consequently, every finite row moment has a factorial bound uniform in $\mu$.
\end{lemma}

\begin{proof}
The functions $h^-_\mu$ and the number $I^-_\mu$ are uniformly bounded because $\g\in L^1(\muunif)$ and $\rho\le M_0$. The logarithmic singularity gives $e^{a|\g|}\in L^1(\muunif)$ for every $a<\d$. This proves \eqref{rep:eq:row-exp}. For $f\in L_0^2(\mu)$,
\begin{equation}\label{rep:eq:fourier-positive}
 \langle f,T_\mu f\rangle
 =\sum_{n\ne0}\widehat \g(n)|\widehat{\rho f}(n)|^2\ge0.
\end{equation}
The uniform second-moment bound for the kernel rows implies the Hilbert--Schmidt property. Canonicality is immediate from \eqref{rep:eq:qmu}.
\end{proof}

Let $H_\vep=e^{\vep^2\Delta}$ be the periodic heat semigroup, set $\g_\vep=H_\vep \g$, and let $q^-_{\mu,\vep}$, $\psi^-_{\mu,\vep}$ be obtained by the same centering formulas as above. The operator with kernel $q^-_{\mu,\vep}$ is nonnegative and trace class. All Gaussian identities below may first be written after a finite Fourier truncation and then passed to the trace-class heat cutoff.

\begin{lemma}[Heat-cutoff estimates]\label{rep:lem:cutoff}
For every finite $r\ge1$, there are $\alpha_r>0$ and $C_r<\infty$ such that, uniformly under \eqref{rep:eq:mu-class},
\begin{align}
 &\|q^-_{\mu,\vep}-q^-_\mu\|_{L^r(\mu^{\otimes2})}
 +\|\psi^-_{\mu,\vep}-\psi^-_\mu\|_{L^r(\mu)}
 \le C_r\vep^{\alpha_r}(1+|\log\vep|),\label{rep:eq:cutoff-Lr}\\
 &\sup_x q^-_{\mu,\vep}(x,x)
 \le C(1+|\log\vep|),\label{rep:eq:cutoff-diagonal}\\
 &\operatorname*{ess\,sup}_x\int e^{a_0|q^-_{\mu,\vep}(x,y)|}\,\dd\mu(y)\le C.\label{rep:eq:cutoff-row}
\end{align}
One may take any $\alpha_r<\min\{1,\d/r\}$.
\end{lemma}

\begin{proof}
The translation modulus of $-\log|x|$ on the torus satisfies
\begin{equation}\label{rep:eq:translation-modulus}
 \|\g(\cdot-z)-\g\|_{L^r(\muunif)}
 \le C_{r,\alpha}|z|^\alpha(1+|\log|z||),
 \qquad \alpha<\min\{1,\d/r\}.
\end{equation}
To prove it, split the torus into the union of the two balls of radius $2|z|$ centered at $0$ and $z$, and their complement. On the balls, integrate the $r$th power of the logarithm. On the complement, use the mean-value estimate $|\g(x-z)-\g(x)|\le C|z|/\operatorname{dist}(x,\{0,z\})$ and integrate in polar coordinates. Averaging \eqref{rep:eq:translation-modulus} against the heat kernel proves the corresponding estimate for $\g_\vep-\g$. Centering is bounded on every finite $L^r$ because $\rho\le M_0$, proving \eqref{rep:eq:cutoff-Lr}. The diagonal estimate follows from the heat average of the logarithm at spatial scale $\vep$. Finally, $|H_\vep \g|\le H_\vep|\g|$ and Jensen's inequality imply the uniform exponential row estimate; the centering terms are bounded.
\end{proof}

\subsection{The exact oscillatory marked identity and integrated Wick coefficients}

Fix a cutoff and suppress $\vep$ temporarily. Put
\begin{equation}\label{rep:eq:c-s}
 c=\be q^-_{\mu,\vep},\qquad s=\be\psi^-_{\mu,\vep},
\end{equation}
and let $\phi$ be a centered real Gaussian field with covariance $c$. Canonicality gives
\begin{equation}\label{rep:eq:field-centered}
 \int_{\T^{\d}}\phi(x)\,\dd\mu(x)=0\qquad\text{almost surely}.
\end{equation}
Write $D(x)=c(x,x)$ and, for a centered deterministic source $u$, define
\begin{equation}\label{rep:eq:Aminus}
 A^-_u(t)=\int_{\T^{\d}}
 \exp\left\{\mathrm{i} t\phi(y)+\frac{t^2}{2}D(y)+t^2u(y)\right\}\,\dd\mu(y).
\end{equation}

For a marked configuration $x=(x_1,\ldots,x_k)$, set
\begin{align}
 u_x^-(y)&=s(y)-\sum_{a=1}^kc(y,x_a),\label{rep:eq:ux-minus}\\
 J_k^-(x)&=\sum_{a=1}^ks(x_a)-\sum_{1\le a<b\le k}c(x_a,x_b).\label{rep:eq:J-minus}
\end{align}
The lower bound in \eqref{rep:eq:row-exp} implies
\begin{equation}\label{rep:eq:u-upper}
 \operatorname*{ess\,sup}_{x,y}u_x^-(y)\le C_{\be_0,k}.
\end{equation}

\begin{lemma}[Exact oscillatory marked identity]\label{rep:lem:marked-identity}
Let $t=N^{-1/2}$. At every cutoff and for $N\ge k$,
\begin{equation}\label{rep:eq:marked-identity}
 P_{N,\be}^{-,\vep}(\mu)
 F_{N,\be,\mu}^{-,(k),\vep}(x)
 =e^{t^2J_k^-(x)}
 \E_\phi\left[A^-_{u_x^-}(t)^{N-k}\right]
\end{equation}
for $\mu^{\otimes k}$-almost every $x$.
\end{lemma}

\begin{proof}
Choose a feature representation $c(x,y)=v(x)\cdot v(y)$, first in finite dimension. The oscillatory Gaussian identity gives
\begin{equation}\label{rep:eq:pre-shift}
 \E\left[e^{\mathrm{i} t\sum_i\phi(x_i)}\right]
 =\exp\left\{-\frac{t^2}{2}\sum_{i,j}c(x_i,x_j)\right\}.
\end{equation}
The factors $e^{t^2D(x_i)/2}$ remove the diagonal terms. Before integrating the marked variables, the right side of \eqref{rep:eq:marked-identity} is therefore
\begin{equation}\label{rep:eq:before-complex-shift}
 \E\left[A^-_s(t)^{N-k}\prod_{a=1}^k
 e^{\mathrm{i} t\phi(x_a)+t^2D(x_a)/2+t^2s(x_a)}\right].
\end{equation}
For a standard Gaussian vector $Z$ and $a=\sum_av(x_a)$,
\begin{equation}\label{rep:eq:complex-CM}
 \E\left[e^{\mathrm{i} t a\cdot Z}G(Z)\right]
 =e^{-t^2|a|^2/2}\E[G(Z+\mathrm{i} t a)].
\end{equation}
Under this shift, $\mathrm{i} t\phi(y)$ becomes $\mathrm{i} t\phi(y)-t^2\sum_ac(y,x_a)$. The self-diagonal factors cancel, while the cross terms give $-t^2\sum_{a<b}c(x_a,x_b)$. The sources at the marks give $t^2\sum_as(x_a)$. This proves \eqref{rep:eq:marked-identity}. Degenerate and trace-class covariances follow by approximation on the range of the covariance operator.
\end{proof}

Let $:\!\phi^r\!:_c$ denote Wick ordering relative to $c$, and put
\begin{equation}\label{rep:eq:weighted-Wick}
 \mathcal W_{r,\ell}(u)=\frac{\mathrm{i}^r}{r!\,\ell!}
 \int_{\T^{\d}}u(y)^\ell:\!\phi(y)^r\!:_c\,\dd\mu(y).
\end{equation}
Then
\begin{equation}\label{rep:eq:A-Wick-series}
 A^-_u(t)=1+\sum_{r+2\ell\ge1}t^{r+2\ell}\mathcal W_{r,\ell}(u)
\end{equation}
in every fixed Gaussian $L^v$ for $t$ sufficiently small.

\begin{lemma}[Cutoff-uniform weighted Wick bounds]\label{rep:lem:wick-bounds}
For every fixed $k$ and $1\le v<\infty$, there is $C_{k,v}$ such that, uniformly in the cutoff, reference, inverse temperature, and marked configuration,
\begin{equation}\label{rep:eq:Wick-geometric}
 \operatorname*{ess\,sup}_x\left\|\mathcal W_{r,\ell}(u_x^-)\right\|_{L^v(\phi)}
 \le C_{k,v}^{\,r+2\ell}.
\end{equation}
The same estimate holds with $u_x^-$ replaced by $s$.
\end{lemma}

Cutoff removal will be performed on the Gaussian expectations of coefficient polynomials. The proof of \Cref{rep:prop:moving-cutoff} establishes their convergence, including the quadratic tilt, by comparing deterministic covariance and resolvent kernels.

\begin{proof}
The Gaussian covariance identity gives
\begin{multline}\label{rep:eq:wick-L2}
 \E\left|\int u(y)^\ell:\!\phi(y)^r\!:_c\,\dd\mu(y)\right|^2\\
 =r!\iint u(y)^\ell u(z)^\ell c(y,z)^r\,\dd\mu(y)\dd\mu(z).
\end{multline}
We estimate the absolute value of the last integral.  Write
$u_x^-=s-\sum_{a=1}^kc(\cdot,x_a)$.  By \eqref{rep:eq:cutoff-row}, after decreasing the exponential parameter by a factor depending only on $k$, H\"older's inequality gives numbers $\delta_k,C_k>0$ such that
\begin{equation}\label{rep:eq:three-factor-exponential}
 \operatorname*{ess\,sup}_x\iint
 \exp\!\left\{\delta_k\bigl(
 |u_x^-(y)|+|u_x^-(z)|+|c(y,z)|\bigr)\right\}
 \,\dd\mu(y)\dd\mu(z)\le C_k.
\end{equation}
Indeed, for fixed $y$, apply H\"older in $z$ to the $k$ row factors
$e^{\delta_k|c(z,x_a)|}$ and the additional row factor
$e^{\delta_k|c(z,y)|}$; then apply the same argument in $y$.

Put $n=r+2\ell$.  The coefficient of
$A^\ell B^\ell C^r$ in the exponential series for
$e^{\delta_k(A+B+C)}$ yields, for $A,B,C\ge0$,
\begin{equation}\label{rep:eq:factorial-monomial}
 A^\ell B^\ell C^r
 \le \ell!^2r!\,\delta_k^{-n}e^{\delta_k(A+B+C)}.
\end{equation}
Equations \eqref{rep:eq:wick-L2}--\eqref{rep:eq:factorial-monomial} therefore imply
\begin{equation}\label{rep:eq:factorial-bound}
 \left\|\int u_x^-(y)^\ell:\!\phi(y)^r\!:_c\,\dd\mu(y)\right\|_{L^2(\phi)}
 \le C_k^{1/2}\delta_k^{-n/2}r!\,\ell!.
\end{equation}
After division by $r!\ell!$, this is geometric in $n=r+2\ell$.
Gaussian hypercontractivity on the $r$th homogeneous chaos \cite[Proposition~2.6]{NourdinPeccatiReinert2010} multiplies the
$L^2$ norm by at most $(v-1)^{r/2}$ for $v\ge2$; the case $1\le v<2$ follows by monotonicity of Gaussian $L^v$ norms.  This proves \eqref{rep:eq:Wick-geometric}.

For untilted polynomial moments, no convergence of covariance square roots is required.  After the Gaussian expectation is taken, every fixed product of weighted Wick coefficients is a finite sum of Wick-contraction diagrams.  Each such diagram is an integral of a finite product of factors drawn from $c$, $u_x^-$, and their cutoff analogues.  Replacing one factor at a time, \eqref{rep:eq:cutoff-Lr}, H\"older's inequality, and the exponential product domination in \eqref{rep:eq:three-factor-exponential} give convergence in every required spatial $L^p$.  The additional resolvent kernels created by the quadratic tilt are treated explicitly in the proof of \Cref{rep:prop:moving-cutoff}.
\end{proof}

The first coefficients of \eqref{rep:eq:A-Wick-series} are
\begin{align}
 A_2&=-\frac12\int:\!\phi^2\!:_c\,\dd\mu,\label{rep:eq:A2minus}\\
 A_3(u)&=\mathrm{i}\left(\int u\phi\,\dd\mu-\frac16\int:\!\phi^3\!:_c\,\dd\mu\right),\label{rep:eq:A3minus}\\
 A_4(u)&=\frac1{24}\int:\!\phi^4\!:_c\,\dd\mu
 -\frac12\int u:\!\phi^2\!:_c\,\dd\mu
 +\frac12\int u^2\,\dd\mu.\label{rep:eq:A4minus}
\end{align}
Since $A^-_u(t)-1=O_{L^v}(t^2)$ uniformly, its logarithm has a cutoff-uniform finite expansion on the event $|A^-_u(t)-1|\le1/2$. The coefficient recursion is the usual finite logarithmic and Bell-polynomial recursion; no infinite formal series is used.

\subsection{Oscillatory expansion with a moving heat cutoff}

The coefficient estimates are uniform in the heat cutoff, whereas a direct modulus bound for the exact oscillatory integral has a polynomial cutoff cost. We choose the cutoff as a power of $N$ and expand far enough to absorb that cost. A class-uniform repulsive partition bound then permits removal of the cutoff. This argument does not use the signed-gas theorem of Appendix~\ref{app:signed-log-heat-kernel}.

\subsubsection{Polynomial control of the exact oscillatory factor}

\begin{lemma}[Deterministic cutoff cost]\label{rep:lem:deterministic-cost}
For $N\ge2k$, $t=N^{-1/2}$, and $0<\vep\le1/2$,
\begin{equation}\label{rep:eq:deterministic-cost}
 \operatorname*{ess\,sup}_x\left|A^-_{u_x^-}(t)\right|^{N-k}
 \le C_{\be_0,k}\vep^{-C_{\be_0}}.
\end{equation}
\end{lemma}

\begin{proof}
By \eqref{rep:eq:u-upper},
\begin{equation}\label{rep:eq:A-modulus}
 |A^-_{u_x^-}(t)|
 \le\int e^{t^2(D/2+u_x^-)}\,\dd\mu.
\end{equation}
For $0<t^2\le1$, concavity of $r\mapsto r^{t^2}$ gives
\begin{equation}\label{rep:eq:concavity-cost}
 \int e^{t^2(D/2+u_x^-)}\,\dd\mu
 \le\left(\int e^{D/2+u_x^-}\,\dd\mu\right)^{t^2}.
\end{equation}
Raise this to the power $N-k\le N$ and use \eqref{rep:eq:cutoff-diagonal} and \eqref{rep:eq:u-upper}.
\end{proof}

\subsubsection{Overexpansion absorbs the cutoff cost}

\begin{prop}[Moving-cutoff expansion with fixed marked coordinates]\label{rep:prop:moving-cutoff}
Fix $M,k,p$ and $a>0$, and set $\vep_N=N^{-a}$. There are coefficient functions $U^-_{j,k,\vep_N}$, obtained from finite Gaussian polynomial expectations, such that
\begin{equation}\label{rep:eq:moving-cutoff-expansion}
 \left\|e^{J_k^-/N}\E\left[A^-_{u_x^-}(N^{-1/2})^{N-k}\right]
 -D^-_{\be,\mu,\vep_N}
 \sum_{j=0}^{M}\frac{U^-_{j,k,\vep_N}(x)}{N^j}
 \right\|_{L^p(\mu^{\otimes k})}
 \le C_{M,k,p,a}N^{-M-1},
\end{equation}
where
\begin{equation}\label{rep:eq:D-epsilon}
 D^-_{\be,\mu,\vep}
 =\det{}_2(\Id+\be T_{\mu,\vep})^{-1/2}.
\end{equation}
For each fixed $j\le M$, $U^-_{j,k,\vep_N}$ converges in every finite $L^p$ to a cutoff-free coefficient as $N\to\infty$.
\end{prop}

\begin{proof}
Write $t=N^{-1/2}$ and
$E_{0,t,x}=\{|A^-_{u_x^-}(t)-1|\le1/2\}$.  By \Cref{rep:lem:wick-bounds}, for every $v<\infty$,
\begin{equation}\label{rep:eq:bad-probability}
 \sup_x\mathbb P(E_{0,t,x}^c)\le C_vt^{2v}.
\end{equation}
Fix an auxiliary integer $L$, to be chosen after the cutoff exponent $a$ and the desired order $M$.  On $E_{0,t,x}$, finite Taylor expansion of the logarithm, combined with the cutoff-uniform expansion \eqref{rep:eq:A-Wick-series}, gives
\begin{equation}\label{rep:eq:log-overexpansion}
 (t^{-2}-k)\log A^-_{u_x^-}(t)
 =A_2+Y_{L,t,x}+\mathcal R_{L,t,x},
 \qquad
 Y_{L,t,x}=\sum_{j=1}^{L}t^jL_{j,k,x},
\end{equation}
where every $L_{j,k,x}$ is a fixed polynomial in the weighted Wick coefficients and
\begin{equation}\label{rep:eq:log-remainder-moment}
 \sup_x\|\mathbf1_{E_{0,t,x}}\mathcal R_{L,t,x}\|_{L^v(\phi)}
 \le C_{L,k,v}t^{L+1},
 \qquad
 \sup_x\|Y_{L,t,x}\|_{L^v(\phi)}\le C_{L,k,v}t.
\end{equation}

Under the quadratic weight $e^{A_2}$, the Gaussian covariance is
\begin{equation}\label{rep:eq:tilted-covariance}
 \kappa_{\vep}
 =\be T_{\mu,\vep}(\Id+\be T_{\mu,\vep})^{-1}.
\end{equation}
For every fixed $r>0$,
\begin{equation}\label{rep:eq:quadratic-exp-bound}
 \E[e^{rA_2}]
 =\det{}_2(\Id+r\be T_{\mu,\vep})^{-1/2}
 \le\exp\left\{\frac{r^2\be_0^2}{4}\|T_\mu\|_{\HS}^2\right\}.
\end{equation}
Thus, every fixed polynomial in the weighted Wick coefficients, multiplied by $e^{A_2}$, has cutoff-uniform finite moments.

We next record the cutoff continuity of the kernels generated by this quadratic tilt.  Restore the cutoff subscript and write
\begin{equation}\label{rep:eq:tilt-resolvent-notation}
 C_\vep=\be T_{\mu,\vep},\qquad
 R_\vep=(\Id+C_\vep)^{-1},\qquad
 \kappa_\vep=C_\vep R_\vep=\Id-R_\vep.
\end{equation}
Since $C_\vep\ge0$, one has $\|R_\vep\|_{\op}\le1$, and the resolvent identity gives
\begin{equation}\label{rep:eq:kappa-HS-continuity}
 \|\kappa_\vep-\kappa_0\|_{\HS}
 =\|R_\vep(C_\vep-C_0)R_0\|_{\HS}
 \le \|C_\vep-C_0\|_{\HS}
 =\be\|q^-_{\mu,\vep}-q^-_\mu\|_{L^2(\mu^{\otimes2})}.
\end{equation}
For $c_{\vep,x}=c_\vep(x,\cdot)$, the kernel rows satisfy
$(\kappa_\vep)_x=R_\vep c_{\vep,x}$ and therefore
\begin{equation}\label{rep:eq:kappa-row-L2}
 \operatorname*{ess\,sup}_x\|(\kappa_\vep)_x\|_{L^2(\mu)}
 \le \operatorname*{ess\,sup}_x\|c_{\vep,x}\|_{L^2(\mu)}\le C.
\end{equation}
Moreover, the operator identity $\kappa_\vep=C_\vep-C_\vep\kappa_\vep$ and Cauchy--Schwarz show that the kernel of $C_\vep\kappa_\vep$ is uniformly bounded.  Hence
$|\kappa_\vep(x,y)|\le |c_\vep(x,y)|+C$, so $\kappa_\vep$ inherits the cutoff-uniform exponential row bound and all finite row moments.

The self-contraction created by changing Wick order is
\begin{equation}\label{rep:eq:omega-tilt-definition}
 \omega_\vep(x)
 =\langle c_{\vep,x},R_\vep c_{\vep,x}\rangle_{L^2(\mu)}.
\end{equation}
It is uniformly bounded, and at a finite cutoff it equals
$c_\vep(x,x)-\kappa_\vep(x,x)$.  Combining
\eqref{rep:eq:kappa-HS-continuity}, the resolvent identity, and
\eqref{rep:eq:cutoff-Lr}, and then interpolating with the uniform row moments, yields: for every finite $p$, there are $\gamma_p>0$ and $C_p<\infty$ such that
\begin{equation}\label{rep:eq:tilt-kernel-cutoff-rate}
 \|\kappa_\vep-\kappa_0\|_{L^p(\mu^{\otimes2})}
 +\|\omega_\vep-\omega_0\|_{L^p(\mu)}
 \le C_p\vep^{\gamma_p}(1+|\log\vep|)^{C_p}.
\end{equation}
Indeed, for the second term one expands the difference of the two quadratic forms in \eqref{rep:eq:omega-tilt-definition}; the two row differences are controlled by \eqref{rep:eq:cutoff-Lr}, and the middle resolvent difference is controlled by \eqref{rep:eq:kappa-HS-continuity}.  Consequently, after division by $\E[e^{A_2}]$, every tilted Wick expectation is a finite contraction diagram with factors drawn from $c_\vep$, $s_\vep$, $\kappa_\vep$, and $\omega_\vep$.  This proves cutoff convergence of the tilted deterministic coefficients without choosing, or comparing, covariance square roots.

Introduce the smaller good event
\begin{equation}\label{rep:eq:very-good}
 E_{t,x}=E_{0,t,x}\cap\{|Y_{L,t,x}|\le1\}\cap\{|\mathcal R_{L,t,x}|\le1\}.
\end{equation}
Equations \eqref{rep:eq:bad-probability} and \eqref{rep:eq:log-remainder-moment}, with arbitrarily large moment exponents, imply
$\sup_x\mathbb P(E_{t,x}^c)\le C_Qt^Q$ for every prescribed $Q$, after increasing the logarithmic expansion order if necessary.  On $E_{t,x}$, Taylor's theorem applied to
$e^{Y_{L,t,x}+\mathcal R_{L,t,x}}$ gives, after retaining all monomials of total $t$-degree at most $2M+1$,
\begin{equation}\label{rep:eq:good-overexpand}
 e^{(t^{-2}-k)\log A^-_{u_x^-}(t)}
 =e^{A_2}\sum_{j=0}^{2M+1}t^jP_{j,k,x}
 +\mathcal E_{M,L,t,x},
\end{equation}
where the $P_{j,k,x}$ are fixed polynomials and
\begin{equation}\label{rep:eq:good-error-moment}
 \sup_x\|\mathbf1_{E_{t,x}}\mathcal E_{M,L,t,x}\|_{L^1(\phi)}
 \le C\bigl(t^{2M+2}+t^{L+1}\bigr).
\end{equation}
This estimate uses H\"older, \eqref{rep:eq:quadratic-exp-bound}, and the polynomial moment bounds; no exponential moment of a cubic or higher Gaussian polynomial is used.

On $E_{t,x}^c$, the exact term is bounded by the deterministic bound \eqref{rep:eq:deterministic-cost}, whereas every retained polynomial--quadratic term is controlled by \eqref{rep:eq:quadratic-exp-bound}.  Therefore, for every $Q$,
\begin{equation}\label{rep:eq:bad-overexpand}
 \sup_x\E\left[\mathbf1_{E_{t,x}^c}
 \left(
 |A^-_{u_x^-}(t)|^{N-k}
 +e^{A_2}\sum_{j=0}^{2M+1}t^j|P_{j,k,x}|
 \right)\right]
 \le C\vep_N^{-C_0}t^Q.
\end{equation}
Choose $L$ and $Q$ so large that
\begin{equation}\label{rep:eq:choose-Lv}
 \frac{L+1}{2}-aC_0\ge M+1,
 \qquad \frac Q2-aC_0\ge M+1.
\end{equation}
Equations \eqref{rep:eq:good-error-moment}--\eqref{rep:eq:choose-Lv} prove an expansion with remainder $O(N^{-M-1})$.

It remains to include the deterministic marked factor.  The row-exponential estimate gives every finite spatial moment of $J_k^-$ uniformly in the cutoff, and $J_k^-$ is bounded above because $s$ is bounded and $c$ is bounded below.  Taylor's theorem on the half-line $(-\infty,C]$ therefore gives
\begin{equation}\label{rep:eq:exterior-factor-expansion}
 \left\|e^{t^2J_k^-}
 -\sum_{\ell=0}^{M}\frac{t^{2\ell}(J_k^-)^\ell}{\ell!}\right\|_{L^p(\mu^{\otimes k})}
 \le C_{M,k,p}t^{2M+2}.
\end{equation}
Multiplying this finite even polynomial into \eqref{rep:eq:good-overexpand}, retaining total $t$-degree at most $2M+1$, and using H\"older's inequality preserves the same remainder order.

Under $\phi\mapsto-\phi$, a coefficient of total $t$-degree $j$ is multiplied by $(-1)^j$, while $e^{A_2}$ is invariant.  Hence every odd coefficient expectation vanishes.  Since \eqref{rep:eq:exterior-factor-expansion} contains only even powers, retaining the even powers proves \eqref{rep:eq:moving-cutoff-expansion}.

After normalization by $D^-_{\be,\mu,\vep}=\E[e^{A_2}]$, every fixed retained coefficient is a finite sum of tilted contraction diagrams with factors drawn from $c_\vep$, $s_\vep$, $\kappa_\vep$, and $\omega_\vep$.  In a sufficiently large finite $L^r$ product space, each such formula is locally Lipschitz in these factors.  The row-exponential bounds supply uniform product domination, while \eqref{rep:eq:cutoff-Lr} and \eqref{rep:eq:tilt-kernel-cutoff-rate} control their differences.  Consequently, for some $\alpha>0$ depending on the fixed coefficient,
\begin{equation}\label{rep:eq:coefficient-cutoff-rate}
 \|U^-_{j,k,\vep}-U^-_{j,k,0}\|_{L^p}
 \le C\vep^\alpha(1+|\log\vep|)^C.
\end{equation}
This follows directly by replacing one cutoff factor at a time in each finite tilted diagram.  Local Lipschitz continuity of $A\mapsto\log\det{}_2(\Id+A)$ on bounded nonnegative Hilbert--Schmidt sets, together with \eqref{rep:eq:kappa-HS-continuity}, gives the same type of algebraic rate for the determinant factor.  Thus, only the exact bad-event term pays the polynomial cutoff cost, and every retained coefficient and the determinant converge at algebraic cutoff rates.
\end{proof}

\subsubsection{Uniform periodic partition bounds and cutoff removal}

The next proposition makes the reference dependence in the repulsive partition estimate of Delgadino, Gvalani, and Rosenzweig \cite{DelgadinoGvalaniRosenzweig2026} uniform over \eqref{rep:eq:mu-class}. This uniformity is needed when comparing a reference measure with its heat regularization.

For $0\le\theta\le1$ and $0<\vep\le1$, write
\begin{equation}\label{rep:eq:interpolated-kernel-routeI}
 \g_{\vep,\theta}=(1-\theta)\g+\theta \g_\vep,
\end{equation}
and define the associated centered kernel, one-body source, and centered partition function by the definitions \eqref{rep:eq:normalization-crosswalk}, \eqref{rep:eq:qmu}, and \eqref{rep:eq:centered-P}.

\begin{prop}[Uniform repulsive partition bounds on the torus]\label{rep:prop:DGR-class-uniform}
Fix $\d\ge2$, $M_0<\infty$, and $\overline\be<\infty$.  There is a constant
\begin{equation}\label{rep:eq:DGR-uniform-constant}
 C_{\mathrm{pf}}=C_{\mathrm{pf}}(\d,M_0,\overline\be)<\infty
\end{equation}
such that
\begin{equation}\label{rep:eq:DGR-class-uniform}
 \sup_{\substack{N\ge1,\ 0\le\be\le\overline\be\\
 0<\vep\le1,\ 0\le\theta\le1\\
 \mu=\rho\muunif,\ 0\le\rho\le M_0,\ \int\rho\,\dd\muunif=1}}
 P^{-,\vep,\theta}_{N,\be}(\mu)
 \le C_{\mathrm{pf}}.
\end{equation}
The same estimate holds for the unregularized endpoint $\vep=0$.
\end{prop}

\begin{proof}
The cited partition theorem has a reference-dependent constant. We first bound its inputs uniformly over the density class.  On the volume-one torus, the periodic logarithmic kernel belongs to $L^r(\muunif)$ for every finite $r$.  Hence, uniformly over $0\le\rho\le M_0$,
\begin{align}
 \|h^-_\mu\|_\infty+|I^-_\mu|+\|\psi^-_\mu\|_\infty&\le C(\d,M_0),\label{rep:eq:uniform-onebody-package}\\
 \iint |\log\operatorname{dist}(x,y)|\,\dd\mu(x)\dd\mu(y)
 +\|q^-_\mu\|_{L^2(\mu^{\otimes2})}^2&\le C(\d,M_0),\label{rep:eq:uniform-log-package}\\
 \sup_{0\le\be\le\overline\be}
 \int e^{-4\be h^-_\mu}\,\dd\mu&\le C(\d,M_0,\overline\be).
 \label{rep:eq:uniform-exp-package}
\end{align}
For the dependence of the imported bounds, we use \cite[Section~4.2, Lemma~4.3(ii), and Section~5.1.2]{DelgadinoGvalaniRosenzweig2026}, with the periodic interpretation specified in its Remark~2.1. The high-frequency $L^2(\mu^{\otimes2})$, row-$L^2(\mu)$, and operator estimates depend on the reference only through its upper density bound. The low-frequency diagonal estimate is independent of the reference. In the logarithmic argument there, the low-frequency interaction is combined with the one-body term before removing the infrared cutoff; equations~(4.28) and~(4.31)--(4.39) identify this cancellation. The high-frequency kernel has an integrable logarithmic majorant, so convolution against $\rho\le M_0$ gives a bound depending only on $\d$, $M_0$, and the fixed frequency cutoff. Equations \eqref{rep:eq:uniform-onebody-package}--\eqref{rep:eq:uniform-exp-package} control the other reference-dependent quantities. Thus the frequency cutoff $K_{\be,\mu}$, the lower tail-integration level $U_0$, and the spatial truncation scale $\delta_{N,u}$ in equations~(5.25)--(5.27) of that proof, as well as its lower bound on $N$, can be chosen using only $\d$, $M_0$, and the fixed upper coupling $\overline\be$.

For the finitely many particle numbers below that uniform threshold, the lower bound $q^-_\mu\ge-C(\d,M_0)$ and \eqref{rep:eq:uniform-onebody-package} give a uniform elementary bound on the exponent.  This closes the finite-$N$ remainder in the layer-cake argument and proves the unregularized endpoint uniformly over the class.

We next compare the heat-regularized interaction with the singular interaction.  Let $P_t=e^{t\Delta}$ and put $t=\vep^2/2$ and $\mu_t=(P_t)^*\mu$.  If $X_i\sim\mu$ and, conditionally on $X_i$, the variables $Y_i$ are obtained by independent periodic Brownian increments of time $t$, then $Y_i\sim\mu_t$ and $0\le \frac{\dd\mu_t}{\dd\muunif}\le M_0$.  Since $2t=\vep^2$, independence of the increments and self-adjointness of the heat semigroup give, for $i\ne j$,
\begin{align}\label{rep:eq:heat-component-identities}
 \E[\g(Y_i-Y_j)\mid X_N]&=P_{2t}\g(X_i-X_j)=\g_\vep(X_i-X_j),\\
 \E[h_{\mu_t}^{-,\g}(Y_i)\mid X_i]&=P_t(\g*P_t\mu)(X_i)=(P_{2t}\g*\mu)(X_i)=h_{\mu}^{-,\g_\vep}(X_i),\\
 I_{\mu_t}^{-,\g}&=\frac12\langle P_t\mu,\g*P_t\mu\rangle
 =\frac12\langle\mu,P_{2t}\g*\mu\rangle=I_{\mu}^{-,\g_\vep}.
\end{align}
The first identity is used only for $i\ne j$, so no diagonal term is introduced.  Substitution into the pair, one-body, and energy definitions yields
\begin{equation}\label{rep:eq:heat-conditional-energy}
 \mathbb E\left[\Fr_N^{-,\g}(Y_N,\mu_t)\,\middle|\,X_N\right]
 =\Fr_N^{-,\g_\vep}(X_N,\mu).
\end{equation}
The map $r\mapsto e^{-\be Nr}$ is convex.  Conditional Jensen therefore gives
\begin{equation}\label{rep:eq:heat-jensen-partition}
 K_{N,\be}^{-,\g_\vep}(\mu)
 \le K_{N,\be}^{-,\g}(\mu_t),
\end{equation}
which is uniformly bounded by the unregularized result.  Since $|I^-_\mu|$ and $|I_{\mu}^{-,\g_\vep}|$ are uniformly bounded, the same conclusion holds for the centered partitions $P^-$.  Finally, for the linear interpolation in \eqref{rep:eq:interpolated-kernel-routeI}, the centered exponent is the convex combination of the two endpoint exponents, so H\"older's inequality gives
\begin{equation}\label{rep:eq:interpolation-holder-routeI}
 P^{-,\vep,\theta}_{N,\be}
 \le
 \bigl(P^-_{N,\be}\bigr)^{1-\theta}
 \bigl(P^{-,\vep}_{N,\be}\bigr)^\theta.
\end{equation}
Log-convexity in $\be$, together with the value one at $\be=0$, reduces the compact interval to the endpoint $\overline\be$.  This proves \eqref{rep:eq:DGR-class-uniform}.
\end{proof}

\begin{lemma}[Cutoff comparison for normalized marginals]\label{rep:lem:cutoff-comparison}
Fix $\d\ge2$, $M_0<\infty$, $\be_0<\infty$, $k\ge0$, and $1\le p<\infty$.  For every sufficiently large finite $r$, there is $C=C(\d,M_0,\be_0,k,p,r)<\infty$ such that, uniformly over \eqref{rep:eq:mu-class},
\begin{equation}\label{rep:eq:marginal-cutoff-comparison}
 \sup_{0\le\be\le\be_0}
 \|F_{N,\be,\mu}^{-,(k)}-F_{N,\be,\mu}^{-,(k),\vep}\|_{L^p(\mu^{\otimes k})}
 \le C\left(
 \|q^-_{\mu,\vep}-q^-_\mu\|_{L^r(\mu^{\otimes2})}
 +\|\psi^-_{\mu,\vep}-\psi^-_\mu\|_{L^r(\mu)}
 \right).
\end{equation}
The same estimate holds for the centered partition functions.
\end{lemma}

\begin{proof}
Interpolate linearly between the cutoff and uncut exponents and denote the normalized full density by $G_\theta$. Differentiation gives
\begin{equation}\label{rep:eq:density-derivative}
 \partial_\theta G_\theta
 =G_\theta\left(\Delta_N-\E_{G_\theta\mu^{\otimes N}}[\Delta_N]\right),
\end{equation}
where
\begin{equation}\label{rep:eq:DeltaN}
 \Delta_N=-\frac\be N\sum_{i<j}(q^-_\mu-q^-_{\mu,\vep})(x_i,x_j)
 +\frac\be N\sum_i(\psi^-_\mu-\psi^-_{\mu,\vep})(x_i).
\end{equation}
For a canonical kernel $h$ and every fixed $r\ge2$, we use the canonical moment estimate
\begin{equation}\label{rep:eq:canonical-moment}
 \left\|\frac1N\sum_{i<j}h(X_i,X_j)\right\|_{L^r(\mu^{\otimes N})}
 \le C_r\|h\|_{L^r(\mu^{\otimes2})}.
\end{equation}
To obtain it, apply the decoupling theorem \cite[Theorem~1]{delaPenaMontgomerySmith1995} and integrate its tail inequality to compare $L^r$ norms. Canonicality then allows symmetrization separately in the two independent samples. Two applications of the Khintchine inequality bound the resulting signed sum by $C_r$ times the $L^r$ norm of the square root of the sum of the squared kernel values. Minkowski's inequality bounds this norm by $N\|h\|_{L^r(\mu^{\otimes2})}$, which cancels the denominator $N$. The same argument allows deterministic coefficients of absolute value at most one on the pair terms.
The one-body term is bounded by its $L^r$ norm.  Choose $r>\max\{p,2\}$ and $s<\infty$ by $1/p=1/r+1/s$.  If $H_\theta$ denotes the centered exponent and $P_\theta(\be)=\E[e^{H_\theta}]$, then Jensen gives $P_\theta(\be)\ge1$, and for every fixed $b\ge1$,
\begin{equation}\label{rep:eq:normalized-density-amplification}
 \|G_\theta\|_{L^b(\mu^{\otimes N})}
 =\frac{P_\theta(b\be)^{1/b}}{P_\theta(\be)}
 \le P_\theta(b\be)^{1/b}.
\end{equation}
Thus, Proposition~\ref{rep:prop:DGR-class-uniform} on the single enlarged window
$0\le\be\le A_{p,r}\be_0$, where
$A_{p,r}=\max\{s,p,r/(r-1)\}$, controls the three norms
$\|G_\theta\|_{L^s}$, $\|G_\theta\|_{L^p}$, and
$\|G_\theta\|_{L^{r/(r-1)}}$ uniformly over the full density class.  H\"older in \eqref{rep:eq:density-derivative}, first for $G_\theta\Delta_N$ and then for
$\E[G_\theta\Delta_N]G_\theta$, therefore gives
\begin{equation}\label{rep:eq:full-density-comparison}
 \|G_1-G_0\|_{L^p(\mu^{\otimes N})}
 \le C\big(\|q-q_\vep\|_{L^r}+\|\psi-\psi_\vep\|_{L^r}\big).
\end{equation}
Conditional expectation onto the first $k$ coordinates is an $L^p$ contraction, proving \eqref{rep:eq:marginal-cutoff-comparison}. The partition estimate follows by differentiating the unnormalized expectation before division.
\end{proof}

\begin{proof}[First, class-uniform proof of \Cref{rep:thm:main}]
Fix the requested order $M$ (and, for a marginal, $k$ and $p$).  Choose $r$ in \Cref{rep:lem:cutoff-comparison}.  Among the finitely many scalar and marked diagrams retained through order $M$, let $\gamma_{M,k,p}>0$ be the minimum of the exponents supplied by \eqref{rep:eq:coefficient-cutoff-rate} and by the Hilbert--Schmidt determinant estimate.  Then choose $a$ so large that
\begin{equation}\label{rep:eq:choose-a}
 a\min\{\alpha_r,\gamma_{M,k,p}\}>M+2.
\end{equation}
At the moving cutoff $\vep_N=N^{-a}$, \Cref{rep:prop:moving-cutoff} gives the all-order marked expansion; after this choice of $a$, the auxiliary orders $L$ and $Q$ are selected as in \eqref{rep:eq:choose-Lv}.  Equations \eqref{rep:eq:cutoff-Lr} and \eqref{rep:eq:marginal-cutoff-comparison} make the difference between cutoff and singular marginals $O(N^{-M-1})$, with the strict inequality in \eqref{rep:eq:choose-a} absorbing the logarithms.  Equation \eqref{rep:eq:coefficient-cutoff-rate} and determinant continuity give the same order after replacing every retained cutoff coefficient by its limit, including the leading determinant.  Divide the resulting unnormalized marginal series by the scalar series. This proves \eqref{rep:eq:partition-all-order} and \eqref{rep:eq:marginal-all-order} uniformly over the full density class. The explicit coefficients are computed in \Cref{rep:sec:first-coefficients}. The $k$-bounds are proved in \Cref{rep:sec:consequences}.  This proof of the partition and marginal expansions uses no signed-gas estimate from \Cref{rep:sec:heat-route}.
\end{proof}

\subsection{Explicit first coefficients}\label{rep:sec:first-coefficients}

We compute at a finite cutoff and then remove it. Let
\begin{equation}\label{rep:eq:C-R-kappa}
 C=\be T_{\mu,\vep},\qquad
 R=(\Id+C)^{-1},\qquad
 \kappa=CR,
\end{equation}
so that $B^-=R-\Id=-\kappa$. Under the normalized quadratic tilt $e^{A_2}$, the field covariance is $\kappa$. Define
\begin{equation}\label{rep:eq:omega-finite}
 \omega(x)=c(x,x)-\kappa(x,x)
 =\langle c(x,\cdot),R c(x,\cdot)\rangle.
\end{equation}
This is nonnegative and converges to \eqref{rep:eq:omega-minus}.

Changing Wick order from $c$ to $\kappa$ gives
\begin{equation}\label{rep:eq:wick-change-minus}
 :\!\phi^2\!:_c=:\!\phi^2\!:_\kappa-\omega,
 \qquad
 :\!\phi^3\!:_c=:\!\phi^3\!:_\kappa-3\omega\phi.
\end{equation}
Let $\E_\kappa$ denote expectation under the tilted Gaussian law. From \eqref{rep:eq:A2minus}--\eqref{rep:eq:A4minus},
\begin{align}
 \E_\kappa[A_4(u)]
 &=\frac18\int\omega^2+\frac12\int u\omega+\frac12\int u^2,\label{rep:eq:A4-expect-minus}\\
 \E_\kappa[A_2^2]
 &=\frac14\left[\left(\int\omega\right)^2+2\Tr(\kappa^2)\right],\label{rep:eq:A2-expect-minus}\\
 \E_\kappa[A_3(u)^2]
 &=-\frac16\iint\kappa(x,y)^3\,\dd\mu(x)\dd\mu(y)
 -\left\langle u+\frac12\omega^\circ,
 \kappa\left(u+\frac12\omega^\circ\right)\right\rangle.\label{rep:eq:A3-expect-minus}
\end{align}
The minus signs in \eqref{rep:eq:A3-expect-minus} are the factor $\mathrm{i}^2=-1$. Therefore,
\begin{equation}\label{rep:eq:z1-source-minus}
 \E_\kappa\left[A_4(u)-\frac12A_2^2+\frac12A_3(u)^2\right]
 =\frac12\left\langle u+\frac12\omega^\circ,
 R\left(u+\frac12\omega^\circ\right)\right\rangle
 -\frac14\Tr(\kappa^2)+\frac1{12}\iint(-\kappa)^3.
\end{equation}
Taking $u=s=\be\psi^-_{\mu,\vep}$ and removing the cutoff proves \eqref{rep:eq:z1-minus}.

For the marked coefficient, the order-$N^{-1}$ term is
\begin{equation}\label{rep:eq:profile-before-simplify}
 f^-_{1,k}(x)
 =J_k^-(x)+z^-_1(u_x^-)-z^-_1(s)-k\E_\kappa[A_2].
\end{equation}
Here, $\E_\kappa[A_2]=\frac12\int\omega$. Put $d=s+\omega^\circ/2$. Since $u_x^--s=-\sum_ac(x_a,\cdot)$,
\begin{equation}\label{rep:eq:z-difference-minus}
 z^-_1(u_x^-)-z^-_1(s)
 =-\sum_a\langle c(x_a,\cdot),Rd\rangle
 +\frac12\sum_{a,b}\langle c(x_a,\cdot),Rc(x_b,\cdot)\rangle.
\end{equation}
The singleton contribution is
\begin{equation}\label{rep:eq:singleton-simplify}
 s(x)+\frac12\omega^\circ(x)-(CRd)(x)
 =d(x)-\kappa d(x)=Rd(x),
\end{equation}
while, for two distinct marks,
\begin{equation}\label{rep:eq:pair-simplify}
 -c(x,y)+\langle c(x,\cdot),Rc(y,\cdot)\rangle
 =-\kappa(x,y)=b^-_{\be,\mu,\vep}(x,y).
\end{equation}
This proves \eqref{rep:eq:first-profile-minus} and, in particular, the positive sign of $\be\psi^-_\mu$ in \eqref{rep:eq:a-minus}.

\subsection{Structural consequences and dependence on the marginal size}\label{rep:sec:consequences}

We record the deductions used in \Cref{rep:cor:consequences}. They are algebraic once \eqref{rep:eq:marginal-all-order} holds in every finite $L^p$.

For a coefficient order $j$, define its canonical kernels by subset-lattice M\"obius inversion:
\begin{equation}\label{rep:eq:hoeffding-coeff}
 h^-_{j,r}(x_1,\ldots,x_r)
 =\sum_{A\subseteq[r]}(-1)^{r-|A|}f^-_{j,|A|}(x_A).
\end{equation}
The marked Gaussian recursion shows that $h^-_{j,r}=0$ for $r>2j$, and
\begin{equation}\label{rep:eq:coefficient-hoeffding-sum}
 f^-_{j,k}
 =\sum_{r=1}^{\min\{2j,k\}}\ \sum_{\substack{A\subseteq[k]\\|A|=r}}h^-_{j,r}(x_A).
\end{equation}
For an $r$-variable canonical kernel, every nonzero term in an even moment of the sum over $r$-subsets uses each particle index at least twice. Consequently,
\begin{equation}\label{rep:eq:canonical-subset-sum}
 \left\|\sum_{|A|=r}h(x_A)\right\|_{L^p(\mu^{\otimes k})}
 \le C_{r,p}k^{r/2}\|h\|_{L^{p'}(\mu^{\otimes r})}
\end{equation}
for a fixed $p'$ depending on $p$ and $r$. Equations \eqref{rep:eq:coefficient-hoeffding-sum} and $r\le2j$ prove \eqref{rep:eq:k-polynomial}.  In the exact remainder argument, the pointwise upper bound for the marked source costs at most $e^{Ck}$, while every retained weighted Wick polynomial has degree bounded in terms of $M$ and hence grows only polynomially in $k$.  The same count in the bad-event estimate proves \eqref{rep:eq:k-remainder}.

The exact $r$th Hoeffding projection of the marginal hierarchy is
\begin{equation}\label{rep:eq:exact-hoeffding}
 H^-_{N,r}=\sum_{A\subseteq[r]}(-1)^{r-|A|}F_{N}^{-,(|A|)}(x_A).
\end{equation}
Thus,
\begin{equation}\label{rep:eq:hoeffding-all-order}
 H^-_{N,r}
 =\sum_{j=\lceil r/2\rceil}^{M}\frac{h^-_{j,r}}{N^j}
 +O_{r,M,p}(N^{-M-1})
\end{equation}
in every finite $L^p$.

The connected correlation density of the labeled marginals is obtained by the cumulant formula, equivalently partition-lattice M\"obius inversion:
\begin{equation}\label{rep:eq:ursell}
 U^-_{N,r}
 =\sum_{\pi\in\Pi_r}(-1)^{|\pi|-1}(|\pi|-1)!
 \prod_{B\in\pi}F_N^{-,(|B|)}(x_B).
\end{equation}
Substitution of \eqref{rep:eq:marginal-all-order} gives an all-order expansion.  At a strictly positive cutoff parameter, use the coefficient recursion \eqref{eq:prescribed-ell-recursion} and the Bell-polynomial formula \eqref{eq:prescribed-Bell-polynomial}, followed by Gaussian expectation using Wick's rule.  Each coefficient is a finite sum of dressed contraction diagrams.  A
\emph{dressed} two-point line may contain the resolvent or $\kappa$; for the sole
purpose of the $N$-power count it can be viewed as the resummation of chains of
pair bonds and degree-two unmarked vertices.  Every such insertion adds one
bond and one summed vertex and therefore preserves the excess $e-v$; no
Neumann-series convergence is being asserted.

After division by the scalar series, components containing no marked vertex (vacuum components) cancel.  Partition-
lattice M\"obius inversion then cancels every remaining diagram whose marked
vertices do not lie in one connected component.  In this particle-index count, vertices carry particle labels, pair bonds are occurrences of two-body interaction factors, and a one-body insertion comes from the source term with coefficient $N^{-1}$. For a surviving diagram let
$r$ be the number of marked vertices, $v$ the number of distinct summed
unmarked vertices, $e$ the number of pair-bond occurrences in its undressed
topological count, and $\ell$ the number of weak one-body insertions.  The
mean-field scaling is
\begin{equation}
 N^{v-e-\ell}.
\end{equation}
The pair-bond multigraph spanning the relevant vertices is connected, hence
\begin{equation}
 e\ge r+v-1,
 \qquad
 N^{v-e-\ell}\le N^{-(r-1)-\ell}.
\end{equation}
One-body insertions cannot connect components and only lower the order.  Thus
the first possible order of the normalized $r$-point connected density is
$N^{-(r-1)}$.  The established finite-diagram cutoff convergence passes the
M\"obius cancellation and this edge-and-vertex count to the singular
kernel coefficient by coefficient.

For the divergences, write $F_N^{-,(k)}=1+u_N$. The global scalar remainder
\begin{equation}\label{rep:eq:entropy-global-remainder}
 \left|(1+u)\log(1+u)-u-
 \sum_{r=2}^{L}\frac{(-1)^r}{r(r-1)}u^r\right|
 \le C_L|u|^{L+1},\qquad u\ge-1,
\end{equation}
combined with $\|u_N\|_{L^{L+1}}=O_k(N^{-1})$ proves every finite-order relative-entropy expansion. Equations \eqref{rep:eq:chi-leading}--\eqref{rep:eq:tv-leading} follow from the first coefficient and Hoeffding orthogonality.

For bounded linear statistics, use \Cref{rep:lem:linear-source-control}, whose exponential bound follows from the signed-gas theorem in Appendix~\ref{app:signed-log-heat-kernel}. Apply it to the insertion $N^{-1/2}\sum_i\sum_a z_a\varphi_a(x_i)$. The resulting expansion is uniform on a fixed complex polydisc, so Cauchy differentiation gives every fixed joint moment and cumulant expansion. Odd half-powers may occur because the external source breaks field parity. Completing the square under the quadratic Gaussian tilt, of covariance $\be T_\mu(\Id+\be T_\mu)^{-1}$, gives \eqref{rep:eq:linear-covariance}.




\section{A heat-kernel Sine--Gordon proof for signed logarithmic gases}
\label{app:signed-log-heat-kernel}

We prove a partition estimate uniform over arbitrary deterministic charge vectors, and then derive marked expansions for the signed logarithmic gas with a common prescribed background. The analytic estimate uses a heat-semigroup decomposition, oscillatory cancellation for the long-range field, and a partial-matching argument for the short-range interaction. Together with the finite-cutoff identities from Appendix~\ref{app:repulsive-marked-oscillatory}, it also gives a second proof of the one-component partition and marginal expansions without the repulsive partition theorem of \cite{DelgadinoGvalaniRosenzweig2026}.

\subsection{Common-background deterministic-charge theorem}
\label{rep:subsec:signed-theorem}
\begin{thm}[Arbitrary deterministic charge imbalances]\label{rep:thm:arbitrary-charge}
Fix $\d\ge2$, $M_0<\infty$, and $\be_0<\infty$.  For each $N$, let
\begin{equation}\label{rep:eq:charge-vector-intro}
 \boldsymbol\sigma=(\sigma_1,\ldots,\sigma_N)\in\{-1,1\}^N,
 \qquad
 \vartheta_N=\frac1N\sum_{i=1}^N\sigma_i.
\end{equation}
The \emph{charge imbalance} $\vartheta_N$ is the fraction of positive charges minus the fraction of negative charges; it is not either species fraction.
Define the diagonal-free signed modulated energy around the prescribed charge profile $\vartheta_N\mu$ by
\begin{equation}\label{rep:eq:signed-modulated-energy-intro}
 \Fr_N^{\boldsymbol\sigma}(X_N,\mu)
 =\frac1{2N^2}\sum_{i\ne j}\sigma_i\sigma_j\g(x_i-x_j)
 -\frac{\vartheta_N}{N}\sum_{i=1}^N\sigma_i h^-_\mu(x_i)
 +\vartheta_N^2I^-_\mu.
\end{equation}
Then
\begin{equation}\label{rep:eq:signed-exact-reduction-intro}
 -\be N\Fr_N^{\boldsymbol\sigma}
 =-\frac\be N\sum_{i<j}\sigma_i\sigma_jq^-_\mu(x_i,x_j)
 +\frac\be N\sum_{i=1}^N\psi^-_\mu(x_i)+\be I^-_\mu.
\end{equation}

For every fixed expansion order $M$, the associated partition function satisfies
\begin{equation}\label{rep:eq:signed-partition-all-order-intro}
 \mathcal K_{N,\be}^{\boldsymbol\sigma}(\mu)
 =e^{\be I^-_\mu}\det{}_2(\Id+\be T_\mu)^{-1/2}
 \left(\sum_{j=0}^{M}\frac{z_j(\be,\mu,\vartheta_N)}{N^j}
 +O_M(N^{-M-1})\right),
 \qquad z_0=1.
\end{equation}
The error is uniform over all references satisfying \eqref{rep:eq:mu-class}, all $0\le\be\le\be_0$, and all charge vectors.  For each fixed $j$, the coefficient $z_j(\be,\mu,\vartheta)$ is a polynomial in $\vartheta$.

Fix a marked sign pattern $\boldsymbol\tau=(\tau_1,\ldots,\tau_k)\in\{-1,1\}^k$.  Whenever the full charge vector contains these marked signs, the corresponding species-resolved marginal density relative to $\mu^{\otimes k}$ has, for every fixed $M$ and finite $p$, an expansion
\begin{equation}\label{rep:eq:signed-marginal-all-order-intro}
 \left\|F_{N,\be,\mu}^{\boldsymbol\sigma,\boldsymbol\tau,(k)}
 -1-\sum_{j=1}^{M}\frac{f_{j,k,\boldsymbol\tau}(\be,\mu,\vartheta_N)}{N^j}
 \right\|_{L^p(\mu^{\otimes k})}
 \le C_{M,k,p}N^{-M-1}.
\end{equation}
The estimate holds for all $N\ge N_0(\d,M_0,\be_0,M,k,p)$; this lower bound ensures the $L^p$ integrability of every opposite-charge collision.  The coefficients are polynomials in the exact finite-$N$ charge imbalance $\vartheta_N$.

Set
\begin{equation}\label{rep:eq:kappa-intro}
 \kappa_{\be,\mu}=\be T_\mu R^-_{\be,\mu}=\Id-R^-_{\be,\mu}.
\end{equation}
For a marked sign $\tau\in\{-1,1\}$ and a charge imbalance $\vartheta\in[-1,1]$, define
\begin{equation}\label{rep:eq:species-singleton-intro}
 a_{\tau,\vartheta}
 =\left(\Id-\tau\vartheta\,\kappa_{\be,\mu}\right)d^-_{\be,\mu}.
\end{equation}
Then the first marginal coefficient is
\begin{equation}\label{rep:eq:signed-first-profile-intro}
 f_{1,k,\boldsymbol\tau}(x_1,\ldots,x_k;\vartheta)
 =\sum_{a=1}^k a_{\tau_a,\vartheta}(x_a)
 +\sum_{a<b}\tau_a\tau_b b^-_{\be,\mu}(x_a,x_b).
\end{equation}
The first scalar correction is
\begin{multline}\label{rep:eq:signed-z1-intro}
 z_1(\be,\mu,\vartheta)
 =\frac12\left\langle d^-_{\be,\mu},
 (\Id-\vartheta^2\kappa_{\be,\mu})d^-_{\be,\mu}\right\rangle
 -\frac14\Tr(\kappa_{\be,\mu}^2)\\
 +\frac{\vartheta^2}{12}
 \iint b^-_{\be,\mu}(x,y)^3\,\dd\mu(x)\dd\mu(y).
\end{multline}
At $\vartheta=1$ with positive marked charges, \eqref{rep:eq:signed-first-profile-intro} reduces to the one-component coefficient \eqref{rep:eq:first-profile-minus}; at $\vartheta=0$, it gives the signed two-component profile.
\end{thm}

By uniqueness of asymptotic coefficients, the specialization $\vartheta_N=1$ with all marked charges positive recovers the one-component expansion of \Cref{rep:thm:main}; we therefore do not restate its one-component corollaries in this appendix.

\begin{remark}[Exact charge imbalance and rounding]\label{rep:rem:exact-fraction}
The coefficients in \eqref{rep:eq:signed-partition-all-order-intro}--\eqref{rep:eq:signed-marginal-all-order-intro} are evaluated at the exact charge imbalance $\vartheta_N$.  An expansion with coefficients depending only on a limit $\vartheta$ requires an inverse-$N$ expansion of $\vartheta_N$ itself.  The usual integer rounding of the species counts can otherwise create a bounded, oscillatory order-$N^{-1}$ correction.
\end{remark}

\subsection{The heat-kernel partition estimate}\label{rep:sec:heat-route}

We now prove a cutoff-uniform exponential estimate using oscillatory cancellation.  The argument follows the field-theoretic organization of Grotto and Romito \cite{GrottoRomito2020CLT,GrottoRomito2020Correlation} but uses the periodic heat semigroup rather than a Yukawa/Bessel-potential decomposition.  The resulting estimate is uniform over every deterministic charge vector.

\subsubsection{The heat-semigroup decomposition}

Let
\begin{equation}\label{rep:eq:periodic-heat-kernel}
 p_t(x)=\sum_{n\in\Z^{\d}}e^{-4\pi^2t|n|^2}e^{2\pi\mathrm{i} n\cdot x}
\end{equation}
be the periodic heat kernel and set
\begin{equation}\label{rep:eq:ad-heat}
 a_{\d}=\frac{c_{\d}}{\Gamma(\d/2)}=2^{\d-1}\pi^{\d/2}.
\end{equation}
Fourier inversion gives
\begin{equation}\label{rep:eq:heat-representation-g}
 \g(x)=a_{\d}\int_0^\infty\bigl(p_t(x)-1\bigr)t^{\d/2-1}\,\dd t.
\end{equation}
For $m\ge2$, define
\begin{align}
 S_m(x)&=a_{\d}\int_0^{m^{-2}}p_t(x)t^{\d/2-1}\,\dd t,\label{rep:eq:Sm-def}\\
 \kappa_m&=a_{\d}\int_0^{m^{-2}}t^{\d/2-1}\,\dd t
 =\frac{2a_{\d}}{\d}m^{-\d},\label{rep:eq:kappam-def}\\
 V_m(x)&=a_{\d}\int_{m^{-2}}^\infty\bigl(p_t(x)-1\bigr)t^{\d/2-1}\,\dd t.\label{rep:eq:Vm-heat-def}
\end{align}
Then
\begin{equation}\label{rep:eq:heat-split}
 \g=V_m+S_m-\kappa_m,
 \qquad
 q^-_\mu=q_{V_m,\mu}+q_{S_m,\mu}.
\end{equation}
The constant disappears under Hoeffding centering.

\begin{lemma}[Heat-kernel bounds]\label{rep:lem:heat-split-bounds}
The kernel $S_m$ is nonnegative, $V_m$ has nonnegative nonzero Fourier coefficients, and $q_{V_m,\mu}$ is a positive trace-class covariance.  Moreover,
\begin{align}
 \int_{\T^{\d}}S_m\,\dd\muunif&=\kappa_m,\label{rep:eq:Sm-mass}\\
 V_m(0)&=\log m+O_{\d}(1),\label{rep:eq:Vm-diagonal}\\
 \|S_m*\mu\|_\infty+|I_{S_m,\mu}|&\le C_{\d,M_0}m^{-\d},\label{rep:eq:Sm-centering}\\
 \|q_{V_m,\mu}\|_{\HS}&\le C_{\d,M_0},\label{rep:eq:Vm-HS}\\
 q_{V_m,\mu}(x,x)&=\log m+O_{\d,M_0}(1),\label{rep:eq:Vm-centered-diagonal}\\
 \left\|q_{V_m,\mu}(\cdot,\cdot)\big|_{x=x}
 -\int q_{V_m,\mu}(y,y)\,\dd\mu(y)\right\|_{L^\infty_x}
 &\le C_{\d,M_0}.
 \label{rep:eq:Vm-diagonal-oscillation}
\end{align}
Finally, with $E_1(s)=\int_s^\infty e^{-u}\,\dd u/u$, for $x\ne0$,
\begin{equation}\label{rep:eq:Sm-E1}
 S_m(x)=\frac12\sum_{z\in\Z^{\d}}
 E_1\left(\frac{m^2|x+z|^2}{4}\right),
 \qquad S_m(0)=+\infty.
\end{equation}
Consequently, if $r=\operatorname{dist}(x,0)$, then
\begin{equation}\label{rep:eq:Sm-pointwise}
 S_m(x)\le
 \begin{cases}
 -\log(mr)+C_{\d},&0<mr\le1,\\
 C_{\d}e^{-c_{\d}m^2r^2}+C_{\d}e^{-c_{\d}m^2},&mr\ge1.
 \end{cases}
\end{equation}
\end{lemma}

\begin{proof}
The Fourier coefficient of the right side of \eqref{rep:eq:heat-representation-g} at $n\ne0$ is
\begin{equation}\label{rep:eq:heat-fourier-check}
 a_{\d}\int_0^\infty e^{-4\pi^2t|n|^2}t^{\d/2-1}\,\dd t
 =\frac{c_{\d}}{(2\pi|n|)^{\d}},
\end{equation}
which proves \eqref{rep:eq:heat-representation-g}.  Positivity and \eqref{rep:eq:Sm-mass} follow directly from the heat kernel.  The nonzero Fourier coefficients of $V_m$ are positive and decay exponentially at high frequency, so the centered operator is positive and trace class.

Poisson summation gives
\begin{equation}\label{rep:eq:heat-poisson}
 p_t(x)=\sum_{z\in\Z^{\d}}(4\pi t)^{-\d/2}
 e^{-|x+z|^2/(4t)}.
\end{equation}
Since $a_{\d}(4\pi)^{-\d/2}=1/2$, the substitution $u=|x+z|^2/(4t)$ proves \eqref{rep:eq:Sm-E1}.  Choose a nearest lift $x_0\in[-1/2,1/2]^{\d}$, so $|x_0|=r$.  Every other lift has distance at least $1/2$, and the lattice sum of their contributions is $O_{\d}(e^{-c_{\d}m^2})$.  The standard estimates
\begin{equation}\label{rep:eq:E1-estimates}
 \frac12E_1(s)\le-\frac12\log s+C\quad(0<s\le1/4),
 \qquad
 E_1(s)\le Ce^{-s/2}\quad(s\ge1)
\end{equation}
(and a change of constants on the compact interval $[1/4,1]$) yield \eqref{rep:eq:Sm-pointwise}, with the exact logarithmic coefficient one in the first line.  Integrating $p_t$ proves \eqref{rep:eq:Sm-mass}, and hence \eqref{rep:eq:Sm-centering} under $\rho\le M_0$.

For \eqref{rep:eq:Vm-diagonal}, use \eqref{rep:eq:heat-poisson} on $m^{-2}\le t\le1$: the zero image contributes
\begin{equation}\label{rep:eq:Vm-log}
 a_{\d}(4\pi)^{-\d/2}\int_{m^{-2}}^1\frac{\dd t}{t}
 =\log m,
\end{equation}
while the nonzero images and the range $t\ge1$ are uniformly integrable.  Since $V_m=\g-S_m+\kappa_m$, its $L^1$ norm is uniformly bounded.  Therefore, $V_m*\mu$ and $I_{V_m,\mu}$ are uniformly bounded, which proves \eqref{rep:eq:Vm-centered-diagonal}--\eqref{rep:eq:Vm-diagonal-oscillation}.  Finally, the nonzero Fourier multiplier of $V_m$ is bounded by that of $\g$, and centering is bounded on $L^2(\mu^{\otimes2})$ under $\rho\le M_0$; this proves \eqref{rep:eq:Vm-HS}.
\end{proof}

\subsubsection{A heat-kernel dipole inequality}

\begin{lemma}[Dipole difference]\label{rep:lem:heat-dipole}
There are constants $c_{\d},C_{\d}>0$ such that, for every $m\ge2$ and $x,z,y\in\T^{\d}$ satisfying
\begin{equation}\label{rep:eq:dipole-assumption}
 \delta=\operatorname{dist}(x,z)\le\operatorname{dist}(z,y),
\end{equation}
one has
\begin{equation}\label{rep:eq:heat-dipole-bound}
 \bigl[S_m(z-y)-S_m(x-y)\bigr]_+
 \le C_{\d}e^{-c_{\d}m^2\delta^2}+C_{\d}e^{-c_{\d}m^2}.
\end{equation}
\end{lemma}

\begin{proof}
We use the convention that the left side is zero when $x=z$, including the common singular case $x=z=y$; this is the value obtained by cancelling the identical terms, and the exceptional collision set is null under every admissible product measure.  Thus assume $\delta>0$.  By translation, take $y=0$.  Choose a nearest lift $z_0\in[-1/2,1/2]^{\d}$ of $z$, with $|z_0|=r=\operatorname{dist}(z,0)$, and choose a minimal displacement $h$ from $z$ to $x$, so $|h|=\delta$ and $x_0=z_0+h$ is a lift of $x$.  Then $|x_0|\le r+\delta\le2r$.

For every nonzero $k\in\Z^{\d}$, $|z_0+k|\ge1/2$; hence, by the large-argument bound for $E_1$ and a Gaussian lattice-sum estimate,
\begin{equation}\label{rep:eq:dipole-image-tail}
 \sum_{k\ne0}E_1\left(\frac{m^2|z_0+k|^2}{4}\right)
 \le C_{\d}e^{-c_{\d}m^2}.
\end{equation}
In the positive difference, discard every nonnegative term of $S_m(x)$ except the image corresponding to $x_0$.  If $x_0=0$, the resulting difference is nonpositive.  Otherwise, with $s=m^2r^2/4$ and using that $E_1$ is decreasing,
\begin{align}
 S_m(z)-S_m(x)
 &\le\frac12\left[E_1(s)-E_1\left(\frac{m^2|x_0|^2}{4}\right)\right]
 +C_{\d}e^{-c_{\d}m^2}\notag\\
 &\le\frac12\bigl[E_1(s)-E_1(4s)\bigr]+C_{\d}e^{-c_{\d}m^2}.
 \label{rep:eq:dipole-E1-step}
\end{align}
(The first bracket is negative when $|x_0|<r$.)  Finally,
\begin{equation}\label{rep:eq:E1-fourfold}
 E_1(s)-E_1(4s)=\int_s^{4s}\frac{e^{-u}}u\,\dd u
 \le\min\{\log4,Ce^{-cs}\}.
\end{equation}
Since $r\ge\delta$, \eqref{rep:eq:heat-dipole-bound} follows.
\end{proof}

\subsubsection{The arbitrary-sign partition bound}

For $L\ge2$ and $\sigma=(\sigma_1,\ldots,\sigma_L)\in\{-1,1\}^L$, define
\begin{equation}\label{rep:eq:signed-Z}
 \mathcal Z_{L,\lambda}^{\sigma}(\mu)
 =\int_{(\T^{\d})^L}\exp\left\{-\frac\lambda L
 \sum_{a<b}\sigma_a\sigma_bq^-_\mu(x_a,x_b)\right\}
 \,\dd\mu^{\otimes L}.
\end{equation}

\begin{thm}[Uniform signed-gas bound]\label{rep:thm:signed-bound}
For every $\d\ge2$, $M_0<\infty$, and $\Lambda<\infty$, there are $L_0$ and $C<\infty$ such that
\begin{equation}\label{rep:eq:signed-uniform}
 \sup_{\substack{L\ge L_0,\ 0\le\lambda\le\Lambda\\
 \sigma\in\{-1,1\}^L\\
 \mu=\rho\muunif,\ 0\le\rho\le M_0,\ \int_{\T^{\d}}\rho\,\dd\muunif=1}}
 \mathcal Z_{L,\lambda}^{\sigma}(\mu)\le C.
\end{equation}
The same bound holds for every conditional-expectation cutoff $q^-_{\mu,j}=Q_jT_\mu Q_j$ associated with a finite measurable partition.
\end{thm}

\begin{proof}
By \eqref{rep:eq:heat-split} and Cauchy--Schwarz, it suffices to bound the signed partition functions with kernels $q_{V_m,\mu}$ and $q_{S_m,\mu}$ at coupling $2\lambda$.  In the two factor estimates below we relabel this factor coupling as $\lambda$; thus it ranges over $[0,2\Lambda]$, and every choice of constants, expansion depth, and $L_0$ is made on that doubled window.

\step{Long-range factor} Let $F_m$ be the centered real Gaussian field with covariance $q_{V_m,\mu}$ and write
\begin{equation}\label{rep:eq:Dm-def}
 D_m(x)=q_{V_m,\mu}(x,x),\qquad
 \overline D_m=\int D_m\,\dd\mu,\qquad
 \delta D_m=D_m-\overline D_m.
\end{equation}
Set $t=(\lambda/L)^{1/2}$, with the case $\lambda=0$ understood by continuity, and define
\begin{equation}\label{rep:eq:Ahat-heat}
 \widehat A_\sigma
 =\int e^{\mathrm{i}\sigma tF_m(x)+t^2\delta D_m(x)/2}\,\dd\mu(x),
 \qquad
 E_m=e^{-t^2\|F_m\|_{L^2(\mu)}^2/2}.
\end{equation}
The oscillatory Hubbard--Stratonovich transformation gives
\begin{equation}\label{rep:eq:long-field-heat}
 \mathcal Z_{L,\lambda}^{\sigma,V_m}
 =e^{\lambda\overline D_m/2}\E\prod_{a=1}^L\widehat A_{\sigma_a}.
\end{equation}
The two factors in \eqref{rep:eq:Ahat-heat} have the same constant and quadratic Taylor terms, while their linear terms vanish.  Taylor's formula with integral remainder, \eqref{rep:eq:Vm-diagonal-oscillation}, and $|e^{\mathrm{i} r}|=1$ give
\begin{equation}\label{rep:eq:Ahat-error-heat}
 |\widehat A_\sigma-E_m|
 \le C_\Lambda t^3\left(1+\|F_m\|_{L^4(\mu)}^4\right),
 \qquad
 |\widehat A_\sigma|\le e^{Ct^2},
 \qquad 0<E_m\le1.
\end{equation}
For every fixed $r,s<\infty$,
\begin{equation}\label{rep:eq:Fm-moments-heat}
 \E\|F_m\|_{L^r(\mu)}^s
 \le C_{r,s,\d,M_0}(1+\log m)^{s/2},
\end{equation}
by \eqref{rep:eq:Vm-centered-diagonal} and Gaussian moments.

Put $\Delta_a=\widehat A_{\sigma_a}-E_m$.  Taylor's theorem with integral remainder for the polynomial $z\mapsto\prod_a(E_m+z\Delta_a)$ yields, for every fixed $p\ge1$,
\begin{equation}\label{rep:eq:long-product-expansion}
 \prod_{a=1}^L\widehat A_{\sigma_a}
 =\sum_{j=0}^{p-1}E_m^{L-j}\sum_{|A|=j}\prod_{a\in A}\Delta_a
 +\mathcal R_p,
\end{equation}
and
\begin{equation}\label{rep:eq:long-product-remainder}
 |\mathcal R_p|
 \le\binom Lp\left[Ct^3(1+\|F_m\|_4^4)\right]^p e^C.
\end{equation}
For the $j$th retained term, set $a_j=\lambda(1-j/L)$.  Tilting the Gaussian law by $e^{-a_j\|F_m\|_2^2/2}$ replaces its covariance by
\begin{equation}\label{rep:eq:long-tilted-covariance}
 T_{V_m,\mu}(\Id+a_jT_{V_m,\mu})^{-1}\le T_{V_m,\mu}.
\end{equation}
Moreover,
\begin{align}
 e^{\lambda\overline D_m/2}
 \E e^{-a_j\|F_m\|_2^2/2}
 &=e^{(\lambda-a_j)\overline D_m/2}
 \det{}_2(\Id+a_jT_{V_m,\mu})^{-1/2}\notag\\
 &\le C_{\Lambda,\d,M_0}
 \exp\left\{C_{\Lambda,j}\frac{\log m}{L}\right\},
 \label{rep:eq:long-det-bound}
\end{align}
where we used \eqref{rep:eq:Vm-HS} and $x-\log(1+x)\le x^2/2$.  It follows from \eqref{rep:eq:Ahat-error-heat}--\eqref{rep:eq:long-det-bound} that the sum of all terms with a fixed $j\ge1$ is bounded by
\begin{equation}\label{rep:eq:long-j-bound}
 C_{\Lambda,j}L^{-j/2}(1+\log m)^{2j}
 \exp\left\{C_{\Lambda,j}\frac{\log m}{L}\right\}.
\end{equation}
The $j=0$ term is the uniformly bounded regularized determinant.  The remainder satisfies
\begin{equation}\label{rep:eq:long-rem-bound}
 e^{\lambda\overline D_m/2}\E|\mathcal R_p|
 \le C_{\Lambda,p}L^{-p/2}m^{\Lambda+C/L}(1+\log m)^{2p}.
\end{equation}
Choose $m=L^a$ for a fixed $a>0$, and then choose an integer $p>2a\Lambda+2$ and enlarge $L_0$ so that $L_0>p$.  Equations \eqref{rep:eq:long-j-bound}--\eqref{rep:eq:long-rem-bound} give a uniform long-range bound on the doubled coupling window.

\step{Short-range factor} Let $P$ and $Q$ be the numbers of positive and negative charges, and assume $Q\le P$ after a global reversal if necessary.  Greedily choose a closest remaining positive--negative pair, breaking ties by a fixed lexicographic rule, and label the selected pairs $(y_i,z_i)$, $1\le i\le Q$; write $y_{Q+1},\ldots,y_P$ for the unpaired positive charges.  If $\delta_i=\operatorname{dist}(y_i,z_i)$, then, at the time the $i$th pair is selected,
\begin{equation}\label{rep:eq:greedy-heat}
 \delta_i\le\operatorname{dist}(z_i,y_j),\qquad
 \delta_i\le\operatorname{dist}(y_i,z_j)
 \qquad(j>i).
\end{equation}
The raw signed short-range Hamiltonian has the exact decomposition
\begin{multline}\label{rep:eq:short-dipole-decomposition}
 \sum_{a<b}\sigma_a\sigma_bS_m(x_a-x_b)
 =-\sum_{i=1}^Q S_m(y_i-z_i)\\
 +\sum_{1\le i<j\le Q}
 \bigl[S_m(y_i-y_j)-S_m(z_i-y_j)
 +S_m(z_i-z_j)-S_m(y_i-z_j)\bigr]\\
 +\sum_{1\le i\le Q<j\le P}
 \bigl[S_m(y_i-y_j)-S_m(z_i-y_j)\bigr]
 +\sum_{Q<i<j\le P}S_m(y_i-y_j).
\end{multline}
The last term is nonnegative.  Applying \Cref{rep:lem:heat-dipole} to every difference in which the $i$th dipole meets a later charge gives
\begin{equation}\label{rep:eq:short-raw-lower}
 \sum_{a<b}\sigma_a\sigma_bS_m(x_a-x_b)
 \ge-\sum_{i=1}^QS_m(y_i-z_i)
 -CL\sum_{i=1}^Qe^{-cm^2\delta_i^2}-CL^2e^{-cm^2}.
\end{equation}
Writing $s_\sigma=\sum_a\sigma_a$, the centered Hamiltonian is
\begin{multline}\label{rep:eq:centered-signed-short}
 \sum_{a<b}\sigma_a\sigma_bq_{S_m,\mu}(x_a,x_b)
 =\sum_{a<b}\sigma_a\sigma_bS_m(x_a-x_b)
 -s_\sigma\sum_a\sigma_a(S_m*\mu)(x_a)\\
 +\sum_a(S_m*\mu)(x_a)+I_{S_m,\mu}(s_\sigma^2-L).
\end{multline}
Equation \eqref{rep:eq:Sm-centering} therefore changes the right side of \eqref{rep:eq:short-raw-lower} by at most $CL^2m^{-\d}$.

Define
\begin{equation}\label{rep:eq:heat-edge-weight}
 w_m(y,z)=
 \exp\left\{\frac\lambda L S_m(y-z)
 +C_\Lambda e^{-cm^2\operatorname{dist}(y,z)^2}+C_\Lambda e^{-cm^2}\right\}-1.
\end{equation}
For $L\ge L_0(\d,\Lambda)$,
\begin{equation}\label{rep:eq:heat-edge-integral}
 \iint w_m(y,z)\,\dd\mu(y)\dd\mu(z)
 \le C_{\d,M_0,\Lambda}m^{-\d}.
\end{equation}
Indeed, enlarge $L_0$ so that $\lambda/L\le \d/2$ throughout the doubled window.  On $r=\operatorname{dist}(y,z)\le m^{-1}$, the first line of \eqref{rep:eq:Sm-pointwise} gives
\begin{equation}
 w_m(y,z)\le C(mr)^{-\lambda/L},
\end{equation}
and therefore, using $\rho\le M_0$ and polar coordinates, its integral over this region is at most
$Cm^{-\d}\int_0^1s^{\d-1-\lambda/L}\,\dd s\le Cm^{-\d}$.  On the complement the exponent defining $w_m$ is uniformly bounded, so $e^u-1\le C u$; the Gaussian tail in \eqref{rep:eq:Sm-pointwise}, together with $e^{-cm^2}\le Cm^{-\d}$, again has total mass $O(m^{-\d})$.

By \eqref{rep:eq:short-raw-lower}--\eqref{rep:eq:centered-signed-short}, the short-range Boltzmann factor is at most
\begin{equation}\label{rep:eq:selected-matching-product}
 e^{C_\Lambda Lm^{-\d}}
 \prod_{i=1}^Q\bigl(1+w_m(y_i,z_i)\bigr).
\end{equation}
Although the greedy matching depends on the configuration, expansion of the selected product is bounded pointwise by the sum of the products over all partial bipartite matchings.  A $j$-edge matching has disjoint endpoints, so its integral is the $j$th power of the left side of \eqref{rep:eq:heat-edge-integral}; the number of such matchings is at most $L^{2j}/j!$.  Hence
\begin{equation}\label{rep:eq:heat-short-final}
 \mathcal Z_{L,\lambda}^{\sigma,S_m}
 \le\exp\left\{C_{\d,M_0,\Lambda}(L^2+L)m^{-\d}\right\}.
\end{equation}
Choose the exponent $a$ above so that $ad>2$.  Then \eqref{rep:eq:heat-short-final} is uniform.  Together with the long-range estimate and Cauchy--Schwarz, this proves \eqref{rep:eq:signed-uniform}.

For a finite measurable partition, let $\mathcal G_j$ be the sigma-field generated by the cell labels of all particles.  The kernel of $Q_jT_\mu Q_j$ is the two-cell conditional average of $q^-_\mu$, and hence, for
$H=\sum_{a<b}\sigma_a\sigma_bq^-_\mu(X_a,X_b)$,
\begin{equation}
 H_j=\sum_{a<b}\sigma_a\sigma_bq^-_{\mu,j}(X_a,X_b)
 =\E[H\mid\mathcal G_j].
\end{equation}
Since $H\in L^1(\mu^{\otimes L})$ and is finite almost surely, conditional Jensen in the extended-expectation sense therefore gives
\begin{equation}\label{rep:eq:conditional-Jensen-signed}
 \int e^{-\lambda H_j/L}\,\dd\mu^{\otimes L}
 \le\int e^{-\lambda H/L}\,\dd\mu^{\otimes L},
\end{equation}
which proves the cutoff assertion.
\end{proof}

\subsubsection{A cutoff-uniform expansion with finitely many marked coordinates}

Use conditional-expectation cutoffs $q^-_{\mu,j}$ and $s_j=Q_j(\be\psi^-_\mu)$.  Let $A^-_{j,u}$ be \eqref{rep:eq:Aminus} with covariance $\be q^-_{\mu,j}$.

\begin{prop}[Cutoff-uniform expansion with fixed marked coordinates]\label{rep:prop:exact-remaining}
Fix $k,M,p$, $\be_0<\infty$, and the bounded-density class \eqref{rep:eq:mu-class}.  For all sufficiently large $N$,
\begin{equation}\label{rep:eq:exact-exponential-moment}
 \sup_{j,\mu,\,0\le\be\le\be_0}
 \operatorname*{ess\,sup}_{x\in(\T^{\d})^k}
 \E\left[\left|A^-_{j,u_x^-}(N^{-1/2})\right|^{2(N-k)}\right]
 <\infty.
\end{equation}
Consequently, the marked expectation in \eqref{rep:eq:marked-identity} admits, uniformly in the cutoff, an expansion through order $N^{-M}$ with an $O(N^{-M-1})$ remainder in $L^p(\mu^{\otimes k})$.  Every deterministic coefficient converges in the required finite $L^p$ spaces as the partition is refined.
\end{prop}

\begin{proof}
Put $L_N=N-k$.  Expanding the even absolute moment introduces $L_N$ positive and $L_N$ negative auxiliary charges.  Wick's formula gives
\begin{multline}\label{rep:eq:even-moment-signed}
 \E|A^-_{j,u_x^-}(N^{-1/2})|^{2L_N}\\
 =\int\exp\left\{-\frac\be N
 \sum_{a<b}\sigma_a\sigma_bq^-_{\mu,j}(z_a,z_b)
 +\frac1N\sum_{a=1}^{2L_N}u_x^-(z_a)\right\}
 \,\dd\mu^{\otimes2L_N}(z),
\end{multline}
with exactly $L_N$ signs of each type.  By \eqref{rep:eq:u-upper}, the source is bounded above.  The interaction is \eqref{rep:eq:signed-Z} at
\begin{equation}\label{rep:eq:effective-lambda}
 \lambda_N=2\be\frac{L_N}{N}\le2\be_0.
\end{equation}
The signed-gas theorem proves \eqref{rep:eq:exact-exponential-moment}.

The good-event logarithmic expansion is identical to the proof of \Cref{rep:prop:moving-cutoff}, except that no deterministic cutoff cost occurs.  The conditional-expectation kernels inherit the row-exponential bound and the uniform lower bound from $q^-_\mu$ by Jensen, so all weighted Wick estimates used there remain uniform in $j$.  Let $E_{t,x}$ be the smaller event defined by the three conditions in \eqref{rep:eq:very-good}, with the conditional-expectation cutoff $j$ used here in place of the heat cutoff. Then, for every $Q$,
$\sup_x\mathbb P(E_{t,x}^c)\le C_Qt^Q$, and Cauchy--Schwarz gives the explicit replacement for \eqref{rep:eq:bad-overexpand}:
\begin{equation}
 \E\big[\mathbf1_{E_{t,x}^c}|A^-_{j,u_x^-}(t)|^{N-k}\big]
 \le \mathbb P(E_{t,x}^c)^{1/2}\,
 \Big(\E|A^-_{j,u_x^-}(t)|^{2(N-k)}\Big)^{1/2}.
\end{equation}
The second factor is uniformly bounded by \eqref{rep:eq:exact-exponential-moment}; the retained polynomial--quadratic terms are handled by \eqref{rep:eq:quadratic-exp-bound}.  Choosing $Q$ and the finite logarithmic expansion depth as in \Cref{rep:prop:moving-cutoff}, expanding through total degree $2M+1$, and using invariance under $\phi\mapsto-\phi$ leave an $O(N^{-M-1})$ remainder and only integer powers of $N^{-1}$.

For cutoff removal, choose refining partitions whose sigma-fields generate the Borel sigma-field.  Then $q^-_{\mu,j}\to q^-_\mu$ in every finite $L^r(\mu^{\otimes2})$ and $s_j\to\be\psi^-_\mu$ in every finite $L^r(\mu)$, so the full cutoff densities converge in measure.  Let $W_j$ be an unnormalized full $N$-particle density and let $G_j=W_j/\int W_j$.  Jensen gives $\int W_j\ge1$.  For every fixed $r>1$, the one-body source is bounded and \Cref{rep:thm:signed-bound}, on the amplified window $r\be_0$, gives
\begin{equation}
 \sup_j\int G_j^r\,\dd\mu^{\otimes N}\le C_r.
\end{equation}
Every fixed marginal is a conditional expectation of $G_j$, hence has the same uniform $L^r$ bound.  Taking $r>p$, Vitali's theorem yields $L^p$ convergence of the exact marginals.  The deterministic coefficient convergence follows from the finite Wick-contraction diagrams, as explained after \eqref{rep:eq:factorial-bound}.
\end{proof}

\begin{proof}[Second proof of \Cref{rep:thm:main}]
Apply \Cref{rep:prop:exact-remaining} to the scalar and marked identities.  The scalar leading Gaussian integral is
\begin{equation}\label{rep:eq:leading-determinant}
 \E\exp\left\{-\frac12\int:\!\phi^2\!:_c\,\dd\mu\right\}
 =\det{}_2(\Id+\be T_\mu)^{-1/2}.
\end{equation}
Finite logarithmic and Bell-polynomial recursion gives coefficients independent of the requested stopping order.  Division by the scalar series gives the normalized marginal coefficients.  The explicit first profiles follow from \Cref{rep:sec:first-coefficients}.  This route does not invoke \Cref{rep:prop:DGR-class-uniform}.
\end{proof}

\subsection{Arbitrary deterministic charge vectors}\label{rep:sec:arbitrary-charge}

We prove \Cref{rep:thm:arbitrary-charge}.  The common-background centering gives a one-body insertion independent of the charge vector; the pair interaction carries all the signs.

\subsubsection{Exact reduction}

Let $S_N=\sum_i\sigma_i=N\vartheta_N$.  Expanding \eqref{rep:eq:qmu} gives
\begin{align}
 \sum_{i<j}\sigma_i\sigma_j\g(x_i-x_j)
 ={}&\sum_{i<j}\sigma_i\sigma_jq^-_\mu(x_i,x_j)
 +S_N\sum_i\sigma_i h^-_\mu(x_i)-\sum_i h^-_\mu(x_i)\notag\\
 &\quad-I^-_\mu(S_N^2-N).
 \label{rep:eq:signed-pair-decomposition}
\end{align}
Substitution into \eqref{rep:eq:signed-modulated-energy-intro} proves \eqref{rep:eq:signed-exact-reduction-intro}.  In particular, the weak one-body insertion is independent of the charge vector.

\begin{remark}[Modulated versus uncentered nonneutral gases]\label{rep:rem:modulated-versus-physical}
The theorem centers the signed empirical measure around $\vartheta_N\mu$.  If one starts instead from the uncentered physical energy with one common external potential, a term of order $N$ proportional to $\vartheta_N\sum_i\sigma_i h^-_\mu(x_i)$ remains when $\vartheta_N\ne0$.  The two species then generally have different Euler--Lagrange equations.  That problem requires a coupled equilibrium problem for the two species or species-dependent backgrounds and is not covered here.
\end{remark}

\subsubsection{Scalar and marked field identities}

At a positive finite cutoff, let $c=\be q^-_{\mu,\vep}$, $s=\be\psi^-_{\mu,\vep}$, and let $\phi$ have covariance $c$.  For $\sigma\in\{-1,1\}$ and centered $u$, set
\begin{equation}\label{rep:eq:A-sigma-arbitrary}
 A_{\sigma,u}(t)=\int_{\T^{\d}}
 e^{\mathrm{i}\sigma t\phi(y)+t^2D(y)/2+t^2u(y)}\,\dd\mu(y).
\end{equation}
If $N_+$ and $N_-$ are the charge counts, then
\begin{equation}\label{rep:eq:arbitrary-scalar-field}
 P_{N,\be}^{\boldsymbol\sigma,\vep}
 =\E\left[A_{+,s}(N^{-1/2})^{N_+}
 A_{-,s}(N^{-1/2})^{N_-}\right].
\end{equation}
Since $s$ is real, $A_{-,s}=\overline{A_{+,s}}$; hence the modulus of the exact scalar integrand is $|A_{+,s}|^N$, independently of the imbalance.  Its square is controlled by \Cref{rep:thm:signed-bound}.

For marked signs $\tau_1,\ldots,\tau_k$, let $h$ and $\ell$ be the numbers of positive and negative marks and put
\begin{equation}\label{rep:eq:cx-arbitrary}
 c_x(y)=\sum_{a=1}^k\tau_a c(y,x_a).
\end{equation}
The complex Cameron--Martin identity yields
\begin{multline}\label{rep:eq:arbitrary-marked-field}
 P_{N,\be}^{\boldsymbol\sigma,\vep}
 F_{N,\be,\mu}^{\boldsymbol\sigma,\boldsymbol\tau,(k),\vep}(x)
 =\exp\left\{\frac1N\left[
 \sum_{a=1}^ks(x_a)-\sum_{a<b}\tau_a\tau_bc(x_a,x_b)
 \right]\right\}\\
 \times\E\left[
 A_{+,s-c_x}(N^{-1/2})^{N_+-h}
 A_{-,s+c_x}(N^{-1/2})^{N_--\ell}
 \right].
\end{multline}
The weighted integrated Wick estimates are unchanged by the signs.  We give the cutoff and removal details because the sources $s\mp c_x$ need not have a cutoff-uniform pointwise upper bound.

At the heat cutoff $\vep$, write $D_\vep(x)=c(x,x)$.  Positivity of the covariance and \eqref{rep:eq:cutoff-diagonal} imply
\begin{equation}
 |c(x,y)|\le D_\vep(x)^{1/2}D_\vep(y)^{1/2}
 \le C_{\be_0}(1+|\log\vep|).
\end{equation}
Thus, for $u_{+,x}=s-c_x$ and $u_{-,x}=s+c_x$,
\begin{equation}
 \sup_{x,y}\bigl(D_\vep(y)/2+u_{\pm,x}(y)\bigr)
 \le C_{k,\be_0}(1+|\log\vep|).
\end{equation}
If $n_+=N_+-h$, $n_-=N_--\ell$, and $t=N^{-1/2}$, concavity of $r\mapsto r^{t^2}$ gives, uniformly in the charge counts and marked configuration,
\begin{equation}\label{rep:eq:arbitrary-deterministic-cutoff-cost}
 |A_{+,u_{+,x}}(t)|^{n_+}|A_{-,u_{-,x}}(t)|^{n_-}
 \le C_{k,\be_0}\vep^{-C_{k,\be_0}}.
\end{equation}
Write
\begin{equation}
 J_k^{\boldsymbol\tau}(x)
 =\sum_{a=1}^k s(x_a)-\sum_{a<b}\tau_a\tau_b c(x_a,x_b).
\end{equation}
The exterior marked factor requires an $L^p$ argument rather than the
one-component pointwise upper bound.  Since $s$ is uniformly bounded and
$c=\be q^-_{\mu,\vep}$ is uniformly bounded below, there is a constant
$C=C(k,\be_0)$ such that
\begin{equation}
 (J_k^{\boldsymbol\tau})_+
 \le C+\sum_{\substack{a<b\\ \tau_a\tau_b=-1}}|c(x_a,x_b)|.
\end{equation}
Let $E_k=\binom{k}{2}$.  The row-exponential estimate for $q$ and generalized
H\"older over the at most $E_k$ pair factors show that, for every fixed
$p<\infty$, there is $N_0=N_0(k,p,\be_0)$ such that
\begin{equation}\label{rep:eq:signed-exterior-exponential}
 \sup_{N\ge N_0}\sup_{\vep,\mu,\be,\boldsymbol\tau}
 \left\|\exp\{(J_k^{\boldsymbol\tau})_+/N\}\right\|_{L^{2p}(\mu^{\otimes k})}
 \le C_{k,p,\be_0}.
\end{equation}
The same row estimate gives
\begin{equation}\label{rep:eq:signed-exterior-moments}
 \sup_{\vep,\mu,\be,\boldsymbol\tau}
 \|J_k^{\boldsymbol\tau}\|_{L^r(\mu^{\otimes k})}<\infty
 \qquad (r<\infty).
\end{equation}
Taylor's integral remainder, in the form
\begin{equation}
 \left|e^x-\sum_{j=0}^{M}\frac{x^j}{j!}\right|
 \le C_M |x|^{M+1}e^{x_+},
\end{equation}
therefore yields
\begin{equation}\label{rep:eq:signed-exterior-expansion}
 \left\|e^{J_k^{\boldsymbol\tau}/N}
 -\sum_{j=0}^{M}\frac{(J_k^{\boldsymbol\tau})^j}{j!N^j}
 \right\|_{L^p(\mu^{\otimes k})}
 \le C_{M,k,p,\be_0}N^{-M-1}.
\end{equation}

Apply the finite logarithmic expansion simultaneously to
$n_+\log A_{+,u_{+,x}}+n_-\log A_{-,u_{-,x}}$.  Since
\begin{equation}
 n_+=\frac N2(1+\vartheta_N)-h,
 \qquad
 n_-=\frac N2(1-\vartheta_N)-\ell,
\end{equation}
every retained coefficient is a polynomial in the exact fraction
$\vartheta_N$.  On the good event, multiply the finite logarithmic expansion
by the finite polynomial in \eqref{rep:eq:signed-exterior-expansion}.  On the bad
event, combine \eqref{rep:eq:signed-exterior-exponential} with H\"older, the
polynomial bound in the inverse cutoff \eqref{rep:eq:arbitrary-deterministic-cutoff-cost}, and
arbitrarily high moment depth.  With a moving cutoff
$\vep_N=N^{-b}$, choose the auxiliary logarithmic depth and bad-event
moment exponent only after $b$ and the requested order.  This gives the
species-resolved expansion at the moving cutoff, uniformly over all charge
vectors, with an $O(N^{-M-1})$ remainder in every fixed finite $L^p$.

To remove the cutoff, we use the signed partition bound and independent heat increments.  Put $\tau=\vep^2/2$, let $\mu_\tau=(e^{\tau\Delta})^*\mu$, and obtain $Y_i$ from $X_i$ by independent periodic Brownian increments of time $\tau$.  Directly from the centered-kernel definitions,
\begin{equation}
 q^-_{\mu,\vep}(X_i,X_j)
 =\E\bigl[q^-_{\mu_\tau}(Y_i,Y_j)\mid X_i,X_j\bigr],
 \qquad
 \psi^-_{\mu,\vep}(X_i)
 =\E\bigl[\psi^-_{\mu_\tau}(Y_i)\mid X_i\bigr].
\end{equation}
Consequently, conditional Jensen and \Cref{rep:thm:signed-bound} give a signed heat-cutoff bound uniform in the charge vector; the density of $\mu_\tau$ is still bounded by $M_0$.  H\"older's inequality gives the same bound for every linear interpolation between the cutoff and uncut exponents.

Differentiate the normalized full density along this interpolation.  The derivative has the form \eqref{rep:eq:density-derivative}, with
$\sigma_i\sigma_j(q^-_\mu-q^-_{\mu,\vep})$ in the pair term.  The signs do not affect the canonical moment estimate \eqref{rep:eq:canonical-moment}; H\"older and the signed partition bound on a fixed amplified coupling window therefore imply, for sufficiently large finite $r$,
\begin{equation}\label{rep:eq:signed-marginal-cutoff-comparison}
 \|F_{N,\be,\mu}^{\boldsymbol\sigma,\boldsymbol\tau,(k)}
 -F_{N,\be,\mu}^{\boldsymbol\sigma,\boldsymbol\tau,(k),\vep}\|_{L^p}
 \le C\bigl(\|q^-_\mu-q^-_{\mu,\vep}\|_{L^r}
 +\|\psi^-_\mu-\psi^-_{\mu,\vep}\|_{L^r}\bigr).
\end{equation}
Here conditional expectation onto the marked coordinates is an $L^p$ contraction, and the normalizing partition function is at least one by Jensen.  Fix the spatial exponent $r$ used in
\eqref{rep:eq:signed-marginal-cutoff-comparison}.  Among the finitely many signed
scalar and marked diagrams retained through order $M$ (including the
polynomials generated by $J_k^{\boldsymbol\tau}$), and the determinant, let
$\gamma_{M,k,p}>0$ be the minimum of their algebraic cutoff exponents.  Choose
$b$ so large that
\begin{equation}
 b\min\{\alpha_r,\gamma_{M,k,p}\}>M+2.
\end{equation}
Only after this choice select the auxiliary logarithmic and bad-event moment
depths.  Then \eqref{rep:eq:cutoff-Lr} and
\eqref{rep:eq:signed-marginal-cutoff-comparison} make the difference of the exact
cutoff and singular marginals $O(N^{-M-1})$, while one-factor-at-a-time
replacement in the finite signed contraction diagrams and determinant
continuity replace every retained cutoff coefficient by its singular limit at
the same order.  The strict inequality absorbs all logarithmic losses.  This proves \eqref{rep:eq:signed-marginal-all-order-intro}.  Locally, an opposite-charge marked pair contributes $|x-y|^{-p\be/N}$, so the sharp two-particle integrability condition is $p\be/N<\d$.  More explicitly, a colliding subcluster with $P$ positive and $Q$ negative charges, $K=P+Q\ge2$, has net adverse logarithmic degree
\begin{equation}\label{rep:eq:collision-cluster-degree}
 D(P,Q)=PQ-\binom P2-\binom Q2
 =\frac{K-(P-Q)^2}{2}\le K-1.
\end{equation}
The relative collision dimension is $\d(K-1)$, so the subcluster condition
$(p\be/N)D(P,Q)<\d(K-1)$ follows from the opposite-pair condition $p\be/N<\d$.  Thus, $p\be/N<\d$ controls every nested collision stratum: by the same calculation, every colliding subcluster has adverse degree at most its relative-coordinate multiplicity $K-1$.  Hence no larger cluster is worse, and the amplified signed partition bound supplies the stated sufficiently-large-$N$ threshold.

For the scalar logarithm, the charge counts give
\begin{equation}\label{rep:eq:arbitrary-bulk-log}
 A_2+\vartheta_NN^{-1/2}A_3(s)
 +N^{-1}\left(A_4(s)-\frac12A_2^2\right)+\cdots.
\end{equation}
A coefficient of total degree $r$ in $N^{-1/2}$ has Gaussian parity $r$.  Thus, odd total degrees vanish after expectation even when $\vartheta_N\ne0$, and the scalar and fixed-marginal expansions contain integer powers of $N^{-1}$.  The dependence on the charge counts is polynomial, proving the all-order assertions of \Cref{rep:thm:arbitrary-charge}.

\subsubsection{First coefficients}

At finite cutoff, put
\begin{equation}\label{rep:eq:finite-arbitrary-operators}
 C=\be T_{\mu,\vep},\qquad
 R=(\Id+C)^{-1},\qquad
 \kappa=CR=\Id-R,
\end{equation}
and let $b^-$ be the kernel of $-\kappa$.  Define
\begin{equation}\label{rep:eq:finite-arbitrary-d}
 \omega(x)=\langle c(x,\cdot),Rc(x,\cdot)\rangle,
 \qquad
 d=s+\frac12\omega^\circ.
\end{equation}
The scalar exponent \eqref{rep:eq:arbitrary-bulk-log} shows that the relative order-$N^{-1}$ correction is
\begin{equation}\label{rep:eq:arbitrary-z1-start}
 \E_\kappa\left[A_4(s)-\frac12A_2^2
 +\frac{\vartheta^2}{2}A_3(s)^2\right].
\end{equation}
The contractions in \eqref{rep:eq:A4-expect-minus}--\eqref{rep:eq:A3-expect-minus} give
\begin{equation}\label{rep:eq:arbitrary-z1-finite}
 z_1(\vartheta)
 =\frac12\langle d,(\Id-\vartheta^2\kappa)d\rangle
 -\frac14\Tr(\kappa^2)
 +\frac{\vartheta^2}{12}\iint b^-(x,y)^3\,\dd\mu(x)\dd\mu(y).
\end{equation}
Removing the cutoff proves \eqref{rep:eq:signed-z1-intro}.

For the marginal, expand \eqref{rep:eq:arbitrary-marked-field} through order $N^{-1}$.  With $L_x=\int c_x\phi\,\dd\mu$, the normalized coefficient is
\begin{multline}\label{rep:eq:arbitrary-f1-pre}
 J_k-k\E_\kappa[A_2]-\vartheta\E_\kappa[\mathcal L(c_x)]
 +\frac12\int c_x^2\,\dd\mu
 -\mathrm{i}\vartheta\E_\kappa[A_3(s)L_x]
 -\frac12\E_\kappa[L_x^2],
\end{multline}
where
\begin{equation}\label{rep:eq:arbitrary-L-J}
 \mathcal L(c_x)=-\frac12\int c_x:\!\phi^2\!:_c\,\dd\mu
 +\int sc_x\,\dd\mu,
 \qquad
 J_k=\sum_as(x_a)-\sum_{a<b}\tau_a\tau_bc(x_a,x_b).
\end{equation}
Changing Wick order from $c$ to $\kappa$ gives
\begin{align}
 \E_\kappa[A_2]&=\frac12\int\omega\,\dd\mu,\label{rep:eq:arbitrary-EA2}\\
 \E_\kappa[\mathcal L(c_x)]&=\langle c_x,d\rangle,\label{rep:eq:arbitrary-EL}\\
 -\mathrm{i}\E_\kappa[A_3(s)L_x]&=\langle c_x,\kappa d\rangle,\label{rep:eq:arbitrary-cross}\\
 \frac12\left(\int c_x^2\,\dd\mu-\E_\kappa[L_x^2]\right)
 &=\frac12\langle c_x,Rc_x\rangle.
 \label{rep:eq:arbitrary-quadratic-mark}
\end{align}
Since $\Id-\kappa=R$ and $\langle c(x,\cdot),Rd\rangle=(\kappa d)(x)$, the singleton contribution is $(\Id-\tau\vartheta\kappa)d$.  The off-diagonal contribution is
\begin{equation}\label{rep:eq:arbitrary-pair-profile}
 -c(x,y)+\langle c(x,\cdot),Rc(y,\cdot)\rangle
 =-\kappa(x,y)=b^-(x,y).
\end{equation}
This proves \eqref{rep:eq:signed-first-profile-intro}.  At $\vartheta=1$, $\tau=1$, one has $(\Id-\kappa)d=Rd$, and at $\vartheta=0$ the singleton is $d$.


\subsubsection{Bounded linear-statistic sources}
\label{rep:subsubsec:signed-linear-sources}
\begin{lemma}[Uniform control with bounded linear-statistic sources]\label{rep:lem:linear-source-control}
Fix centered functions $\varphi_1,\ldots,\varphi_q\in L^\infty(\mu)$ and a compact set $K\subset\mathbb C^q$.  Put
\begin{equation}\label{rep:eq:Phi-z}
 \Phi_z=\sum_{a=1}^qz_a\varphi_a,
 \qquad
 A^-_{j,u;z}(t)=\int
 e^{\mathrm{i} t\phi_j+t\Phi_z+t^2(D_j/2+u)}\,\dd\mu.
\end{equation}
The bounded test source $\Phi_z$ generates the linear-statistic transform and enters with coefficient $t=N^{-1/2}$. It is distinct from the one-body and fixed-coordinate source $u$, whose coefficient is $t^2=N^{-1}$.
For every fixed number of marks and every $r<\infty$, there is $C<\infty$ such that, for $t=N^{-1/2}$ and all sufficiently large $N$,
\begin{equation}\label{rep:eq:linear-source-moment}
 \sup_{z\in K}\sup_j\operatorname*{ess\,sup}_x
 \E\left[|A^-_{j,u_x^-;z}(t)|^{2(N-k)}\right]\le C.
\end{equation}
Consequently, the scalar and marked Laplace transforms with the insertion
$\exp\{N^{-1/2}\sum_i\Phi_z(x_i)\}$ admit cutoff-uniform expansions to every prescribed order in $N^{-1/2}$, uniformly for $z\in K$.
\end{lemma}

\begin{proof}
The weighted Wick estimate is unchanged after adjoining the bounded factor $e^{t\Phi_z}$: every coefficient is a finite sum of the coefficients in \eqref{rep:eq:weighted-Wick} multiplied by bounded powers of the functions $\varphi_a$.  Centering of both $\phi$ and $\Phi_z$ removes the nominal order-$t$ term, but
the leading order-$t^2$ exponent is changed.  Uniformly for $z$ in the fixed
compact set,
\begin{equation}\label{rep:eq:linear-source-leading-tilt}
 A^-_{j,u;z}(t)
 =1+t^2 A_{2,z}+O_{L^v}(t^3),
 \qquad
 A_{2,z}=A_2+\mathrm{i}\langle\Phi_z,\phi\rangle
 +\frac12\int\Phi_z^2\,\dd\mu.
\end{equation}
For every fixed polynomial $P$ in the weighted Wick coefficients and every
fixed $q<\infty$,
\begin{equation}\label{rep:eq:linear-source-polynomial-tilt-bound}
 \sup_{z\in K}\sup_j
 \E\left[|P|^q e^{q\operatorname{Re}A_{2,z}}\right]<\infty.
\end{equation}
Indeed, the deterministic term in \eqref{rep:eq:linear-source-leading-tilt} is
uniformly bounded on $K$; Cauchy--Schwarz separates the quadratic weight from
the real Gaussian linear functional; \eqref{rep:eq:quadratic-exp-bound}, the Wick
polynomial moment bounds, and the Gaussian linear moment-generating function
then give \eqref{rep:eq:linear-source-polynomial-tilt-bound}.  Completing the
square under the $A_2$-tilted Gaussian law gives the exact leading transform
\begin{equation}\label{rep:eq:linear-source-leading-transform}
 \E e^{A_{2,z}}
 =\det{}_2(\Id+\be T_{\mu,j})^{-1/2}
 \exp\left\{\frac12\langle\Phi_z,
 (\Id+\be T_{\mu,j})^{-1}\Phi_z\rangle_{L^2(\mu)}\right\}.
\end{equation}
The pairings here are complex bilinear, as appropriate for the holomorphic
Laplace transform.  It remains to control the exact exponential on the bad event.

Expand the even absolute moment in \eqref{rep:eq:linear-source-moment}.  Apart from the order-$t^2$ marked source, which is bounded above as in \eqref{rep:eq:u-upper}, its absolute value is bounded by
\begin{equation}\label{rep:eq:linear-source-even-moment}
 \int e^{-\frac\be N\sum_{a<b}\sigma_a\sigma_bq^-_{\mu,j}(z_a,z_b)}
 \exp\left\{t\sum_{a=1}^{2(N-k)}\Re\Phi_z(z_a)\right\}
 \,\dd\mu^{\otimes2(N-k)}(z).
\end{equation}
Choose conjugate H\"older exponents $r_0>1$ and $r_0'$.  The first factor is bounded by \Cref{rep:thm:signed-bound} at the fixed amplified window $2r_0\be_0$: with $L=2(N-k)$, its signed-gas coupling is exactly
\begin{equation}\label{rep:eq:linear-source-signed-window}
 \lambda=\frac{r_0\be L}{N}
 =2r_0\be\frac{N-k}{N}
 \le 2r_0\be_0.
\end{equation}
Thus the required signed-gas window is exactly $[0,2r_0\be_0]$.  Since $\int\Re\Phi_z\,\dd\mu=0$ and $\Re\Phi_z$ is uniformly bounded for $z\in K$,
\begin{equation}\label{rep:eq:linear-source-product-bound}
 \left(\int e^{r_0't\Re\Phi_z}\,\dd\mu\right)^{2(N-k)/r_0'}
 \le \exp\{C_K(N-k)t^2\}\le C_K.
\end{equation}
H\"older therefore proves \eqref{rep:eq:linear-source-moment}.  The good-event logarithmic expansion now uses
\eqref{rep:eq:linear-source-leading-tilt} and
\eqref{rep:eq:linear-source-polynomial-tilt-bound}; the complementary-event
argument uses \eqref{rep:eq:linear-source-moment}.  Expanding through the requested
finite total degree therefore gives the cutoff-uniform half-integer series.
The deterministic order-$t$ insertion breaks field parity, so both odd and
even powers of $N^{-1/2}$ may occur.  Uniformity on $K$ follows from the same
bounds, and coefficientwise holomorphy permits Cauchy differentiation on
smaller polydiscs.
\end{proof}


\subsection{Comparison of the two proofs and scope}
\label{rep:subsec:signed-scope}

For the one-component partition and marginal expansions, the two proofs share the finite-cutoff oscillatory identity and Wick-contraction algebra, but use different exponential estimates. A direct modulus bound costs a power of the cutoff because the diagonal covariance grows logarithmically. The moving-cutoff proof absorbs that cost by expanding to a higher order and uses the class-uniform repulsive partition bound. The heat-kernel proof instead converts an even absolute moment into a signed-gas integral and bounds it by the long-range Gaussian and short-range matching estimates. The bounded-linear-statistic consequence uses this signed-gas estimate through \Cref{rep:lem:linear-source-control}.

The signed theorem applies to a common prescribed background, arbitrary deterministic charge vectors, and their exact finite-particle charge imbalances. The one-component cases and exact neutrality are included. An uncentered nonneutral physical gas or distinct species backgrounds require a different macroscopic equilibrium, as explained in \Cref{rep:rem:modulated-versus-physical}. The fixed-order, fixed-marginal expansions do not assert convergence of the full series or uniformity for growing expansion order.

For repulsive Riesz kernels, the Hilbert--Schmidt partition estimate of \cite{DelgadinoGvalaniRosenzweig2026} remains available. An all-order marginal expansion would additionally require order- and $L^p$-dependent integrability estimates or a separate analysis near collisions. The signed argument here is specifically logarithmic: for $s>0$ and $c>0$, an opposite-charge factor $e^{c|x-y|^{-s}/N}$ is not locally integrable.

\section{Proof of the planar confined estimates}
\label{app:planar-confined-static-package}

We prove \Cref{thm:planar-confined-static-package}. The equilibrium estimates first give a strict global entropy--energy margin. A finite spatial partition then gives a larger-coupling partition bound at that same equilibrium. Finally, a Bernoulli replacement identity transfers the centered bound to the dominated reference measures in \eqref{pcsp:eq:tube}.

\subsection{The equilibrium family and uniform spatial bounds}

We use the subcritical stationary Keller--Segel theory in the normalization specified below; see \cite[Section~1, equation~(1.10)]{CamposDolbeault2014}. The subsequent global-margin argument uses the sharp planar logarithmic HLS inequality of \cite{CarlenLoss1992}. We treat the zero-mean constraint imposed by conservation of total mass and the passage from small relative entropy to the corresponding spectral estimate explicitly.

\begin{prop}[Equilibrium family]\label{pcsp:prop:equilibrium-family}
For each $0\leq\be<4$, the functional \eqref{pcsp:eq:mean-field-functional} has a unique minimizer $m_\be$.  It is smooth, strictly positive, and radial, and satisfies \eqref{pcsp:eq:euler-lagrange}.  Uniformly for $0\leq\be\leq\be_0$,
\begin{equation}\label{pcsp:eq:potential-asymptotic}
 \abs{U_{m_\be}(x)+\log(1+|x|)}\leq C.
\end{equation}
Moreover, if $\be_n\to\be$, then
\begin{align}
 m_{\be_n}&\longrightarrow m_\be
 &&\text{in }L^1(\R^2),\label{pcsp:eq:m-continuity}\\
 U_{m_{\be_n}}&\longrightarrow U_{m_\be}
 &&\text{locally uniformly on }\R^2,\label{pcsp:eq:U-continuity}\\
 \Gcal(m_{\be_n})&\longrightarrow\Gcal(m_\be).\label{pcsp:eq:G-continuity}
\end{align}
\end{prop}

\begin{proof}
For $\be>0$, existence, smoothness, radiality, and uniqueness follow from the stationary Keller--Segel result recalled in \cite[Section~1, equation~(1.10)]{CamposDolbeault2014}, after the change of variables
\begin{equation}\label{pcsp:eq:KS-scaling}
 z=\sqrt2x,
 \qquad
 M=2\pi\be,
 \qquad
 n_M(z)=\frac M2m_\be(z/\sqrt2).
\end{equation}
For $\be>0$, the condition $\be<4$ is exactly $M<8\pi$. At $\be=0$, the unique minimizer is the Gaussian measure \eqref{pcsp:eq:gaussian-reference}. The Euler--Lagrange equation for \eqref{pcsp:eq:mean-field-functional} is \eqref{pcsp:eq:euler-lagrange}.

Let
\begin{equation}\label{pcsp:eq:radial-distribution}
 M_\be(r)=\int_{|x|\leq r}m_\be(x)\,\dd x.
\end{equation}
Radial Newton's theorem gives
\begin{align}
 U_{m_\be}(r)
 &=-\int_0^\infty\log(r\vee s)\,\dd M_\be(s),
 \label{pcsp:eq:newton-potential}\\
 U_{m_\be}'(r)&=-\frac{M_\be(r)}r.
 \label{pcsp:eq:newton-derivative}
\end{align}
The uniform subcritical equilibrium estimates, together with \eqref{pcsp:eq:newton-potential}, give \eqref{pcsp:eq:potential-asymptotic}.  Substitution in \eqref{pcsp:eq:euler-lagrange} yields
\begin{equation}\label{pcsp:eq:raw-envelope}
 C^{-1}e^{-|x|^2}(1+|x|)^{-\be}
 \leq m_\be(x)
 \leq Ce^{-|x|^2}(1+|x|)^{-\be}.
\end{equation}
The elementary comparison
\begin{equation}\label{pcsp:eq:weight-comparison}
 (1+|x|^2)^{1/2}\leq1+|x|
 \leq\sqrt2(1+|x|^2)^{1/2}
\end{equation}
shows that \eqref{pcsp:eq:raw-envelope} is equivalent, with constants uniform for $\be\leq\be_0$, to \eqref{pcsp:eq:equilibrium-envelope} with $R=1$ after normalization.

For completeness, the derivative bounds follow directly from the radial formulas.  Writing $\phi_\be=\log m_\be$, one has
\begin{equation}\label{pcsp:eq:score}
 \nabla\phi_\be(x)=-2x+\be\nabla U_{m_\be}(x).
\end{equation}
The uniform $L^\infty$ bound in \eqref{pcsp:eq:raw-envelope} implies $M_\be(r)\leq Cr^2$ for $r\leq1$.  Equations \eqref{pcsp:eq:newton-derivative} and \eqref{pcsp:eq:radial-distribution} therefore give
\begin{equation}\label{pcsp:eq:potential-first-derivative}
 \abs{\nabla U_{m_\be}(x)}\leq C\min\{|x|,|x|^{-1}\}
 \quad\text{away from }x=0,
\end{equation}
with the continuous value $0$ at the origin.  Differentiating \eqref{pcsp:eq:newton-derivative} and using
\begin{equation}\label{pcsp:eq:M-prime}
 M_\be'(r)=2\pi r m_\be(r)
\end{equation}
gives a uniform bound on $D^2U_{m_\be}$.  Differentiating once more, and using
\begin{equation}\label{pcsp:eq:m-radial-derivative}
 m_\be'(r)=m_\be(r)\left(-2r-\be\frac{M_\be(r)}r\right),
\end{equation}
gives a uniform bound on $D^3U_{m_\be}$ on bounded sets and the bound $O(|x|^{-3})$ outside the unit ball.  The identities
\begin{align}
 Dm_\be&=m_\be D\phi_\be,\label{pcsp:eq:derivative-one}\\
 D^2m_\be&=m_\be\left(D^2\phi_\be+D\phi_\be^{\otimes2}\right),\label{pcsp:eq:derivative-two}\\
 D^3m_\be&=m_\be\left(D^3\phi_\be
 +3\operatorname{Sym}(D^2\phi_\be\otimes D\phi_\be)
 +D\phi_\be^{\otimes3}\right)
 \label{pcsp:eq:derivative-three}
\end{align}
then prove \eqref{pcsp:eq:equilibrium-derivatives}.

Finally, the uniform two-sided spatial bounds give tightness and uniform local integrability of the logarithmic singularity.  Along any sequence $\be_n\to\be$, a weakly convergent subsequence has potentials converging locally uniformly after splitting the convolution into a small diagonal ball, a compact annulus, and a Gaussian tail.  Passing to the normalized Euler--Lagrange equation identifies the limit with the unique $m_\be$, which proves \eqref{pcsp:eq:m-continuity}--\eqref{pcsp:eq:G-continuity}.
\end{proof}

\subsection{The strict global entropy--energy margin}

For a probability density $\rho$, set
\begin{equation}\label{pcsp:eq:centered-energy}
 \Ecal_\be(\rho)
 =\iint\g(x-y)(\rho-m_\be)(x)(\rho-m_\be)(y)\,\dd x\dd y.
\end{equation}
For zero-mass signed measures in its natural finite-energy class, logarithmic energy is a nonnegative quadratic form.  In particular, it has the Cauchy--Schwarz property
\begin{equation}\label{pcsp:eq:energy-CS}
 \abs{\Ecal_\g(\sigma,\tau)}^2
 \leq\Ecal_\g(\sigma)\Ecal_\g(\tau).
\end{equation}
The equilibrium identity gives
\begin{equation}\label{pcsp:eq:free-energy-defect}
 \Ent(\rho\mid m_\be)-\frac\be2\Ecal_\be(\rho)
 =\mathcal I_\be(\rho)-\mathcal I_\be(m_\be)\geq0,
\end{equation}
with equality only for $\rho=m_\be$. The global entropy--energy margin means that $\be$ in the coefficient of \eqref{pcsp:eq:free-energy-defect} can be increased by a fixed positive amount uniformly over the stated family; see \eqref{pcsp:eq:global-gap}. This concerns all admissible probability densities. The spectral estimate below concerns only zero-mean perturbations.

\subsubsection{Spectral stability on the zero-mean space}

Let $T_\be$ be the centered logarithmic operator on
\begin{equation}\label{pcsp:eq:centered-space}
 L_0^2(m_\be)=\left\{f\in L^2(m_\be):\int f\,\dd m_\be=0\right\},
\end{equation}
which represents signed perturbations $fm_\be$ of zero total mass. Its quadratic form is
\begin{equation}\label{pcsp:eq:T-quadratic-form}
 \pcspip{f}{T_\be f}_{L^2(m_\be)}
 =\Ecal_\g(fm_\be).
\end{equation}
The centered kernel is Hilbert--Schmidt under \eqref{pcsp:eq:raw-envelope}; hence $T_\be$ is compact and self-adjoint.  The linearized equilibrium inequality gives
\begin{equation}\label{pcsp:eq:linearized-nonstrict}
 0\leq\be T_\be\leq\Id.
\end{equation}

\begin{lemma}[Strict stability on the zero-mean space]\label{pcsp:lem:tangent-reduction}
For every $0<\be<4$, the endpoint $1$ is not an eigenvalue of $\be T_\be$ on $L_0^2(m_\be)$.
\end{lemma}

\begin{proof}
Suppose $f\neq0$ and $\be T_\be f=f$.  Define the uncentered potential operator
\begin{equation}\label{pcsp:eq:uncentered-operator}
 K_\be f(x)=\int\g(x-y)f(y)\,m_\be(y)\,\dd y,
 \qquad
 C_f=\int K_\be f\,\dd m_\be.
\end{equation}
Then
\begin{equation}\label{pcsp:eq:full-potential-eigenrelation}
 \be K_\be f=f+\be C_f.
\end{equation}
This full potential relation, rather than only the equation obtained after applying $-\Delta$, is essential.

Use the scaling \eqref{pcsp:eq:KS-scaling} and set
\begin{equation}\label{pcsp:eq:F-scaled}
 F(z)=f(z/\sqrt2),
 \qquad
 G_2(z)=-\frac1{2\pi}\log|z|.
\end{equation}
Because $f$ has zero $m_\be$-mean, the additive $\log\sqrt2$ term cancels, and \eqref{pcsp:eq:full-potential-eigenrelation} becomes
\begin{equation}\label{pcsp:eq:KS-potential-relation}
 F=G_2*(Fn_M)-\be C_f.
\end{equation}
The linearized Keller--Segel operator has the form
\begin{equation}\label{pcsp:eq:KS-linearized-operator}
 L_MF=\frac1{n_M}\nabla\cdot\left[n_M\nabla
 \left(F-G_2*(Fn_M)\right)\right].
\end{equation}
Thus \eqref{pcsp:eq:KS-potential-relation} implies $L_MF=0$.  The mass mode, obtained by differentiating the equilibrium family with respect to total mass, is identified in \cite[Proposition~9]{CamposDolbeault2014} and has the form
\begin{equation}\label{pcsp:eq:KS-kernel}
 F=a\,\partial_M\log n_M
\end{equation}
for a scalar $a$.  This mode is excluded by the zero-mean constraint: on the one hand,
\begin{equation}\label{pcsp:eq:F-zero-mass}
 \int F n_M\,\dd z=M\int f\,\dd m_\be=0.
\end{equation}
On the other hand,
\begin{equation}\label{pcsp:eq:mass-mode-nonzero}
 \int(\partial_M\log n_M)n_M\,\dd z
 =\partial_M\int n_M\,\dd z=1.
\end{equation}
Hence a mass-mode vector satisfying the zero-mean constraint has $a=0$.
We exclude the endpoint eigenvector using the full potential relation: it supplies the boundary behavior needed for the following ground-state identities.

Write $m=m_\be$, $M(r)=M_\be(r)$, and $V(r)=2\pi\be m(r)$.
The Gaussian upper bound, Cauchy--Schwarz, and $\int fm=0$ imply that
$K_\be f$ is locally bounded and is $O(r^{-1})$ at infinity.
Indeed, subtract $-\log r$ in its convolution, split at $|y|=r/2$,
and use the first moment on the inner region and the Gaussian-weighted
$L^2$ bound on the outer region, including its integrable logarithmic
singularity.  The full potential relation therefore makes $f$ bounded,
with $f+\be C_f=O(r^{-1})$.  Elliptic regularity now makes $f$ smooth;
since $fm$ is bounded with Gaussian decay and has zero mass,
differentiating its potential gives $\nabla f=O(r^{-2})$ at infinity.

Set
\begin{equation}
 w(r)=-(\log m)'(r)=2r+\be M(r)/r>0.
\end{equation}
The radial equilibrium equations give
\begin{equation}
 -w''-r^{-1}w'+(r^{-2}-V)w=0,
 \qquad w(r)\sim(2+\pi\be m(0))r\quad(r\downarrow0),
 \qquad w(r)\sim2r\quad(r\to\infty).
\end{equation}
For each real sine or cosine Fourier coefficient $u_\ell$, $\ell\geq1$,
of $f$, the equation $-\Delta f=Vf$ and the ground-state identity give
\begin{equation}
 \begin{split}
 0&=\int_0^\infty\bigl[(u_\ell')^2+
              (\ell^2/r^2-V)u_\ell^2\bigr]r\,\dd r\\
  &=\int_0^\infty r w^2\bigl[(u_\ell/w)'\bigr]^2\,\dd r
       +(\ell^2-1)\int_0^\infty u_\ell^2/r\,\dd r.
 \end{split}
\end{equation}
There are no boundary terms: smoothness gives $u_\ell=O(r^\ell)$
at zero, while the full potential relation gives $u_\ell=O(r^{-1})$
and $u_\ell'=O(r^{-2})$ at infinity.  Thus $u_\ell=0$ for $\ell\geq2$.
For $\ell=1$ the identity gives $u_1=cw$, and boundedness forces $c=0$.
Consequently $f$ is radial.

For the radial coefficient, put
\begin{equation}
 Q(r)=2\pi\int_0^r t m(t)f(t)\,\dd t,
 \qquad p(r)=\frac1{r m(r)},\qquad v(r)=r^2m(r).
\end{equation}
Here $Q(0)=Q(\infty)=0$, and the full potential relation and Newton's
formula give $f'=-\be Q/r$.  Direct differentiation yields
\begin{equation}
 (pQ')'+\frac{2\pi\be}{r}Q=0,
 \qquad
 (pv')'+\frac{2\pi\be}{r}v=-4pv.
\end{equation}
The corresponding ground-state identity is
\begin{equation}
 0=\int_0^\infty
       \left[p(Q')^2-\frac{2\pi\be}{r}Q^2\right]\,\dd r
   =\int_0^\infty pv^2\bigl[(Q/v)'\bigr]^2\,\dd r
       +4\int_0^\infty pQ^2\,\dd r.
\end{equation}
The boundary terms again vanish: $Q=O(r^2)$ at zero and
$Q(r)=-2\pi\int_r^\infty t m(t)f(t)\,\dd t
 =O(e^{-r^2}r^{-\be})$ at infinity, by zero mass and boundedness of $f$.
Hence $Q=0$ and then $f=0$, the desired contradiction.
\end{proof}

\begin{cor}[Uniform spectral gap on the zero-mean space]\label{pcsp:cor:tangent-gap}
There is $a_0>0$ such that, for every $0\leq\be\leq\be_0$ and every $f\in L_0^2(m_\be)$,
\begin{equation}\label{pcsp:eq:uniform-tangent-gap}
 \int f^2\,\dd m_\be
 \geq(\be+a_0)\Ecal_\g(fm_\be).
\end{equation}
\end{cor}

\begin{proof}
Conjugate $T_\be$ to $L^2(\dd x)$ using multiplication by $\sqrt{m_\be}$.  The resulting centered Hilbert--Schmidt kernels depend continuously on $\be$ by \cref{pcsp:prop:equilibrium-family}.  Therefore $t_\be=\norm{T_\be}_{\op}$ is continuous.  It is positive for every $\be$, since logarithmic energy is strictly positive on a nonzero smooth compactly supported zero-mass density.  By \cref{pcsp:lem:tangent-reduction} and \eqref{pcsp:eq:linearized-nonstrict},
\begin{equation}\label{pcsp:eq:strict-operator-norm}
 \be t_\be<1
 \quad\text{for }0<\be\leq\be_0.
\end{equation}
At $\be=0$, $t_0$ is finite and positive.  Consequently,
\begin{equation}\label{pcsp:eq:a0-definition}
 a_0=\inf_{0\leq\be\leq\be_0}\left(t_\be^{-1}-\be\right)>0,
\end{equation}
which is \eqref{pcsp:eq:uniform-tangent-gap}.
\end{proof}

\subsubsection{Entropy-small perturbations}

The spectral inequality on the zero-mean space must be extended to perturbations that are small in relative entropy but not necessarily in $L^2$.

\begin{lemma}[Uniform entropy-small passage]\label{pcsp:lem:entropy-small}
Suppose $\be_n\in[0,\be_0]$ and $\rho_n$ are probability densities such that
\begin{equation}\label{pcsp:eq:entropy-small-assumption}
 H_n:=\Ent(\rho_n\mid m_{\be_n})\longrightarrow0.
\end{equation}
Then
\begin{equation}\label{pcsp:eq:entropy-small-conclusion}
 \liminf_{n\to\infty}
 \frac{2H_n}{\Ecal_{\be_n}(\rho_n)}
 \geq
 \liminf_{n\to\infty}t_{\be_n}^{-1},
\end{equation}
with the convention that the quotient is $+\infty$ when the denominator vanishes.
\end{lemma}

\begin{proof}
Write
\begin{equation}\label{pcsp:eq:f-perturbation}
 \rho_n=(1+f_n)m_{\be_n},
 \qquad
 \int f_n\,\dd m_{\be_n}=0,
\end{equation}
and set
\begin{equation}\label{pcsp:eq:Psi}
 \Psi(1+s)=(1+s)\log(1+s)-s.
\end{equation}
Fix $0<\varepsilon<1/2$ and let $B_n=\{|f_n|>\varepsilon\}$.  Define
\begin{equation}\label{pcsp:eq:bad-parts}
 \sigma_{n,+}=(f_n)_+\one_{B_n}m_{\be_n},
 \quad
 \sigma_{n,-}=(f_n)_-\one_{B_n}m_{\be_n},
 \quad
 a_{n,\pm}=\sigma_{n,\pm}(\R^2).
\end{equation}
Since $\Psi(1+s)\geq c_\varepsilon|s|$ on $|s|>\varepsilon$,
\begin{equation}\label{pcsp:eq:bad-masses}
 a_{n,+}+a_{n,-}\leq C_\varepsilon H_n.
\end{equation}

For a positive subprobability $\sigma=a\bar\sigma$, the sharp logarithmic HLS inequality \cite{CarlenLoss1992}, applied to $\bar\sigma$, gives
\begin{equation}\label{pcsp:eq:scaled-logHLS}
 \Ecal_\g(\sigma)
 \leq\frac a2\Ent(\sigma)-\frac{a^2}{2}\log a+Ca^2.
\end{equation}
On $\{f_n>\varepsilon\}$,
\begin{equation}\label{pcsp:eq:f-log-f}
 f_n\log f_n\leq C_\varepsilon\Psi(1+f_n)+C_\varepsilon f_n.
\end{equation}
The upper bound on $\log m_{\be_n}$ therefore gives
\begin{equation}\label{pcsp:eq:positive-subentropy}
 \Ent(\sigma_{n,+})\leq C_\varepsilon(H_n+a_{n,+}),
\end{equation}
and
\begin{equation}\label{pcsp:eq:positive-relative-entropy}
 a_{n,+}\Ent(a_{n,+}^{-1}\sigma_{n,+}\mid m_{\be_n})
 \leq C_\varepsilon H_n+C_\varepsilon a_{n,+}
 +a_{n,+}|\log a_{n,+}|.
\end{equation}
By \eqref{pcsp:eq:potential-asymptotic}, for every fixed $\lambda>0$,
\begin{equation}\label{pcsp:eq:potential-exponential}
 \sup_{0\leq\be\leq\be_0}
 \int e^{\lambda(-U_{m_\be})_+}\,\dd m_\be<\infty.
\end{equation}
The Donsker--Varadhan lemma applied to \eqref{pcsp:eq:positive-relative-entropy} yields
\begin{equation}\label{pcsp:eq:positive-potential-bound}
 \int(-U_{m_{\be_n}})_+\,\dd\sigma_{n,+}
 \leq C_\varepsilon H_n
 +C_\varepsilon a_{n,+}(1+|\log a_{n,+}|).
\end{equation}
Expanding the centered energy and using \eqref{pcsp:eq:scaled-logHLS} gives
\begin{equation}\label{pcsp:eq:positive-bad-energy}
 \Ecal_\g(\sigma_{n,+}-a_{n,+}m_{\be_n})
 \leq C_\varepsilon\left[
 a_{n,+}H_n+a_{n,+}^2(1+|\log a_{n,+}|)
 \right].
\end{equation}

For the negative part, write $\sigma_{n,-}=u_nm_{\be_n}$ with $0\leq u_n\leq1$.  Then
\begin{equation}\label{pcsp:eq:negative-entropy-bounds}
 \Ent(\sigma_{n,-})\leq Ca_{n,-},
 \qquad
 \Ent(a_{n,-}^{-1}\sigma_{n,-}\mid m_{\be_n})
 \leq|\log a_{n,-}|.
\end{equation}
The same argument gives
\begin{equation}\label{pcsp:eq:negative-bad-energy}
 \Ecal_\g(\sigma_{n,-}-a_{n,-}m_{\be_n})
 \leq Ca_{n,-}^2(1+|\log a_{n,-}|).
\end{equation}
Set
\begin{equation}\label{pcsp:eq:bad-centered}
 b_n=a_{n,+}-a_{n,-},
 \qquad
 \eta_n=f_n\one_{B_n}m_{\be_n}-b_nm_{\be_n}.
\end{equation}
By \eqref{pcsp:eq:energy-CS}, \eqref{pcsp:eq:bad-masses}, \eqref{pcsp:eq:positive-bad-energy}, and \eqref{pcsp:eq:negative-bad-energy},
\begin{equation}\label{pcsp:eq:bad-oH}
 \Ecal_\g(\eta_n)
 \leq C_\varepsilon H_n^2(1+|\log H_n|)=o_\varepsilon(H_n).
\end{equation}

Define the centered good perturbation
\begin{equation}\label{pcsp:eq:good-centered}
 \widetilde f_n=f_n\one_{B_n^c}+b_n.
\end{equation}
Then $\int\widetilde f_n\,\dd m_{\be_n}=0$ and
\begin{equation}\label{pcsp:eq:good-L2-identity}
 \int\widetilde f_n^2\,\dd m_{\be_n}
 =\int_{B_n^c}f_n^2\,\dd m_{\be_n}-b_n^2.
\end{equation}
Taylor's inequality gives
\begin{equation}\label{pcsp:eq:Psi-quadratic}
 \Psi(1+s)\geq\frac{1-C\varepsilon}{2}s^2
 \quad\text{for }|s|\leq\varepsilon.
\end{equation}
The spectral inequality on the zero-mean space therefore implies
\begin{equation}\label{pcsp:eq:good-energy-bound}
 \Ecal_\g(\widetilde f_nm_{\be_n})
 \leq\frac{2H_n}{(1-C\varepsilon)t_{\be_n}^{-1}}.
\end{equation}
Since
\begin{equation}\label{pcsp:eq:full-decomposition}
 (\rho_n-m_{\be_n})
 =\widetilde f_nm_{\be_n}+\eta_n,
\end{equation}
energy Cauchy--Schwarz and \eqref{pcsp:eq:bad-oH} show
\begin{equation}\label{pcsp:eq:full-energy-local}
 \Ecal_{\be_n}(\rho_n)
 \leq\frac{2H_n}{(1-C\varepsilon)t_{\be_n}^{-1}}+o_\varepsilon(H_n).
\end{equation}
First let $n\to\infty$, then $\varepsilon\downarrow0$, to obtain \eqref{pcsp:eq:entropy-small-conclusion}.
\end{proof}

\subsubsection{Energy compactness and the global entropy--energy margin}

\begin{lemma}[Uniform logarithmic-energy continuity]\label{pcsp:lem:energy-continuity}
Suppose $\be_n\to\be\in[0,\be_0]$, $\rho_n\rightharpoonup\rho$, and
\begin{equation}\label{pcsp:eq:entropy-bounded}
 \sup_n\Ent(\rho_n\mid m_{\be_n})<\infty.
\end{equation}
Then
\begin{equation}\label{pcsp:eq:energy-continuity}
 \Ecal_{\be_n}(\rho_n)\longrightarrow\Ecal_\be(\rho).
\end{equation}
\end{lemma}

\begin{proof}
The Donsker--Varadhan lemma against a sufficiently small positive multiple of $|x|^2$, together with the uniform Gaussian upper bound, gives a uniform second-moment bound.  Binary data processing on any square $Q$ of side length $r$ gives
\begin{equation}\label{pcsp:eq:square-mass}
 \sup_{n,Q}\rho_n(Q)\leq\frac C{\log(1/r)}=:b(r),
 \qquad b(r)\downarrow0.
\end{equation}
Cover the plane by squares of side $r$ and use fixed finite enlargements $Q^*$.  Apply the scaled inequality \eqref{pcsp:eq:scaled-logHLS} to the restrictions $\rho_n|_{Q^*}$.  Finite overlap and the uniform positive-entropy bound imply
\begin{equation}\label{pcsp:eq:near-diagonal-UI}
 \sup_n\iint_{|x-y|<r}
 \log\frac1{|x-y|}\,\rho_n(\dd x)\rho_n(\dd y)
 \leq C\left[b(r)+b(r)|\log b(r)|\right].
\end{equation}
The right side tends to zero.  The negative logarithmic tail is uniformly integrable because
\begin{equation}\label{pcsp:eq:negative-log-tail}
 \log_+|x-y|
 \leq\log(1+|x|)+\log(1+|y|),
\end{equation}
and the second moments are uniformly bounded.  Truncating near the diagonal and at spatial infinity reduces the interaction to a bounded continuous kernel, for which weak convergence applies.  Finally, \eqref{pcsp:eq:U-continuity}--\eqref{pcsp:eq:G-continuity} treat the centering terms in \eqref{pcsp:eq:centered-energy}.
\end{proof}

\begin{prop}[Uniform global entropy--energy margin]\label{pcsp:prop:global-gap}
There is $\eta_0>0$ such that, for every $0\leq\be\leq\be_0$ and every admissible probability density $\rho$,
\begin{equation}\label{pcsp:eq:global-gap}
 \Ent(\rho\mid m_\be)
 \geq\frac{\be+\eta_0}{2}\Ecal_\be(\rho).
\end{equation}
\end{prop}

\begin{proof}
The sharp planar logarithmic HLS inequality and \eqref{pcsp:eq:potential-asymptotic} give the following bound with the sharp planar log-HLS coefficient
\begin{equation}\label{pcsp:eq:critical-defect}
 \Ent(\rho\mid m_\be)\geq2\Ecal_\be(\rho)-C_{\be_0}.
\end{equation}
Indeed, after using \eqref{pcsp:eq:euler-lagrange}, the difference between the two sides is the sharp log-HLS functional plus
\begin{equation}\label{pcsp:eq:effective-potential}
 \int\left[|x|^2+(4-\be)U_{m_\be}(x)\right]\rho(x)\,\dd x
\end{equation}
up to a bounded constant, and the bracket is bounded below uniformly on $\be\leq\be_0$.

Assume that no uniform $\eta_0$ exists.  Then there are $\be_n\to\be$ and probability densities $\rho_n$ with $\Ecal_{\be_n}(\rho_n)>0$ such that
\begin{equation}\label{pcsp:eq:quotient-losing}
 \frac{2\Ent(\rho_n\mid m_{\be_n})}
 {\Ecal_{\be_n}(\rho_n)}-\be_n\longrightarrow0.
\end{equation}
Since $\be_0<4$, \eqref{pcsp:eq:critical-defect} and \eqref{pcsp:eq:quotient-losing} bound the energies and entropies uniformly.  Extract a weakly convergent subsequence.  By \cref{pcsp:lem:energy-continuity}, the centered energies converge.

If their limit is positive, the defect in \eqref{pcsp:eq:free-energy-defect} tends to zero.  Varying-reference lower semicontinuity of entropy, energy continuity, and uniqueness imply that the limit is $m_\be$, whose centered energy is zero, a contradiction.

If the energies tend to zero, \eqref{pcsp:eq:quotient-losing} forces the relative entropies to zero.  Lemma~\ref{pcsp:lem:entropy-small} and \eqref{pcsp:eq:a0-definition} then give
\begin{equation}\label{pcsp:eq:local-contradiction}
 \liminf_{n\to\infty}
 \left[
 \frac{2\Ent(\rho_n\mid m_{\be_n})}
 {\Ecal_{\be_n}(\rho_n)}-\be_n
 \right]\geq a_0>0,
\end{equation}
again a contradiction.
\end{proof}

\subsection{Partition bound at a larger coupling with the equilibrium fixed}

We next pass from the global margin to a particle-uniform partition bound. The noncompact step is the treatment of the tail and its adjacent cells in a single canonical array, controlled by their total occupancy.

\begin{prop}[Partition bound at the fixed equilibrium]\label{pcsp:prop:equilibrium-buffer}
There are $\delta_0>0$ and $C_0<\infty$, depending only on $\be_0$, such that
\begin{equation}\label{pcsp:eq:equilibrium-buffer}
 \sup_{\substack{N\geq2,\ 0\leq\be\leq\be_0\\
                  0\leq s\leq\be+\delta_0}}
 \K_N(s;m_\be)\leq C_0.
\end{equation}
\end{prop}

The reference in \eqref{pcsp:eq:equilibrium-buffer} is fixed at $m_\be$ while $s$ is increased above $\be$.

\begin{proof}
We prove the finite-partition implication needed here. Consider a family of nonnegative couplings $\alpha_\be$ satisfying
\begin{equation}\label{pcsp:eq:source-buffer-hypotheses}
 \sup_{0\leq\be\leq\be_0}\alpha_\be<4,
 \qquad
 \inf_{0\leq\be\leq\be_0}
 \bigl(\Lambda_\be-\alpha_\be\bigr)>0,
\end{equation}
where
\begin{equation}\label{pcsp:eq:Lambda-definition}
 \Lambda_\be
 =\inf_{0<\Ecal_\be(\rho)<\infty}
 \frac{2\Ent(\rho\mid m_\be)}{\Ecal_\be(\rho)}.
\end{equation}
By \cref{pcsp:prop:global-gap}, $\Lambda_\be\geq\be+\eta_0$ uniformly. Choose
\begin{equation}\label{pcsp:eq:delta0-choice}
 0<\delta_0<\frac14\min\{\eta_0,4-\be_0\}.
\end{equation}
For any $s\in[0,\be+\delta_0]$, take $\alpha_\be=s$ at the chosen value of $\be$. Then
\begin{equation}\label{pcsp:eq:source-buffer-margins}
 s\leq\be_0+\delta_0<4,
 \qquad
 \Lambda_\be-s\geq\eta_0-\delta_0>0.
\end{equation}
The proof below gives a constant depending only on these strict margins and the common equilibrium bounds. It therefore applies uniformly to all pairs $(\be,s)$ in \eqref{pcsp:eq:equilibrium-buffer}.

\step{1. A preliminary linear bound on the log partition function.}
Write $h_\be=\widetilde\g_{m_\be}$ and
\begin{equation}\label{pcsp:eq:partition-linear-notation}
 Z_{N,b,\be}
 =\E_{m_\be^{\otimes N}}\left[\exp\left\{\frac bN\sum_{i<j}h_\be(X_i,X_j)\right\}\right].
\end{equation}
The inequality $|x-y|\leq(1+|x|)(1+|y|)$ and \eqref{pcsp:eq:potential-asymptotic} give $h_\be\geq-C$ uniformly. For each fixed integer $l\geq2$ and $b_*<4$, the integrals $Z_{l,b,\be}$ are bounded uniformly for $0\leq b\leq b_*$ and $0\leq\be\leq\be_0$. Indeed, the full-cluster condition is $b<4$, and the induction of \Cref{lem:nested-cluster-power-count} treats all partial diagonals. The centering terms grow at most logarithmically in each coordinate, so the common Gaussian bound \eqref{pcsp:eq:raw-envelope} controls the spatial tails.

The subset-marginal variational argument of \Cref{lem:subset-marginal-partition}, now under the product reference $m_\be^{\otimes N}$, gives
\begin{equation}\label{pcsp:eq:partition-linear-subset}
 \log Z_{N,b,\be}
 \leq \frac Nl\log Z_{l,\,bl(N-1)/(N(l-1)),\,\be}
 \qquad(N\geq l).
\end{equation}
Choose $l$ with $b_*l/(l-1)<4$, and include the finitely many $N<l$ by the preceding fixed-particle bound. Then
\begin{equation}\label{pcsp:eq:partition-linear-bound}
 \sup_{\substack{0\leq\be\leq\be_0\\0\leq b\leq b_*}}
 \log Z_{N,b,\be}\leq C_{b_*,\be_0}N.
\end{equation}
This bound will control the exceptional occupancy event before the covariance decomposition is used.

\step{2. A finite partition with uniform conditional moments.}
Fix a near range $L\geq4$. Choose a radius $R_0$ large compared with $L$, and then a mesh size $h\leq c/R_0$. Partition the disk $B_{R_0}$ into a polar mesh with radial thickness comparable to $h$ and angular arc length comparable to $h$; use one central disk and choose $R_0/h$ integral to avoid arbitrarily thin outer cells. Every core cell has diameter at most $Ch$, area at least $ch^2$, and uniformly bounded aspect ratio. Add the tail cell $A_\infty=\mathbb R^2\setminus B_{R_0}$. Core cells at distance at most $Lh$ are called near. Their adjacency graph has degree at most $C L^2$ by packing.

Let $\nu_a$ be the normalized restriction of $m_\be$ to a cell $A_a$. For $x\in A_a$, $y\in A_b$, define
\begin{align}\label{pcsp:eq:partition-double-centering}
 r_\be(x,y)={}&\g(x-y)-\int\g(u-y)\,\dd\nu_a(u)
                 -\int\g(x-v)\,\dd\nu_b(v)\notag\\
 &+\iint\g(u-v)\,\dd\nu_a(u)\dd\nu_b(v).
\end{align}
This kernel is canonical on every cell pair. On the core, \eqref{pcsp:eq:score} gives $|\nabla\log m_\be|\leq C(1+R_0)$. Since $hR_0$ is bounded, the conditional densities on the rescaled cells are uniformly comparable to Lebesgue measure. For near core cells, the constant $-\log h$ cancels from \eqref{pcsp:eq:partition-double-centering}. The logarithmic integral on the rescaled cells therefore gives, for every integer $p\geq1$,
\begin{equation}\label{pcsp:eq:partition-near-moments}
 \sup_y\int|r_\be(x,y)|^p\,\dd\nu_a(x)
 +\sup_x\int|r_\be(x,y)|^p\,\dd\nu_b(y)
 \leq 2 C_L^p p!.
\end{equation}
Here $C_L$ is independent of $R_0,h,\be$.

Let the collar consist of core cells at distance at most $(L+2)h$ from $\partial B_{R_0}$, and let $E$ be its union with the tail. All tail--tail and tail--collar pairs will be estimated together. Their conditional moment constant is independent of the number of collar cells, as we now show.

The radial density $w_\be(r)=2\pi r m_\be(r)$ satisfies
\begin{equation}\label{pcsp:eq:partition-tail-radial}
 \frac{\dd}{\dd r}\log w_\be(r)
 =r^{-1}-2r-\be M_\be(r)/r.
\end{equation}
For sufficiently large $R_0$, this derivative lies between $-3r$ and $-r$ on $[R_0,\infty)$, uniformly in $\be$. Integration on the full ray for the upper bound and on $[R_0,R_0+R_0^{-1}]$ for the lower bound gives
\begin{equation}\label{pcsp:eq:partition-tail-normalization}
 c\frac{w_\be(R_0)}{R_0}
 \leq\int_{R_0}^\infty w_\be(r)\,\dd r
 \leq C\frac{w_\be(R_0)}{R_0}.
\end{equation}
Consequently, under the conditional tail law,
\begin{equation}\label{pcsp:eq:partition-tail-exponential}
 \Pb\{R_0(|X|-R_0)>u\mid |X|>R_0\}\leq C e^{-cu}
 \qquad(u\geq0).
\end{equation}
The angle is independent and uniform. For $y=qe^{i\varphi}$ in the tail or collar, the exact identity
\begin{equation}\label{pcsp:eq:partition-angular-log}
 -\log|re^{i\vartheta}-qe^{i\varphi}|
 =-\log(r\vee q)
  -\log\left|e^{i(\vartheta-\varphi)}-\frac{r\wedge q}{r\vee q}\right|
\end{equation}
and the bound
\begin{equation}\label{pcsp:eq:partition-angular-moment}
 \sup_{0\leq a\leq1}\frac1{2\pi}\int_0^{2\pi}
       |\log|e^{i\theta}-a||^p\,\dd\theta\leq C^p p!
\end{equation}
control the angular part. To verify \eqref{pcsp:eq:partition-angular-moment}, the positive logarithm is bounded by $\log2$; for the negative logarithm, $|e^{i\theta}-a|^2=(1-a)^2+2a(1-\cos\theta)$ reduces the case $a\geq1/2$ to the integrable one-dimensional logarithm, while $a<1/2$ is bounded. The angular mean of the last term in \eqref{pcsp:eq:partition-angular-log} is zero.

The fluctuation of $\log(r\vee q)$ is bounded in $L^p$ by twice that of $\log(r/R_0)$, since the maximum is a contraction as a function of $\log r$. Equation~\eqref{pcsp:eq:partition-tail-exponential} bounds this last norm by $Cp$. Thus
\begin{equation}\label{pcsp:eq:partition-tail-centered-moment}
 \sup_{y\in E}\|\g(X-y)-\E[\g(X-y)]\|_{L^p(\nu_\infty)}\leq Cp.
\end{equation}
Conversely, for $Y$ in any collar cell, uniformly in $x$, the same bound holds for $\g(x-Y)-\E[\g(x-Y)]$. If $x$ is within $O(h)$ of the cell, rescaling removes $-\log h$ and gives the bounded-cell logarithmic estimate. Otherwise the oscillation is at most $Ch/\operatorname{dist}(x,A_a)$, with the intermediate range covered by the same rescaling. Applying Jensen to these centered estimates proves \eqref{pcsp:eq:partition-near-moments} for tail--tail and tail--collar pairs with one constant $C_{\rm exc}$ independent of $L,R_0,h$ once $R_0$ is large and $Lh\leq R_0/4$.

For non-near core cells, the mixed second difference gives $|r_\be|\leq C/L^2$. For a tail--core pair outside the collar, center first in the core cell. Its distance from the tail is at least $Lh$, so this first difference is at most $C/L$; centering in the tail multiplies that bound by at most two. The full remaining far kernel is therefore globally canonical and satisfies
\begin{equation}\label{pcsp:eq:partition-far-bound}
 \|r_{\be,\mathrm{far}}\|_\infty\leq C/L.
\end{equation}
Finally, the Gaussian bound gives
\begin{equation}\label{pcsp:eq:partition-cell-masses}
 \sup_\be m_\be(E)\leq C e^{-R_0^2/2}(1+R_0Lh),
 \qquad
 \sup_{\be,a\ne\infty}m_\be(A_a)\leq Ch^2.
\end{equation}
The conditional moment constants were fixed before these masses are made small.

\step{3. Coarse covariance and its diagonal correction.}
Let $Q$ be conditional expectation onto the cell-constant functions in $L^2(m_\be)$, and use its restriction to the mean-zero subspace. Define
\begin{equation}\label{pcsp:eq:partition-operators}
 \begin{aligned}
 B&=QT_\be Q, & A&=T_\be QB^{-1}QT_\be,\\
 R&=(\Id-Q)T_\be(\Id-Q), &
 D&=(\Id-Q)T_\be QB^{-1}QT_\be(\Id-Q).
 \end{aligned}
\end{equation}
The inverse is on the cell-constant mean-zero space. The logarithmic energy is strictly positive on its nonzero elements, so $B$ is invertible. The block calculation of \Cref{prop:regression-covariance-identity} gives
\begin{equation}\label{pcsp:eq:regression-decomposition}
 T_\be=A+R-D,
 \qquad
 0\leq A\leq T_\be,
 \qquad
 D\geq0.
\end{equation}
The centered kernel of $R$ is exactly \eqref{pcsp:eq:partition-double-centering}: the one-body centering terms cancel.

The functions spanning the range of $T_\be Q$ are potentials of zero-mass densities bounded by a common Gaussian, up to a bounded constant. Their potentials are bounded: on bounded sets, use the integrable logarithmic singularity and the uniform density bound; at infinity, subtract $-\log|x|$ and split the convolution at $|y|=|x|/2$. The inner part is controlled by the first moment and the outer part by Gaussian decay. Hence $A$ has bounded continuous features, and $D$ has bounded features. After the finite partition is fixed, continuity in $\be$ and the positive cell masses give a uniform bound on $B^{-1}$ and on these features. In particular, with $d(x,y)=\mathbf d(x)\cdot\mathbf d(y)$,
\begin{equation}\label{pcsp:eq:partition-diagonal-cost}
 -\frac{\alpha_\be}{N}\sum_{i<j}d(X_i,X_j)
 \leq \frac{\alpha_\be}{2}\sup_x d(x,x)\leq C_{\mathscr P}.
\end{equation}

Write $a(x,y)=\mathbf a(x)\cdot\mathbf a(y)$. The entropy--energy inequality and $A\leq T_\be$ yield
\begin{equation}\label{pcsp:eq:partition-coarse-entropy}
 \Ent(\rho\mid m_\be)
 \geq \frac{\Lambda_\be}{2}
       \left|\int\mathbf a\,\dd\rho\right|^2.
\end{equation}
Initially this follows for bounded densities relative to $m_\be$. For a general finite-entropy density, truncate that relative density and renormalize. The entropies converge, and the bounded feature averages converge, which proves \eqref{pcsp:eq:partition-coarse-entropy} on the full variational domain. Entropy duality and the Hubbard--Stratonovich transformation, as in \Cref{prop:coarse-partition-estimate}, give
\begin{equation}\label{pcsp:eq:partition-coarse-bound}
 \E\left[\exp\left\{\frac\gamma N\sum_{i<j}a(X_i,X_j)\right\}\right]
 \leq(1-\gamma/\Lambda_\be)^{-\operatorname{rank}(A)/2}
 \qquad(0\leq\gamma<\Lambda_\be).
\end{equation}

\step{4. Conditional residual estimate.}
Choose $p>1$ with $p\alpha_\be<\Lambda_\be$ uniformly, and write $q=p/(p-1)$. Choose also $r>1$ with $r\sup_\be\alpha_\be<4$, and let $r'$ be conjugate to $r$. These choices use only \eqref{pcsp:eq:source-buffer-hypotheses}. For the residual at coefficient $q\alpha_\be$, use H\"older with exponent three to separate the near-core, exceptional, and far parts.

First choose $L$ so large that \eqref{pcsp:eq:partition-far-bound} lies within the canonical-array exponential range of \cref{pcsp:lem:canonical-array-nonidentical} at coefficient $3q\sup_\be\alpha_\be$. The near-core graph then has a fixed edge-coloring number $\chi_L$. Choose $\delta,\delta_{\rm exc}>0$ so small that
\begin{equation}\label{pcsp:eq:partition-local-temperatures}
 6q\chi_L\sup_\be\alpha_\be\, C_{\rm can}C_L\delta<\tfrac12,
 \qquad
 3q\sup_\be\alpha_\be\, C_{\rm can}C_{\rm exc}\delta_{\rm exc}<\tfrac12.
\end{equation}
Consider the event
\begin{equation}\label{pcsp:eq:partition-good-event}
 \mathcal G=\{\max_{a\ne\infty}n_a\leq\delta N,
                     \ n_E\leq\delta_{\rm exc}N\}.
\end{equation}
Condition on the complete cell-label vector. The particle positions are independent with their conditional cell laws. Within a near-edge color class the arrays use disjoint cells. Each edge involves at most $2\delta N$ variables, and its effective coefficient is at most $6q\chi_L\alpha_\be\delta$. The first inequality in \eqref{pcsp:eq:partition-local-temperatures} and \eqref{pcsp:eq:partition-near-moments} therefore bound every edge factor.

For the exceptional part, use one array on all $n_E$ particles, setting its unused pair kernels to zero. Its effective coefficient is $3q\alpha_\be n_E/N\leq3q\alpha_\be\delta_{\rm exc}$. The second inequality in \eqref{pcsp:eq:partition-local-temperatures} bounds this entire array at once. In particular, there is no factor proportional to the number of collar cells. The far kernel is handled on all $N$ variables by \eqref{pcsp:eq:partition-far-bound}. After averaging over labels,
\begin{equation}\label{pcsp:eq:partition-residual-good}
 \sup_{N,\be}\E\left[\mathbf1_{\mathcal G}
       \exp\left\{\frac{q\alpha_\be}{N}\sum_{i<j}r_\be(X_i,X_j)\right\}\right]
 <\infty.
\end{equation}
The number of near edges affects the constant, but the partition is fixed before $N$ varies.

\step{5. Occupancy tails and assembly.}
Take $K>C_{r\sup\alpha_\be,\be_0}r'/r$, with the constant from \eqref{pcsp:eq:partition-linear-bound}. By \eqref{pcsp:eq:partition-cell-masses}, first increasing $R_0$ and then decreasing $h$ makes
\begin{equation}\label{pcsp:eq:partition-Chernoff-choice}
 \mathsf D\!\left(\delta_{\rm exc}\,\middle\Vert\,
              \sup_\be m_\be(E)\right)>K,
 \qquad
 \mathsf D\!\left(\delta\,\middle\Vert\,
              \sup_{\be,a\ne\infty}m_\be(A_a)\right)>K,
\end{equation}
where $\mathsf D$ is the binary relative entropy \eqref{eq:binary-relative-entropy}. Neither choice changes the conditional moment constants. The binomial Chernoff estimate and a union bound give $\Pb(\mathcal G^c)\leq C_{\mathscr P}e^{-KN}$. Thus the full, undecomposed exponent satisfies
\begin{equation}\label{pcsp:eq:partition-bad-bound}
 \E\left[\mathbf1_{\mathcal G^c}e^{\alpha_\be N^{-1}\sum_{i<j}h_\be(X_i,X_j)}\right]
 \leq C_{\mathscr P}^{1/r'}
       \exp\left\{\left(\frac{C_{r\sup\alpha_\be,\be_0}}r
                              -\frac K{r'}\right)N\right\}.
\end{equation}
On $\mathcal G$, use \eqref{pcsp:eq:regression-decomposition}, \eqref{pcsp:eq:partition-diagonal-cost}, and H\"older with $p,q$. Equations \eqref{pcsp:eq:partition-coarse-bound} and \eqref{pcsp:eq:partition-residual-good} bound this part uniformly as well. This proves the order-one bound under \eqref{pcsp:eq:source-buffer-hypotheses}. Since all choices depend only on the common equilibrium bounds and the strict margins, \eqref{pcsp:eq:source-buffer-margins} proves \eqref{pcsp:eq:equilibrium-buffer} uniformly in both $\be$ and $s$.
\end{proof}

\begin{remark}[Equilibrium-centered finite-partition input]\label{pcsp:rem:partition-dependency}
The coarse estimate uses the global entropy--energy margin. The near and exceptional arrays instead use their small local particle fractions; their global H\"older coefficients need not lie below the collision threshold. Only the independent linear bound on the log partition function for the bad occupancy event requires a global coupling below $4$.
\end{remark}

\begin{lemma}[Canonical arrays under possibly different laws]\label{pcsp:lem:canonical-array-nonidentical}
Let $X_1,\ldots,X_m$ be independent, with no identical-distribution
assumption.  For each $i<j$, let $r_{ij}(X_i,X_j)$ be a real kernel
centered in each variable under the corresponding marginal law.  Suppose
that, for every integer $p\geq1$,
\begin{equation}
 \mathop{\mathrm{ess\,sup}}_y\E|r_{ij}(X_i,y)|^p
 +\mathop{\mathrm{ess\,sup}}_x\E|r_{ij}(x,X_j)|^p
 \leq 2L^p p!.
\end{equation}
Zero kernels are permitted.  With $S=\sum_{i<j}r_{ij}(X_i,X_j)$ and
$C_{\rm can}=16e$, one has, for every integer $k\geq1$,
\begin{equation}
 \E|S|^{2k}\leq (C_{\rm can}Lm)^{2k}(2k)!.
\end{equation}
Consequently, for $m\geq2$ and $|\vartheta|C_{\rm can}L<1$,
\begin{equation}
 \E\exp\!\left(\frac{|\vartheta S|}{m}\right)
 \leq \frac{2}{1-(C_{\rm can}L\vartheta)^2}.
\end{equation}
For $m=0,1$ the pair sum is zero.
\end{lemma}

\begin{proof}
Put $q=2k$.  Expand $\E S^q$ by edge multiplicities
$a_{ij}\geq0$, $\sum_{i<j}a_{ij}=q$.  Every monomial with a vertex of
total degree one has expectation zero, by independence and centering in
that variable.  Thus a surviving monomial has $v\leq\min(q,m)$ active
vertices.  For each fixed monomial, eliminate the active vertices in
increasing index order and put $d_i=\sum_{j>i}a_{ij}$.  At a step with
$d_i>0$, conditional H\"older gives
\begin{equation}
 \E_i\prod_{j>i}|r_{ij}(X_i,x_j)|^{a_{ij}}
 \leq\prod_{j:a_{ij}>0}
       \bigl(\E_i|r_{ij}(X_i,x_j)|^{d_i}\bigr)^{a_{ij}/d_i}
 \leq 2L^{d_i}d_i!.
\end{equation}
The estimate is uniform in all remaining variables.  A step with $d_i=0$
costs one.  Iterating proves absolute integrability of each monomial and
bounds its absolute expectation by $(2L)^q\prod_i d_i!$.
In particular, the preceding uses of centering and Fubini are justified.

After division by $q!$, the multiplicity coefficient is
$\prod_i d_i!/\prod_{i<j}a_{ij}!$.  For a fixed active set and fixed
outdegrees $d_i$, this counts ordered outgoing neighbor lists.  There
are at most $v^q$ such lists; allowing earlier neighbors and loops only
increases this upper bound.  There are
$\binom{q+v-1}{v-1}$ possible outdegree lists.  Dropping every remaining
restriction except $v\leq q$ therefore yields
\begin{equation}
 \frac{\E S^q}{q!}
 \leq (2L)^q\sum_{v=1}^{\min(q,m)}
       \binom mv\binom{q+v-1}{v-1}v^q.
\end{equation}
For $v\leq\min(q,m)$,
\begin{equation}
 \binom mv v^q
 \leq e^v m^q(v/m)^{q-v}\leq e^q m^q,
 \qquad
 \binom{q+v-1}{v-1}\leq4^q.
\end{equation}
There are at most $q\leq2^q$ summands.  This gives the asserted moment
bound with $C_{\rm can}=16e$.  Finally, $e^{|u|}\leq2\cosh u$ and Tonelli's
theorem for the nonnegative even-power series give the displayed
exponential bound.  This proof uses neither identical marginal laws,
translation invariance, nor a decoupling theorem.
\end{proof}

\subsection{Exact transfer from equilibrium to a dominated reference}

A direct product-density comparison between $\nu^{\otimes N}$ and $m^{\otimes N}$ would cost a factor exponential in $N$ and is therefore unusable.  The following replacement-channel identity is exact and avoids that loss.

\begin{lemma}[Bernoulli replacement channel]\label{pcsp:lem:replacement-channel}
Let $m$ and $\nu$ be probability measures on a measurable space, and let $0<\theta\leq1$ satisfy
\begin{equation}\label{pcsp:eq:domination}
 \nu\leq\theta^{-1}m.
\end{equation}
Let $h_m$ be a symmetric integrable kernel canonical under $m$:
\begin{equation}\label{pcsp:eq:hm-canonical}
 \int h_m(x,y)\,m(\dd y)=0
 \quad\text{for $m$-almost every $x$}.
\end{equation}
For $\theta<1$, define
\begin{equation}\label{pcsp:eq:lambda}
 \lambda=\frac{m-\theta\nu}{1-\theta};
\end{equation}
for $\theta=1$, set $\lambda=m$.  Given $X\sim\nu$, let $Y$ have conditional law
\begin{equation}\label{pcsp:eq:channel}
 \mathsf K(x,\dd y)=\theta\delta_x(\dd y)+(1-\theta)\lambda(\dd y).
\end{equation}
Then $Y\sim m$.  For two conditionally independent outputs $Y_1,Y_2$ given $X_1=x,X_2=y$,
\begin{equation}\label{pcsp:eq:conditional-kernel-identity}
 \E[h_m(Y_1,Y_2)\mid X_1=x,X_2=y]
 =\theta^2h_\nu(x,y),
\end{equation}
where
\begin{equation}\label{pcsp:eq:hnu-double-centered}
 h_\nu(x,y)
 =h_m(x,y)-\int h_m(x,z)\,\nu(\dd z)
 -\int h_m(z,y)\,\nu(\dd z)
 +\iint h_m(z,w)\,\nu(\dd z)\nu(\dd w).
\end{equation}
\end{lemma}

\begin{proof}
The domination \eqref{pcsp:eq:domination} makes \eqref{pcsp:eq:lambda} a probability measure and gives
\begin{equation}\label{pcsp:eq:mixture}
 m=\theta\nu+(1-\theta)\lambda.
\end{equation}
Hence the output of \eqref{pcsp:eq:channel} has law $m$.

Set
\begin{equation}\label{pcsp:eq:A-nu-lambda}
 A_\nu(x)=\int h_m(x,z)\,\nu(\dd z),
 \qquad
 A_\lambda(x)=\int h_m(x,z)\,\lambda(\dd z),
\end{equation}
and
\begin{equation}\label{pcsp:eq:c-nu}
 c_\nu=\int A_\nu\,\dd\nu=\iint h_m\,\dd\nu\dd\nu.
\end{equation}
Canonicality under $m$ and \eqref{pcsp:eq:mixture} imply
\begin{equation}\label{pcsp:eq:A-relation}
 \theta A_\nu+(1-\theta)A_\lambda=0.
\end{equation}
Consequently,
\begin{align}
 \iint h_m\,\dd\nu\dd\lambda
 &=-\frac\theta{1-\theta}c_\nu,
 \label{pcsp:eq:cross-mixture}\\
 \iint h_m\,\dd\lambda\dd\lambda
 &=\frac{\theta^2}{(1-\theta)^2}c_\nu.
 \label{pcsp:eq:lambda-square}
\end{align}
Expanding the four outcomes of the two independent Bernoulli replacements and using \eqref{pcsp:eq:A-relation}--\eqref{pcsp:eq:lambda-square} gives
\begin{align}
 \E[h_m(Y_1,Y_2)\mid X_1=x,X_2=y]
 &=\theta^2\left[h_m(x,y)-A_\nu(x)-A_\nu(y)+c_\nu\right]
 \label{pcsp:eq:channel-expansion}\\
 &=\theta^2h_\nu(x,y),
 \label{pcsp:eq:channel-expansion-final}
\end{align}
which is \eqref{pcsp:eq:conditional-kernel-identity}.  The case $\theta=1$ is immediate.
\end{proof}

\begin{prop}[Dominated-reference Laplace transfer]\label{pcsp:prop:laplace-transfer}
Under the assumptions of \cref{pcsp:lem:replacement-channel}, let $X_1,\ldots,X_N$ be independent with law $\nu$, and, conditionally on them, let $Y_1,\ldots,Y_N$ be independent outputs of \eqref{pcsp:eq:channel}.  For every $s\geq0$ for which the right side is finite,
\begin{equation}\label{pcsp:eq:abstract-laplace-transfer}
 \E_{\nu^{\otimes N}}\exp\left\{
 \frac sN\sum_{i<j}h_\nu(X_i,X_j)\right\}
 \leq
 \E_{m^{\otimes N}}\exp\left\{
 \frac{s}{\theta^2N}\sum_{i<j}h_m(Y_i,Y_j)\right\}.
\end{equation}
\end{prop}

\begin{proof}
Summing \eqref{pcsp:eq:conditional-kernel-identity} over unordered pairs gives
\begin{equation}\label{pcsp:eq:sum-conditional}
 \sum_{i<j}h_\nu(X_i,X_j)
 =\theta^{-2}\E\left[
 \sum_{i<j}h_m(Y_i,Y_j)\,\bigg|\,X_1,\ldots,X_N\right].
\end{equation}
Conditional Jensen applied to the exponential, followed by integration, yields \eqref{pcsp:eq:abstract-laplace-transfer}.  This applies directly to the unbounded kernel: integrability of $h_m$ under $m\otimes m$ and the independent output law $m^{\otimes N}$ make the output pair sum integrable.  Domination also gives integrability under every input and replacement product law; conditional expectations and identities are understood almost surely.  Jensen is valid with an extended-valued right side, and the finite equilibrium exponential moment makes that side finite in the application.  No clipping of a canonical kernel is needed, since clipping without recentering would generally destroy canonicality.
\end{proof}

\begin{cor}[Centered logarithmic transfer]\label{pcsp:cor:log-transfer}
For the kernel \eqref{pcsp:eq:kernel}, the double centering of $\widetilde\g_m$ under $\nu$ is exactly $\widetilde\g_\nu$.  Consequently, if \eqref{pcsp:eq:domination} holds, then
\begin{equation}\label{pcsp:eq:log-laplace-transfer}
 \K_N(s;\nu)\leq\K_N(s/\theta^2;m).
\end{equation}
\end{cor}

\begin{proof}
Set
\begin{equation}\label{pcsp:eq:a-cross-potential}
 a=\int U_m\,\dd\nu.
\end{equation}
Then
\begin{equation}\label{pcsp:eq:A-explicit}
 \int\widetilde\g_m(x,z)\,\nu(\dd z)
 =U_\nu(x)-U_m(x)-a+\Gcal(m),
\end{equation}
and
\begin{equation}\label{pcsp:eq:c-explicit}
 \iint\widetilde\g_m(z,w)\,\nu(\dd z)\nu(\dd w)
 =\Gcal(\nu)-2a+\Gcal(m).
\end{equation}
Substitution in \eqref{pcsp:eq:hnu-double-centered} gives \eqref{pcsp:eq:centered-kernel} with $\mu=\nu$.  Proposition~\ref{pcsp:prop:laplace-transfer} then gives \eqref{pcsp:eq:log-laplace-transfer}.
\end{proof}

\subsection{Uniform constants for the reference-neighborhood bound}

Let $\eta_0$, $\delta_0$, and $C_0$ be supplied by \cref{pcsp:prop:global-gap,pcsp:prop:equilibrium-buffer}.  Define
\begin{equation}\label{pcsp:eq:eta-choice}
 \eta_{\rm eq}
 =\min\left\{\frac{\eta_0}{2},
               \frac{4-\be_0}{2},
               \frac{2\delta_0}{9}\right\}.
\end{equation}
Then \eqref{pcsp:eq:F-EQ} follows from \eqref{pcsp:eq:global-gap}.  Since $\be_0>0$, the minimum in \eqref{pcsp:eq:frozen-radius} is attained by the second term, so
\begin{equation}\label{pcsp:eq:r-explicit}
 r_*=\frac{\eta_{\rm eq}}{4(\eta_{\rm eq}+\be_0)}<\frac14.
\end{equation}
Set
\begin{equation}\label{pcsp:eq:delta-choice}
 \delta
 =\min\left\{\frac{4-\be_0}{2},
 \frac{\delta_0}{2(1+r_*)^2}\right\}.
\end{equation}
This is positive and belongs to $(0,4-\be_0)$.

Suppose \eqref{pcsp:eq:tube} holds.  Then
\begin{equation}\label{pcsp:eq:upper-domination-tube}
 \nu\leq(1+r_*)m_\be.
\end{equation}
Apply \cref{pcsp:cor:log-transfer} with
\begin{equation}\label{pcsp:eq:theta-choice}
 \theta=(1+r_*)^{-1}.
\end{equation}
It remains only to check that the amplified coupling lies within the interval of the equilibrium partition bound.  Since $r_*\leq1/4$,
\begin{equation}\label{pcsp:eq:amplification-elementary}
 (1+r_*)^2-1=2r_*+r_*^2\leq\frac94r_*.
\end{equation}
Also, by \eqref{pcsp:eq:r-explicit},
\begin{equation}\label{pcsp:eq:r-upper-eta}
 r_*\leq\frac{\eta_{\rm eq}}{4\be_0}.
\end{equation}
Therefore
\begin{equation}\label{pcsp:eq:beta-amplification-cost}
 \be\left[(1+r_*)^2-1\right]
 \leq\be_0\left[(1+r_*)^2-1\right]
 \leq\frac{9\eta_{\rm eq}}{16}
 \leq\frac{\delta_0}{8}.
\end{equation}
By \eqref{pcsp:eq:delta-choice},
\begin{equation}\label{pcsp:eq:delta-amplification-cost}
 \delta(1+r_*)^2\leq\frac{\delta_0}{2}.
\end{equation}
Combining \eqref{pcsp:eq:beta-amplification-cost} and \eqref{pcsp:eq:delta-amplification-cost},
\begin{equation}\label{pcsp:eq:amplified-coupling}
 (\be+\delta)(1+r_*)^2
 \leq\be+\frac{5\delta_0}{8}
 \leq\be+\delta_0.
\end{equation}
Hence \eqref{pcsp:eq:log-laplace-transfer} and \eqref{pcsp:eq:equilibrium-buffer} give
\begin{align}
 \K_N(\be+\delta;\nu)
 &\leq\K_N((\be+\delta)(1+r_*)^2;m_\be)
 \label{pcsp:eq:transfer-application}\\
 &\leq C_0.
 \label{pcsp:eq:transfer-application-final}
\end{align}
Thus \eqref{pcsp:eq:F-BUF} holds with
\begin{equation}\label{pcsp:eq:Cbuf-choice}
 C_{\rm buf}=C_0.
\end{equation}
Together with \cref{pcsp:prop:equilibrium-family,pcsp:prop:global-gap}, this proves \cref{thm:planar-confined-static-package}.

\begin{remark}[One-sided domination suffices for the moment transfer]\label{pcsp:rem:one-sided}
The proof of \eqref{pcsp:eq:F-BUF} uses only $\nu\leq(1+r_*)m_\be$.  The lower bound in the two-sided neighborhood also controls $m_\be/\nu$ by $M_*=(1-r_*)^{-1}$, as required for the dynamical applications in Part~II.
\end{remark}

\begin{remark}[Why no exponential-in-$N$ loss occurs]\label{pcsp:rem:no-exp-loss}
The product density ratio satisfies only the crude bound
\begin{equation}\label{pcsp:eq:bad-direct-comparison}
 \frac{\dd\nu^{\otimes N}}{\dd m_\be^{\otimes N}}
 \leq(1+r_*)^N,
\end{equation}
which grows exponentially with $N$.  The replacement channel instead transfers the \emph{centered kernel itself}; the only cost is the coupling amplification $(1+r_*)^2$, independent of $N$.
\end{remark}

\subsection{Normalization of the static input to Part~II}

For comparison with the conventions of Part~II, note that Part~II uses
\begin{equation}\label{pcsp:eq:W-convention}
 W(x,y)=-\g(x-y),
 \qquad
 h_\nu=W*\nu=-U_\nu,
 \qquad
 I_\nu=\frac12\iint W\,\dd\nu\dd\nu=-\frac12\Gcal(\nu).
\end{equation}
Its centered interaction is
\begin{equation}\label{pcsp:eq:W-centered}
 \widetilde W_\nu(x,y)
 =W(x,y)-h_\nu(x)-h_\nu(y)+2I_\nu
 =-\widetilde\g_\nu(x,y).
\end{equation}
Therefore the centered moment hypothesis in Part~II
\begin{equation}\label{pcsp:eq:manuscript-centered-hypothesis}
 \E_{\nu^{\otimes N}}\exp\left\{
 -\frac{\be+\delta}{N}\sum_{i<j}\widetilde W_\nu(x_i,x_j)
 \right\}\leq C_{\rm buf}
\end{equation}
is exactly \eqref{pcsp:eq:F-BUF}.  No factor of two is introduced: both formulas use unordered pairs $i<j$, and the coefficient is $(\be+\delta)/N$.

The scaled equilibrium equation in Part~II is \eqref{pcsp:eq:euler-lagrange}.  Its comparison density is \eqref{pcsp:eq:comparison-density}, and \cref{pcsp:prop:equilibrium-family} supplies it with $R=1$.  Thus the constants in \cref{thm:planar-confined-static-package} give the confined equilibrium and reference estimates used in Part~II.


\bibliographystyle{amsalpha}
\bibliography{references}

\end{document}
